% !TeX program = lualatex
\documentclass[11pt,leqno]{article}
\usepackage{amsmath,amssymb}
\usepackage{fontspec}
\usepackage{unicode-math}
\setmainfont{STIXTwoText}[Extension=.otf,UprightFont=*-Regular,ItalicFont=*-Italic,BoldFont=*-Bold,BoldItalicFont=*-BoldItalic]
\setmathfont{STIXTwoMath-Regular.otf}
\usepackage[main=french]{babel}
\usepackage{tikz-cd}
\usetikzlibrary{babel}
\usepackage[a4paper,inner=27mm,outer=23mm,top=25mm,bottom=27mm]{geometry}
\usepackage{xcolor}
\usepackage[hidelinks,unicode]{hyperref}
\setlength{\parindent}{0pt}
\setlength{\parskip}{0.55\baselineskip}
\setlength{\emergencystretch}{3em}
\newcommand{\Page}[1]{\ifvmode\noindent\else\unskip\ \fi{\footnotesize\textcolor{gray}{[p.\,#1]}}\ \ignorespaces}
\title{Théorie de Hodge, III}
\author{Pierre Deligne}
\date{}
\begin{document}
\begin{center}\small\itshape Version corrigée : les 53 coquilles de l'imprimé relevées dans corrections/Deligne\_1974\_Theorie\_de\_Hodge\_III\_misprints\_corrected.csv y sont corrigées. L'édition principale est le texte tel qu'imprimé.\end{center}

\maketitle
% ---- 005-012 ----
\Page{5}

\begin{center}
{\Large THÉORIE DE HODGE, III\par}
\medskip
{\large \textit{par} \textsc{Pierre} DELIGNE\par}
\end{center}

\section*{\centering\normalfont SOMMAIRE}

\begingroup
\small
\noindent\textsc{Introduction}\dotfill 6\par
\medskip
\noindent\textsc{Terminologie et notations}\dotfill 7\par
\medskip
\noindent 5. \textbf{Descente cohomologique}\dotfill 8\par
\noindent\hspace*{1.5em}1. Espaces topologiques simpliciaux\dotfill 8\par
\noindent\hspace*{1.5em}2. Cohomologie des espaces topologiques simpliciaux\dotfill 12\par
\noindent\hspace*{1.5em}3. Descente cohomologique\dotfill 14\par
\medskip
\noindent 6. \textbf{Exemples d’espaces topologiques simpliciaux}\dotfill 16\par
\noindent\hspace*{1.5em}1. Espaces classifiants\dotfill 16\par
\noindent\hspace*{1.5em}2. Construction d’hyperrecouvrements\dotfill 19\par
\noindent\hspace*{1.5em}3. Cohomologie relative\dotfill 21\par
\noindent\hspace*{1.5em}4. Espaces multisimpliciaux\dotfill 22\par
\medskip
\noindent 7. \textbf{Compléments au § 1}\dotfill 23\par
\noindent\hspace*{1.5em}1. Catégorie dérivée filtrée\dotfill 23\par
\noindent\hspace*{1.5em}2. Compléments au lemme des deux filtrations\dotfill 25\par
\medskip
\noindent 8. \textbf{Théorie de Hodge des espaces algébriques}\dotfill 28\par
\noindent\hspace*{1.5em}1. Complexes de Hodge\dotfill 28\par
\noindent\hspace*{1.5em}2. Espaces algébriques séparés\dotfill 38\par
\noindent\hspace*{1.5em}3. Théorie de Hodge des schémas simpliciaux\dotfill 41\par
\medskip
\noindent 9. \textbf{Exemples et applications}\dotfill 45\par
\noindent\hspace*{1.5em}1. Cohomologie des groupes et des classifiants\dotfill 45\par
\noindent\hspace*{1.5em}2. Théorie de Hodge des hypersurfaces lisses, d’après Griffiths\dotfill 48\par
\noindent\hspace*{1.5em}3. Construction de complexes d’opérateurs différentiels du 1\textsuperscript{er} ordre\dotfill 50\par
\medskip
\noindent 10. \textbf{\boldmath Théorie de Hodge en niveau \(\leq 1\)}\dotfill 53\par
\noindent\hspace*{1.5em}1. 1-motifs\dotfill 53\par
\noindent\hspace*{1.5em}2. 1-motifs et biextensions\dotfill 60\par
\noindent\hspace*{1.5em}3. Interprétation algébrique du \(\mathrm{H}^{1}\) mixte : cas des courbes\dotfill 69\par
\noindent\hspace*{1.5em}4. Traduction d’un théorème de Picard\dotfill 75\par
\medskip
\noindent\textsc{Bibliographie}\dotfill 77\par
\endgroup

\Page{6}

\section*{\centering Introduction}

Cet article fait suite à [4] et [5], cités I et II. La numérotation de ses paragraphes prolonge celle de II (qui contient les §§ 1 à 4). Pour renvoyer à un résultat de II, on mentionne simplement son numéro. Par exemple : (3.2.17).

Dans I, nous avons exposé le yoga qui sous-tend II et le présent article III. Toutefois, II et III sont logiquement indépendants de I. Ils ne contiennent pas tous les résultats qui sont annoncés dans I.

Dans II, nous avons exposé la théorie de Hodge des variétés algébriques non singulières (non nécessairement complètes). On traite ici le cas de singularités quelconques. A partir du § 7, un usage essentiel sera fait des résultats des §§ 1 à 3 de II.

Dans le cas d’une variété algébrique singulière \textit{complète} \(\mathrm{X}\), l’idée fondamentale est d’utiliser la résolution des singularités pour « remplacer » \(\mathrm{X}\) par un système simplicial de schémas projectifs et lisses \(\mathrm{X}_{\bullet}\)
\[
\begin{tikzcd}[column sep=2.6em]
\ldotp\,\ldotp\,\ldotp\,\ldotp\,\ldotp\ \mathrm{X}_2
  \arrow[r, shift left=3.2ex]
  \arrow[r]
  \arrow[r, shift right=3.2ex]
&
\mathrm{X}_1
  \arrow[l, shift right=1.6ex, shorten <=1.1em]
  \arrow[l, shift left=1.6ex, shorten <=1.1em]
  \arrow[r, shift left=1.6ex]
  \arrow[r, shift right=1.6ex]
&
\mathrm{X}_0
  \arrow[l, shorten <=1.1em]
\end{tikzcd}
\]
(ou, comme nous dirons, par un schéma simplicial projectif et lisse). Les résultats des §§ 5 et 6 permettent de trouver de tels schémas simpliciaux \(\mathrm{X}_{\bullet}\) ayant, en un sens convenable, même cohomologie que \(\mathrm{X}\). On peut alors « exprimer », \textit{via} une suite spectrale, la cohomologie de \(\mathrm{X}\) en terme de celle des \(\mathrm{X}_n\). En conformité avec les principes généraux de I, on peut alors, sur chacun des \(\mathrm{X}_n\), utiliser la théorie de Hodge classique, et en déduire une structure de Hodge mixte sur la cohomologie de \(\mathrm{X}\).

Dans le cas d’une variété algébrique quelconque \(\mathrm{X}\), on « remplace » \(\mathrm{X}\) par un schéma simplicial lisse \(\mathrm{X}_{\bullet}\), compactifié par un schéma simplicial projectif et lisse \(\overline{\mathrm{X}}_{\bullet}\). On se débrouille de plus pour que \(\overline{\mathrm{X}}_n-\mathrm{X}_n\) soit un diviseur à croisements normaux, réunion de diviseurs lisses \(\mathrm{D}_{n,i}\). On « exprime » alors \textit{via} une suite spectrale la cohomologie de \(\mathrm{X}\) en terme de la cohomologie des intersections \(p\) à \(p\) des \(\mathrm{D}_{n,i}\) (\(n\geq 0,\ p\geq 0\)), et on en déduit une structure de Hodge mixte sur \(\mathrm{H}^{\bullet}(\mathrm{X})\).

Les §§ 5 et 6 constituent un résumé (sans démonstrations) de la théorie de la « descente cohomologique ». Cette théorie, dans le cadre plus général des topos, est exposée par B. Saint-Donat dans \textit{S.G.A.}4, V \textit{bis}.

Au § 7 n\textsuperscript{o} 1, on reprend dans le langage des catégories dérivées filtrées certains des résultats obtenus antérieurement. Le n\textsuperscript{o} 2 contient une démonstration du lemme des deux filtrations plus simple que celle du § 1, n\textsuperscript{o} 3.

Le rebutant n\textsuperscript{o} 1 du § 8 est le cœur de ce travail. Le résultat essentiel est le théorème (8.1.15). Dans sa démonstration, pour pouvoir appliquer le lemme des deux
\Page{7}
filtrations, on utilise que tout morphisme de structures de Hodge mixtes est \textit{strictement} compatible aux filtrations \(\mathrm{W}\) et \(\mathrm{F}\) (2.3.5). Au n\textsuperscript{o} 2, on définit la structure de Hodge mixte de \(\mathrm{H}^n(\mathrm{X},\mathbf{Z})\) (\(\mathrm{X}\) schéma séparé). On montre que les nombres de Hodge \(h^{pq}\) de \(\mathrm{H}^n(\mathrm{X},\mathbf{Z})\) ne peuvent être non nuls que pour \((p,q)\in[0,n]\times[0,n]\). Pour \(n>\dim\mathrm{X}\), on a un résultat plus précis (8.2.4). Si \(\mathrm{X}\) est lisse (resp. complet) ils ne peuvent être non nuls que si de plus \(p+q\geq n\) (resp. \(p+q\leq n\)). Si \(\mathrm{X}\) est un sous-schéma complet de complément \(\mathrm{U}\) d’un schéma \(\mathrm{Z}\), complet et lisse de dimension \(n\), on peut, d’après N. Katz, interpréter cette « dualité » entre les cas complets et lisses comme provenant de la « dualité d’Alexander »
\[
\ldots\ \mathrm{H}^i(\mathrm{Z},\mathbf{Q})
\to \mathrm{H}^i(\mathrm{X},\mathbf{Q})
\to \bigl(\mathrm{H}^{2n-i-1}(\mathrm{U},\mathbf{Q})\bigr)^*
\to \mathrm{H}^{i+1}(\mathrm{Z},\mathbf{Q})\ \ldots
\]

Des propriétés de fonctorialité de la théorie construite, on déduit les utiles corollaires (8.2.5) à (8.2.8), qui s’énoncent en termes indépendants de toute théorie de Hodge.

Au § 8, n\textsuperscript{o} 3, on munit la cohomologie de tout schéma simplicial \(\mathrm{X}_{\bullet}\) d’une structure de Hodge mixte. Cette généralité n’est pas illusoire.

\textit{a)} On peut interpréter des espaces de cohomologie relative comme étant la cohomologie de schémas simpliciaux convenables.

\textit{b)} Soit \(\mathrm{B}_{\mathrm{G}}\) l’espace classifiant du groupe de Lie sous-jacent à un groupe algébrique \(\mathrm{G}\). On peut interpréter la cohomologie de \(\mathrm{B}_{\mathrm{G}}\) comme étant celle d’un schéma simplicial convenable.

Au § 9, n\textsuperscript{o} 1, après avoir calculé la structure de Hodge mixte de la cohomologie de \(\mathrm{B}_{\mathrm{G}}\) (pour \(\mathrm{G}\) linéaire), on en déduit celle de \(\mathrm{G}\). Comme corollaire, on trouve que si un groupe linéaire \(\mathrm{G}\) agit sur une variété complète non singulière \(\mathrm{X}\), alors l’application
\[
\mathrm{H}^{\bullet}(\mathrm{X},\mathbf{Q})
\to \mathrm{H}^{\bullet}(\mathrm{X},\mathbf{Q})
\otimes \mathrm{H}^{\bullet}(\mathrm{G},\mathbf{Q})
\]
se factorise par \(\mathrm{H}^{\bullet}(\mathrm{X},\mathbf{Q})\otimes\mathrm{H}^{0}(\mathrm{G},\mathbf{Q})\subset\mathrm{H}^{\bullet}(\mathrm{X},\mathbf{Q})\otimes\mathrm{H}^{\bullet}(\mathrm{G},\mathbf{Q})\).

Au § 10, on interprète en terme de schémas abéliens les structures de Hodge mixtes purement de type \(\{(-1,-1),(-1,0),(0,-1),(0,0)\}\) (III.0.5), et on considère en détail le \(\mathrm{H}^{1}\) des courbes.

La théorie développée jusqu’ici est une théorie absolue (on n’y étudie pas les foncteurs \(\mathrm{R}f_*\)), et ne concerne que des coefficients constants. Je conjecture que si \(\mathscr{H}\) est une variation de structures de Hodge polarisable (au sens de Griffiths [6]) sur un schéma \(\mathrm{X}\), alors la cohomologie de \(\mathrm{X}\) à coefficients dans le système local \(\mathscr{H}\) est munie d’une structure de Hodge mixte naturelle. Je ne puis le prouver que pour \(\mathrm{X}\) complet.

\section*{\centering Terminologie et notations}

\textbf{(III.0.1)} Soit \(u:\mathrm{X}\to\mathrm{Y}\) une application continue entre espaces topologiques. Nous dirons que \(u\) est \textit{propre} si elle est propre au sens de Bourbaki (i.e. universellement fermée) et de plus séparée (i.e. si la diagonale de \(\mathrm{X}\times_{\mathrm{Y}}\mathrm{X}\) est fermée).

\Page{8}

\textbf{(III.0.2)} Suivant en cela Gabriel et Zisman, nous dirons \textit{simplicial} là où autrefois on disait semi-simplicial.

\textbf{(III.0.3)} Nous désignerons par \(\mathrm{A}\) un sous-anneau noethérien de \(\mathbf{R}\) tel que \(\mathrm{A}\otimes\mathbf{Q}\) soit un corps. Les cas utiles sont \(\mathrm{A}=\mathbf{Z}\), \(\mathbf{Q}\) ou \(\mathbf{R}\).

\textbf{(III.0.4)} Une \(\mathrm{A}\)-\textit{structure de Hodge mixte} consiste en un \(\mathrm{A}\)-module de type fini \(\mathrm{H}_{\mathrm{A}}\), une filtration finie croissante \(\mathrm{W}\) du \(\mathrm{A}\otimes\mathbf{Q}\)-vectoriel \(\mathrm{H}_{\mathrm{A}\otimes\mathbf{Q}}=\mathrm{H}_{\mathrm{A}}\otimes\mathbf{Q}\), et une filtration finie décroissante \(\mathrm{F}\) du \(\mathbf{C}\)-vectoriel \(\mathrm{H}_{\mathbf{C}}=\mathrm{H}_{\mathrm{A}}\otimes_{\mathrm{A}}\mathbf{C}\). On exige que les \((\operatorname{Gr}^{\mathrm{W}}_n(\mathrm{H}_{\mathrm{A}\otimes\mathbf{Q}}),\operatorname{Gr}^{\mathrm{W}}_n(\mathrm{F}))\) soient des \(\mathrm{A}\otimes\mathbf{Q}\)-structures de Hodge. Pour \(\mathrm{A}=\mathbf{Z}\) (resp. \(\mathbf{Q}\)) on retrouve (2.3.1) (resp. 2.3.8); les résultats de (2.3) se transposent tels quels.

\textbf{(III.0.5)} Soit \(\mathscr{E}\) une partie de \(\mathbf{Z}\times\mathbf{Z}\). Une structure de Hodge mixte \(\mathrm{H}\) est de \textit{type} \(\mathscr{E}\) si ses nombres de Hodge \(h^{pq}\) sont nuls pour \((p,q)\notin\mathscr{E}\).

\textbf{(III.0.6)} Nous noterons désormais \(\Omega_{\mathrm{X}}^{\bullet}(\log\mathrm{D})\) ce que nous notions \(\Omega_{\mathrm{X}}^*\langle\mathrm{D}\rangle\) (3.1.2).

\textbf{(III.0.7)} Sauf mention explicite du contraire, \textit{schéma} signifie « schéma de type fini sur \(\mathbf{C}\) », et un \textit{faisceau} sur un schéma \(\mathrm{X}\) est un faisceau sur l’espace topologique sous-jacent à \(\mathrm{X}_{\mathrm{an}}\).

\section*{\centering 5. DESCENTE COHOMOLOGIQUE}

\section*{5.1. Espaces topologiques simpliciaux.}

Ce numéro commence par des rappels pour lesquels on peut renvoyer au séminaire homotopique de Strasbourg, 63/64.

\textbf{(5.1.1)} On utilisera les notations suivantes.

\noindent \(\mathscr{A}^{0}=\) catégorie opposée d’une catégorie \(\mathscr{A}\).\par
\noindent \(\mathit{Hom}(\mathscr{A},\mathscr{B})=\) catégorie des foncteurs de \(\mathscr{A}\) dans \(\mathscr{B}\).\par

Désignons par \(n\) et \(k\) des entiers \(\geq-1\).

\noindent \(\Delta_n=\) l’ensemble fini totalement ordonné \([0,n]\).\par
\noindent \(\delta_i:\Delta_n\to\Delta_{n+1}=\) l’injection croissante telle que \(i\notin\delta_i(\Delta_n)\) (\(0\leq i\leq n+1\)).\par
\noindent \(s_i:\Delta_{n+1}\to\Delta_n=\) la surjection croissante telle que \(s_i(i)=s_i(i+1)\) (\(0\leq i\leq n\)).\par
\noindent \(\varepsilon:\Delta_{-1}\to\Delta_n=\) l’unique application de \(\Delta_{-1}\) dans \(\Delta_n\).\par
\noindent \((\Delta^+)=\) la catégorie ayant pour objets les \(\Delta_n\) (\(n\geq-1\)), et pour flèches les applications croissantes entre les \(\Delta_n\).\par
\noindent \((\Delta)=\) la sous-catégorie pleine de \((\Delta^+)\) d’objets les \(\Delta_n\) (\(n\geq0\)).\par
\noindent \((\Delta^+)_k=\) la sous-catégorie pleine de \((\Delta^+)\) d’objets les \(\Delta_n\) (\(k\geq n\)).\par
\noindent \((\Delta)_k=\) la sous-catégorie pleine de \((\Delta^+)\) d’objets les \(\Delta_n\) (\(k\geq n\geq0\)).\par

\Page{9}

Pour toute catégorie \(\mathscr{C}\), un \textit{objet simplicial} (resp. \textit{objet simplicial \(k\)-tronqué}) de \(\mathscr{C}\) est un objet de \(\mathit{Hom}((\Delta)^0,\mathscr{C})\) (resp. \(\mathit{Hom}((\Delta)_k^0,\mathscr{C})\)). De même, un \textit{objet cosimplicial} (resp. …) de \(\mathscr{C}\) est un objet de \(\mathit{Hom}((\Delta),\mathscr{C})\) (resp. …).

Le foncteur \textit{squelette} \(\operatorname{sq}_k\) est, pour \(k\geq-1\), le foncteur de restriction
\[
\operatorname{sq}_k:\mathit{Hom}((\Delta)^0,\mathscr{C})
\to\mathit{Hom}((\Delta)_k^0,\mathscr{C}).
\]
Le foncteur \textit{cosquelette} \(\operatorname{cosq}_k\) est le foncteur adjoint à droite à \(\operatorname{sq}_k\)
\[
\operatorname{cosq}_k:\mathit{Hom}((\Delta)_k^0,\mathscr{C})
\to\mathit{Hom}((\Delta)^0,\mathscr{C}).
\]
Soit \(\mathrm{Y}\) un objet simplicial de \(\mathscr{C}\). On appelle encore \textit{squelette} le foncteur
\[
\operatorname{sq}_k:\mathit{Hom}((\Delta)^0,\mathscr{C})/\mathrm{Y}
\to\mathit{Hom}((\Delta)_k^0,\mathscr{C})/\operatorname{sq}_k(\mathrm{Y}).
\]
Son adjoint à droite se note \(\operatorname{cosq}_k^{\mathrm{Y}}\) (cosquelette relatif à \(\mathrm{Y}\))
\[
\operatorname{cosq}_k^{\mathrm{Y}}:
\mathit{Hom}((\Delta)_k^0,\mathscr{C})/\operatorname{sq}_k\mathrm{Y}
\to\mathit{Hom}((\Delta)^0,\mathscr{C})/\mathrm{Y}.
\]

Les foncteurs cosquelettes sont définis si les limites projectives finies existent dans \(\mathscr{C}\). Les foncteurs \(\operatorname{cosq}_k^{\mathrm{Y}}\) existent dès qu’existent les produits fibrés. On a
\[
\operatorname{cosq}_k^{\mathrm{Y}}(\mathrm{X})
=\operatorname{cosq}_k(\mathrm{X})
\times_{\operatorname{cosq}_k\operatorname{sq}_k(\mathrm{Y})}\mathrm{Y}.
\]

\textbf{(5.1.2)} Si \(\mathrm{X}_{\bullet}:(\Delta)^0\to\mathscr{C}\) est un objet simplicial de \(\mathscr{C}\), on pose \(\mathrm{X}_n=\mathrm{X}_{\bullet}(\Delta_n)\), \(\delta_i=\mathrm{X}_{\bullet}(\delta_i:\Delta_n\to\Delta_{n+1}):\mathrm{X}_{n+1}\to\mathrm{X}_n\) et \(s_i=\mathrm{X}_{\bullet}(s_i:\Delta_{n+1}\to\Delta_n):\mathrm{X}_n\to\mathrm{X}_{n+1}\).
\[
\begin{tikzcd}[column sep=2.8em]
\ldotp\,\ldotp\,\ldotp\,\ldotp\,\ldotp\ \mathrm{X}_2
  \arrow[r, shift left=4.4ex, "\delta_0"{near end}]
  \arrow[r, "\delta_1"{near end}]
  \arrow[r, shift right=4.4ex, "\delta_2"{near end}]
&
\mathrm{X}_1
  \arrow[l, shift right=2.2ex, shorten <=1.2em, "s_0"'{near end}]
  \arrow[l, shift left=2.2ex, shorten <=1.2em, "s_1"'{near end}]
  \arrow[r, shift left=2.2ex, "\delta_0"{near end}]
  \arrow[r, shift right=2.2ex, "\delta_1"{near end}]
&
\mathrm{X}_0
  \arrow[l, shorten <=1.2em, "s_0"'{near end}]
\end{tikzcd}
\tag*{\(\mathrm{X}_{\bullet}:\)}
\]

\textbf{(5.1.3)} Soit \(\mathrm{S}\in\operatorname{Ob}\mathscr{C}\). L’\textit{objet simplicial constant} \(\mathrm{S}_{\bullet}\) est l’objet simplicial pour lequel \(\mathrm{S}_n=\mathrm{S}\) et \(\delta_i=s_i=\operatorname{Id}_{\mathrm{S}}\). Un objet simplicial (resp. simplicial \(k\)-tronqué) de \(\mathscr{C}\), \textit{augmenté vers} \(\mathrm{S}\), est un morphisme \(a:\mathrm{X}_{\bullet}\to\mathrm{S}_{\bullet}\) (resp. \(a:\mathrm{X}_{\bullet}\to\operatorname{sq}_k(\mathrm{S}_{\bullet})\)). Nous identifierons les objets simpliciaux (resp. simpliciaux \(k\)-tronqués) de \(\mathscr{C}\), augmentés vers \(\mathrm{S}\), \(a:\mathrm{X}_{\bullet}\to\mathrm{S}_{\bullet}\) (resp. …) aux objets \(\mathrm{X}_{\bullet}^{+}\) de \(\mathit{Hom}((\Delta^+)^0,\mathscr{C})\) (resp. \(\mathit{Hom}((\Delta^+)_k^0,\mathscr{C})\)), tels que \(\mathrm{X}_{\bullet}^{+}(\Delta_{-1})=\mathrm{S}\). On a \(a_n=\mathrm{X}_{\bullet}^{+}(\varepsilon:\Delta_{-1}\to\Delta_n)\). Ces objets seront désignés par une notation du type \(a:\mathrm{X}_{\bullet}\to\mathrm{S}\). Les cosquelettes relatifs seront surtout utilisés dans ce cadre, et notés \(\operatorname{cosq}_k^{\mathrm{S}}\) ou simplement \(\operatorname{cosq}_k\).

\textbf{(5.1.4)} Le cosquelette de l’objet simplicial 0-tronqué augmenté \(a_0:\mathrm{X}\to\mathrm{S}\) est l’objet simplicial augmenté vers \(\mathrm{S}\) de composantes les puissances cartésiennes de \(\mathrm{X}\) dans \(\mathscr{C}/\mathrm{S}\)
\[
\operatorname{cosq}_0(\mathrm{X}\to\mathrm{S})
=(((\mathrm{X}/\mathrm{S})^{\Delta_n})_{n\geq0}\to\mathrm{S}).
\]

\textbf{(5.1.5)} Soient \(u:\mathrm{X}\to\mathrm{Y}\) une application continue entre espaces topologiques, \(\mathrm{F}\) un faisceau sur \(\mathrm{X}\) et \(\mathrm{G}\) un faisceau sur \(\mathrm{Y}\). L’ensemble \(\operatorname{Hom}_u(\mathrm{G},\mathrm{F})\) des \textit{\(u\)-morphismes} de \(\mathrm{G}\) dans \(\mathrm{F}\) est l’ensemble \(\operatorname{Hom}(u^*\mathrm{G},\mathrm{F})\sim\operatorname{Hom}(\mathrm{G},u_*\mathrm{F})\).

\Page{10}

Se donner un \(u\)-morphisme \(f\) de \(\mathrm{G}\) dans \(\mathrm{F}\) revient à se donner, pour tout couple d’ouverts \(\mathrm{U}\subset\mathrm{X}\) et \(\mathrm{V}\subset\mathrm{Y}\) tel que \(u(\mathrm{U})\subset\mathrm{V}\), une application \(f_{\mathrm{UV}}:\mathrm{G}(\mathrm{V})\to\mathrm{F}(\mathrm{U})\); ces applications doivent vérifier la condition

\noindent (*) Pour \(\mathrm{U}'\subset\mathrm{U}\subset\mathrm{X}\), \(\mathrm{V}'\subset\mathrm{V}\subset\mathrm{Y}\), \(u(\mathrm{U})\subset\mathrm{V}\) et \(u(\mathrm{U}')\subset\mathrm{V}'\), le diagramme
\[
\begin{tikzcd}[column sep=2em, row sep=3em]
\mathrm{G}(\mathrm{V}) \arrow[r] \arrow[d]
  & \mathrm{G}(\mathrm{V}') \arrow[d] \\
\mathrm{F}(\mathrm{U}) \arrow[r]
  & \mathrm{F}(\mathrm{U}')
\end{tikzcd}
\]
est commutatif.

\textbf{(5.1.6)} Un \textit{espace topologique simplicial} est un objet simplicial de la catégorie d’objets les espaces topologiques, et de morphismes les applications continues.

Un \textit{faisceau} \(\mathrm{F}^{\bullet}\) sur un espace topologique simplicial \(\mathrm{X}_{\bullet}\) consiste en

\noindent\textit{a)} une famille \(\mathrm{F}^n\) de faisceaux sur les \(\mathrm{X}_n\);\par
\noindent\textit{b)} pour \(f:\Delta_n\to\Delta_m\), un \(\mathrm{X}_{\bullet}(f)\)-morphisme \(\mathrm{F}^{\bullet}(f)\) de \(\mathrm{F}^n\) dans \(\mathrm{F}^m\).\par

On exige de plus que \(\mathrm{F}^{\bullet}(f\circ g)=\mathrm{F}^{\bullet}(f)\circ\mathrm{F}^{\bullet}(g)\).

Un morphisme \(u\) de \(\mathrm{F}^{\bullet}\) dans \(\mathrm{G}^{\bullet}\) est une famille de morphismes \(u^n:\mathrm{F}^n\to\mathrm{G}^n\) telle que pour \(f\in\operatorname{Hom}_{(\Delta)}(\Delta_n,\Delta_m)\), on ait \(u^m\mathrm{F}^{\bullet}(f)=\mathrm{G}^{\bullet}(f)\,u^n\).

\textbf{(5.1.7)} Pour \(\mathrm{U}\) un ouvert de \(\mathrm{X}_n\), soit \(\mathrm{F}^{\bullet}(\mathrm{U})=\mathrm{F}^n(\mathrm{U})\). Pour \(f:\Delta_n\to\Delta_m\), \(\mathrm{U}\subset\mathrm{X}_n\) et \(\mathrm{V}\) dans \(\mathrm{X}_m\) tels que \(\mathrm{X}_{\bullet}(f)(\mathrm{V})\subset\mathrm{U}\), soit \(\mathrm{F}^{\bullet}(f,\mathrm{V},\mathrm{U}):\mathrm{F}^{\bullet}(\mathrm{U})\to\mathrm{F}^{\bullet}(\mathrm{V})\) l’application induite par \(\mathrm{F}^{\bullet}(f)\). On vérifie aussitôt que le système des ensembles \(\mathrm{F}^{\bullet}(\mathrm{U})\) (indexés par \(n\) et \(\mathrm{U}\subset\mathrm{X}_n\)) et des applications \(\mathrm{F}^{\bullet}(f,\mathrm{V},\mathrm{U})\) détermine \(\mathrm{F}^{\bullet}\) uniquement (cf. (5.1.5)). Pour qu’un tel système provienne d’un faisceau sur \(\mathrm{X}_{\bullet}\), il faut et il suffit que :

\noindent\textit{a)} \(\mathrm{F}^{\bullet}(fg,\mathrm{U},\mathrm{W})=\mathrm{F}^{\bullet}(f,\mathrm{U},\mathrm{V})\,\mathrm{F}^{\bullet}(g,\mathrm{V},\mathrm{W})\) chaque fois que cela a un sens;\par
\noindent\textit{b)} quel que soit \(n\), les \(\mathrm{F}^{\bullet}(\mathrm{U})\) pour \(\mathrm{U}\subset\mathrm{X}_n\) forment un faisceau sur \(\mathrm{X}_n\).\par

\textbf{(5.1.8)} Ceci montre que les faisceaux sur \(\mathrm{X}_{\bullet}\) peuvent s’interpréter comme les faisceaux sur un site convenable, et en particulier forment un topos \((\mathrm{X}_{\bullet})^{\sim}\). On utilisera librement, pour les faisceaux sur \(\mathrm{X}_{\bullet}\), la terminologie en usage pour les faisceaux sur un site. Ainsi, un faisceau en groupes abéliens (resp. un faisceau d’anneaux, …) \(\mathrm{F}^{\bullet}\) est un système de faisceaux en groupes abéliens (resp. de faisceaux d’anneaux, …) \(\mathrm{F}^n\) sur les \(\mathrm{X}_n\), et de morphismes \(\mathrm{F}^{\bullet}(f)\).

\textbf{(5.1.9)} \textit{Exemples.} (I) Soit \(\mathrm{X}_{\bullet}\) un espace analytique simplicial. Les faisceaux structuraux \(\mathscr{O}_{\mathrm{X}_n}\) forment un faisceau d’anneaux sur \(\mathrm{X}_{\bullet}\).

(II) Soit \(\mathrm{X}_{\bullet}\) un espace analytique simplicial augmenté vers \(\mathrm{S}\). Les faisceaux \(\Omega^1_{\mathrm{X}_n/\mathrm{S}}\) forment un faisceau de \(\mathscr{O}\)-modules sur \(\mathrm{X}_{\bullet}\). Sa puissance extérieure \(i\)-ème est le faisceau
\Page{11}
de \(\mathscr{O}\)-modules \((\Omega^i_{\mathrm{X}_n/\mathrm{S}})_{n\geq0}\), noté \(\Omega^i_{\mathrm{X}_{\bullet}/\mathrm{S}}\). Les complexes de De Rham \((\Omega^{\bullet}_{\mathrm{X}_n/\mathrm{S}})_{n\geq0}\) forment un complexe de faisceaux sur \(\mathrm{X}_{\bullet}\).

(III) Soit \(\mathrm{F}^{\bullet}\) un faisceau de groupes abéliens sur \(\mathrm{X}_{\bullet}\). Les résolutions flasques canoniques de Godement \(\mathscr{C}^{\bullet}(\mathrm{F}^n)\) forment un complexe de faisceaux sur \(\mathrm{X}_{\bullet}\), et une résolution de \(\mathrm{F}^{\bullet}\).

(IV) Soit \(\mathrm{S}\) un espace topologique. Les faisceaux \(\mathrm{F}^{\bullet}\) sur l’espace simplicial constant \(\mathrm{S}_{\bullet}\) s’identifient aux faisceaux cosimpliciaux sur \(\mathrm{S}\). En particulier, un faisceau abélien \(\mathrm{F}^{\bullet}\) sur \(\mathrm{S}_{\bullet}\) définit un complexe différentiel \((\mathrm{F}^n,d=\sum\limits_i(-1)^i\delta_i)\). Un complexe de faisceaux abéliens \(\mathrm{K}\) sur \(\mathrm{S}_{\bullet}\) définit un double complexe \((\mathrm{K}^{n,m})\) encore noté \(\mathrm{K}\) (\(m=\) degré cosimplicial) de complexe simple associé \(\boldsymbol{s}\mathrm{K}\) :
\[
(\boldsymbol{s}\mathrm{K})^n
=\bigoplus_{p+q=n}\mathrm{K}^{pq}
\tag*{\textbf{(5.1.9.1)}}
\]
et
\[
d(x^{pq})=d_{\mathrm{K}}(x^{pq})
+(-1)^p\sum_i(-1)^i\delta_i x^{pq}.
\tag*{\textbf{(5.1.9.2)}}
\]
On notera \(\mathrm{L}\) la seconde filtration de \(\boldsymbol{s}\mathrm{K}\) :
\[
\mathrm{L}^r(\boldsymbol{s}\mathrm{K})
=\bigoplus_{q\geq r}\mathrm{K}^{pq}.
\tag*{\textbf{(5.1.9.3)}}
\]

\textbf{(5.1.10)} Soit \(u:\mathrm{X}_{\bullet}\to\mathrm{Y}_{\bullet}\) un morphisme d’espaces topologiques simpliciaux, de composantes \(u_n:\mathrm{X}_n\to\mathrm{Y}_n\). Si \(\mathrm{G}\) (resp. \(\mathrm{F}\)) est un faisceau sur \(\mathrm{Y}_{\bullet}\) (resp. sur \(\mathrm{X}_{\bullet}\)), alors \((u_n^*\mathrm{G}^n)_{n\geq0}\) (resp. \((u_{n*}\mathrm{F}^n)_{n\geq0}\)) est un faisceau sur \(\mathrm{X}_{\bullet}\) (resp. sur \(\mathrm{Y}_{\bullet}\)); on le note \(u^*\mathrm{G}\) (resp. \(u_*\mathrm{F}\)). Les foncteurs \(u_*\) et \(u^*\) sont adjoints; ce sont les morphismes image directe et réciproque d’un morphisme de topos \(u:(\mathrm{X}_{\bullet})^{\sim}\to(\mathrm{Y}_{\bullet})^{\sim}\).

\textbf{(5.1.11)} Soit \(a:\mathrm{X}_{\bullet}\to\mathrm{S}\) un espace topologique simplicial augmenté vers \(\mathrm{S}\). Si \(\mathrm{F}\) est un faisceau sur \(\mathrm{S}\), \(a^*\mathrm{F}=(a_n^*\mathrm{F})_{n\geq0}\) « est » un faisceau sur \(\mathrm{X}_{\bullet}\). Le foncteur \(a^*\) a un adjoint à droite
\[
a_*:\mathrm{F}^{\bullet}\mapsto
\operatorname{Ker}\bigl(
\begin{tikzcd}[ampersand replacement=\&, cramped, column sep=1.7em, baseline=-0.5ex]
a_{0*}\mathrm{F}^0
  \arrow[r, shift left=0.5ex, "\delta_0^*"]
  \arrow[r, shift right=0.5ex, "\delta_1^*"']
\& a_{1*}\mathrm{F}^1
\end{tikzcd}
\bigr).
\tag*{\textbf{(5.1.11.1)}}
\]
Les foncteurs \(a_*\) et \(a^*\) sont les foncteurs image directe et réciproque d’un morphisme de topos \(a:(\mathrm{X}_{\bullet})^{\sim}\to(\mathrm{S})^{\sim}\).

\textbf{(5.1.12)} Soient \(\mathrm{S}_{\bullet}\) l’espace simplicial constant associé à \(\mathrm{S}\) et \(a_{\bullet}:\mathrm{X}_{\bullet}\to\mathrm{S}_{\bullet}\) le morphisme défini par \(a\). Pour \(\mathrm{F}^{\bullet}\) un faisceau abélien sur \(\mathrm{X}_{\bullet}\), on identifie \(a_{\bullet *}(\mathrm{F}^{\bullet})\) à un faisceau cosimplicial sur \(\mathrm{S}\) ((5.1.9) (IV)). Pour \(\mathrm{K}\) un complexe de faisceaux abéliens sur \(\mathrm{X}_{\bullet}\), si \(\mathrm{K}^{pq}\) est la composante de degré \(p\) de la restriction de \(\mathrm{K}\) à \(\mathrm{X}_q\), les composantes du complexe simple associé au double complexe défini par \(a_{\bullet *}\mathrm{K}\) sont
\[
(\boldsymbol{s}a_{\bullet *}\mathrm{K})^n
=\bigoplus_{p+q=n}a_{q*}\mathrm{K}^{pq}.
\tag*{\textbf{(5.1.12.1)}}
\]
\Page{12}
La suite spectrale définie par la filtration \(\mathrm{L}\) ((5.1.9.3)) de \(\boldsymbol{s}a_{\bullet *}\mathrm{K}\) s’écrit
\[
\mathrm{E}_1^{pq}
=H^q(a_{p*}(\mathrm{K}\mid\mathrm{X}_{p}))
\Rightarrow H^{p+q}(\boldsymbol{s}a_{\bullet *}\mathrm{K}).
\tag*{\textbf{(5.1.12.2)}}
\]

\textbf{(5.1.13)} Faisons \(\mathrm{S}=\mathrm{P}^{t}\) (espace topologique réduit à un point). On trouve que, pour \(\mathrm{F}^{\bullet}\) un faisceau sur \(\mathrm{X}_{\bullet}\), les \(\Gamma(\mathrm{X}_n,\mathrm{F}^n)\) sont les composantes d’un ensemble cosimplicial \(\Gamma^{\bullet}(\mathrm{X}_{\bullet},\mathrm{F}^{\bullet})\). Le foncteur \(\Gamma\) (sections globales) est le foncteur
\[
\Gamma:\mathrm{F}^{\bullet}\mapsto
\operatorname{Ker}\bigl(
\begin{tikzcd}[ampersand replacement=\&, cramped, column sep=1.4em, baseline=-0.5ex]
\Gamma(\mathrm{X}_0,\mathrm{F}^0)
  \arrow[r, shift left=0.5ex]
  \arrow[r, shift right=0.5ex]
\& \Gamma(\mathrm{X}_1,\mathrm{F}^1)
\end{tikzcd}
\bigr).
\tag*{\textbf{(5.1.13.1)}}
\]
Si \(\mathrm{K}\) est un complexe de faisceaux abéliens, on désigne par \(\Gamma^{\bullet}(\mathrm{X}_{\bullet},\mathrm{K})\) le complexe différentiel de composantes les groupes abéliens cosimpliciaux \(\Gamma^{\bullet}(\mathrm{X}_{\bullet},\mathrm{K}^n)\) et par \(\boldsymbol{s}\Gamma^{\bullet}(\mathrm{X}_{\bullet},\mathrm{K})\) le complexe simple associé.

\section*{5.2. Cohomologie des espaces topologiques simpliciaux.}

\textbf{(5.2.1)} A chaque espace topologique simplicial \(\mathrm{X}_{\bullet}\) est attaché un espace topologique usuel \(\lvert\mathrm{X}_{\bullet}\rvert\), appelé sa réalisation géométrique (voir ci-dessous). Dans le cas particulier où les \(\mathrm{X}_n\) sont discrets, on retrouve la notion habituelle de réalisation géométrique d’un ensemble simplicial.

Dans [14], G. Segal définit la cohomologie de \(\mathrm{X}_{\bullet}\) à valeurs dans un groupe abélien \(\mathrm{A}\) comme étant \(\mathrm{H}^{\bullet}(\lvert\mathrm{X}_{\bullet}\rvert,\mathrm{A})\). Sous des hypothèses convenables, la filtration de \(\lvert\mathrm{X}_{\bullet}\rvert\) par les squelettes successifs fournit une suite spectrale.
\[
\mathrm{E}_1^{pq}
=\mathrm{H}^q(\mathrm{X}_p,\mathrm{A})
\Rightarrow\mathrm{H}^{p+q}(\mathrm{X}_{\bullet},\mathrm{A}).
\tag*{\textbf{(5.2.1.1)}}
\]

Un des intérêts de cette définition est qu’elle s’applique aussi bien, par exemple, à la définition de groupes \(\mathrm{K}(\mathrm{X}_{\bullet})\). Nous en adopterons une autre, mieux adaptée aux techniques faisceautiques.

\textit{Réalisation géométrique.} — Pour \(n\) entier \(\geq0\), on désigne par \(\lvert\Delta_n\rvert\) le simplexe dans \(\mathbf{R}^{\Delta_n}\) de sommets l’ensemble des vecteurs de base. On identifie \(\Delta_n\) à l’ensemble des sommets de \(\lvert\Delta_n\rvert\). Toute fonction \(f:\Delta_n\to\Delta_m\) se prolonge par linéarité en \(\lvert f\rvert:\lvert\Delta_n\rvert\to\lvert\Delta_m\rvert\).

Soit \(\mathrm{X}_{\bullet}\) un espace topologique simplicial. Soit \(\mathrm{Y}=\coprod\limits_{n\geq0}(\mathrm{X}_n\times\lvert\Delta_n\rvert)\). Soit \(\mathrm{R}\) la plus fine des relations d’équivalence sur \(\mathrm{Y}\) pour lesquelles, quels que soient \(f:\Delta_n\to\Delta_m\) dans \((\Delta)\), \(x_m\in\mathrm{X}_m\) et \(a\in\lvert\Delta_n\rvert\), on ait
\[
(x_m,\lvert f\rvert(a))\equiv(f(x_m),a)\qquad(\operatorname{mod}\mathrm{R}).
\]
La \textit{réalisation géométrique} de \(\mathrm{X}_{\bullet}\) est par définition
\[
\lvert\mathrm{X}_{\bullet}\rvert=\mathrm{Y}/\mathrm{R}.
\]

\textit{Définition} \textbf{(5.2.2)}. — \textit{Soit} \(\mathrm{X}_{\bullet}\) \textit{un espace topologique simplicial. Les foncteurs} cohomologie à coefficients dans le faisceau abélien \(\mathrm{F}^{\bullet}\) sur \(\mathrm{X}_{\bullet}\), \(\mathrm{H}^i(\mathrm{X}_{\bullet},\mathrm{F}^{\bullet})\), \textit{sont les foncteurs dérivés du foncteur} \(\Gamma\) (5.1.13.1).

Cette définition équivaut à la suivante, parfois plus maniable.

% ---- 013-020 ----
\Page{13}
\textbf{(5.2.3)} Soit \(\mathrm{F}^{\bullet}\) un faisceau abélien sur \(\mathrm{X}_{\bullet}\). On montre (par exemple à l’aide de (5.1.9) (III)) que \(\mathrm{F}^{\bullet}\) admet toujours des résolutions à droite \(\mathrm{K}\) telles que
\[
\mathrm{H}^{r}(\mathrm{X}_{q},\mathrm{K}^{p,q})=0
\qquad \text{pour}\quad p,q\geq 0
\quad\text{et}\quad r>0.
\]
Si \(\mathrm{K}\) est une telle résolution, on a canoniquement
\[
\mathrm{H}^{n}(\mathrm{X}_{\bullet},\mathrm{F}^{\bullet})
\simeq
\mathrm{H}^{n}(\boldsymbol{s}\Gamma^{\bullet}(\mathrm{X}_{\bullet},\mathrm{K})).
\tag*{\textbf{(5.2.3.1)}}
\]
Il est facile de vérifier directement que le membre de droite est indépendant (à isomorphisme unique près) du choix de \(\mathrm{K}\); pour ce qui suit, il serait loisible de \textit{définir} \(\mathrm{H}^{\bullet}(\mathrm{X}_{\bullet},\mathrm{F}^{\bullet})\) par (5.2.3.1). On vérifie de plus que la suite spectrale du complexe \(\boldsymbol{s}\Gamma^{\bullet}(\mathrm{X}_{\bullet},\mathrm{K})\), filtré par \(\mathrm{L}\) ((5.1.9.3) et (5.1.12.2) pour \(\mathrm{S}=\mathrm{P}^{t}\))
\[
\mathrm{E}_{1}^{pq}=\mathrm{H}^{q}(\mathrm{X}_{p},\mathrm{F}^{p})
\Rightarrow
\mathrm{H}^{p+q}(\mathrm{X}_{\bullet},\mathrm{F}^{\bullet})
\tag*{\textbf{(5.2.3.2)}}
\]
est indépendante (à isomorphisme unique près) du choix de \(\mathrm{K}\) (cf. (5.2.1.1)).

\textbf{(5.2.4)} Il y a lieu de préciser la construction qui précède en passant aux catégories dérivées, et d’en donner une variante relative.

Soient \(a:\mathrm{X}_{\bullet}\to\mathrm{S}\), \(\mathrm{S}_{\bullet}\) l’espace simplicial constant défini par \(\mathrm{S}\), \(a_{\bullet}:\mathrm{X}_{\bullet}\to\mathrm{S}_{\bullet}\) déduit de \(a\) et \(\varepsilon:\mathrm{S}_{\bullet}\to\mathrm{S}\) le morphisme d’augmentation. On a
\[
a_*=\varepsilon_*a_{\bullet *}:
(\mathrm{X}_{\bullet})^{\sim}\to(\mathrm{S})^{\sim}.
\tag*{\textbf{(5.2.4.1)}}
\]

Cette formule se dérive en
\[
\mathrm{R}a_*=\mathrm{R}\varepsilon_*\,\mathrm{R}a_{\bullet *}:
\mathrm{D}^{+}(\mathrm{X}_{\bullet})\to\mathrm{D}^{+}(\mathrm{S})
\tag*{\textbf{(5.2.4.2)}}
\]
(\(\mathrm{D}^{+}(\mathrm{X}_{\bullet})=\)la catégorie dérivée bornée inférieurement de la catégorie des faisceaux abéliens sur \(\mathrm{X}_{\bullet}\)). Calculons \(\mathrm{R}a_{\bullet *}\) et \(\mathrm{R}\varepsilon_*\).

\textbf{(5.2.5)} Soit \(u:\mathrm{X}_{\bullet}\to\mathrm{Y}_{\bullet}\) un morphisme d’espaces topologiques simpliciaux. Le foncteur \(\mathrm{R}u_*:\mathrm{D}^{+}(\mathrm{X}_{\bullet})\to\mathrm{D}^{+}(\mathrm{Y}_{\bullet})\) se calcule « composante par composante » : si \(\mathrm{K}\) est un complexe de faisceaux abéliens sur \(\mathrm{X}_{\bullet}\), pour calculer \(\mathrm{R}u_*\mathrm{K}\), on prend une résolution \(\mathrm{K}\xrightarrow{\approx}\mathrm{K}'\) telle que les composantes \(\mathrm{F}\) de \(\mathrm{K}'\) vérifient \(\mathrm{R}^{i}_{u_{p*}}(\mathrm{F}^{p})=0\) pour \(i>0\), \(p\geq 0\) (cf. (5.1.9) (III)); on a alors \(\mathrm{R}u_*\mathrm{K}\approx u_*\mathrm{K}'\).

\textbf{(5.2.6)} Le foncteur \(\boldsymbol{s}:\mathrm{C}^{+}(\mathrm{S}_{\bullet})\to\mathrm{C}^{+}(\mathrm{S})\) transforme complexe acyclique en complexe acyclique. Il se dérive donc trivialement en
\[
\boldsymbol{s}:\mathrm{D}^{+}(\mathrm{S}_{\bullet})\to\mathrm{D}^{+}(\mathrm{S}).
\]

Pour \(\mathrm{F}\) un faisceau injectif sur \(\mathrm{S}_{\bullet}\), le complexe différentiel \(\boldsymbol{s}\mathrm{F}\) est une résolution de \(\varepsilon_*\mathrm{F}\). On en déduit un isomorphisme \(\mathrm{R}\varepsilon_*\xrightarrow{\sim}\boldsymbol{s}\) et (5.2.4.2) devient
\[
\mathrm{R}a_*=\boldsymbol{s}\mathrm{R}a_{\bullet *}.
\tag*{\textbf{(5.2.6.1)}}
\]

\Page{14}
\textbf{(5.2.7)} Jointe à (5.2.5) et spécialisée au cas \(\mathrm{S}=\mathrm{P}^{t}\), cette formule prouve (5.2.3.1). En termes concrets, elle signifie que pour calculer \(\mathrm{R}a_*\mathrm{K}\), on peut procéder en deux étapes :

\textit{a)} on prend une résolution \(\mathrm{K}\xrightarrow{\approx}\mathrm{K}'\), telle que les composantes \(\mathrm{F}\) de \(\mathrm{K}'\) vérifient \(\mathrm{R}^{i}a_{p*}(\mathrm{F}^{p})=0\) pour \(i>0\). Le complexe \(a_{\bullet *}\mathrm{K}'\in\mathrm{D}^{+}(\mathrm{S}_{\bullet})\) (catégorie dérivée de la catégorie des faisceaux abéliens cosimpliciaux sur \(\mathrm{S}\)) s’identifie à \(\mathrm{R}a_{\bullet *}\mathrm{K}\).

\textit{b)} \(\mathrm{R}a_*\mathrm{K}\) est le simple complexe \(\boldsymbol{s}a_{\bullet *}\mathrm{K}'\) associé au double complexe \(a_{\bullet *}\mathrm{K}'\).

La suite spectrale (5.2.3.2) se généralise en une suite spectrale déduite de (5.1.12.2)
\[
\mathrm{E}_{1}^{pq}=\mathrm{R}^{q}a_{p*}(\mathrm{K}\mid\mathrm{X}_{p})
\Rightarrow \mathrm{R}^{p+q}a_*\mathrm{K}.
\tag*{\textbf{(5.2.7.1)}}
\]

\section*{5.3. Descente cohomologique.}

\textbf{(5.3.1)} Soit \(a:\mathrm{X}_{\bullet}\to\mathrm{S}\) un espace topologique simplicial augmenté. Pour tout faisceau \(\mathrm{F}\) sur \(\mathrm{S}\), on dispose d’un morphisme d’adjonction
\[
\varphi:\mathrm{F}\to a_*a^*\mathrm{F}.
\]
Ce morphisme se dérive en un morphisme de foncteurs de \(\mathrm{D}^{+}(\mathrm{S})\) dans \(\mathrm{D}^{+}(\mathrm{S})\) :
\[
\varphi:\mathrm{Id}\to\mathrm{R}a_*a^*.
\tag*{\textbf{(5.3.1.1)}}
\]

\textit{Définition} \textbf{(5.3.2)}. — \textit{On dit que \(a\) est de \textup{descente cohomologique} si pour tout faisceau abélien \(\mathrm{F}\) sur \(\mathrm{S}\), on a}
\[
\mathrm{F}\xrightarrow{\sim}
\operatorname{Ker}(a_{0*}a_{0}^{*}\mathrm{F}\to a_{1*}a_{1}^{*}\mathrm{F})
\]
\[
\makebox[\linewidth]{\textit{et}\hfill
\(\mathrm{R}^{i}a_*a^*\mathrm{F}=0
\qquad\text{\textit{pour}}\quad i>0.\)\hfill}
\]

Il revient au même de demander que (5.3.1.1) soit un isomorphisme.

\textbf{(5.3.3)} Si \(a\) est de descente cohomologique, alors, pour \(\mathrm{K}\in\operatorname{Ob}\mathrm{D}^{+}(\mathrm{S})\), l’application canonique
\[
\mathrm{R}\Gamma(\mathrm{S},\mathrm{K})
\to\mathrm{R}\Gamma(\mathrm{S},\mathrm{R}a_*a^*\mathrm{K})
\sim\mathrm{R}\Gamma(\mathrm{X}_{\bullet},a^*\mathrm{K})
\tag*{\textbf{(5.3.3.1)}}
\]
est un isomorphisme. En particulier, pour \(\mathrm{F}\) un faisceau abélien sur \(\mathrm{S}\), on a une suite spectrale (5.2.3.2)
\[
\mathrm{E}_{1}^{pq}=\mathrm{H}^{q}(\mathrm{X}_{p},a_{p}^{*}\mathrm{F})
\Rightarrow\mathrm{H}^{p+q}(\mathrm{S},\mathrm{F}).
\tag*{\textbf{(5.3.3.2)}}
\]

Pour un complexe, on a encore, en hypercohomologie, une suite spectrale
\[
\mathrm{E}_{1}^{pq}=\mathbf{H}^{q}(\mathrm{X}_{p},a_{p}^{*}\mathrm{K})
\Rightarrow\mathbf{H}^{p+q}(\mathrm{S},\mathrm{K}).
\tag*{\textbf{(5.3.3.3)}}
\]

Dans les deux cas, les \(\mathrm{E}_{1}^{\bullet q}\) (\(q\) fixe) forment un groupe cosimplicial, et
\[
d_{1}=\sum_i(-1)^i\delta_i:
\mathrm{E}_{1}^{p,q}\to\mathrm{E}_{1}^{p+1,q}.
\]

\Page{15}
\textit{Définition} \textbf{(5.3.4)}. — \textit{Une application continue \(a:\mathrm{X}\to\mathrm{S}\) est de \textup{descente cohomologique} si le morphisme d’augmentation de \(\operatorname{cosq}(\mathrm{X}\to\mathrm{S})\) :}
\[
((\mathrm{X}/\mathrm{S})^{\Delta_n})_{n\geq 0}\to\mathrm{S}
\]
\textit{est de descente cohomologique. On dit que \(a\) est de \textup{descente cohomologique universelle} si pour tout \(u:\mathrm{S}'\to\mathrm{S}\), l’application continue \(a':\mathrm{X}\times_{\mathrm{S}}\mathrm{S}'\to\mathrm{S}'\) est de descente cohomologique.}

\textbf{(5.3.5)} Les résultats fondamentaux, démontrés dans \textit{S.G.A.}4, V \textit{bis}, sont les suivants.

(I) Les applications continues de descente cohomologique universelle forment une topologie de Grothendieck sur la catégorie des espaces topologiques.

On l’appelle la \textit{topologie de la descente cohomologique universelle}.

(II) Une application propre (III.0.1) et surjective est de descente cohomologique universelle.

(III) Une application \(a:\mathrm{X}\to\mathrm{S}\) admettant des sections localement sur \(\mathrm{S}\) est de descente cohomologique universelle.

(IV) Soit \(a:\mathrm{X}_{\bullet}\to\mathrm{S}\) un espace simplicial augmenté \(k\)-tronqué \((-1\leq k\leq\infty)\). Pour \(k\geq n\geq-1\), soit \(\varphi_n:\operatorname{cosq}\mathrm{X}_{\bullet}\to\operatorname{cosq}\operatorname{sq}_n\mathrm{X}_{\bullet}\) la flèche évidente. On dit que \(\mathrm{X}_{\bullet}\) est un \textit{hyperrecouvrement \(k\)-tronqué de \(\mathrm{S}\), pour la topologie de la descente cohomologique universelle}, si les applications
\[
(\varphi_n)_{n+1}:\mathrm{X}_{n+1}
\to(\operatorname{cosq}\operatorname{sq}_n\mathrm{X}_{\bullet})_{n+1}
\qquad(-1\leq n\leq k-1)
\tag*{\textbf{(5.3.5.1)}}
\]
sont de descente cohomologique universelle. Si \(\mathrm{X}_{\bullet}\) est un tel hyperrecouvrement, alors l’espace simplicial augmenté vers \(\mathrm{S}\), \(\operatorname{cosq}(\mathrm{X}_{\bullet})\), est de descente cohomologique.

(V) Soit \(a\) un morphisme d’espaces topologiques simpliciaux augmenté vers \(\mathrm{S}\)
\[
\begin{tikzcd}[column sep=2.5em,row sep=2.5em]
\mathrm{X}_{\bullet}\arrow[r,"a"]\arrow[d,"x"]
  & \mathrm{Y}_{\bullet}\arrow[d,"y"]\\
\mathrm{S}\arrow[r,equal] & \mathrm{S}
\end{tikzcd}
\]
On dit que \(a\) est un \textit{hyperrecouvrement} pour la topologie de la descente cohomologique universelle si les applications évidentes \(\mathrm{X}_n\to(\operatorname{cosq}_{n-1}^{\mathrm{Y}_{\bullet}}\operatorname{sq}_{n-1}\mathrm{X}_{\bullet})_n\) sont de descente cohomologique universelle. Si \(a\) est un tel hyperrecouvrement, alors, pour tout \(\mathrm{K}\in\operatorname{Ob}\mathrm{D}^{+}(\mathrm{S})\)
\[
a^*:\mathrm{R}y_*y^*\mathrm{K}\xrightarrow{\sim}\mathrm{R}x_*x^*\mathrm{K}.
\]

\textbf{(5.3.6)} Pour \(k=\infty\), (IV) affirme que \(a:\mathrm{X}_{\bullet}\to\mathrm{S}\) est de descente cohomologique si les \((\varphi_n)_{n+1}\) sont de descente cohomologique universelle. Pour \(n=-1\) ou \(0\), ces applications sont
\[
\begin{aligned}
n=-1&:\quad \mathrm{X}_0\xrightarrow{a}\mathrm{S}\\
n=0&:\quad \mathrm{X}_1\xrightarrow{(\delta_0,\delta_1)}
\mathrm{X}_0\times_{\mathrm{S}}\mathrm{X}_0.
\end{aligned}
\]
\Page{16}
Pour \(n=1\), \((\operatorname{cosq}\operatorname{sq}_1(\mathrm{X}_{\bullet}))_2\) est le sous-espace de \(\mathrm{X}_1\times_{\mathrm{S}}\mathrm{X}_1\times_{\mathrm{S}}\mathrm{X}_1\) formé des triplets \((x,y,z)\) tels que \(\delta_0x=\delta_0y\), \(\delta_1x=\delta_0z\) et \(\delta_1y=\delta_1z\). L’application \((\varphi_1)_2\) est \(x\mapsto(\delta_0x,\delta_1x,\delta_2x)\).

Pour \(k=0\), (IV) est la définition (5.3.4).

\textit{Exemple} \textbf{(5.3.7)}. — Soit \(\mathcal{U}=(\mathrm{U}_i)_{i\in\mathrm{I}}\) un recouvrement ouvert, ou un recouvrement fermé localement fini de \(\mathrm{S}\). Soit \(\mathrm{X}=\coprod\limits_{i\in\mathrm{I}}\mathrm{U}_i\). Alors, \(a:\mathrm{X}\to\mathrm{S}\) est de descente cohomologique. La suite spectrale (5.3.3.2) pour \(\mathrm{X}_{\bullet}=\operatorname{cosq}(\mathrm{X}\to\mathrm{S})\) n’est autre que la suite spectrale de Leray du recouvrement \(\mathcal{U}\).

\textbf{(5.3.8)} Soit \(a:\mathrm{X}_{\bullet}\to\mathrm{S}\) comme en (5.3.5) (IV). On dit que \(\mathrm{X}_{\bullet}\) est un \textit{hyperrecouvrement propre \(k\)-tronqué} de \(\mathrm{S}\) si les flèches (5.3.5.1) sont propres et surjectives. Pour \(k=\infty\), on parle simplement d’hyperrecouvrement propre.

\section*{6. EXEMPLES D’ESPACES TOPOLOGIQUES SIMPLICIAUX}

\section*{6.1. Espaces classifiants.}

\textbf{(6.1.1)} Soit \(u:\mathrm{X}\to\mathrm{S}\) une application continue. Pour tout faisceau \(\mathrm{F}\) sur \(\mathrm{S}\), le faisceau \(u^*\mathrm{F}\) est muni d’une « donnée de descente » relativement à \(u\) : on dispose d’un isomorphisme entre les deux images réciproques de \(u^*\mathrm{F}\) sur \(\mathrm{X}\times_{\mathrm{S}}\mathrm{X}\), et cet isomorphisme vérifie une condition de cocycles. Si, localement sur \(\mathrm{S}\), \(u\) admet une section, cette construction définit une équivalence entre la catégorie des faisceaux sur \(\mathrm{S}\) et celle des faisceaux sur \(\mathrm{X}\), munis d’une donnée de descente.

Prenons pour \(\mathrm{X}\) un espace principal homogène (à gauche) de groupe \(\mathrm{G}\) sur un espace \(\mathrm{S}\). Les faisceaux \(\mathrm{G}\)-équivariants sur \(\mathrm{X}\) ne sont autres que les faisceaux munis d’une donnée de descente : tout faisceau équivariant sur \(\mathrm{X}\) est d’une et d’une seule façon (en tant que faisceau équivariant) image réciproque d’un faisceau sur \(\mathrm{S}=\mathrm{X}/\mathrm{G}\).

\textbf{(6.1.2)} Si un groupe topologique \(\mathrm{G}\) agit sur un espace \(\mathrm{X}\), alors \(\mathrm{G}\) agit sur \(\mathrm{G}^{\Delta_n}\times\mathrm{X}\) par
\[
g\mathbin{.}(g_0,\ldots,g_n,x)
=(g_0g^{-1},\ldots,g_ng^{-1},gx).
\]
On désigne par \([\mathrm{X}/\mathrm{G}]_{\bullet}\) l’espace simplicial
\[
[\mathrm{X}/\mathrm{G}]_{\bullet}
=((\mathrm{G}^{\Delta_n}\times\mathrm{X})/\mathrm{G})_{n\geq 0}.
\tag*{\textbf{(6.1.2.1)}}
\]

\textit{a)} Si \(\mathrm{X}\) est un espace principal homogène de groupe \(\mathrm{G}\) sur \(\mathrm{S}=\mathrm{X}/\mathrm{G}\), alors l’application
\[
\mathrm{G}^{\Delta_n}\times\mathrm{X}\to\mathrm{X}^{\Delta_n}:
(g_0,\ldots,g_n,x)\mapsto(g_0x,\ldots,g_nx)
\]
identifie \([\mathrm{X}/\mathrm{G}]_n\) au produit fibré itéré \((\mathrm{X}/\mathrm{S})^{\Delta_n}\) :
\[
\operatorname{cosq}(\mathrm{X}\to\mathrm{S})
=([\mathrm{X}/\mathrm{G}]_{\bullet}\to\mathrm{S}).
\]
\Page{17}
En particulier, on a (5.3.5 (III))
\[
\mathrm{H}^{\bullet}([\mathrm{X}/\mathrm{G}]_{\bullet})
\simeq\mathrm{H}^{\bullet}(\mathrm{X}/\mathrm{G})
\qquad\text{(pour un espace principal homogène).}
\tag*{\textbf{(6.1.2.2)}}
\]

\textit{b)} Pour tout \(n\), \(\mathrm{G}^{\Delta_n}\times\mathrm{X}\) est un espace principal homogène de groupe \(\mathrm{G}\) sur \([\mathrm{X}/\mathrm{G}]_n\). Pour tout faisceau équivariant \(\mathrm{F}\) sur \(\mathrm{X}\), \(\operatorname{pr}_2^*\mathrm{F}\) est un faisceau équivariant sur \(\mathrm{G}^{\Delta_n}\times\mathrm{X}\); d’après (6.1.1), ce dernier est image réciproque de \(\mathrm{F}^n\) sur \([\mathrm{X}/\mathrm{G}]_n\). Tout faisceau équivariant \(\mathrm{F}\) sur \(\mathrm{X}\) définit ainsi un faisceau sur \([\mathrm{X}/\mathrm{G}]_{\bullet}\). Il est facile de vérifier qu’on obtient ainsi une équivalence de la catégorie des faisceaux équivariants sur \(\mathrm{X}\) avec la catégorie des faisceaux \(\mathrm{F}^{\bullet}\) sur \([\mathrm{X}/\mathrm{G}]_{\bullet}\) qui vérifient

\((\ast)\) Pour tout \(f:\Delta_n\to\Delta_m\), le morphisme structural \(f^*\mathrm{F}^n\to\mathrm{F}^m\) est un isomorphisme.

\textit{c)} La construction \textit{b)} est naturelle en \((\mathrm{G},\mathrm{X},\mathrm{F})\). On pose
\[
\mathrm{H}^{\bullet}(\mathrm{X},\mathrm{G};\mathrm{F})
=\mathrm{H}^{\bullet}([\mathrm{X}/\mathrm{G}]_{\bullet},\mathrm{F}^{\bullet})
\tag*{\textbf{(6.1.2.3)}}
\]
(cohomologie mixte de \(\mathrm{X}\), \(\mathrm{G}\) à coefficients dans \(\mathrm{F}\)). Sous les hypothèses de \textit{a)}, si \(\mathrm{F}\) est image réciproque de \(\mathrm{F}^{-1}\) sur \(\mathrm{S}=\mathrm{X}/\mathrm{G}\), on a
\[
\mathrm{H}^{\bullet}(\mathrm{X},\mathrm{G};\mathrm{F})
=\mathrm{H}^{\bullet}(\mathrm{X}/\mathrm{G},\mathrm{F}^{-1})
\qquad\text{(pour un espace principal homogène).}
\tag*{\textbf{(6.1.2.4)}}
\]
Ceci généralise (6.1.2.2) (qui est le cas \(\mathrm{F}=\mathbf{Z}\)).

\textbf{(6.1.3)} Soit \(\mathrm{P}^{t}\) l’espace topologique réduit à un point. On appelle \textit{espace classifiant simplicial} de \(\mathrm{G}\) et on note \(\mathrm{B}_{\bullet\mathrm{G}}\) l’espace simplicial
\[
\mathrm{B}_{\bullet\mathrm{G}}=[\mathrm{P}^{t}/\mathrm{G}]_{\bullet}.
\]

Soit \(\mathrm{P}\) un espace principal homogène de groupe \(\mathrm{G}\) sur \(\mathrm{S}\). Le morphisme évident
\[
\operatorname{cosq}(\mathrm{P}\to\mathrm{S})
=[\mathrm{P}/\mathrm{G}]_{\bullet}
\to[\mathrm{P}^{t}/\mathrm{G}]_{\bullet}
=\mathrm{B}_{\bullet\mathrm{G}}
\]
définit un morphisme composé
\[
[\mathrm{P}]:
\mathrm{H}^{\bullet}(\mathrm{B}_{\bullet\mathrm{G}})
\longrightarrow
\mathrm{H}^{\bullet}([\mathrm{P}/\mathrm{G}]_{\bullet})
\xleftarrow[\text{\textbf{6.1.2.2}}]{\sim}
\mathrm{H}^{\bullet}(\mathrm{S}).
\tag*{\textbf{(6.1.3.1)}}
\]

On verra ci-dessous que dans les bons cas \(\mathrm{H}^{\bullet}(\mathrm{B}_{\bullet\mathrm{G}})=\mathrm{H}^{\bullet}(\mathrm{B}_{\mathrm{G}})\), et que l’image de \([\mathrm{P}]\) est formée des \textit{classes caractéristiques} de \(\mathrm{P}\).

\textbf{(6.1.4)} Soient \(\mathrm{G}\) un groupe de Lie, \(\mathrm{B}_{\mathrm{G}}\) un espace classifiant pour \(\mathrm{G}\) et \(a:\mathrm{U}_{\mathrm{G}}\to\mathrm{B}_{\mathrm{G}}\) le \(\mathrm{G}\)-espace principal homogène universel. Soit \(\mathrm{X}\) un \(\mathrm{G}\)-espace; \(\mathrm{X}\times\mathrm{U}_{\mathrm{G}}\) est un \(\mathrm{G}\)-espace principal homogène sur \(\mathrm{X}\times\mathrm{U}_{\mathrm{G}}/\mathrm{G}\), de sorte que pour tout faisceau équivariant \(\mathrm{F}\) sur \(\mathrm{X}\), \(\operatorname{pr}_1^*\mathrm{F}\) est image réciproque d’un faisceau \(\mathrm{F}^{\mathrm{G}}\) sur \(\mathrm{X}\times\mathrm{U}_{\mathrm{G}}/\mathrm{G}\). Puisque \(\mathrm{U}_{\mathrm{G}}\) est contractile et d’après (6.1.2.2), on a
\[
\mathrm{H}^{\bullet}(\mathrm{X},\mathrm{G};\mathrm{F})
\xrightarrow{\sim}
\mathrm{H}^{\bullet}(\mathrm{X}\times\mathrm{U}_{\mathrm{G}},\mathrm{G};
\operatorname{pr}_1^*\mathrm{F})
\xleftarrow{\sim}
\mathrm{H}^{\bullet}(\mathrm{X}\times\mathrm{U}_{\mathrm{G}}/\mathrm{G},\mathrm{F}^{\mathrm{G}}).
\tag*{\textbf{(6.1.4.1)}}
\]

En particulier, pour \(\mathrm{X}=\mathrm{P}^{t}\)
\[
\mathrm{H}^{\bullet}(\mathrm{B}_{\bullet\mathrm{G}})
=\mathrm{H}^{\bullet}(\mathrm{B}_{\mathrm{G}}).
\tag*{\textbf{(6.1.4.2)}}
\]

\Page{18}
On vérifie que l’isomorphisme (6.1.4.2) est le cas particulier de (6.1.3.1) pour \(\mathrm{S}=\mathrm{B}_{\mathrm{G}}\), \(\mathrm{P}=\mathrm{U}_{\mathrm{G}}\).

\textbf{(6.1.5)} La suite spectrale (5.2.3.2) pour \(\mathrm{B}_{\bullet\mathrm{G}}\)
\[
\mathrm{E}_{1}^{pq}=\mathrm{H}^{q}(\mathrm{G}^{\Delta_p}/\mathrm{G})
\Rightarrow
\mathrm{H}^{p+q}(\mathrm{B}_{\bullet\mathrm{G}})
=\mathrm{H}^{p+q}(\mathrm{B}_{\mathrm{G}})
\]
est essentiellement la \textit{suite spectrale d’Eilenberg-Moore}. Rappelons brièvement comment elle permet de relier les cohomologies rationnelles de \(\mathrm{G}\) et de \(\mathrm{B}_{\mathrm{G}}\), pour \(\mathrm{G}\) connexe.

\textit{a)} L’algèbre \(\mathrm{H}^{\bullet}(\mathrm{G},\mathbf{Q})\) est une algèbre de Hopf graduée connexe de dimension finie sur \(\mathbf{Q}\). Si \(\mathrm{P}^{\bullet}(\mathrm{G})\) est le module gradué de ses éléments primitifs, on a donc
\[
\mathrm{H}^{\bullet}(\mathrm{G},\mathbf{Q})
=\mathop{\textstyle\bigwedge}\mathrm{P}^{\bullet}(\mathrm{G}),
\]
et les générateurs de \(\mathrm{P}^{\bullet}(\mathrm{G})\) sont de degrés impairs.

\textit{b)} L’algèbre cosimpliciale \((\mathrm{E}_{1}^{p\bullet})_{p\geq 0}\) est
\[
\mathrm{E}_{1}^{p\bullet}
=\mathop{\textstyle\bigwedge}
(\mathrm{P}^{\bullet}(\mathrm{G})^{\Delta_p}/\mathrm{P}^{\bullet}(\mathrm{G}));
\]
c’est l’algèbre extérieure de la suspension du module cosimplicial constant \(\mathrm{P}^{\bullet}(\mathrm{G})\); on a donc (Quillen [12]) \(\mathrm{E}_{2}^{p\bullet}=\operatorname{Sym}^{p}(\mathrm{P}^{\bullet}(\mathrm{G}))\). Les termes \(\mathrm{E}_{2}^{pq}\) ne sont non nuls que pour \(p+q\) pair; on a donc \(\mathrm{E}_{2}^{pq}=\mathrm{E}_{\infty}^{pq}\) et, pour une filtration convenable, on a canoniquement
\[
\operatorname{Gr}\mathrm{H}^{\bullet}(\mathrm{B}_{\mathrm{G}},\mathbf{Q})
\simeq\operatorname{Sym}^{\bullet}(\mathrm{P}^{\bullet}(\mathrm{G})[-1])
\]
et non canoniquement
\[
\mathrm{H}^{\bullet}(\mathrm{B}_{\mathrm{G}},\mathbf{Q})
\simeq\operatorname{Sym}^{\bullet}(\mathrm{P}^{\bullet}(\mathrm{G})[-1]).
\]

\textbf{(6.1.6)} Prenons pour \(\mathrm{G}\) un groupe algébrique linéaire (complexe).

\textit{a)} Si \(\mathrm{T}\) est un tore maximal de \(\mathrm{G}\), de groupe de Weyl \(\mathrm{W}\), on a
\[
\mathrm{H}^{\bullet}(\mathrm{B}_{\mathrm{G}},\mathbf{Q})
\xrightarrow{\sim}
\mathrm{H}^{\bullet}(\mathrm{B}_{\mathrm{T}},\mathbf{Q})^{\mathrm{W}}.
\tag*{\textbf{(6.1.6.1)}}
\]

\textit{b)} Pour un tore \(\mathrm{T}\) de groupe de caractères \(\mathrm{X}(\mathrm{T})\), on a
\[
\mathrm{H}^{\bullet}(\mathrm{T},\mathbf{Z})
\simeq\overset{\bullet}{\mathop{\textstyle\bigwedge}}\mathrm{X}(\mathrm{T})
\tag*{\textbf{(6.1.6.2)}}
\]
(isomorphisme d’algèbres de Hopf graduées).

Nous n’utiliserons \textit{a)} que sous la forme affaiblie suivante :

\(a')\) L’application \(\mathrm{H}^{\bullet}(\mathrm{B}_{\mathrm{G}},\mathbf{Q})\to\mathrm{H}^{\bullet}(\mathrm{B}_{\mathrm{T}},\mathbf{Q})\) est injective \textit{(splitting principle)}.

Pour être complet, rappelons une démonstration de \(a')\).

Si \(\mathrm{B}\) est un sous-groupe de Borel de \(\mathrm{G}\), le fibré \(\mathrm{U}_{\mathrm{G}}/\mathrm{B}\) sur \(\mathrm{B}_{\mathrm{G}}\) est un fibré en espaces de drapeaux. D’après ([3] \(2.1+2.6.3\)) mieux expliqué dans P. A. Griffiths, Periods III, \textit{Publ. Math. I.H.E.S.}, \textbf{38}, prop. (3.1), ou d’après [1], la suite spectrale de Leray de \(\mathrm{U}_{\mathrm{G}}/\mathrm{B}\to\mathrm{B}_{\mathrm{G}}\) dégénère en cohomologie rationnelle. On a donc \(\mathrm{H}^{\bullet}(\mathrm{B}_{\mathrm{G}},\mathbf{Q})\hookrightarrow\mathrm{H}^{\bullet}(\mathrm{U}_{\mathrm{G}}/\mathrm{B},\mathbf{Q})\). On conclut en notant que \(\mathrm{U}_{\mathrm{G}}/\mathrm{B}\sim\mathrm{B}_{\mathrm{B}}\sim\mathrm{B}_{\mathrm{T}}\).

\Page{19}
\section*{6.2. Construction d’hyperrecouvrements.}

\textbf{(6.2.1)} Soit \(\mathrm{X}_{\bullet}\) un ensemble simplicial. Désignons par \(\mathrm{D}(\Delta_n,\Delta_m)\) l’ensemble des applications surjectives croissantes de \(\Delta_n\) dans \(\Delta_m\) (opérateurs de dégénérescence), et posons
\[
\mathrm{N}(\mathrm{X}_n)
=\mathrm{X}_n-\bigcup_{s\in\mathrm{D}(\Delta_n,\Delta_{n-1})}
s(\mathrm{X}_{n-1}).
\tag*{\textbf{(6.2.1.1)}}
\]

Rappelons que pour tout \(n\), l’application
\[
\coprod s:
\coprod_{m\leq n}\coprod_{s\in\mathrm{D}(\Delta_n,\Delta_m)}
\mathrm{N}(\mathrm{X}_m)\to\mathrm{X}_n
\tag*{\textbf{(6.2.1.2)}}
\]
est bijective.

\textit{Définition} \textbf{(6.2.2)}. — \textit{Un espace topologique simplicial est dit \textup{\(s\)-scindé} si les applications \textup{(6.2.1.2)} sont des homéomorphismes.}

Soit \(\mathrm{X}_{\bullet}\) un ensemble simplicial \(k\)-tronqué. Pour \(n\leq k\), on définit encore \(\mathrm{N}(\mathrm{X}_n)\) par (6.2.1.1), et (6.2.1.2) est une bijection. Un espace topologique simplicial \(k\)-tronqué est dit \textit{\(s\)-scindé} si (6.2.1.2) est un homéomorphisme pour \(n\leq k\).

\textbf{(6.2.3)} Pour \(\mathrm{X}\) un espace topologique simplicial \((n+1)\)-tronqué \(s\)-scindé, augmenté vers \(\mathrm{S}\), soit \(\alpha(\mathrm{X})\) le triplet consistant en \(\operatorname{sq}_n(\mathrm{X})\), \(\mathrm{N}\mathrm{X}_{n+1}\) et en l’application évidente de \(\mathrm{N}\mathrm{X}_{n+1}\) dans \((\operatorname{cosq}\operatorname{sq}_n\mathrm{X})_{n+1}\).

Ce triplet \((\mathrm{Y},\mathrm{N},\beta)\) vérifie

\((\ast)\) \(\mathrm{Y}\) est un espace topologique simplicial \(n\)-tronqué \(s\)-scindé augmenté vers \(\mathrm{S}\), et \(\beta\) est une application continue de \(\mathrm{N}\) dans \((\operatorname{cosq}\mathrm{Y})_{n+1}\).

\textit{Proposition} \textbf{(6.2.4)} (\textit{S.G.A.}4, V \textit{bis} (5.1.3)). — \textit{Soit \((\mathrm{Y},\mathrm{N},\beta)\) vérifiant \((\ast)\) ci-dessus.}

(i) \textit{A isomorphisme unique près, il existe un et un seul espace topologique \(\mathrm{X}\), \((n+1)\)-tronqué \(s\)-scindé augmenté vers \(\mathrm{S}\), avec \(\alpha(\mathrm{X})\simeq(\mathrm{Y},\mathrm{N},\beta)\).}

(ii) \textit{Il revient au même de se donner \(f:\mathrm{X}\to\mathrm{Z}\) ou de se donner :}

a) \textit{un morphisme \(f':\mathrm{Y}\to\operatorname{sq}_n(\mathrm{Z})\),}

b) \textit{un morphisme \(f'':\mathrm{N}\to\mathrm{Z}_{n+1}\) rendant commutatif le diagramme}
\[
\begin{tikzcd}[column sep=3em,row sep=3em]
\mathrm{N}\arrow[r,"\beta"]\arrow[d,"f''"']
  & (\operatorname{cosq}\mathrm{Y})_{n+1}\arrow[d,"f'"']\\
\mathrm{Z}_{n+1}\arrow[r]
  & (\operatorname{cosq}\operatorname{sq}_n\mathrm{Z})_{n+1}
\end{tikzcd}
\]

Cette proposition (6.2.4) s’applique aussi aux objets simpliciaux d’autres catégories \(\mathcal{C}\) que celle des espaces topologiques. Il suffit que \(\mathcal{C}\) vérifie
\Page{20}

\textbf{(6.2.4.1)} Les limites projectives finies existent. Les sommes finies existent, sont disjointes et universelles.

\textbf{(6.2.5)} Cette proposition permet de construire comme suit, par induction, des hyperrecouvrements propres de \(\mathrm{S}\).

0) On prend \(f_0:\mathrm{X}_0\to\mathrm{S}\), propre et surjectif. \(\{\mathrm{X}_0\}\) est un hyperrecouvrement propre \(0\)-tronqué de \(\mathrm{S}\) (5.3.8), il est \(s\)-scindé.

1) On prend \(f_1:\mathrm{N}_1\to\operatorname{cosq}(\{\mathrm{X}_0\})_1\), i.e. \(f_1:\mathrm{N}_1\to\mathrm{X}_0\times_{\mathrm{S}}\mathrm{X}_0\). Appliquant (6.2.4), on associe à \(f_1\) l’espace topologique simplicial \(1\)-tronqué augmenté \(s\)-scindé
\[
\begin{tikzcd}[column sep=2em]
{}_{1}\mathrm{X}_{\bullet}:\ \mathrm{N}_1\amalg\mathrm{X}_0
  \arrow[r,shift left=1.1ex]\arrow[r,shift right=1.1ex]
  & \mathrm{X}_0\arrow[l]\arrow[r] & \mathrm{S}.
\end{tikzcd}
\]
On suppose \(f_1\) choisi de telle sorte que
\[
f'_1:\mathrm{N}_1\amalg\mathrm{X}_0
\to\mathrm{X}_0\times_{\mathrm{S}}\mathrm{X}_0
\]
soit propre et surjectif (par exemple \(f_1\) propre et surjectif). Alors, \({}_{1}\mathrm{X}_{\bullet}\) est un hyperrecouvrement propre \(1\)-tronqué \(s\)-scindé de \(\mathrm{S}\).

\noindent\(\ldotp\;\ldotp\;\ldotp\;\ldotp\;\ldotp\;\ldotp\;\ldotp\;\ldotp\;\ldotp\;\ldotp\;\ldotp\;\ldotp\;\ldotp\)

\(k+1)\) Soit déjà construit un hyperrecouvrement propre \(k\)-tronqué \(s\)-scindé \({}_{k}\mathrm{X}_{\bullet}\to\mathrm{S}\). On prend \(f_{k+1}:\mathrm{N}_{k+1}\to(\operatorname{cosq}({}_{k}\mathrm{X}_{\bullet}))_{k+1}\), et, appliquant (6.2.4), on associe à \(f_{k+1}\) un espace topologique semi-simplicial \(k+1\)-tronqué augmenté \(s\)-scindé \({}_{k+1}\mathrm{X}_{\bullet}\). On suppose que \(f_{k+1}\) est tel que
\[
f'_{k+1}:{}_{k+1}\mathrm{X}_{k+1}
\to\operatorname{cosq}({}_{k}\mathrm{X}_{\bullet})_{k+1}
\]
soit propre et surjectif (par exemple \(f_{k+1}\) propre et surjectif). Alors, \({}_{k+1}\mathrm{X}_{\bullet}\) est un hyperrecouvrement propre \(k+1\)-tronqué \(s\)-scindé de \(\mathrm{S}\).

\noindent\(\ldotp\;\ldotp\;\ldotp\;\ldotp\;\ldotp\;\ldotp\;\ldotp\;\ldotp\;\ldotp\;\ldotp\;\ldotp\;\ldotp\;\ldotp\)

\(\infty)\) Les \({}_{k}\mathrm{X}_{\bullet}\) ainsi construits sont les squelettes successifs d’un hyperrecouvrement propre \(s\)-scindé de \(\mathrm{S}\).

\textbf{(6.2.6)} Un schéma simplicial \(\mathrm{X}_{\bullet}\) sur \(\mathbf{C}\) sera dit \textit{lisse} si les \(\mathrm{X}_n\) sont lisses; il sera dit \textit{propre} si les \(\mathrm{X}_n\) sont compacts. Un \textit{diviseur à croisements normaux} \(\mathrm{D}_{\bullet}\) de \(\mathrm{X}_{\bullet}\), supposé lisse, est une famille \(\mathrm{D}_n\subset\mathrm{X}_n\) de diviseurs à croisements normaux (3.1.2) telle que les \(\mathrm{U}_n=\mathrm{X}_n-\mathrm{D}_n\) forment un sous-schéma simplicial \(\mathrm{U}_{\bullet}\) de \(\mathrm{X}_{\bullet}\). Cette définition est justifiée par le lemme suivant.

\textit{Lemme} \textbf{(6.2.7)}. — \textit{Si \(\mathrm{D}_{\bullet}\) est un diviseur à croisements normaux de \(\mathrm{X}_{\bullet}\), alors les complexes de De Rham logarithmiques \((\Omega_{\mathrm{X}_n}^{\bullet}(\log\mathrm{D}_n))_{n\geq 0}\), munis de la filtration par le poids \textup{(3.1.5)}, forment un complexe filtré sur \(\mathrm{X}_{\bullet}\).}

Résulte de (3.1.3 (ii)). Le complexe \((\Omega_{\mathrm{X}_n}^{\bullet}(\log\mathrm{D}_n))_{n\geq 0}\) se notera \(\Omega_{\mathrm{X}_{\bullet}}^{\bullet}(\log\mathrm{D}_{\bullet})\).

\textbf{(6.2.8)} En utilisant (6.2.5), on montre que pour tout schéma séparé \(\mathrm{S}\) sur \(\mathbf{C}\), il existe :

% ---- 021-028 ----
\Page{21}
\noindent
\(\alpha)\) un schéma simplicial \(\mathrm{X}_{\bullet}\) sur \(\mathbf{C}\), propre et lisse, qu’on peut prendre \(s\)-scindé;

\noindent
\(\beta)\) un diviseur à croisements normaux \(\mathrm{D}_{\bullet}\) dans \(\mathrm{X}_{\bullet}\); posons \(\mathrm{U}_{\bullet}=\mathrm{X}_{\bullet}-\mathrm{D}_{\bullet}\);

\noindent
\(\gamma)\) une augmentation \(a:\mathrm{U}_{\bullet}\to\mathrm{S}\) qui fasse de \(\mathrm{U}_{\bullet}^{\mathrm{an}}\) un hyperrecouvrement propre de \(\mathrm{S}^{\mathrm{an}}\).

De plus, deux tels systèmes sont coiffés par un même troisième, un morphisme \(u:\mathrm{S}\to\mathrm{T}\) peut être coiffé par un morphisme
\[
\begin{tikzcd}[column sep=1.3em,row sep=1.8em]
& \mathrm{X}_{\bullet} \arrow[rr] & & \mathrm{Y}_{\bullet} \\
\mathrm{U}_{\bullet} \arrow[ur,hook] \arrow[rr] \arrow[d,"a"']
& & \mathrm{V}_{\bullet} \arrow[ur,hook] \arrow[d,"b"'] & \\
\mathrm{S} \arrow[rr,"u"] & & \mathrm{T} &
\end{tikzcd}
\]
de tels systèmes... (voir \textit{S.G.A.}4, V \textit{bis}).

\section*{6.3. Cohomologie relative.}

\textbf{(6.3.1)} La construction \textit{cône d’application} pour les morphismes d’ensembles simpliciaux se transpose telle quelle aux objets simpliciaux de toute catégorie \(\mathscr{C}\) ayant un objet final \(e\) et des sommes finies. Pour \(u:\mathrm{Y}_{\bullet}\to\mathrm{X}_{\bullet}\), le cône \(\mathrm{C}(u)\) vérifie
\[
\mathrm{C}(u)_n
=\mathrm{X}_n\amalg\coprod_{i<n}\mathrm{Y}_i\amalg e.
\]

Nous prendrons pour \(\mathscr{C}\) :

\noindent
\textit{a)} la catégorie des espaces topologiques et applications continues (objet final : \(\mathrm{P}^{t}\));

\noindent
\textit{b)} la catégorie des couples \((\mathrm{X},\mathrm{F})\) formés d’un espace topologique \(\mathrm{X}\) et d’un faisceau abélien \(\mathrm{F}\) sur \(\mathrm{X}\), une flèche \((u,f):(\mathrm{Y},\mathrm{F})\to(\mathrm{X},\mathrm{G})\) étant une application continue \(u:\mathrm{Y}\to\mathrm{X}\), plus un \(u\)-morphisme (5.1.5) \(f:\mathrm{G}\to\mathrm{F}\) (objet final : \((\mathrm{P}^{t},0)\)).

\textbf{(6.3.2)} Soient \(u:\mathrm{Y}_{\bullet}\to\mathrm{X}_{\bullet}\) un morphisme d’espaces topologiques simpliciaux, de cône \(\mathrm{C}(u)\), \(\mathrm{F}\) un faisceau abélien sur \(\mathrm{X}_{\bullet}\), \(\mathrm{G}\) un faisceau abélien sur \(\mathrm{Y}_{\bullet}\) et \(f:\mathrm{F}\to\mathrm{G}\) un \(u\)-morphisme. Le cône \(\mathrm{C}(f)\) de \(f\) est un faisceau abélien sur \(\mathrm{C}(u)\), et on pose
\[
\mathrm{H}^{n}(\mathrm{C}(u),\mathrm{C}(f))
=\mathrm{H}^{n}(\mathrm{X}_{\bullet}\operatorname{mod}\mathrm{Y}_{\bullet},
\mathrm{F}\operatorname{mod}\mathrm{G}).
\tag*{\textbf{(6.3.2.1)}}
\]
Ce sont les \textit{groupes de cohomologie relative}. On vérifie facilement qu’ils s’insèrent dans une suite exacte longue
\[
\ldots\to
\mathrm{H}^{i}(\mathrm{X}_{\bullet}\operatorname{mod}\mathrm{Y}_{\bullet},
\mathrm{F}\operatorname{mod}\mathrm{G})
\to\mathrm{H}^{i}(\mathrm{X}_{\bullet},\mathrm{F})
\to\mathrm{H}^{i}(\mathrm{Y}_{\bullet},\mathrm{G})
\to\ldots
\tag*{\textbf{(6.3.2.2)}}
\]

\textbf{(6.3.3)} Plus généralement, soient \(\mathrm{L}\) et \(\mathrm{K}\) des complexes bornés inférieurement de faisceaux abéliens sur \(\mathrm{Y}_{\bullet}\) et \(\mathrm{X}_{\bullet}\), et soit un \(u\)-morphisme \(f:\mathrm{K}\to\mathrm{L}\). On en déduit un complexe \(\mathrm{C}(f)\) sur \(\mathrm{C}(u)\). On pose encore en hypercohomologie
\[
\mathbf{H}^{n}(\mathrm{C}(u),\mathrm{C}(f))
=\mathbf{H}^{n}(\mathrm{X}_{\bullet}\operatorname{mod}\mathrm{Y}_{\bullet},
\mathrm{K}\operatorname{mod}\mathrm{L}).
\]
\Page{22}
Ces groupes figurent dans une suite exacte analogue à (6.3.2.2), provenant, dans la catégorie dérivée convenable, d’un triangle distingué
\[
\begin{tikzcd}[column sep=0.5em,row sep=2em]
& \mathrm{R}\Gamma(\mathrm{Y}_{\bullet},\mathrm{L})
  \arrow[dl,"\partial"'] & \\
\mathrm{R}\Gamma(\mathrm{C}(u),\mathrm{C}(f)) \arrow[rr]
& & \mathrm{R}\Gamma(\mathrm{X}_{\bullet},\mathrm{K}) \arrow[ul]
\end{tikzcd}
\tag*{\textbf{(6.3.3.1)}}
\]

\textbf{(6.3.4)} La construction présentée plus haut n’est pas la seule possible. Elle a l’inconvénient que, même si on part de vrais espaces topologiques \(\mathrm{X}\) et \(\mathrm{Y}\) (i.e. d’espaces simpliciaux constants), on est conduit à considérer des espaces simpliciaux non constants.

\section*{6.4. Espaces multisimpliciaux.}

\textbf{(6.4.1)} Soit \(r\) un entier \(\geq 0\). Un \textit{objet \(r\)-simplicial} \(\mathrm{Z}_{\bullet}\) d’une catégorie \(\mathscr{C}\) est un foncteur contravariant de \((\Delta)^r\) dans \(\mathscr{C}\). L’objet simplicial \textit{diagonal} \(\delta\mathrm{Z}_{\bullet}\) est le foncteur composé \((\Delta)\to(\Delta)^r\to\mathscr{C}\).

\textbf{(6.4.2)} On définit comme en (5.1.6) le topos des \textit{faisceaux} sur un espace topologique \(r\)-simplicial. Soit \(\Gamma^{\bullet}\) le foncteur
\[
\mathrm{F}\mapsto
(\Gamma(\mathrm{X}_{n_1\ldots n_r},\mathrm{F}^{n_1,\ldots,n_r}))
\]
des faisceaux sur \(\mathrm{X}_{\bullet}\) dans les ensembles \(r\)-cosimpliciaux. Pour \(r\) petit, on le notera souvent plutôt \(\Gamma^{\bullet\cdots\bullet}\) (\(r\) points). On dispose d’une co-augmentation \(\Gamma(\mathrm{X}_{\bullet},\mathrm{F}^{\bullet})\to\Gamma^{\bullet}(\mathrm{X}_{\bullet},\mathrm{F}^{\bullet})\). Les foncteurs \textit{cohomologie à valeur dans un faisceau abélien} \(\mathrm{F}\) sont les dérivés du foncteur « section globale » \(\Gamma\). Ils peuvent se calculer par un procédé parallèle à (5.2.3) :
\[
\mathrm{R}\Gamma=\boldsymbol{s}\mathrm{R}\Gamma^{\bullet}:
\mathrm{D}^{+}(\mathrm{X}_{\bullet})\to\mathrm{D}^{+}((\mathrm{Ab})).
\tag*{\textbf{(6.4.2.1)}}
\]

Un faisceau \(\mathrm{F}\) sur un espace topologique \(r\)-simplicial \(\mathrm{X}_{\bullet}\) induit un faisceau \(\delta\mathrm{F}\) sur l’espace simplicial diagonal \(\delta\mathrm{X}_{\bullet}\). Il résulte du théorème de Cartier-Eilenberg-Zilber que
\[
\mathrm{R}\Gamma(\mathrm{X}_{\bullet},\mathrm{K})
\xrightarrow{\sim}
\mathrm{R}\Gamma(\delta\mathrm{X}_{\bullet},\delta\mathrm{K}).
\tag*{\textbf{(6.4.2.2)}}
\]

\textbf{(6.4.3)} Limitons-nous au cas \(r=2\). Un objet bisimplicial 1-augmenté vers un objet simplicial \(\mathrm{S}_{\bullet}\) est un foncteur contravariant de \((\Delta^{+})\times(\Delta)\) dans \(\mathscr{C}\), tel que \(\mathrm{S}_{\bullet}\) soit le foncteur composé
\[
(\Delta)\xrightarrow{(\Delta_{-1},\bullet)}
(\Delta^{+})\times(\Delta)\longrightarrow\mathscr{C}.
\]
Pour désigner un objet bisimplicial 1-augmenté vers \(\mathrm{S}_{\bullet}\), d’objet bisimplicial sous-jacent \(\mathrm{X}_{\bullet\bullet}\), on utilisera une notation du type
\[
a:\mathrm{X}_{\bullet\bullet}\to\mathrm{S}_{\bullet}.
\]
\Page{23}
Pour \(n\geq 0\), \(a:(\mathrm{X}_{\bullet n})\to\mathrm{S}_n\) est un objet simplicial augmenté vers \(\mathrm{S}_n\). Si \(\mathrm{F}\) est un faisceau sur \(\mathrm{X}_{\bullet\bullet}\), les \((a_{\ast}(\mathrm{F}\mid\mathrm{X}_{\bullet n})\text{ sur }\mathrm{S}_n)_{n\geq 0}\) forment un faisceau sur \(\mathrm{S}_{\bullet}\); on définit ainsi un morphisme de topos
\[
a:\mathrm{X}_{\bullet\bullet}^{\sim}\to\mathrm{S}_{\bullet}^{\sim}.
\]

On explique dans \textit{S.G.A.}4, V \textit{bis} que \(\mathrm{R}a_{\ast}\) se calcule « composante par composante » :
\[
\mathrm{R}a_{\ast}\mathrm{K}\mid\mathrm{S}_n
=\mathrm{R}a_{\ast}(\mathrm{K}\mid\mathrm{X}_{\bullet n}).
\]
Il en résulte que si, pour chaque \(n\), \(a:(\mathrm{X}_{\bullet n})\to\mathrm{S}_n\) est de descente cohomologique, alors \(a:\mathrm{X}_{\bullet\bullet}\to\mathrm{S}_{\bullet}\) est de descente cohomologique : pour tout complexe de faisceaux abéliens \(\mathrm{K}\in\mathrm{D}^{+}(\mathrm{S}_{\bullet})\), on a
\[
\mathrm{K}\xrightarrow{\approx}\mathrm{R}a_{\ast}a^{\ast}\mathrm{K}.
\]

\textbf{(6.4.4)} Dans \textit{S.G.A.}4, V \textit{bis}, on montre que pour tout schéma simplicial séparé \(\mathrm{S}_{\bullet}\), il existe un schéma bisimplicial \(\mathrm{X}_{\bullet\bullet}\) 1-augmenté vers \(\mathrm{S}_{\bullet}\), et \(i:\mathrm{X}_{\bullet\bullet}\hookrightarrow\overline{\mathrm{X}}_{\bullet\bullet}\) tels que :

\textit{a)} Les \(\overline{\mathrm{X}}_{nm}\) sont projectifs et lisses; \(\mathrm{X}_{nm}\) est le complément d’un diviseur à croisements normaux \(\mathrm{D}_{nm}\) dans \(\overline{\mathrm{X}}_{nm}\). On peut supposer \(\mathrm{D}_{nm}\) réunion de diviseurs lisses.

\textit{b)} Pour \(n\geq 0\), \(\mathrm{X}_{\bullet n}\) est un hyperrecouvrement propre de \(\mathrm{S}_n\). On peut le supposer \(s\)-scindé.

La construction procède comme en (6.2), mais la récurrence est plus compliquée. Les assertions d’unicité (6.2.8) restent valables, \textit{mutatis mutandis}.

\section*{7. COMPLÉMENTS AU § 1}

\section*{7.1. Catégorie dérivée filtrée.}

Ce numéro complète le § 1, n\textsuperscript{o} 4.

\textbf{(7.1.1)} Soit \(\mathscr{A}\) une catégorie abélienne. Posons :

\noindent
\(\mathrm{F}\mathscr{A}\) (resp. \(\mathrm{F}_2\mathscr{A}\)) \(=\) la catégorie des objets filtrés (resp. bifiltrés) de filtration(s) finie(s) de \(\mathscr{A}\).

\noindent
\(\mathrm{K}^{+}\mathrm{F}\mathscr{A}\) (resp. \(\mathrm{K}^{+}\mathrm{F}_2\mathscr{A}\)) \(=\) la catégorie des complexes filtrés (resp. bifiltrés) bornés inférieurement d’objets de \(\mathscr{A}\), à homotopie respectant la filtration (resp. les filtrations) près.

\noindent
\(\mathrm{D}^{+}\mathrm{F}\mathscr{A}\) (resp. \(\mathrm{D}^{+}\mathrm{F}_2\mathscr{A}\)) \(=\) la catégorie triangulée déduite de \(\mathrm{K}^{+}\mathrm{F}\mathscr{A}\) (resp. \(\mathrm{K}^{+}\mathrm{F}_2\mathscr{A}\)) en inversant les quasi-isomorphismes filtrés (resp. bifiltrés) (1.3.6). C’est la \textit{catégorie dérivée filtrée} (resp. \textit{bifiltrée}).

\textbf{(7.1.2)} Un quasi-isomorphisme filtré \(u:(\mathrm{K},\mathrm{F})\to(\mathrm{K}',\mathrm{F}')\) induit un isomorphisme de suites spectrales \(u:\mathrm{E}_{\bullet}^{\bullet\bullet}(\mathrm{K},\mathrm{F})\to\mathrm{E}_{\bullet}^{\bullet\bullet}(\mathrm{K}',\mathrm{F}')\). Un objet \(\mathrm{K}\) de \(\mathrm{D}^{+}\mathrm{F}\mathscr{A}\) définit donc
\Page{24}
une suite spectrale \(\mathrm{E}_{\bullet}^{\bullet\bullet}(\mathrm{K})\). De même, un objet \(\mathrm{L}\) de \(\mathrm{D}^{+}\mathrm{F}_2\mathscr{A}\) définit un amoncellement de suites spectrales du type considéré en (1.4.9).

\textbf{(7.1.3)} Soit \(\mathrm{T}\) un foncteur exact à gauche de \(\mathscr{A}\) dans une catégorie abélienne \(\mathscr{B}\). On suppose que tout objet de \(\mathscr{A}\) s’injecte dans un objet injectif. Le foncteur \(\mathrm{T}\) se « dérive » alors en des foncteurs
\[
\mathrm{RT}:\mathrm{D}^{+}(\mathscr{A})\to\mathrm{D}^{+}(\mathscr{B})
\qquad\text{(pour mémoire)}
\tag*{\textbf{(7.1.3.1)}}
\]
\[
\mathrm{RT}:\mathrm{D}^{+}\mathrm{F}(\mathscr{A})
\to\mathrm{D}^{+}\mathrm{F}(\mathscr{B}),
\tag*{\textbf{(7.1.3.2)}}
\]
\[
\mathrm{RT}:\mathrm{D}^{+}\mathrm{F}_2(\mathscr{A})
\to\mathrm{D}^{+}\mathrm{F}_2(\mathscr{B}).
\tag*{\textbf{(7.1.3.3)}}
\]
Ils se calculent ainsi : si \(\mathrm{K}'\) est une résolution (resp. une résolution filtrée, resp. une résolution bifiltrée) \(\mathrm{T}\)-acyclique de \(\mathrm{K}\) ((1.4.5), (1.4.11)), alors \(\mathrm{RT}(\mathrm{K})=\mathrm{T}(\mathrm{K}')\).

La suite spectrale d’hypercohomologie (pour \(\mathrm{T}\)) de \(\mathrm{K}\in\operatorname{Ob}\mathrm{D}^{+}\mathrm{F}(\mathscr{A})\) est la suite spectrale de \(\mathrm{RTK}\in\operatorname{Ob}\mathrm{D}^{+}\mathrm{F}(\mathscr{B})\) (cf. (1.4.5)).

\textbf{(7.1.4)} Nous aurons besoin de résultats plus précis pour les foncteurs \(\mathrm{R}a_{\ast}\), \(a:\mathrm{X}_{\bullet}\to\mathrm{S}\) étant une augmentation d’un espace topologique simplicial. Le cas \(\mathrm{S}=\mathrm{P}^{t}\), \(\mathrm{R}a_{\ast}=\mathrm{R}\Gamma\) nous suffirait.

Reprenons les notations de (5.1.12), et rappelons la formule (5.2.6.1) :
\[
\mathrm{R}a_{\ast}=\boldsymbol{s}\mathrm{R}a_{\bullet\ast}.
\]
Pour tout complexe \(\mathrm{K}\in\operatorname{Ob}\mathrm{C}^{+}(\mathrm{S}_{\bullet})\), le simple complexe \(\boldsymbol{s}\mathrm{K}\) est muni d’une filtration naturelle \(\mathrm{L}\) (5.1.9.3). Un quasi-isomorphisme \(u:\mathrm{K}'\xrightarrow{\approx}\mathrm{K}''\) induit un quasi-isomorphisme filtré \(u:(\boldsymbol{s}\mathrm{K}',\mathrm{L})\to(\boldsymbol{s}\mathrm{K}'',\mathrm{L})\). Dès lors, \(\boldsymbol{s}\) se factorise par
\[
\boldsymbol{s}:\mathrm{D}^{+}(\mathrm{S}_{\bullet})
\to\mathrm{D}^{+}\mathrm{F}(\mathrm{S})
\]
et \(\mathrm{R}a_{\ast}\) se factorise par
\[
\mathrm{R}a_{\ast}:\mathrm{D}^{+}(\mathrm{X}_{\bullet})
\to\mathrm{D}^{+}\mathrm{F}(\mathrm{S}).
\]
La suite spectrale du complexe filtré \((\mathrm{R}a_{\ast}\mathrm{K},\mathrm{L})\in\mathrm{D}^{+}\mathrm{F}(\mathrm{S})\) n’est autre que (5.2.7.1).

\textbf{(7.1.5)} Si \(\mathrm{K}\) est filtré (resp. bifiltré), alors \(\mathrm{R}a_{\bullet\ast}\mathrm{K}\) est filtré (bifiltré) : on dispose de
\[
\mathrm{R}a_{\bullet\ast}:
\mathrm{D}^{+}\mathrm{F}(\mathrm{X}_{\bullet})
\to\mathrm{D}^{+}\mathrm{F}(\mathrm{S}_{\bullet})
\tag*{\textbf{(7.1.5.1)}}
\]
\[
\mathrm{R}a_{\bullet\ast}:
\mathrm{D}^{+}\mathrm{F}_2(\mathrm{X}_{\bullet})
\to\mathrm{D}^{+}\mathrm{F}_2(\mathrm{S}_{\bullet}).
\tag*{\textbf{(7.1.5.2)}}
\]

\textbf{(7.1.6)} Soit \(\mathrm{K}\) un complexe de faisceaux cosimpliciaux sur \(\mathrm{S}\), muni d’une filtration croissante \(\mathrm{W}\). On appelle \textit{filtration diagonale} \(\delta(\mathrm{W},\mathrm{L})\) de \(\mathrm{W}\) et \(\mathrm{L}\) la filtration croissante suivante de \(\boldsymbol{s}\mathrm{K}\)
\begin{align*}
\delta(\mathrm{W},\mathrm{L})_n(\boldsymbol{s}\mathrm{K})
&=\bigoplus_{p,q}\mathrm{W}_{n+p}(\mathrm{K}^{p,q})
\tag*{\textbf{(7.1.6.1)}}\\
&=\sum_p\boldsymbol{s}(\mathrm{W}_{n+p}(\mathrm{K}))
\cap\mathrm{L}^{p}(\boldsymbol{s}\mathrm{K}).
\end{align*}

\Page{25}
On a
\[
\operatorname{Gr}_{n}^{\delta(\mathrm{W},\mathrm{L})}
(\boldsymbol{s}\mathrm{K})
\simeq
\bigoplus_p\operatorname{Gr}_{n+p}^{\mathrm{W}}
(\mathrm{K}^{\bullet p})[-p].
\tag*{\textbf{(7.1.6.2)}}
\]

Le foncteur \((\mathrm{K},\mathrm{W})\mapsto
(\boldsymbol{s}\mathrm{K},\delta(\mathrm{W},\mathrm{L}))\) transforme quasi-isomorphismes filtrés en quasi-isomorphismes filtrés et définit
\[
(\boldsymbol{s},\delta):
(\mathrm{K},\mathrm{W})\mapsto
(\boldsymbol{s}\mathrm{K},\delta(\mathrm{W},\mathrm{L})):
\mathrm{D}^{+}\mathrm{F}(\mathrm{S}_{\bullet})
\to\mathrm{D}^{+}\mathrm{F}(\mathrm{S})
\tag*{\textbf{(7.1.6.3)}}
\]
d’où par composition avec (7.1.5.1)
\[
(\mathrm{R}\Gamma,\delta(\, ,\mathrm{L})):
(\mathrm{K},\mathrm{W})\mapsto
(\mathrm{R}\Gamma\mathrm{K},\delta(\mathrm{W},\mathrm{L})):
\mathrm{D}^{+}\mathrm{F}(\mathrm{X}_{\bullet})
\to\mathrm{D}^{+}\mathrm{F}(\mathrm{S}).
\tag*{\textbf{(7.1.6.4)}}
\]
On tire de (7.1.6.2) que
\[
\operatorname{Gr}_{n}^{\delta(\mathrm{W},\mathrm{L})}
(\mathrm{R}a_{\ast}\mathrm{K})
=\bigoplus_p\mathrm{R}a_{p\ast}
(\operatorname{Gr}_{n+p}^{\mathrm{W}}\mathrm{K})[-p].
\tag*{\textbf{(7.1.6.5)}}
\]

\textbf{(7.1.7)} Si \((\mathrm{K},\mathrm{W},\mathrm{F})\) est un complexe de faisceaux cosimpliciaux bifiltré, alors \(\boldsymbol{s}\mathrm{K}\) est muni des trois filtrations \(\mathrm{W}\), \(\mathrm{F}\) et \(\mathrm{L}\) et définit divers complexes bifiltrés.

Par exemple, pour \(\mathrm{W}\) croissante, le foncteur \((\mathrm{K},\mathrm{W},\mathrm{F})\mapsto
(\boldsymbol{s}\mathrm{K},\delta(\mathrm{W},\mathrm{L}),\mathrm{F})\) transforme quasi-isomorphisme bifiltré en quasi-isomorphisme bifiltré, donc définit
\[
(\mathrm{K},\mathrm{W},\mathrm{F})\mapsto
(\boldsymbol{s}\mathrm{K},\delta(\mathrm{W},\mathrm{L}),\mathrm{F}):
\mathrm{D}^{+}\mathrm{F}_2(\mathrm{S}_{\bullet})
\to\mathrm{D}^{+}\mathrm{F}_2(\mathrm{S}).
\tag*{\textbf{(7.1.7.1)}}
\]
Par composition avec (7.1.5.2), on en déduit
\[
(\mathrm{K},\mathrm{W},\mathrm{F})\mapsto
(\mathrm{R}\Gamma\mathrm{K},\delta(\mathrm{W},\mathrm{L}),\mathrm{F}):
\mathrm{D}^{+}\mathrm{F}_2(\mathrm{X}_{\bullet})
\to\mathrm{D}^{+}\mathrm{F}_2(\mathrm{S}),
\tag*{\textbf{(7.1.7.2)}}
\]
et on a
\[
\operatorname{Gr}_{n}^{\delta(\mathrm{W},\mathrm{L})}
(\mathrm{R}\Gamma\mathrm{K},\mathrm{F})
=\bigoplus_p\mathrm{R}a_{p\ast}
(\operatorname{Gr}_{n+p}^{\mathrm{W}}\mathrm{K},\mathrm{F})[-p]
\tag*{\textbf{(7.1.7.3)}}
\]
dans \(\mathrm{D}^{+}\mathrm{F}(\mathrm{S})\).

\section*{7.2. Compléments au lemme des deux filtrations.}

Dans ce numéro, nous donnons une nouvelle démonstration du lemme des deux filtrations (1.3.16) et quelques compléments.

\textbf{(7.2.1)} Soit \((\mathrm{K},\mathrm{F},\mathrm{W})\) un complexe bifiltré borné inférieurement d’objets d’une catégorie abélienne \(\mathscr{A}\). La filtration \(\mathrm{F}\) est supposée birégulière.

On dit que \((\mathrm{K},\mathrm{F},\mathrm{W})\) est \(\mathrm{F}\)-\textit{scindable} si le complexe filtré \((\mathrm{K},\mathrm{W})\) peut se représenter comme une somme de complexes filtrés \((\mathrm{K}_n,\mathrm{W}_n)_{n\in\mathbf{Z}}\)
\[
(\mathrm{K},\mathrm{W})=\bigoplus_n(\mathrm{K}_n,\mathrm{W}_n)
\]
avec
\[
\mathrm{F}^{n}\mathrm{K}=\bigoplus_{n'\geq n}\mathrm{K}_{n'}.
\]

\Page{26}
Soit \(r_0\) un entier \(\geq 0\), ou \(+\infty\). La condition suivante a été considérée en (1.3.16), (1.3.17) :

\textbf{(7.2.2; \(r_0\))} Pour tout entier positif ou nul \(r<r_0\), les différentielles \(d_r\) du complexe gradué \(\mathrm{E}_r(\mathrm{K},\mathrm{W})\) sont strictement compatibles à la filtration récurrente définie par \(\mathrm{F}\).

\textbf{(7.2.3)} Il est clair que si \((\mathrm{K},\mathrm{F},\mathrm{W})\) est \(\mathrm{F}\)-scindable, la condition (7.2.2; \(\infty\)) est vérifiée. Réciproquement, il semble que lorsque la condition (7.2.2; \(\infty\)) est vérifiée, tout se passe comme si le foncteur \(\operatorname{Gr}_{\mathrm{F}}\) était exact. Par exemple, nous montrerons que le \(\operatorname{Gr}_{\mathrm{F}}\) de la suite spectrale \(\mathrm{E}(\mathrm{K},\mathrm{W})\) s’identifie alors à la suite spectrale \(\mathrm{E}(\operatorname{Gr}_{\mathrm{F}}\mathrm{K},\mathrm{W})\), et que la suite spectrale \(\mathrm{E}(\mathrm{K},\mathrm{F})\) dégénère \((\mathrm{E}_1=\mathrm{E}_{\infty})\).

\textbf{(7.2.4)} On déduit aussitôt de la définition (1.3.8) que la première filtration directe de \(\mathrm{E}_r(\mathrm{K},\mathrm{W})\) est la filtration \(\mathrm{F}_d\) de \(\mathrm{E}_r(\mathrm{K},\mathrm{W})\) par les images
\[
\mathrm{F}_d^p(\mathrm{E}_r(\mathrm{K},\mathrm{W}))
=\operatorname{Im}(\mathrm{E}_r(\mathrm{F}^p\mathrm{K},\mathrm{W})
\to\mathrm{E}_r(\mathrm{K},\mathrm{W})).
\]
Dualement, la seconde filtration directe (1.3.9) de \(\mathrm{E}_r(\mathrm{K},\mathrm{W})\) est la filtration \(\mathrm{F}_{d^{*}}\) de \(\mathrm{E}_r(\mathrm{K},\mathrm{W})\) par les noyaux
\[
\mathrm{F}_{d^{*}}^p(\mathrm{E}_r(\mathrm{K},\mathrm{W}))
=\operatorname{Ker}(\mathrm{E}_r(\mathrm{K},\mathrm{W})
\to\mathrm{E}_r(\mathrm{K}/\mathrm{F}^p\mathrm{K},\mathrm{W})).
\]
La filtration récurrente \(\mathrm{F}_{\mathrm{rec}}\) de \(\mathrm{E}_r(\mathrm{K},\mathrm{W})\) est intermédiaire entre ces deux filtrations (1.3.13 (iii)).

\textit{Proposition} \textbf{(7.2.5)}. — \textit{Supposons que \((\mathrm{K},\mathrm{F},\mathrm{W})\) vérifie \((7.2.2;\,r_0)\). Alors}

\noindent
(i) \textit{\((\mathrm{F}^a\mathrm{K}/\mathrm{F}^b\mathrm{K},\mathrm{F},\mathrm{W})\) vérifie encore \((7.2.2;\,r_0)\).}

\noindent
(ii) \textit{Pour \(r\leq r_0\), la suite}
\[
0\to\mathrm{E}_r(\mathrm{F}^p\mathrm{K},\mathrm{W})
\to\mathrm{E}_r(\mathrm{K},\mathrm{W})
\to\mathrm{E}_r(\mathrm{K}/\mathrm{F}^p\mathrm{K},\mathrm{W})\to 0
\]
\textit{est exacte; pour \(r=r_0+1\), la suite}
\[
\mathrm{E}_r(\mathrm{F}^p\mathrm{K},\mathrm{W})
\to\mathrm{E}_r(\mathrm{K},\mathrm{W})
\to\mathrm{E}_r(\mathrm{K}/\mathrm{F}^p\mathrm{K},\mathrm{W})
\]
\textit{est exacte. En particulier pour \(r\leq r_0+1\), les deux filtrations directes et la filtration récurrente de \(\mathrm{E}_r(\mathrm{K},\mathrm{W})\) coïncident.}

Fixons un entier \(p\) et prouvons par récurrence sur \(r\) l’assertion

\((\ast)_r\) Si \(r\leq r_0\), \(\mathrm{E}_r(\mathrm{F}^p\mathrm{K},\mathrm{W})\) s’injecte dans \(\mathrm{E}_r(\mathrm{K},\mathrm{W})\); si \(r\leq r_0+1\), son image est \(\mathrm{F}_{\mathrm{rec}}^p(\mathrm{E}_r(\mathrm{K},\mathrm{W}))\).

L’assertion \((\ast)_0\) est toujours vraie. Admettons \((\ast)_r\) et prouvons \((\ast)_{r+1}\). On peut supposer que \(0\leq r\leq r_0\). On dispose d’un diagramme
\[
\begin{tikzcd}[column sep=2em,row sep=3em]
\mathrm{E}_r(\mathrm{K},\mathrm{W}) \arrow[r,"d_r"]
& \mathrm{E}_r(\mathrm{K},\mathrm{W}) \arrow[r,"d_r"]
& \mathrm{E}_r(\mathrm{K},\mathrm{W}) \\
\mathrm{E}_r(\mathrm{F}^p\mathrm{K},\mathrm{W}) \arrow[u,hook] \arrow[r]
& \mathrm{E}_r(\mathrm{F}^p\mathrm{K},\mathrm{W}) \arrow[u,hook] \arrow[r]
& \mathrm{E}_r(\mathrm{F}^p\mathrm{K},\mathrm{W}), \arrow[u,hook,"i"']
\end{tikzcd}
\]
\Page{27}
et l’image des inclusions verticales est \(\mathrm{F}_{\mathrm{rec}}^p(\mathrm{E}_r(\mathrm{K},\mathrm{W}))\). Parce que \(i\) est injectif
\begin{align*}
\mathrm{F}_{\mathrm{rec}}^p(\mathrm{E}_{r+1}(\mathrm{K},\mathrm{W}))
&=\operatorname{Im}(\operatorname{Ker}(
\mathrm{F}_{\mathrm{rec}}^p\mathrm{E}_r(\mathrm{K},\mathrm{W})
\xrightarrow{d_r}\mathrm{E}_r(\mathrm{K},\mathrm{W}))
\to\mathrm{E}_{r+1}(\mathrm{K},\mathrm{W}))\\
&=\operatorname{Im}(\operatorname{Ker}(
\mathrm{E}_r(\mathrm{F}^p\mathrm{K},\mathrm{W})
\xrightarrow{d_r}\mathrm{E}_r(\mathrm{F}^p\mathrm{K},\mathrm{W}))
\to\mathrm{E}_{r+1}(\mathrm{K},\mathrm{W}))\\
&=\operatorname{Im}(
\mathrm{E}_{r+1}(\mathrm{F}^p\mathrm{K},\mathrm{W})
\to\mathrm{E}_{r+1}(\mathrm{K},\mathrm{W})).
\end{align*}
Si \(r<r_0\), \(d_r\) est strictement compatible à \(\mathrm{F}_{\mathrm{rec}}\), d’où
\[
d_r\mathrm{E}_r(\mathrm{K},\mathrm{W})
\cap\mathrm{E}_r(\mathrm{F}^p\mathrm{K},\mathrm{W})
=d_r\mathrm{E}_r(\mathrm{F}^p\mathrm{K},\mathrm{W}),
\]
et \(\mathrm{E}_{r+1}(\mathrm{F}^p\mathrm{K},\mathrm{W})\) s’injecte donc dans \(\mathrm{E}_{r+1}(\mathrm{K},\mathrm{W})\). Ceci prouve \((\ast)_{r+1}\).

Les assertions \((\ast)_r\), jointes aux assertions duales, impliquent (ii).

Il résulte de (ii) que pour \(r<r_0\) et \(a\leq b\), \(\mathrm{E}_r(\mathrm{F}^b\mathrm{K},\mathrm{W})\) s’injecte dans \(\mathrm{E}_r(\mathrm{F}^a\mathrm{K},\mathrm{W})\) et que la différentielle \(d_r\) de \(\mathrm{E}_r(\mathrm{F}^a\mathrm{K},\mathrm{W})\) est strictement compatible à la première filtration directe de \(\mathrm{E}_r(\mathrm{F}^a\mathrm{K},\mathrm{W})\). On en déduit par récurrence sur \(r<r_0\) que la première filtration directe de \(\mathrm{E}_r(\mathrm{F}^a\mathrm{K},\mathrm{W})\) coïncide avec la filtration récurrente, et (i) en résulte.

\textbf{(7.2.6)} Si la condition (7.2.2; \(r_0\)) est vérifiée, et que \(r\leq r_0+1\), on désignera par \(\mathrm{F}\) la filtration \(\mathrm{F}_d=\mathrm{F}_{d^{*}}=\mathrm{F}_{\mathrm{rec}}\) de \(\mathrm{E}_r(\mathrm{K},\mathrm{W})\). Les suites exactes (7.2.5) (ii)
\[
\begin{tikzcd}[column sep=0.6em,row sep=1.8em]
0
& \mathrm{E}_r(\mathrm{K}/\mathrm{F}^p\mathrm{K},\mathrm{W}) \arrow[l]
& & \mathrm{E}_r(\mathrm{K}/\mathrm{F}^{p+1}\mathrm{K},\mathrm{W}) \arrow[ll]
& & 0 \\
& & \mathrm{E}_r(\mathrm{K},\mathrm{W}) \arrow[ul] \arrow[ur]
& & \mathrm{E}_r(\operatorname{Gr}_{\mathrm{F}}^p\mathrm{K},\mathrm{W})
    \arrow[ul] \arrow[ur]
& \\
0 \arrow[r]
& \mathrm{E}_r(\mathrm{F}^{p+1}\mathrm{K},\mathrm{W}) \arrow[ur] \arrow[rr]
& & \mathrm{E}_r(\mathrm{F}^p\mathrm{K},\mathrm{W}) \arrow[ul] \arrow[ur]
& & 0 \arrow[ul]
\end{tikzcd}
\]
définissent pour \(r\leq r_0\) un isomorphisme (autodual) compatible aux différentielles \(d_r\)
\[
\operatorname{Gr}_{\mathrm{F}}^p(\mathrm{E}_r(\mathrm{K},\mathrm{W}))
\simeq\mathrm{E}_r(\operatorname{Gr}_{\mathrm{F}}^p\mathrm{K},\mathrm{W})
\qquad(r\leq r_0).
\tag*{\textbf{(7.2.6.1)}}
\]

\textbf{(7.2.7)} Si la condition (7.2.2; \(\infty\)) est vérifiée, les suites
\[
0\to\mathrm{E}_r(\mathrm{F}^p\mathrm{K},\mathrm{W})
\to\mathrm{E}_r(\mathrm{K},\mathrm{W})
\to\mathrm{E}_r(\mathrm{K}/\mathrm{F}^p\mathrm{K},\mathrm{W})\to 0
\]
sont exactes pour tout \(r\). Si la filtration \(\mathrm{W}\) est birégulière, la suite
\[
0\to\mathrm{E}_{\infty}(\mathrm{F}^p\mathrm{K},\mathrm{W})
\to\mathrm{E}_{\infty}(\mathrm{K},\mathrm{W})
\to\mathrm{E}_{\infty}(\mathrm{K}/\mathrm{F}^p\mathrm{K},\mathrm{W})\to 0
\]
est donc exacte. On peut réécrire cette suite
\[
0\to\operatorname{Gr}_{\mathrm{W}}(\mathrm{H}(\mathrm{F}^p\mathrm{K}))
\to\operatorname{Gr}_{\mathrm{W}}(\mathrm{H}(\mathrm{K}))
\to\operatorname{Gr}_{\mathrm{W}}(\mathrm{H}(\mathrm{K}/\mathrm{F}^p\mathrm{K}))
\to 0.
\]

\textit{Lemme} \textbf{(7.2.7.1)}. — \textit{Soit \(\mathrm{L}\) un complexe muni d’une filtration birégulière \(\mathrm{G}\). Pour que \(\operatorname{Gr}_{\mathrm{G}}(\mathrm{L})\) soit acyclique, il faut et il suffit que \(\mathrm{L}\) soit acyclique et que les différentielles de \(\mathrm{L}\) soient strictement compatibles à la filtration \(\mathrm{G}\).}

C’est un cas particulier de (1.3.2).

\Page{28}
Appliquons ce lemme, en prenant pour \(\mathrm{L}\) le complexe
\[
0\to\mathrm{H}(\mathrm{F}^p\mathrm{K})
\to\mathrm{H}(\mathrm{K})
\to\mathrm{H}(\mathrm{K}/\mathrm{F}^p\mathrm{K})\to 0.
\]
On trouve que :

\noindent
\textit{a)} cette suite est exacte, i.e. les différentielles de \(\mathrm{K}\) sont strictement compatibles à la filtration \(\mathrm{F}\);

\noindent
\textit{b)} la filtration de \(\mathrm{H}(\mathrm{F}^p\mathrm{K})\), déduite de la filtration \(\mathrm{W}\) de \(\mathrm{F}^p\mathrm{K}\), coïncide avec la filtration induite par la filtration \(\mathrm{W}\) de \(\mathrm{H}(\mathrm{K})\);

\noindent
\textit{c)} un énoncé analogue vaut pour \(\mathrm{K}/\mathrm{F}^p\mathrm{K}\).

En conclusion :

\textit{Proposition} \textbf{(7.2.8)}. — \textit{Si \((\mathrm{K},\mathrm{F},\mathrm{W})\) vérifie \((7.2.2;\,\infty)\), la suite spectrale \(\mathrm{E}(\mathrm{K},\mathrm{F})\) dégénère \((\mathrm{E}_1=\mathrm{E}_{\infty})\). On a de plus un isomorphisme de suites spectrales}
\[
\operatorname{Gr}_{\mathrm{F}}^p(\mathrm{E}_{\bullet}(\mathrm{K},\mathrm{W}))
\simeq\mathrm{E}_{\bullet}(\operatorname{Gr}_{\mathrm{F}}^p\mathrm{K},\mathrm{W}).
\]

\section*{8. THÉORIE DE HODGE DES ESPACES ALGÉBRIQUES}

\section*{8.1. Complexes de Hodge.}

\textbf{(8.1.1)} Soit \(\mathrm{A}\) comme en (III.0.3). Un \(\mathrm{A}\)-\textit{complexe de Hodge} \(\mathrm{K}\) \textit{de poids \(n\)} consiste en :

\noindent
\(\alpha)\) un complexe \(\mathrm{K}_{\mathrm{A}}\in\operatorname{Ob}\mathrm{D}^{+}(\mathrm{A})\), tel que \(\mathrm{H}^{k}(\mathrm{K}_{\mathrm{A}})\) soit un \(\mathrm{A}\)-module de type fini pour tout \(k\);

\noindent
\(\beta)\) une filtration \(\mathrm{F}\) sur \(\mathrm{K}_{\mathrm{A}}\otimes\mathbf{C}\), i.e. un complexe filtré \((\mathrm{K}_{\mathbf{C}},\mathrm{F})\in\operatorname{Ob}\mathrm{D}^{+}\mathrm{F}(\mathbf{C})\), et un isomorphisme \(\alpha:\mathrm{K}_{\mathbf{C}}\simeq\mathrm{K}_{\mathrm{A}}\otimes\mathbf{C}\) dans \(\mathrm{D}^{+}(\mathbf{C})\).

Les axiomes suivants doivent être vérifiés :

\noindent
(CH 1) la différentielle \(d\) de \(\mathrm{K}_{\mathbf{C}}\) est strictement compatible à la filtration \(\mathrm{F}\); en d’autres termes, la suite spectrale définie par \((\mathrm{K}_{\mathbf{C}},\mathrm{F})\) dégénère en \(\mathrm{E}_1\) \((\mathrm{E}_1=\mathrm{E}_{\infty})\) (1.3.2);

\noindent
(CH 2) pour tout \(k\), la filtration \(\mathrm{F}\) sur \(\mathrm{H}^{k}(\mathrm{K}_{\mathbf{C}})\simeq\mathrm{H}^{k}(\mathrm{K}_{\mathrm{A}})\otimes\mathbf{C}\) définit une \(\mathrm{A}\)-structure de Hodge de poids \(n+k\) sur \(\mathrm{H}^{k}(\mathrm{K}_{\mathrm{A}})\) : la filtration \(\mathrm{F}\) est \((n+k)\)-opposée à sa complexe conjuguée (cela a un sens car \(\mathrm{A}\subset\mathbf{R}\)).

\textbf{(8.1.2)} Soit \(\mathrm{X}\) un espace topologique. Un \(\mathrm{A}\)-\textit{complexe de Hodge cohomologique de poids \(n\)} sur \(\mathrm{X}\), \(\mathrm{K}\), consiste en :

\noindent
\(\alpha)\) un complexe \(\mathrm{K}_{\mathrm{A}}\in\operatorname{Ob}\mathrm{D}^{+}(\mathrm{X},\mathrm{A})\),

\noindent
\(\beta)\) un complexe filtré \((\mathrm{K}_{\mathbf{C}},\mathrm{F})\in\operatorname{Ob}\mathrm{D}^{+}\mathrm{F}(\mathrm{X},\mathbf{C})\),

\noindent
\(\gamma)\) un isomorphisme \(\alpha:\mathrm{K}_{\mathbf{C}}\simeq\mathrm{K}_{\mathrm{A}}\otimes\mathbf{C}\) dans \(\mathrm{D}^{+}(\mathrm{X},\mathbf{C})\).

L’axiome suivant doit être vérifié :

\noindent
(CHC) Le triplet \((\mathrm{R}\Gamma(\mathrm{K}_{\mathrm{A}}),\mathrm{R}\Gamma((\mathrm{K}_{\mathbf{C}},\mathrm{F})),\mathrm{R}\Gamma(\alpha))\) est un complexe de Hodge de poids \(n\).

% ---- 029-036 ----
\Page{29}
Lorsque \(\mathrm{A}=\mathbf{Z}\), on parlera simplement de complexe de Hodge et de complexe de Hodge cohomologique.

L’énoncé suivant reformule une partie de la théorie de Hodge classique (cf. (2.2.1)).

\textit{Scholie} \textbf{(8.1.3)}. — \textit{Soit \(\mathrm{X}\) une variété kählérienne compacte. Soient \(\mathrm{K}_{\mathbf{Z}}\) le complexe \(\mathbf{Z}[0]\) réduit au faisceau constant \(\mathbf{Z}\), \(\mathrm{K}_{\mathbf{C}}\) le complexe de De Rham holomorphe \(\Omega_{\mathrm{X}}^\bullet\) et \(\mathrm{F}\) sa filtration bête. On dispose de \(\alpha:\mathrm{K}_{\mathbf{Z}}\otimes\mathbf{C}\simeq\mathbf{C}[0]\xrightarrow{\approx}\Omega_{\mathrm{X}}^\bullet\) (lemme de Poincaré), et \((\mathrm{K}_{\mathbf{Z}},(\mathrm{K}_{\mathbf{C}},\mathrm{F}),\alpha)\) est un complexe de Hodge cohomologique de poids \(0\).}

\textit{Remarque} \textbf{(8.1.4)}. — Si \(\mathrm{K}\) est un complexe de Hodge (resp. de Hodge cohomologique) de poids \(n\), alors \(\mathrm{K}[m]\) est un complexe de Hodge (resp. de Hodge cohomologique) de poids \(n+m\).

\textbf{(8.1.5)} Un \(\mathrm{A}\)-\textit{complexe de Hodge mixte} \(\mathrm{K}\) consiste en :

\(\alpha)\) un complexe \(\mathrm{K}_{\mathrm{A}}\in\operatorname{Ob}\mathrm{D}^{+}(\mathrm{A})\) tel que \(\mathrm{H}^{k}(\mathrm{K}_{\mathrm{A}})\) soit un \(\mathrm{A}\)-module de type fini pour tout \(k\);

\(\beta)\) un complexe filtré \((\mathrm{K}_{\mathrm{A}\otimes\mathbf{Q}},\mathrm{W})\in\operatorname{Ob}\mathrm{D}^{+}\mathrm{F}(\mathrm{A}\otimes\mathbf{Q})\) (\(\mathrm{W}\) filtration \textit{croissante}) et un isomorphisme \(\mathrm{K}_{\mathrm{A}\otimes\mathbf{Q}}\simeq\mathrm{K}_{\mathrm{A}}\otimes\mathbf{Q}\) dans \(\mathrm{D}^{+}(\mathrm{A}\otimes\mathbf{Q})\);

\(\gamma)\) un complexe bifiltré \((\mathrm{K}_{\mathbf{C}},\mathrm{W},\mathrm{F})\in\operatorname{Ob}\mathrm{D}^{+}\mathrm{F}_{2}(\mathbf{C})\) (\(\mathrm{W}\) filtration \textit{croissante} et \(\mathrm{F}\) filtration \textit{décroissante}) et un isomorphisme \(\alpha:(\mathrm{K}_{\mathbf{C}},\mathrm{W})\simeq(\mathrm{K}_{\mathrm{A}\otimes\mathbf{Q}},\mathrm{W})\otimes\mathbf{C}\) dans \(\mathrm{D}^{+}\mathrm{F}(\mathbf{C})\).

L’axiome suivant doit être vérifié :

(CHM) Pour tout \(n\), le système formé du complexe \(\operatorname{Gr}_{n}^{\mathrm{W}}(\mathrm{K}_{\mathrm{A}\otimes\mathbf{Q}})\in\operatorname{Ob}\mathrm{D}^{+}(\mathrm{A}\otimes\mathbf{Q})\), du complexe \(\operatorname{Gr}_{n}^{\mathrm{W}}(\mathrm{K}_{\mathbf{C}})\in\operatorname{Ob}\mathrm{D}^{+}\mathrm{F}(\mathbf{C})\) filtré par \(\mathrm{F}\) et de l’isomorphisme
\[
\operatorname{Gr}_{n}^{\mathrm{W}}(\alpha):
\operatorname{Gr}_{n}^{\mathrm{W}}(\mathrm{K}_{\mathrm{A}\otimes\mathbf{Q}})\otimes\mathbf{C}
\simeq
\operatorname{Gr}_{n}^{\mathrm{W}}(\mathrm{K}_{\mathbf{C}}),
\]
est un \(\mathrm{A}\otimes\mathbf{Q}\)-complexe de Hodge de poids \(n\).

\textbf{(8.1.6)} Un \(\mathrm{A}\)-\textit{complexe de Hodge mixte cohomologique} \(\mathrm{K}\) sur un espace topologique \(\mathrm{X}\) consiste en :

\(\alpha)\) un complexe \(\mathrm{K}_{\mathrm{A}}\in\operatorname{Ob}\mathrm{D}^{+}(\mathrm{X},\mathrm{A})\), tel que les \(\mathbf{H}^{k}(\mathrm{X},\mathrm{K}_{\mathrm{A}})\) soient des \(\mathrm{A}\)-modules de type fini, et que \(\mathbf{H}^{\bullet}(\mathrm{X},\mathrm{K}_{\mathrm{A}})\otimes\mathbf{Q}\xrightarrow{\sim}\mathbf{H}^{\bullet}(\mathrm{X},\mathrm{K}_{\mathrm{A}}\otimes\mathbf{Q})\);

\(\beta)\) un complexe filtré \((\mathrm{K}_{\mathrm{A}\otimes\mathbf{Q}},\mathrm{W})\in\operatorname{Ob}\mathrm{D}^{+}\mathrm{F}(\mathrm{X},\mathrm{A}\otimes\mathbf{Q})\) (\(\mathrm{W}\) une filtration \textit{croissante}) et un isomorphisme \(\mathrm{K}_{\mathrm{A}\otimes\mathbf{Q}}\simeq\mathrm{K}_{\mathrm{A}}\otimes\mathbf{Q}\) dans \(\mathrm{D}^{+}(\mathrm{X},\mathrm{A}\otimes\mathbf{Q})\);

\(\gamma)\) un complexe bifiltré \((\mathrm{K}_{\mathbf{C}},\mathrm{W},\mathrm{F})\in\operatorname{Ob}\mathrm{D}^{+}\mathrm{F}_{2}(\mathrm{X},\mathbf{C})\) (\(\mathrm{W}\) une filtration \textit{croissante} et \(\mathrm{F}\) une filtration \textit{décroissante}) et un isomorphisme \(\alpha:(\mathrm{K}_{\mathbf{C}},\mathrm{W})\simeq(\mathrm{K}_{\mathrm{A}\otimes\mathbf{Q}},\mathrm{W})\otimes\mathbf{C}\) dans \(\mathrm{D}^{+}\mathrm{F}(\mathrm{X},\mathbf{C})\).

L’axiome suivant doit être vérifié :

(CHMC) Pour tout \(n\), le système formé du complexe \(\operatorname{Gr}_{n}^{\mathrm{W}}(\mathrm{K}_{\mathrm{A}\otimes\mathbf{Q}})\in\operatorname{Ob}\mathrm{D}^{+}(\mathrm{X},\mathrm{A}\otimes\mathbf{Q})\), du complexe \(\operatorname{Gr}_{n}^{\mathrm{W}}(\mathrm{K}_{\mathbf{C}})\in\operatorname{Ob}\mathrm{D}^{+}\mathrm{F}(\mathrm{X},\mathbf{C})\) filtré par \(\mathrm{F}\) et de l’isomorphisme
\[
\operatorname{Gr}_{n}^{\mathrm{W}}(\alpha):
\operatorname{Gr}_{n}^{\mathrm{W}}(\mathrm{K}_{\mathbf{C}})
\simeq
\operatorname{Gr}_{n}^{\mathrm{W}}(\mathrm{K}_{\mathrm{A}\otimes\mathbf{Q}})\otimes\mathbf{C}
\]
est un \(\mathrm{A}\otimes\mathbf{Q}\)-complexe de Hodge cohomologique de poids \(n\).

\Page{30}
On vérifie trivialement :

\textit{Proposition} \textbf{(8.1.7)}. — \textit{Si \(\mathrm{K}=(\mathrm{K}_{\mathrm{A}},(\mathrm{K}_{\mathrm{A}\otimes\mathbf{Q}},\mathrm{W}),(\mathrm{K}_{\mathbf{C}},\mathrm{W},\mathrm{F}))\) est un \(\mathrm{A}\)-complexe de Hodge mixte cohomologique, alors \(\mathrm{R}\Gamma\mathrm{K}=(\mathrm{R}\Gamma\mathrm{K}_{\mathrm{A}},\mathrm{R}\Gamma(\mathrm{K}_{\mathrm{A}\otimes\mathbf{Q}},\mathrm{W}),\mathrm{R}\Gamma(\mathrm{K}_{\mathbf{C}},\mathrm{W},\mathrm{F}))\) est un \(\mathrm{A}\)-complexe de Hodge mixte.}

\textit{Scholie} \textbf{(8.1.8)}. — \textit{Soient \(\mathrm{U}\) le complément, dans un schéma propre et lisse \(\mathrm{X}\), d’un diviseur à croisements normaux \(\mathrm{D}\), et \(j:\mathrm{U}\hookrightarrow\mathrm{X}\) le morphisme d’inclusion. Soit \(\mathrm{W}\) la filtration canonique de \(\mathrm{R}j_{\ast}\mathbf{Q}=\mathrm{R}j_{\ast}\mathbf{Z}\otimes\mathbf{Q}\) : \(\mathrm{W}_{n}(\mathrm{R}j_{\ast}\mathbf{Q})=\tau_{\leq n}(\mathrm{R}j_{\ast}\mathbf{Q})\). On a (3.1.8)}
\[
(\mathrm{R}j_{\ast}\mathbf{Q})\otimes\mathbf{C}
\xrightarrow{\approx}\mathrm{R}j_{\ast}\mathbf{C}
\xrightarrow{\approx}j_{\ast}\Omega_{\mathrm{U}}^\bullet
\xleftarrow{\approx}\Omega_{\mathrm{X}}^\bullet(\log\mathrm{D}).
\]
\textit{Soient \(\mathrm{W}\) la filtration par le poids de \(\Omega_{\mathrm{X}}^\bullet(\log\mathrm{D})\) et \(\mathrm{F}\) sa filtration de Hodge (par les \(\sigma_{\geq p}\)). Alors, (3.1.8) fournit dans \(\mathrm{D}^{+}\mathrm{F}(\mathrm{X},\mathbf{C})\) un isomorphisme \((\mathrm{R}j_{\ast}\mathbf{Q},\mathrm{W})\otimes\mathbf{C}\simeq(\Omega_{\mathrm{X}}^\bullet(\log\mathrm{D}),\mathrm{W})\) et, d’après (3.1.5) et (3.1.9)}
\[
(\mathrm{R}j_{\ast}\mathbf{Z},(\mathrm{R}j_{\ast}\mathbf{Q},\mathrm{W}),
(\Omega_{\mathrm{X}}^\bullet(\log\mathrm{D}),\mathrm{W},\mathrm{F}))
\]
\textit{est un complexe de Hodge mixte cohomologique sur \(\mathrm{X}\).}

Le résultat suivant est une version « abstraite » du numéro (3.2).

\textit{Scholie} \textbf{(8.1.9)}. — \textit{Soit \(\mathrm{K}\) un \(\mathrm{A}\)-complexe de Hodge mixte.}

(i) \textit{Sur les termes \({}_{\mathrm{W}}\mathrm{E}_{r}^{pq}\) de la suite spectrale de \((\mathrm{K}_{\mathbf{C}},\mathrm{W})\), la filtration récurrente et les deux filtrations directes, définies par \(\mathrm{F}\) coïncident.}

(ii) \textit{La filtration \(\mathrm{W}[n]\) (1.1.2) de \(\mathrm{H}^{n}(\mathrm{K}_{\mathrm{A}\otimes\mathbf{Q}})\simeq\mathrm{H}^{n}(\mathrm{K}_{\mathrm{A}})\otimes\mathbf{Q}\) et la filtration \(\mathrm{F}\) de \(\mathrm{H}^{n}(\mathrm{K}_{\mathbf{C}})\simeq\mathrm{H}^{n}(\mathrm{K}_{\mathrm{A}})\otimes_{\mathrm{A}}\mathbf{C}\) définissent sur \(\mathrm{H}^{n}(\mathrm{K}_{\mathrm{A}})\) une \(\mathrm{A}\)-structure de Hodge mixte ((2.3.1) amplifié par (III.0.4)).}

(iii) \textit{Les morphismes \({}_{\mathrm{W}}d_{1}:{}_{\mathrm{W}}\mathrm{E}_{1}^{pq}\to{}_{\mathrm{W}}\mathrm{E}_{1}^{p+1,q}\) sont compatibles à la bigraduation de Hodge; en particulier, ils sont strictement compatibles à la filtration de Hodge.}

(iv) \textit{La suite spectrale de \((\mathrm{K}_{\mathrm{A}\otimes\mathbf{Q}},\mathrm{W})\) dégénère en \(\mathrm{E}_{2}\) : \({}_{\mathrm{W}}\mathrm{E}_{2}={}_{\mathrm{W}}\mathrm{E}_{\infty}\).}

(v) \textit{La suite spectrale de \((\mathrm{K}_{\mathbf{C}},\mathrm{F})\) dégénère en \(\mathrm{E}_{1}\) : \({}_{\mathrm{F}}\mathrm{E}_{1}={}_{\mathrm{F}}\mathrm{E}_{\infty}\).}

(vi) \textit{La suite spectrale du complexe \(\operatorname{Gr}_{\mathrm{F}}^{p}(\mathrm{K}_{\mathbf{C}})\), muni de la filtration induite par \(\mathrm{W}\), dégénère en \(\mathrm{E}_{2}\).}

La démonstration de (i), (ii), (iii), (iv) est parallèle à celle de (3.2.5). L’analogue des lemmes (3.2.6) et (3.2.7) a été pris pour axiome (CH 1 et CH 2). Les démonstrations de (3.2.8), (3.2.9) et (3.2.10) se transposent alors telles quelles, et comme en (3.2.10), on en déduit (i), (ii), (iii), (iv). Les assertions (v), (vi) résultent alors de (7.2.8).

\textbf{(8.1.10)} Nous abrégerons « différentiel gradué » en \(\mathrm{DG}\) et « différentiel gradué à degrés bornés inférieurement » en \(\mathrm{DG}^{+}\). Un complexe de \(\mathrm{A}\)-modules \(\mathrm{DG}\) peut être vu comme un double complexe de \(\mathrm{A}\)-modules; le \textit{premier} degré est celui du complexe, le \textit{second} est celui défini par la graduation des modules \(\mathrm{DG}\). On désigne par \(\mathrm{D}^{+}((\mathrm{A}\text{-mod }\mathrm{DG})^{+})\)
\Page{31}
la catégorie dérivée de la catégorie des complexes bornés inférieurement de \(\mathrm{A}\)-modules \(\mathrm{DG}\), à degrés uniformément bornés inférieurement.

Un \(\mathrm{A}\)-\textit{complexe de Hodge} \(\mathrm{K}\) \textit{mixte} \(\mathrm{DG}\) consiste en :

\textit{a)} un complexe \(\mathrm{K}_{\mathrm{A}}\in\operatorname{Ob}\mathrm{D}^{+}((\mathrm{A}\text{-mod }\mathrm{DG})^{+})\);

\textit{b)} un complexe filtré \((\mathrm{K}_{\mathrm{A}\otimes\mathbf{Q}},\mathrm{W})\in\operatorname{Ob}\mathrm{D}^{+}\mathrm{F}((\mathrm{A}\otimes\mathbf{Q}\text{-mod }\mathrm{DG})^{+})\), et un isomorphisme \(\mathrm{K}_{\mathrm{A}}\otimes\mathbf{Q}\simeq\mathrm{K}_{\mathrm{A}\otimes\mathbf{Q}}\) dans \(\mathrm{D}^{+}\);

\textit{c)} un complexe bifiltré \((\mathrm{K}_{\mathbf{C}},\mathrm{F},\mathrm{W})\in\operatorname{Ob}\mathrm{D}^{+}\mathrm{F}_{2}((\mathbf{C}\text{-mod }\mathrm{DG})^{+})\), et un isomorphisme \(\mathrm{K}_{\mathrm{A}\otimes\mathbf{Q}}\otimes\mathbf{C}\simeq\mathrm{K}_{\mathbf{C}}\) dans \(\mathrm{D}^{+}\mathrm{F}\).

On exige que, pour chaque \(n\), la composante \((\mathrm{K}^{\bullet n},\mathrm{F},\mathrm{W})\) de second degré \(n\) de \(\mathrm{K}\) soit un \(\mathrm{A}\)-complexe de Hodge mixte.

On désigne par \(\mathrm{L}\) la filtration par le second degré de \(\mathbf{s}\mathrm{K}\). On désigne par \(\operatorname{Dec}_{1}\mathrm{W}\) la filtration de \(\mathrm{K}_{\mathrm{A}\otimes\mathbf{Q}}\) déduite par décalage (1.3.3) de \(\mathrm{W}\). On ne confondra pas la filtration qu’elle induit sur \(\mathbf{s}\mathrm{K}_{\mathrm{A}\otimes\mathbf{Q}}\) (encore notée \(\operatorname{Dec}_{1}\mathrm{W}\)) avec la filtration décalée de la filtration induite sur \(\mathbf{s}\mathrm{K}_{\mathrm{A}\otimes\mathbf{Q}}\) par \(\mathrm{W}\).

\textit{Variante} \textbf{(8.1.10.1)}. — On définit de même les \(\mathrm{A}\)-\textit{complexes de Hodge mixtes cosimpliciaux} (resp. \textit{\(r\)-cosimpliciaux}) en remplaçant partout \(\mathrm{DG}\) par cosimplicial (resp. \(r\)-cosimplicial). Le foncteur (\(\mathrm{A}\)-module cosimplicial) \(\to\) (\(\mathrm{A}\)-module \(\mathrm{DG}\)) transforme \(\mathrm{A}\)-complexes de Hodge mixtes cosimpliciaux en \(\mathrm{A}\)-complexes de Hodge mixtes \(\mathrm{DG}\). De même pour \(r\)-cosimplicial.

\textbf{(8.1.11)} Soit \(\mathrm{X}_{\bullet}\) un espace topologique simplicial. Un \(\mathrm{A}\)-\textit{complexe de Hodge mixte cohomologique} \(\mathrm{K}\) sur \(\mathrm{X}_{\bullet}\) consiste en :

\(\alpha)\) un complexe \(\mathrm{K}_{\mathrm{A}}\in\operatorname{Ob}\mathrm{D}^{+}(\mathrm{X}_{\bullet},\mathrm{A})\);

\(\beta)\) un complexe filtré \((\mathrm{K}_{\mathrm{A}\otimes\mathbf{Q}},\mathrm{W})\in\operatorname{Ob}\mathrm{D}^{+}\mathrm{F}(\mathrm{X}_{\bullet},\mathrm{A}\otimes\mathbf{Q})\), et un isomorphisme \(\mathrm{K}_{\mathrm{A}\otimes\mathbf{Q}}\simeq\mathrm{K}_{\mathrm{A}}\otimes\mathbf{Q}\) dans \(\mathrm{D}^{+}(\mathrm{X}_{\bullet},\mathrm{A}\otimes\mathbf{Q})\);

\(\gamma)\) un complexe bifiltré \((\mathrm{K}_{\mathbf{C}},\mathrm{W},\mathrm{F})\in\operatorname{Ob}\mathrm{D}^{+}\mathrm{F}_{2}(\mathrm{X}_{\bullet},\mathbf{C})\), et un isomorphisme
\[
(\mathrm{K}_{\mathbf{C}},\mathrm{W})
\simeq
(\mathrm{K}_{\mathrm{A}\otimes\mathbf{Q}},\mathrm{W})\otimes\mathbf{C}
\]
dans \(\mathrm{D}^{+}\mathrm{F}(\mathrm{X}_{\bullet},\mathbf{C})\).

L’axiome suivant doit être vérifié :

\((\mathrm{CHM}_{\bullet})\) La restriction de \(\mathrm{K}\) à chacun des \(\mathrm{X}_{n}\) est un \(\mathrm{A}\)-complexe de Hodge mixte cohomologique.

Pour \(\mathrm{A}=\mathbf{Z}\), on parlera simplement de complexe de Hodge mixte cohomologique.

\textit{Exemple} \textbf{(8.1.12)}. — Soit \(\mathrm{X}_{\bullet}\) un schéma simplicial propre et lisse (sur \(\mathbf{C}\)) et soit \(\mathrm{Y}_{\bullet}\) un diviseur à croisements normaux dans \(\mathrm{X}_{\bullet}\). On pose \(\mathrm{U}_{\bullet}=\mathrm{X}_{\bullet}-\mathrm{Y}_{\bullet}\) et on désigne par \(j\) l’inclusion de \(\mathrm{U}_{\bullet}\) dans \(\mathrm{X}_{\bullet}\). Les constructions (8.1.8) peuvent être effectuées « unifor-%
\Page{32}mément en \(\mathrm{X}_{n}\) » comme (5.2.5) et (6.2.7) le montrent. Elles fournissent donc un complexe de Hodge mixte cohomologique
\[
(\mathrm{R}j_{\ast}\mathbf{Z},(\mathrm{R}j_{\ast}\mathbf{Q},\tau_{\leq}),
(\Omega_{\mathrm{X}_{\bullet}}^\bullet(\log\mathrm{Y}_{\bullet}),\mathrm{W},\mathrm{F}))
\]
sur \(\mathrm{X}_{\bullet}\).

\textbf{(8.1.13)} Soit \(\mathrm{K}\) un \(\mathrm{A}\)-complexe de Hodge mixte cohomologique sur \(\mathrm{X}_{\bullet}\). Appliquons à \(\mathrm{K}\) le foncteur \(\mathrm{R}\Gamma^\bullet\). On obtient

\textit{a)} un complexe cosimplicial \(\mathrm{R}\Gamma^\bullet\mathrm{K}_{\mathrm{A}}\in\operatorname{Ob}\mathrm{D}^{+}((\mathrm{A}\text{-modules cosimpliciaux}))\);

\textit{b)} un complexe filtré cosimplicial \(\mathrm{R}\Gamma^\bullet(\mathrm{K}_{\mathrm{A}\otimes\mathbf{Q}},\mathrm{W})\in\operatorname{Ob}\mathrm{D}^{+}\mathrm{F}((\mathrm{A}\otimes\mathbf{Q}\text{-modules cosimpliciaux}))\);

\textit{c)} un complexe bifiltré cosimplicial
\[
\mathrm{R}\Gamma^\bullet(\mathrm{K}_{\mathbf{C}},\mathrm{W},\mathrm{F})
\in\operatorname{Ob}\mathrm{D}^{+}\mathrm{F}_{2}((\mathbf{C}\text{-vectoriels cosimpliciaux}));
\]

\textit{d)} des isomorphismes (8.1.10.\textit{b)}, \textit{c)}) entre ces objets.

Il est clair que \(\mathrm{R}\Gamma^\bullet(\mathrm{K},\mathrm{W},\mathrm{F})\) décrit ci-dessus est un \(\mathrm{A}\)-complexe de Hodge mixte cosimplicial. Il définit un \(\mathrm{A}\)-complexe de Hodge mixte \(\mathrm{DG}\) (8.1.10.1).

\textbf{(8.1.14)} Soit \(\mathrm{K}\) un \(\mathrm{A}\)-complexe de Hodge mixte \(\mathrm{DG}\) et considérons la suite spectrale d’aboutissement la cohomologie de \(\mathbf{s}\mathrm{K}\) qui est définie par la filtration \(\mathrm{L}\). Par hypothèse, les termes \(\mathrm{E}_{1}\) de cette suite spectrale sont munis de \(\mathrm{A}\)-structures de Hodge mixtes. Le contenu du théorème ci-dessous est que toute la suite spectrale, y compris son aboutissement, est munie d’une \(\mathrm{A}\)-structure de Hodge mixte. En d’autres termes, c’est une suite spectrale dans la catégorie abélienne des \(\mathrm{A}\)-structures de Hodge mixtes. La structure sur l’aboutissement correspond à une « structure » de \(\mathrm{A}\)-complexe de Hodge mixte naturelle sur \(\mathbf{s}\mathrm{K}\).

\textit{Théorème} \textbf{(8.1.15)}. — \textit{Soit \(\mathrm{K}\) un \(\mathrm{A}\)-complexe de Hodge mixte \(\mathrm{DG}\).}

(i) \textit{Avec les notations de (8.1.10), (8.1.10.1) et (7.1.6), \((\mathbf{s}\mathrm{K},\delta(\mathrm{W},\mathrm{L}),\mathrm{F})\) est un \(\mathrm{A}\)-complexe de Hodge mixte.}

(ii) \textit{On a}
\[
{}_{\delta(\mathrm{W},\mathrm{L})}\mathrm{E}_{1}^{ab}
(\mathbf{s}\mathrm{K}\otimes\mathbf{Q})
=
\sum_{\substack{a=\gamma-\beta\\b=\alpha+\beta}}
\mathrm{H}^{\alpha}
(\operatorname{Gr}_{\beta}^{\mathrm{W}}\mathrm{K}^{\bullet,\gamma});
\tag*{\textbf{(8.1.15.1)}}
\]
\textit{le complexe \(({}_{\delta(\mathrm{W},\mathrm{L})}\mathrm{E}_{1}^{\bullet b},d_{1})\) est le complexe simple associé au complexe double de \(\mathrm{A}\otimes\mathbf{Q}\)-structures de Hodge de poids \(b\) suivant}
\[
\begin{tikzcd}[ampersand replacement=\&, column sep=1em, row sep=3em]
\mathrm{H}^{b-(\beta+1)}
(\operatorname{Gr}_{\beta+1}^{\mathrm{W}}\mathrm{K}^{\bullet,\gamma+1})
\arrow[r, "\partial"]
\&
\mathrm{H}^{b-\beta}
(\operatorname{Gr}_{\beta}^{\mathrm{W}}\mathrm{K}^{\bullet,\gamma+1})
\arrow[r, "\partial"]
\&
\mathrm{H}^{b-(\beta-1)}
(\operatorname{Gr}_{\beta-1}^{\mathrm{W}}\mathrm{K}^{\bullet,\gamma+1})
\\
\&
\mathrm{H}^{b-\beta}
(\operatorname{Gr}_{\beta}^{\mathrm{W}}\mathrm{K}^{\bullet,\gamma})
\arrow[u, "d''"]
\arrow[r, "\partial"]
\&
\mathrm{H}^{b-(\beta-1)}
(\operatorname{Gr}_{\beta-1}^{\mathrm{W}}\mathrm{K}^{\bullet,\gamma})
\arrow[u, "d''"]
\end{tikzcd}
\tag*{\textbf{(8.1.15.2)}}
\]

\Page{33}
(iii) \textit{Sur les termes \({}_{\mathrm{L}}\mathrm{E}_{r}\) de la suite spectrale définie par \((\mathbf{s}\mathrm{K}_{\mathrm{A}\otimes\mathbf{Q}},\mathrm{L})\), la filtration récurrente et les deux filtrations directes définies par \(\operatorname{Dec}_{1}(\mathrm{W})\) coïncident; de même, sur les \({}_{\mathrm{L}}\mathrm{E}_{r}(\mathbf{s}\mathrm{K}_{\mathbf{C}})\), la filtration récurrente et les deux filtrations directes induites par \(\mathrm{F}\) coïncident. On notera ces filtrations \(\mathrm{W}\) et \(\mathrm{F}\).}

(iv) \textit{Pour \(r\geq1\), \(({}_{\mathrm{L}}\mathrm{E}_{r},\mathrm{W},\mathrm{F})\) est une \(\mathrm{A}\)-structure de Hodge mixte (III.0.4). Les différentielles \(d_{r}\) sont des morphismes de \(\mathrm{A}\)-structures de Hodge mixtes.}

(v) \textit{La filtration \(\mathrm{L}\) de \(\mathrm{H}^{\bullet}(\mathbf{s}\mathrm{K})\) est une filtration dans la catégorie des structures de Hodge mixtes, et}
\[
\operatorname{Gr}_{\mathrm{L}}^{p}
(\mathrm{H}^{p+q}(\mathbf{s}\mathrm{K}),
\delta(\mathrm{W},\mathrm{L})[p+q],\mathrm{F})
=
({}_{\mathrm{L}}\mathrm{E}_{\infty}^{pq},\mathrm{W},\mathrm{F}).
\tag*{\textbf{(8.1.15.3)}}
\]

\textit{Remarque} \textbf{(8.1.16)}. — Dans (8.1.15), la filtration \(\operatorname{Dec}_{1}(\mathrm{W})\) (8.1.10) joue un plus grand rôle que \(\mathrm{W}\). On a
\[
\operatorname{Dec}(\delta(\mathrm{W},\mathrm{L}))
=
\operatorname{Dec}_{1}(\mathrm{W})
\quad\text{sur }\mathbf{s}\mathrm{K},
\tag*{\textbf{(8.1.16.1)}}
\]
et (8.1.15.3) se réécrit
\[
\operatorname{Gr}_{\mathrm{L}}
(\mathrm{H}(\mathbf{s}\mathrm{K}),\operatorname{Dec}_{1}\mathrm{W},\mathrm{F})
=
({}_{\mathrm{L}}\mathrm{E}_{\infty},\mathrm{W},\mathrm{F}).
\tag*{\textbf{(8.1.16.2)}}
\]
Dans le même esprit, notons que, avec les notations de (8.1.8)
\[
\mathrm{R}\Gamma
(\mathrm{R}j_{\ast}\mathbf{Z},
(\mathrm{R}j_{\ast}\mathbf{Q},\operatorname{Dec}\mathrm{W}),
(\Omega_{\mathrm{X}}^\bullet(\log\mathrm{D}),
\operatorname{Dec}\mathrm{W},\mathrm{F}))
\]
ne dépend que de \(\mathrm{U}\) et non de la compactification choisie \(\mathrm{X}\).

\textbf{(8.1.17)} Prouvons (8.1.15). Les assertions (i) et (ii) résultent de ce que
\[
\operatorname{Gr}_{n}^{\delta(\mathrm{W},\mathrm{L})}
(\mathbf{s}\mathrm{K}\otimes\mathbf{Q})
=
\sum_{p-q=n}
\operatorname{Gr}_{p}^{\mathrm{W}}\operatorname{Gr}_{\mathrm{L}}^{q}
(\mathbf{s}\mathrm{K}\otimes\mathbf{Q})
=
\sum_{p-q=n}
(\operatorname{Gr}_{p}^{\mathrm{W}}\mathrm{K}^{\bullet q})[-q]
\]
(cet isomorphisme est compatible à \(\mathrm{F}\), et on applique (8.1.4)). La vérification de (8.1.15.2) est laissée au lecteur.

Soient les assertions :

\((a_{r})\) Les morphismes \(d_{r}\) de \({}_{\mathrm{L}}\mathrm{E}_{r}(\mathbf{s}\mathrm{K}_{\mathrm{A}\otimes\mathbf{Q}})\) sont strictement compatibles à la filtration récurrente définie par \(\operatorname{Dec}_{1}(\mathrm{W})\).

\((b_{r})\) Les morphismes \(d_{r}\) de \({}_{\mathrm{L}}\mathrm{E}_{r}(\mathbf{s}\mathrm{K}_{\mathbf{C}})\) sont strictement compatibles à la filtration récurrente définie par \(\mathrm{F}\).

\((c_{r})\) \({}_{\mathrm{L}}\mathrm{E}_{r}\), muni des filtrations considérées en \((a_{r})\), \((b_{r})\), est une \(\mathrm{A}\)-structure de Hodge mixte.

\((d_{r})\) Les \(d_{r}\) sont des morphismes de \(\mathrm{A}\)-structures de Hodge mixtes.

Nous prouverons \((a_{r})\) (\(r\geq0\)), \((b_{r})\) (\(r\geq0\)), \((c_{r})\) (\(r\geq1\)), \((d_{r})\) (\(r\geq1\)) par une récurrence simultanée.

A) \textit{Preuve de} \((a_{0})\) : on a
\[
({}_{\mathrm{L}}\mathrm{E}_{0}^{p\bullet}
(\mathbf{s}\mathrm{K}_{\mathrm{A}\otimes\mathbf{Q}}),
\operatorname{Dec}_{1}\mathrm{W})
=
(\mathrm{K}_{\mathrm{A}\otimes\mathbf{Q}}^{\bullet p},
\operatorname{Dec}\mathrm{W}).
\]
D’après (8.1.9) (iv), la suite spectrale de \((\mathrm{K}_{\mathrm{A}\otimes\mathbf{Q}}^{\bullet p},\mathrm{W})\) dégénère en \(\mathrm{E}_{2}\). D’après (1.3.4), la suite spectrale de \((\mathrm{K}_{\mathrm{A}\otimes\mathbf{Q}}^{\bullet p},\operatorname{Dec}\mathrm{W})\) dégénère en \(\mathrm{E}_{1}\), et ceci prouve \((a_{0})\) (1.3.2).

\Page{34}
B) \textit{Preuve de} \((b_{0})\) : on a
\[
({}_{\mathrm{L}}\mathrm{E}_{0}^{p\bullet}
(\mathbf{s}\mathrm{K}_{\mathbf{C}}),\mathrm{F})
=
(\mathrm{K}_{\mathbf{C}}^{\bullet p},\mathrm{F}).
\]
On applique (8.1.9) (v) et (1.3.2).

C) \textit{Preuve de} \((c_{1})\) : il faut montrer que \(\mathrm{H}^{n}(\mathrm{K}_{\mathrm{A}}^{\bullet p})\), muni des filtrations induites par \(\operatorname{Dec}\mathrm{W}\) sur \(\mathrm{H}^{n}(\mathrm{K}_{\mathrm{A}\otimes\mathbf{Q}}^{\bullet p})=\mathrm{H}^{n}(\mathrm{K}_{\mathrm{A}}^{\bullet p})\otimes\mathbf{Q}\) et \(\mathrm{F}\) sur \(\mathrm{H}^{n}(\mathrm{K}_{\mathbf{C}}^{\bullet p})=\mathrm{H}^{n}(\mathrm{K}_{\mathrm{A}}^{\bullet p})\otimes_{\mathrm{A}}\mathbf{C}\), est une \(\mathrm{A}\)-structure de Hodge mixte. \(\operatorname{Dec}\mathrm{W}\) induit sur \(\mathrm{H}^{n}(\mathrm{K}_{\mathrm{A}\otimes\mathbf{Q}}^{\bullet p})\) la même filtration que \(\mathrm{W}[n]\), et on conclut par (8.1.9) (ii).

D) \textit{Preuve de} \((c_{r})+(d_{r})\Rightarrow(a_{r})+(b_{r})+(c_{r+1})\) : c’est le point clef de la démonstration, et une simple application de (2.3.5) (amplifié par (III.0.4)).

E) \textit{Preuve de} \((a_{r})_{r\leq r_{0}-2}+(b_{r})_{r\leq r_{0}-2}+(c_{r_{0}})\Rightarrow(d_{r_{0}})\) : D’après (1.3.16), la filtration récurrente et les filtrations directes induites par \(\operatorname{Dec}_{1}\mathrm{W}\) (resp. \(\mathrm{F}\)) sur \({}_{\mathrm{L}}\mathrm{E}_{r}\) coïncident pour \(r\leq r_{0}\), et on applique (1.3.13) (i).

Ceci termine la démonstration par récurrence de \((a)\), \((b)\), \((c)\), \((d)\).

L’assertion (iii) a été prouvée dans la partie E de la récurrence et les assertions de (iv) sont \((c_{r})\) et \((d_{r})\). On prouvera (8.1.15.3) sous la forme (8.1.16.2).

Que
\[
\operatorname{Gr}_{\mathrm{L}}
(\mathrm{H}(\mathbf{s}\mathrm{K}_{\mathrm{A}\otimes\mathbf{Q}}),
\operatorname{Dec}_{1}\mathrm{W})
=
({}_{\mathrm{L}}\mathrm{E}_{\infty}
(\mathbf{s}\mathrm{K}_{\mathrm{A}\otimes\mathbf{Q}}),\mathrm{W})
\]
résulte de \((a_{r})\) et (1.3.17). Que
\[
\operatorname{Gr}_{\mathrm{L}}
(\mathrm{H}(\mathbf{s}\mathrm{K}_{\mathbf{C}}),\mathrm{F})
=
({}_{\mathrm{L}}\mathrm{E}_{\infty}
(\mathbf{s}\mathrm{K}_{\mathbf{C}}),\mathrm{F})
\]
résulte de \((b_{r})\) et (1.3.17). On conclut par le lemme suivant, dont la vérification est laissée au lecteur :

\textit{Lemme} \textbf{(8.1.18)}. — \textit{Soit \(\mathrm{H}=(\mathrm{H}_{\mathrm{A}},\mathrm{W},\mathrm{F})\) une \(\mathrm{A}\)-structure de Hodge mixte. Pour qu’une filtration \(\mathrm{L}\) de \(\mathrm{H}_{\mathrm{A}}\) provienne d’une filtration de \(\mathrm{H}\), dans la catégorie abélienne des \(\mathrm{A}\)-structures de Hodge mixtes, il faut et il suffit que pour tout \(n\), \((\operatorname{Gr}_{\mathrm{L}}^{n}(\mathrm{H}_{\mathrm{A}}),\operatorname{Gr}_{\mathrm{L}}^{n}(\mathrm{W}),\operatorname{Gr}_{\mathrm{L}}^{n}(\mathrm{F}))\) soit une \(\mathrm{A}\)-structure de Hodge mixte.}

\textbf{(8.1.19)} Appliquons (8.1.13) et (8.1.15) à l’exemple (8.1.12).

On obtient un complexe de Hodge mixte consistant en \textit{a)}, \textit{b)}, \textit{c)} ci-dessous.

\textit{a)} Le complexe \(\mathrm{R}\Gamma\mathrm{R}j_{\ast}\mathbf{Z}\simeq\mathrm{R}\Gamma(\mathrm{U}_{\bullet},\mathbf{Z})\), ayant pour groupes de cohomologie les \(\mathrm{H}^{n}(\mathrm{U}_{\bullet},\mathbf{Z})\).

\textit{b)} La filtration \(\delta(\mathrm{W},\mathrm{L})\) de \(\mathrm{R}\Gamma\mathrm{R}j_{\ast}\mathbf{Q}\simeq\mathrm{R}\Gamma(\mathrm{U}_{\bullet},\mathbf{Q})\). Avec les notations de (3.1.4), on a
\begin{align*}
\operatorname{Gr}_{n}^{\delta(\mathrm{W},\mathrm{L})}
\mathrm{R}\Gamma(\mathrm{U}_{\bullet},\mathbf{Q})
&\underset{\text{(7.1.6.5)}}{\simeq}
\bigoplus_m
\operatorname{Gr}_{n+m}^{\mathrm{W}}
\mathrm{R}\Gamma(\mathrm{U}_{m},\mathbf{Q})[-m]
\\
&\simeq
\bigoplus_m
\mathrm{R}\Gamma
(\widetilde{\mathrm{Y}}_{m}^{n+m},
\varepsilon_{\mathbf{Q}}^{n+m}[-n-m])[-m]
\\
&\simeq
\bigoplus_m
\mathrm{R}\Gamma
(\widetilde{\mathrm{Y}}_{m}^{n+m},
\varepsilon_{\mathbf{Q}}^{n+m})[-n-2m].
\end{align*}

\Page{35}
Le terme \(\mathrm{E}_{1}\) de la suite spectrale correspondante est somme de groupes du type \(\mathrm{H}^{p}(\widetilde{\mathrm{Y}}_{q}^{r},\varepsilon_{\mathbf{Q}}^{r})\); ce dernier contribue au \(\mathrm{H}^{p+q+r}(\mathrm{U}_{\bullet})\); il correspond au \(\operatorname{Gr}_{a}^{\delta(\mathrm{W},\mathrm{L})}\) pour \(p+2r=(p+q+r)+a\) :
\[
{}_{\delta(\mathrm{W},\mathrm{L})}\mathrm{E}_{1}^{-a,b}
=
\bigoplus_{\substack{p+2r=b\\q-r=-a}}
\mathrm{H}^{p}(\widetilde{\mathrm{Y}}_{q}^{r},\varepsilon_{\mathbf{Q}}^{r})
\Rightarrow
\mathrm{H}^{-a+b}(\mathrm{U}_{\bullet},\mathbf{Q}).
\tag*{\textbf{(8.1.19.1)}}
\]

\textit{c)} La bifiltration \((\delta(\mathrm{W},\mathrm{L}),\mathrm{F})\) de \(\mathrm{R}\Gamma\mathrm{R}j_{\ast}\mathbf{C}\simeq\mathbf{s}\mathrm{R}\Gamma^\bullet\Omega_{\mathrm{X}_{\bullet}}^\bullet(\log\mathrm{Y}_{\bullet})\).

La filtration \(\mathrm{F}\) munit \({}_{\delta(\mathrm{W},\mathrm{L})}\mathrm{E}_{1}^{-a,b}\) d’une structure de Hodge de poids \(b\), et les différentielles \(d_{1}\) de \({}_{\delta(\mathrm{W},\mathrm{L})}\mathrm{E}\) sont des morphismes de \(\mathbf{Q}\)-structures de Hodge. Le complexe \({}_{\delta(\mathrm{W},\mathrm{L})}\mathrm{E}_{1}^{\bullet b}\) est le complexe simple associé au complexe double suivant, de composantes des structures de Hodge de poids \(b\) :

{\small
\[
\begin{tikzcd}[ampersand replacement=\&, column sep=1em, row sep=4em]
\mathrm{H}^{b-2(r+1)}
(\widetilde{\mathrm{Y}}_{q+1}^{r+1},\varepsilon_{\mathbf{Q}}^{r+1})(-r-1)
\arrow[r, "\mathrm{Gysin}"']
\&
\mathrm{H}^{b-2r}
(\widetilde{\mathrm{Y}}_{q+1}^{r},\varepsilon_{\mathbf{Q}}^{r})(-r)
\arrow[r, "\mathrm{Gysin}"']
\&
\mathrm{H}^{b-2(r-1)}
(\widetilde{\mathrm{Y}}_{q+1}^{r-1},\varepsilon_{\mathbf{Q}}^{r-1})(-r+1)
\\
\&
\mathrm{H}^{b-2r}
(\widetilde{\mathrm{Y}}_{q}^{r},\varepsilon_{\mathbf{Q}}^{r})(-r)
\arrow[u, "{\sum\limits_i(-1)^i\delta_i}"']
\arrow[r, "\mathrm{Gysin}"']
\&
\mathrm{H}^{b-2(r-1)}
(\widetilde{\mathrm{Y}}_{q}^{r-1},\varepsilon_{\mathbf{Q}}^{r-1})(-r+1)
\arrow[u, "{\sum\limits_i(-1)^i\delta_i}"']
\end{tikzcd}
\tag*{\textbf{(8.1.19.2)}}
\]
}

Les indices des composantes dans \(\mathrm{E}_{1}^{ab}\) vérifient \(a=q-r\).

En particulier, les \(\mathrm{H}^{n}(\mathrm{U}_{\bullet},\mathbf{Z})\) se trouvent munis de structures de Hodge mixtes (8.1.9).

\textit{Proposition} \textbf{(8.1.20)}. — \textit{Sous les hypothèses précédentes :}

(i) \textit{La structure de Hodge mixte de \(\mathrm{H}^{n}(\mathrm{U}_{\bullet},\mathbf{Z})\) est fonctorielle en le couple \((\mathrm{U}_{\bullet},\mathrm{X}_{\bullet})\).}

(ii) \textit{En cohomologie rationnelle, la suite spectrale (8.1.19.1) dégénère en \(\mathrm{E}_{2}\) \((\mathrm{E}_{2}=\mathrm{E}_{\infty})\) et aboutit à la filtration par le poids de \(\mathrm{H}^{n}(\mathrm{U}_{\bullet},\mathbf{Q})\). La structure de Hodge de \(\mathrm{E}_{2}\) induite par celle de \(\mathrm{E}_{1}\) est aussi celle de \(\operatorname{Gr}_{\mathrm{W}}\mathrm{H}^{n}(\mathrm{U}_{\bullet},\mathbf{Q})\).}

(iii) \textit{Les nombres de Hodge de \(\mathrm{H}^{n}(\mathrm{U}_{\bullet},\mathbf{Z})\) vérifient}
\[
h^{pq}\neq0
\quad\Rightarrow\quad
0\leq p\leq n\ \textit{et}\ 0\leq q\leq n.
\]

(iv) \textit{Pour \(\mathrm{Y}_{\bullet}=\varnothing\), les nombres de Hodge de \(\mathrm{H}^{n}(\mathrm{X}_{\bullet},\mathbf{Z})\) vérifient}
\[
h^{pq}\neq0
\quad\Rightarrow\quad
0\leq p\leq n,\quad 0\leq q\leq n
\quad\textit{et}\quad p+q\leq n.
\]

(i) et (ii) résultent de la théorie générale (8.1.9). Pour prouver (iii), on contemple (8.1.19.2), et on remarque que les structures de Hodge \(\mathrm{H}^{p}(\widetilde{\mathrm{Y}}_{q}^{r},\varepsilon_{\mathbf{Q}}^{r})(-r)\) \((p+q+r=n)\) vérifient (iii). Pour prouver (iv), on remarque que les structures de Hodge \(\mathrm{H}^{p}(\mathrm{X}_{q})\) \((p+q=n)\) vérifient (iv).

\textbf{(8.1.21)} On définit un \(\mathrm{A}\)-\textit{complexe de Hodge mixte cohomologique sur un espace topologique \(r\)-simplicial} en transposant (8.1.11). Si \(\mathrm{K}\) est un tel complexe, on voit comme en (8.1.13) que \(\mathrm{R}\Gamma^\bullet\mathrm{K}\) est un \(\mathrm{A}\)-complexe de Hodge mixte \(r\)-cosimplicial (8.1.10.1).

\Page{36}
Soit \(\mathrm{X}_{\bullet}\) un schéma propre et lisse \(r\)-simplicial. Un diviseur à croisements normaux \(\mathrm{D}_{\bullet}\) dans \(\mathrm{X}_{\bullet}\) est un système de diviseurs à croisements normaux dans les \(\mathrm{X}_{n}\) tel que \(\mathrm{X}_{\bullet}-\mathrm{D}_{\bullet}=\mathrm{U}_{\bullet}\) soit un sous-schéma \(r\)-simplicial de \(\mathrm{X}_{\bullet}\). Soit \(j:\mathrm{U}_{\bullet}\hookrightarrow\mathrm{X}_{\bullet}\). On définit comme en (8.1.12) un complexe de Hodge mixte cohomologique \(\mathrm{R}j_{\ast}\mathbf{Z}\) sur \(\mathrm{X}_{\bullet}\).

On pourrait comme en (8.1.20) en déduire une structure de Hodge mixte sur \(\mathrm{H}^{\bullet}(\mathrm{U}_{\bullet},\mathbf{Z})\), mais, puisque
\[
\mathrm{H}^{\bullet}(\mathrm{U}_{\bullet},\mathbf{Z})
=
\mathrm{H}^{\bullet}(\delta\mathrm{U}_{\bullet},\mathbf{Z}),
\]
ceci ne nous apprendrait rien de neuf (l’égalité précédente serait une identité entre structures de Hodge mixtes).

\textbf{(8.1.22)} Faisons \(r=2\). Pour tout complexe bicosimplicial \(\mathrm{K}=(\mathrm{K}^{p n_{1} n_{2}})\) (\(p=\text{degré}\) différentiel, \(n_i=i\)-ème degré simplicial) borné inférieurement, on désignera par \(\mathbf{s}_{1}\mathrm{K}\) le complexe cosimplicial de degré simplicial \(n_{2}\), donné par
\[
(\mathbf{s}_{1}\mathrm{K})^{\bullet n_{2}}
=
\mathbf{s}\mathrm{K}^{\bullet\bullet n_{2}}.
\]
On désignera par \(\mathrm{L}_{1}\) la filtration de \(\mathbf{s}_{1}\mathrm{K}\) qui, sur chaque \(\mathbf{s}\mathrm{K}^{\bullet\bullet n_{2}}\) induit la filtration par le degré simplicial \(n_{1}\).

Soit \(\mathrm{K}\) un \(\mathrm{A}\)-complexe de Hodge mixte bicosimplicial. Il résulte de (8.1.15) que \((\mathbf{s}_{1}\mathrm{K},\delta(\mathrm{W},\mathrm{L}_{1}),\mathrm{F})\) est un complexe de Hodge mixte cosimplicial. Soit \(\mathrm{L}_{2}\) la filtration de \(\mathbf{s}\mathrm{K}=\mathbf{s}\mathbf{s}_{1}\mathrm{K}\) par le degré simplicial \(n_{2}\). D’après (8.1.14), (8.1.15), la suite spectrale définie par \(\mathrm{L}_{2}\) provient d’une suite spectrale dans la catégorie abélienne des \(\mathrm{A}\)-structures de Hodge mixtes
\[
{}_{\mathrm{L}_{2}}\mathrm{E}_{1}^{pq}
=
\mathrm{H}^{q}(\mathbf{s}\mathrm{K}^{\bullet\bullet p})
\Rightarrow
\mathrm{H}^{p+q}(\mathbf{s}\mathrm{K}).
\tag*{\textbf{(8.1.22.1)}}
\]
Dans cette suite spectrale, la \(\mathrm{A}\)-structure de Hodge mixte sur \(\mathrm{E}_{1}^{pq}\) est la \(\mathrm{A}\)-structure de Hodge mixte ((8.1.9) (ii)) de la cohomologie du \(\mathrm{A}\)-complexe de Hodge mixte \((\mathbf{s}\mathrm{K}^{\bullet\bullet p},\delta(\mathrm{W},\mathrm{L}_{1}),\mathrm{F})\). La \(\mathrm{A}\)-structure de Hodge mixte de l’aboutissement est la \(\mathrm{A}\)-structure de Hodge mixte de la cohomologie du \(\mathrm{A}\)-complexe de Hodge mixte \((\mathbf{s}\mathrm{K},\delta(\delta(\mathrm{W},\mathrm{L}_{1}),\mathrm{L}_{2}),\mathrm{F})\). La filtration \(\delta(\delta(\mathrm{W},\mathrm{L}_{1}),\mathrm{L}_{2})\underset{\mathrm{dfn}}{=}\delta(\mathrm{W},\mathrm{L}_{1},\mathrm{L}_{2})\) est donnée par
\[
\delta(\mathrm{W},\mathrm{L}_{1},\mathrm{L}_{2})_{k}
(\mathrm{K}^{p n_{1} n_{2}})
=
\mathrm{W}_{k+n_{1}+n_{2}}\mathrm{K}^{p n_{1} n_{2}}.
\]

\textbf{(8.1.23)} En particulier, un \(\mathrm{A}\)-complexe de Hodge mixte cohomologique sur un espace topologique bisimplicial \(\mathrm{X}_{\bullet\bullet}\) définit une suite spectrale de \(\mathrm{A}\)-structures de Hodge mixtes
\[
\mathrm{E}_{1}^{pq}
=
\mathrm{H}^{q}(\mathrm{X}_{\bullet p},\mathrm{K})
\Rightarrow
\mathrm{H}^{p+q}(\mathrm{X}_{\bullet\bullet},\mathrm{K}).
\tag*{\textbf{(8.1.23.1)}}
\]
Dans le cas particulier considéré en (8.1.21), pour \(r=2\), on a \(\mathrm{A}=\mathbf{Z}\) et (8.1.23.1) s’écrit
\[
\mathrm{E}_{1}^{pq}
=
\mathrm{H}^{q}(\mathrm{U}_{\bullet p},\mathbf{Z})
\Rightarrow
\mathrm{H}^{p+q}(\mathrm{U}_{\bullet\bullet},\mathbf{Z}).
\tag*{\textbf{(8.1.23.2)}}
\]

% ---- 037-044 ----
\Page{37}
\textbf{(8.1.24)} Si \(\mathrm{K}'\) et \(\mathrm{K}''\) sont deux \(\mathrm{A}\)-complexes de Hodge mixtes, leur produit tensoriel \(\mathrm{K}'\otimes\mathrm{K}''\), défini par
\[
\begin{aligned}
&(\mathrm{K}'\otimes\mathrm{K}'')_{\mathrm{A}}
 =\mathrm{K}'_{\mathrm{A}}\overset{\mathrm{L}}{\otimes}\mathrm{K}''_{\mathrm{A}},\\
&((\mathrm{K}'\otimes\mathrm{K}'')_{\mathrm{A}\otimes\mathbf{Q}},\mathrm{W})
 =(\mathrm{K}'_{\mathrm{A}\otimes\mathbf{Q}},\mathrm{W}')
 \otimes(\mathrm{K}''_{\mathrm{A}\otimes\mathbf{Q}},\mathrm{W}''),\\
&((\mathrm{K}'\otimes\mathrm{K}'')_{\mathbf{C}},\mathrm{W},\mathrm{F})
 =(\mathrm{K}'_{\mathbf{C}},\mathrm{W}',\mathrm{F}')
 \otimes(\mathrm{K}''_{\mathbf{C}},\mathrm{W}'',\mathrm{F}''),
\end{aligned}
\]
est encore un \(\mathrm{A}\)-complexe de Hodge mixte.

Soit \(\mathrm{U}'\) (resp. \(\mathrm{U}''\)) le complément d’un diviseur à croisements normaux \(\mathrm{Y}'\) (resp. \(\mathrm{Y}''\)) dans un schéma propre et lisse \(\mathrm{X}'\) (resp. \(\mathrm{X}''\)). Alors, \(\mathrm{U}=\mathrm{U}'\times\mathrm{U}''\) est le complément d’un diviseur à croisements normaux \(\mathrm{Y}\) dans \(\mathrm{X}=\mathrm{X}'\times\mathrm{X}''\). Soit \(j\) (resp. \(j',j''\)) l’inclusion de \(\mathrm{U}\) (resp. \(\mathrm{U}',\mathrm{U}''\)) dans \(\mathrm{X}\) (resp. \(\mathrm{X}',\mathrm{X}''\)).

On dispose d’un quasi-isomorphisme
\[
\mathrm{R}j'_*\mathbf{Z}\overset{\mathrm{L}}{\boxtimes}\mathrm{R}j''_*\mathbf{Z}
\underset{\mathrm{dfn}}{=}
\operatorname{pr}_1^*\mathrm{R}j'_*\mathbf{Z}
\overset{\mathrm{L}}{\otimes}
\operatorname{pr}_2^*\mathrm{R}j''_*\mathbf{Z}
\xrightarrow{\sim}\mathrm{R}j_*\mathbf{Z},
\]
d’un quasi-isomorphisme filtré
\[
(\mathrm{R}j'_*\mathbf{Q},\tau_{\leq})
\overset{\mathrm{L}}{\boxtimes}
(\mathrm{R}j''_*\mathbf{Q},\tau_{\leq})
\xrightarrow{\sim}
(\mathrm{R}j_*\mathbf{Q},\tau_{\leq}),
\]
et d’un morphisme bifiltré
\[
(\Omega_{\mathrm{X}'}^\bullet(\log\mathrm{Y}'),\mathrm{W},\mathrm{F})
\boxtimes
(\Omega_{\mathrm{X}''}^\bullet(\log\mathrm{Y}''),\mathrm{W},\mathrm{F})
\longrightarrow
(\Omega_{\mathrm{X}}^\bullet(\log\mathrm{Y}),\mathrm{W},\mathrm{F}).
\]

Appliquons le foncteur \(\mathrm{R}\Gamma\); d’après la formule de Künneth, et d’après son analogue en cohomologie des faisceaux analytiques cohérents, on obtient un isomorphisme de complexes de Hodge mixtes
\[
\mathrm{R}\Gamma(\mathrm{U}',\mathbf{Z})
\overset{\mathrm{L}}{\otimes}
\mathrm{R}\Gamma(\mathrm{U}'',\mathbf{Z})
\xrightarrow{\sim}
\mathrm{R}\Gamma(\mathrm{U},\mathbf{Z}),
\]
d’où un isomorphisme de \(\mathbf{Q}\)-structures de Hodge mixtes
\[
\mathrm{H}^\bullet(\mathrm{U}',\mathbf{Q})
\otimes\mathrm{H}^\bullet(\mathrm{U}'',\mathbf{Q})
\xrightarrow{\sim}\mathrm{H}^\bullet(\mathrm{U},\mathbf{Q}).
\]

\textbf{(8.1.25)} Des arguments standards, basés sur le théorème de Cartier-Eilenberg-Zilber, fournissent la variante simpliciale suivante de (8.1.24).

Soit \(\mathrm{U}'_\bullet\) (resp. \(\mathrm{U}''_\bullet\)) le complément d’un diviseur à croisements normaux dans un schéma simplicial propre et lisse \(\mathrm{X}'_\bullet\) (resp. \(\mathrm{X}''_\bullet\)). Soit \(\mathrm{U}_\bullet=\mathrm{U}'_\bullet\times\mathrm{U}''_\bullet\) et \(\mathrm{X}_\bullet=\mathrm{X}'_\bullet\times\mathrm{X}''_\bullet\).

On a un isomorphisme de complexes de Hodge mixtes
\[
\mathrm{R}\Gamma(\mathrm{U}'_\bullet,\mathbf{Z})
\overset{\mathrm{L}}{\otimes}
\mathrm{R}\Gamma(\mathrm{U}''_\bullet,\mathbf{Z})
\xrightarrow{\sim}
\mathrm{R}\Gamma(\mathrm{U}_\bullet,\mathbf{Z}),
\]
et en particulier un isomorphisme de \(\mathbf{Q}\)-structures de Hodge mixtes
\[
\mathrm{H}^\bullet(\mathrm{U}'_\bullet,\mathbf{Q})
\otimes\mathrm{H}^\bullet(\mathrm{U}''_\bullet,\mathbf{Q})
\xrightarrow{\sim}\mathrm{H}^\bullet(\mathrm{U}_\bullet,\mathbf{Q}).
\]

\textbf{(8.1.26)} Soit \(\mathrm{U}'_{\bullet\bullet}\) (resp. \(\mathrm{U}''_{\bullet\bullet}\)) le complément d’un diviseur à croisements normaux dans un schéma bisimplicial propre et lisse \(\mathrm{X}'_{\bullet\bullet}\) (resp. \(\mathrm{X}''_{\bullet\bullet}\)). Soient \(\mathrm{U}_{\bullet\bullet}=\mathrm{U}'_{\bullet\bullet}\times\mathrm{U}''_{\bullet\bullet}\) et
\Page{38}
\(\mathrm{X}_{\bullet\bullet}=\mathrm{X}'_{\bullet\bullet}\times\mathrm{X}''_{\bullet\bullet}\). Des raisonnements standards montrent que, en cohomologie rationnelle, la suite spectrale (8.1.23.2), pour \(\mathrm{U}_{\bullet\bullet}\), est produit tensoriel des suites spectrales (8.1.23.2) pour \(\mathrm{U}'_{\bullet\bullet}\) et \(\mathrm{U}''_{\bullet\bullet}\).

\section*{8.2. Espaces algébriques séparés.}

\textbf{(8.2.1)} Soit \(\mathrm{X}\) un schéma (ou un espace algébrique) de type fini sur \(\mathbf{C}\), supposé séparé. Il existe alors un diagramme
\[
\begin{tikzcd}[column sep=large,row sep=large]
\mathrm{Y}_\bullet \arrow[r,hook] \arrow[d,"\alpha"']
  & \overline{\mathrm{Y}}_\bullet\\
\mathrm{X} &
\end{tikzcd}
\tag*{\textbf{(8.2.1.1)}}
\]
dans lequel \(\mathrm{Y}_\bullet\) est le complément d’un diviseur à croisements normaux dans un schéma simplicial propre et lisse \(\overline{\mathrm{Y}}_\bullet\), et est un hyperrecouvrement propre de \(\mathrm{X}\). On a (5.3.5 (II))
\[
\mathrm{H}^\bullet(\mathrm{X})
\xrightarrow{\sim}\mathrm{H}^\bullet(\mathrm{Y}_\bullet).
\tag*{\textbf{(8.2.1.2)}}
\]

En (8.1.19), on a muni \(\mathrm{H}^\bullet(\mathrm{Y}_\bullet)\) d’une structure de Hodge mixte. Soit \(\mathcal{H}(\alpha)\) la structure de Hodge mixte sur \(\mathrm{H}^\bullet(\mathrm{X})\) qui s’en déduit par (8.2.1.2).

\textit{Proposition} \textbf{(8.2.2)}. — \textit{Avec les notations précédentes, la structure de Hodge mixte \(\mathcal{H}(\alpha)\) sur \(\mathrm{H}^\bullet(\mathrm{X})\) est indépendante du choix de \(\alpha\). On l’appelle la} structure de Hodge mixte de \(\mathrm{H}^\bullet(\mathrm{X},\mathbf{Z})\). \textit{Pour tout morphisme de schémas \(f:\mathrm{X}_1\to\mathrm{X}_2\), \(f^\bullet:\mathrm{H}^\bullet(\mathrm{X}_2)\to\mathrm{H}^\bullet(\mathrm{X}_1)\) est un morphisme de structures de Hodge mixtes.}

La démonstration utilise (6.2.8), et est parallèle à celle de (3.2.11 C)) (cf. aussi (8.3.3)).

\textit{Remarque} \textbf{(8.2.3)}. — Pour \(\mathrm{X}\) lisse, on vérifie en prenant pour \(\overline{\mathrm{Y}}_\bullet\) un schéma simplicial constant que la structure de Hodge mixte (8.2.2) coïncide avec celle définie en (3.2.12).

\textit{Théorème} \textbf{(8.2.4)}. — \textit{Les couples \((p,q)\) tels que le nombre de Hodge \(h^{pq}\) de \(\mathrm{H}^n(\mathrm{X})\) soit non nul vérifient les conditions suivantes}

\noindent\hspace*{\parindent}\makebox[2em][l]{(i)}
\((p,q)\in[0,n]\times[0,n]\).

\noindent\hspace*{\parindent}\makebox[2em][l]{(ii)}
\textit{Si \(\dim\mathrm{X}=\mathrm{N}\), et que \(n\geq\mathrm{N}\), alors}
\[
(p,q)\in[n-\mathrm{N},\mathrm{N}]\times[n-\mathrm{N},\mathrm{N}].
\]

\noindent\hspace*{\parindent}\makebox[2em][l]{(iii)}
\textit{Si \(\mathrm{X}\) est propre, alors \(p+q\leq n\).}

\noindent\hspace*{\parindent}\makebox[2em][l]{(iv)}
\textit{Si \(\mathrm{X}\) est lisse, alors \(p+q\geq n\). La même conclusion vaut pour \(\mathrm{X}\) une « rational homology manifold » équidimensionnelle de dimension \(\mathrm{N}\), \textup{i.e.} si pour tout \(x\in\mathrm{X}\), \(\mathrm{H}^i_{\{x\}}(\mathrm{X},\mathbf{Q})=0\) si \(i\neq2\mathrm{N}\), \(=\mathbf{Q}\) si \(i=2\mathrm{N}\).}

\Page{39}
\[
\begin{tikzcd}[
  baseline=(current bounding box.south),
  row sep=7mm,column sep=7mm,
  cells={nodes={inner sep=0pt,outer sep=0pt,text height=0pt,text depth=0pt}},
  arrows={start anchor=center,end anchor=center}
]
{} & {} & {} & {} & {}\\
{} \arrow[rr,no head,"\text{lisse}"'{pos=.7,yshift=-3mm}]
   \arrow[ddrr,no head]
 & {}
 & {} \arrow[dd,no head,"n"{pos=1,anchor=north,yshift=-2mm}]
 & {} & {}\\
{} & {} & {} & {} & {}\\
{} \arrow[uuu,"q"{pos=1,anchor=east,xshift=-2mm,yshift=-2mm}]
   \arrow[rrrr,"p"{pos=1,anchor=west,xshift=2mm}]
   \arrow[rrrr,phantom,"\text{propre}"{pos=.18,anchor=south,yshift=1mm}]
   \arrow[rrrr,phantom,"n\leq\mathrm{N}"'{pos=.22,anchor=north,yshift=-7mm}]
 & {} & {} & {} & {}
\end{tikzcd}
\qquad\qquad
\begin{tikzcd}[
  baseline=(current bounding box.south),
  row sep=5.5mm,column sep=5.5mm,
  cells={nodes={inner sep=0pt,outer sep=0pt,text height=0pt,text depth=0pt}},
  arrows={start anchor=center,end anchor=center}
]
{} & {} & {} & {} & {} & {} & {}\\
{} \arrow[dr,dashed,no head]
 & {} & {}
 & {} \arrow[d,dashed,no head]
 & {} & {} & {}\\
{} \arrow[r,dashed,no head,"\mathrm{N}"{pos=0,anchor=east,xshift=-2mm}]
 & {} \arrow[rr,no head,"\text{lisse}"'{pos=.7,yshift=-2mm}]
      \arrow[dd,no head]
      \arrow[ddrr,no head]
 & {}
 & {} \arrow[r,dashed,no head] \arrow[dd,no head]
 & {} & {} & {}\\
{} & {} & {} & {} & {} & {} & {}\\
{} & {} \arrow[rr,no head,"\text{propre}"{pos=.43,anchor=south,yshift=.7mm}]
 & {}
 & {} \arrow[dr,dashed,no head,"n"{pos=1,anchor=north,yshift=-2mm}]
      \arrow[d,dashed,no head,"\mathrm{N}"{pos=1,anchor=north,yshift=-2mm}]
 & {} & {} & {}\\
{} \arrow[uuuuu,"q"{pos=1,anchor=east,xshift=-2mm,yshift=-2mm}]
   \arrow[rrrrrr,"p"{pos=1,anchor=west,xshift=2mm}]
   \arrow[rrrrrr,phantom,"n\geq\mathrm{N}"'{pos=.61,anchor=north,yshift=-7mm}]
 & {} & {} & {} & {} & {} & {}
\end{tikzcd}
\]

L’assertion (ii) sera prouvée en (8.3.10). L’assertion (i) résulte de (8.1.20 (iii)). Pour \(\mathrm{X}\) propre, on peut trouver un diagramme (8.2.1.1) dans lequel \(\mathrm{Y}_\bullet=\overline{\mathrm{Y}}_\bullet\) et (iii) résulte de (8.1.20 (iv)). Pour \(\mathrm{X}\) lisse, (iv) résulte de (8.2.3) et (3.2.15). Le cas général en résulte : si \(p:\widetilde{\mathrm{X}}\to\mathrm{X}\) est une résolution des singularités de \(\mathrm{X}\), \(p^*:\mathrm{H}^\bullet(\mathrm{X},\mathbf{Q})\to\mathrm{H}^\bullet(\widetilde{\mathrm{X}},\mathbf{Q})\) est injectif, car le transposé par dualité de Poincaré \(p_!\) de \(p^*:\mathrm{H}_c^\bullet(\mathrm{X},\mathbf{Q})\to\mathrm{H}_c^\bullet(\widetilde{\mathrm{X}},\mathbf{Q})\) en est une rétraction.

\textit{Proposition} \textbf{(8.2.5)}. — \textit{Supposons \(\mathrm{X}\) propre. Si \(u:\mathrm{Y}\to\mathrm{X}\) est un morphisme propre surjectif, avec \(\mathrm{Y}\) lisse, alors le quotient de poids \(n\) de \(\mathrm{H}^n(\mathrm{X},\mathbf{Q})\) est l’image de \(\mathrm{H}^n(\mathrm{X},\mathbf{Q})\) dans \(\mathrm{H}^n(\mathrm{Y},\mathbf{Q})\). La suite}
\[
\mathrm{H}^n(\mathrm{X},\mathbf{Q})
\longrightarrow
\mathrm{H}^n(\mathrm{Y},\mathbf{Q})
\xrightarrow[\operatorname{pr}_1^*-\operatorname{pr}_2^*]{}
\mathrm{H}^n(\mathrm{Y}\times_{\mathrm{X}}\mathrm{Y},\mathbf{Q})
\tag*{\textbf{(8.2.5.1)}}
\]
\textit{est exacte.}

Soit \(\mathrm{Y}_\bullet\) comme en (8.2.1.1), avec \(\mathrm{Y}_0=\mathrm{Y}\) et \(\mathrm{Y}_\bullet=\overline{\mathrm{Y}}_\bullet\). D’après ((8.1.20 (ii))), on a une suite exacte
\[
\begin{tikzcd}[column sep=small]
0 \arrow[r]
 & \mathrm{H}^n(\mathrm{X},\mathbf{Q})/
   \mathrm{W}^{n-1}\mathrm{H}^n(\mathrm{X},\mathbf{Q}) \arrow[r]
 & \mathrm{H}^n(\mathrm{Y}_0,\mathbf{Q})
   \arrow[r,shift left=.4ex,"\delta_0"]
   \arrow[r,shift right=.4ex,"\delta_1"']
 & \mathrm{H}^n(\mathrm{Y}_1,\mathbf{Q})
\end{tikzcd}
\tag*{\textbf{(8.2.5.2)}}
\]
qui implique (8.2.5) (réciproquement, cette suite exacte se déduit aisément de (8.2.5) et (8.2.7) ci-dessous).

\textit{Proposition} \textbf{(8.2.6)}. — \textit{Soient des morphismes de schémas}
\[
\begin{tikzcd}[column sep=small]
\mathrm{Y} \arrow[r,"f"]
 & \mathrm{X} \arrow[r,hook,"j"]
 & \overline{\mathrm{X}}.
\end{tikzcd}
\]
\textit{On suppose que \(\mathrm{Y}\) est propre, que \(\mathrm{X}\) est lisse et que \(\overline{\mathrm{X}}\) est une compactification lisse de \(\mathrm{X}\). Alors, \(\mathrm{H}^\bullet(\overline{\mathrm{X}},\mathbf{Q})\) et \(\mathrm{H}^\bullet(\mathrm{X},\mathbf{Q})\) ont même image dans \(\mathrm{H}^\bullet(\mathrm{Y},\mathbf{Q})\).}

La démonstration est parallèle à celle de (3.2.18). Il suffit de prouver que, quels que soient \(i\) et \(n\) \(\operatorname{Gr}_i^{\mathrm{W}}(f^*)\) et \(\operatorname{Gr}_i^{\mathrm{W}}((jf)^*)\) ont même image dans \(\operatorname{Gr}_i^{\mathrm{W}}\mathrm{H}^n(\mathrm{Y},\mathbf{Q})\).

\textit{a)} Pour \(i>n\), ces images sont nulles car \(\operatorname{Gr}_i^{\mathrm{W}}\mathrm{H}^n(\mathrm{Y},\mathbf{Q})=0\) (8.2.4 (iii)).

\textit{b)} Pour \(i<n\), ces images sont nulles car \(\operatorname{Gr}_i^{\mathrm{W}}\mathrm{H}^n(\mathrm{X},\mathbf{Q})=0\) (8.2.4 (iv)).

\textit{c)} Pour \(i=n\), ces images sont égales car \(\operatorname{Gr}_n^{\mathrm{W}}(j^*)\) est surjectif (3.2.17).

\Page{40}
\textit{Proposition} \textbf{(8.2.7)}. — \textit{Soient des morphismes de schémas}
\[
\begin{tikzcd}[column sep=small]
\widetilde{\mathrm{X}} \arrow[r,"\pi"]
 & \mathrm{X} \arrow[r,"f"]
 & \mathrm{Y}.
\end{tikzcd}
\]
\textit{On suppose que \(\mathrm{Y}\) est lisse, que \(\mathrm{X}\) est propre, que \(\widetilde{\mathrm{X}}\) est propre et lisse et que \(\pi\) est surjectif. Alors, les noyaux de \(f^*\) et de \((f\pi)^*\) dans \(\mathrm{H}^\bullet(\mathrm{Y},\mathbf{Q})\) sont égaux.}

Prouvons que \(\operatorname{Ker}\operatorname{Gr}_i^{\mathrm{W}}(f^*)=\operatorname{Ker}(\operatorname{Gr}_i^{\mathrm{W}}((f\pi)^*))\) dans \(\operatorname{Gr}_i^{\mathrm{W}}(\mathrm{H}^n(\mathrm{Y},\mathbf{Q}))\).

\textit{a)} Pour \(i<n\), ces noyaux sont nuls car \(\operatorname{Gr}_i^{\mathrm{W}}\mathrm{H}^n(\mathrm{Y},\mathbf{Q})=0\) ((8.2.4 (iv))).

\textit{b)} Pour \(i>n\), ces noyaux sont tout \(\operatorname{Gr}_i^{\mathrm{W}}(\mathrm{H}^n(\mathrm{Y},\mathbf{Q}))\) car \(\operatorname{Gr}_i^{\mathrm{W}}\mathrm{H}^n(\mathrm{X},\mathbf{Q})=0\) ((8.2.4 (iii))).

\textit{c)} Pour \(i=n\), ces noyaux sont égaux car \(\operatorname{Gr}_n^{\mathrm{W}}(\pi^*)\) est injectif (8.2.5.2).

\textit{Corollaire} \textbf{(8.2.8)}. — \textit{Soient des morphismes de schémas}
\[
\begin{tikzcd}[column sep=small]
\widetilde{\mathrm{X}} \arrow[r,"\pi"]
 & \mathrm{X} \arrow[r,hook,"f"]
 & \mathrm{Y}.
\end{tikzcd}
\]
\textit{On suppose que \(\mathrm{Y}\) est propre et lisse, que \(\mathrm{X}\) est un sous-schéma fermé de \(\mathrm{Y}\) et que \(\widetilde{\mathrm{X}}\) est une résolution des singularités de \(\mathrm{X}\). Soient \(q=f\pi\) et \(\mathrm{U}=\mathrm{Y}-\mathrm{X}\). Alors, la suite}
\[
\mathrm{H}^\bullet(\widetilde{\mathrm{X}},\mathbf{Q})
\xrightarrow{q_!}\mathrm{H}^\bullet(\mathrm{Y},\mathbf{Q})
\longrightarrow\mathrm{H}^\bullet(\mathrm{U},\mathbf{Q})
\tag*{\textbf{(8.2.8.1)}}
\]
\textit{est exacte.}

D’après (8.2.7), on a
\[
\operatorname{Ker}(f^*:\mathrm{H}^\bullet(\mathrm{Y},\mathbf{Q})
 \to\mathrm{H}^\bullet(\mathrm{X},\mathbf{Q}))
=
\operatorname{Ker}(q^*:\mathrm{H}^\bullet(\mathrm{Y},\mathbf{Q})
 \to\mathrm{H}^\bullet(\widetilde{\mathrm{X}},\mathbf{Q})).
\]
La suite exacte longue de cohomologie du couple \((\mathrm{Y},\mathrm{X})\) fournit donc une suite exacte
\[
\mathrm{H}_c^\bullet(\mathrm{U},\mathbf{Q})
\longrightarrow\mathrm{H}^\bullet(\mathrm{Y},\mathbf{Q})
\longrightarrow\mathrm{H}^\bullet(\widetilde{\mathrm{X}},\mathbf{Q}),
\tag*{\textbf{(8.2.8.2)}}
\]
et (8.2.8.1) est la suite transposée de (8.2.8.2) par dualité de Poincaré.

\textit{Remarque} \textbf{(8.2.9)}. — Le corollaire (8.2.8) est utilisé sous des hypothèses plus générales dans ([8], (9.7), p. 161). Comme l’a fait remarquer Lieberman, la justification donnée dans \textit{loc. cit.} est insuffisante.

On déduit de (8.1.25) que

\textit{Proposition} \textbf{(8.2.10)}. — \textit{Les isomorphismes de Künneth}
\[
\mathrm{H}^\bullet(\mathrm{X}\times\mathrm{Y},\mathbf{Q})
\simeq
\mathrm{H}^\bullet(\mathrm{X},\mathbf{Q})
\otimes\mathrm{H}^\bullet(\mathrm{Y},\mathbf{Q})
\]
\textit{sont des isomorphismes de structures de Hodge mixtes.}

\textit{Corollaire} \textbf{(8.2.11)}. — \textit{Le cup-produit}
\[
\mathrm{H}^\bullet(\mathrm{X})\otimes\mathrm{H}^\bullet(\mathrm{X})
\longrightarrow\mathrm{H}^\bullet(\mathrm{X})
\]
\textit{est un morphisme de structures de Hodge mixtes.}

Résulte de (8.2.10) et de (8.2.2) appliqué à l’application diagonale \(\Delta:\mathrm{X}\to\mathrm{X}\times\mathrm{X}\).

\Page{41}
\section*{8.3. Théorie de Hodge des schémas simpliciaux.}

\textbf{(8.3.1)} Pour \(\mathrm{X}_\bullet\) un schéma (ou espace algébrique) simplicial séparé, nous aurons à considérer des diagrammes de schémas simpliciaux
\[
\begin{tikzcd}[column sep=large,row sep=large]
\mathrm{Z}_\bullet \arrow[r,hook,"j"] \arrow[d,"p"']
 & \overline{\mathrm{Z}}_\bullet\\
\mathrm{X}_\bullet &
\end{tikzcd}
\tag*{\textbf{(8.3.1.1)}}
\]
vérifiant :

\textit{a)} \(j\) est une immersion ouverte d’image dense; \(\overline{\mathrm{Z}}_\bullet\) est propre.

\textit{b)} Les morphismes évidents \(\mathrm{Z}_n\to(\operatorname{cosq}_{n-1}^{\mathrm{X}_\bullet}\operatorname{sq}_{n-1}\mathrm{Z}_\bullet)_n\) sont propres et surjectifs.

Pour \(\overline{\mathrm{Z}}_\bullet\) réduit, on peut regarder \(j\) comme un objet simplicial de la catégorie \(\mathcal{C}\) des couples \((\mathrm{Y},\overline{\mathrm{Y}})\) formés d’un schéma propre réduit \(\overline{\mathrm{Y}}\) et d’un ouvert (de Zariski) dense \(\mathrm{Y}\) de \(\overline{\mathrm{Y}}\). La catégorie \(\mathcal{C}\) vérifie (6.2.4.1), ce qui permet de lui appliquer l’analogue de (6.2.4).

\textit{Proposition} \textbf{(8.3.2)}. — (i) \textit{Pour tout schéma simplicial séparé \(\mathrm{X}_\bullet\), il existe un diagramme \textup{(8.3.1.1)} vérifiant a), b) et tel que \(\overline{\mathrm{Z}}_\bullet\) soit propre et lisse et \(\mathrm{Z}_\bullet\) le complément d’un diviseur à croisements normaux.}

(ii) \textit{Soit \(f:\mathrm{X}_{\bullet1}\to\mathrm{X}_{\bullet2}\) un morphisme de schémas simpliciaux séparés. Soient \((\mathrm{X}_{\bullet i},\mathrm{Z}_{\bullet i},\overline{\mathrm{Z}}_{\bullet i})_{i=1,2}\) deux diagrammes \textup{(8.3.1.1)} vérifiant a) et b). Il existe alors \((\mathrm{X}_{\bullet1},\mathrm{Z}_\bullet,\overline{\mathrm{Z}}_\bullet)\) comme en \textup{(i)} et un diagramme commutatif}
\[
\begin{tikzcd}[column sep=large,row sep=large]
\mathrm{Z}_\bullet
  \arrow[r,hook] \arrow[d] \arrow[drr]
 & \overline{\mathrm{Z}}_\bullet
   \arrow[d,crossing over] \arrow[drr]
 & &\\
\mathrm{Z}_{\bullet1}
  \arrow[r,hook,"j_1"'] \arrow[d,"p_1"']
 & \overline{\mathrm{Z}}_{\bullet1}
 & \mathrm{Z}_{\bullet2}
   \arrow[r,hook,"j_2"'] \arrow[d,"p_2"']
 & \overline{\mathrm{Z}}_{\bullet2}\\
\mathrm{X}_{\bullet1} \arrow[rr,"f"]
 & & \mathrm{X}_{\bullet2} &
\end{tikzcd}
\]

La démonstration, analogue en plus compliqué à (6.2.5), est laissée au lecteur.

\textbf{(8.3.3)} Soit un diagramme (8.3.1.1) comme en (8.3.2 (i)). D’après (5.3.5 (V)), l’application \(\mathrm{H}^\bullet(\mathrm{X}_\bullet,\mathbf{Z})\to\mathrm{H}^\bullet(\mathrm{Z}_\bullet,\mathbf{Z})\) est un isomorphisme. En (8.1.19) et (8.1.20), nous avons muni \(\mathrm{H}^\bullet(\mathrm{Z}_\bullet,\mathbf{Z})\) d’une structure de Hodge mixte, d’où une structure de
\Page{42}
Hodge mixte sur \(\mathrm{H}^\bullet(\mathrm{X}_\bullet,\mathbf{Z})\). Elle est indépendante du choix du diagramme (8.3.1.1) : d’après (8.3.2 (ii)) pour \(f=\mathrm{Id}\), deux systèmes \((\mathrm{Z}_{\bullet i},\overline{\mathrm{Z}}_{\bullet i},p_i)_{i=1,2}\) sont toujours coiffés par un troisième \((\mathrm{Z}_\bullet,\overline{\mathrm{Z}}_\bullet,p)\), d’où un diagramme commutatif d’isomorphismes
\[
\begin{tikzcd}[column sep=large,row sep=huge]
 & \mathrm{H}^\bullet(\mathrm{X}_\bullet,\mathbf{Z})
   \arrow[dl] \arrow[d] \arrow[dr] &\\
\mathrm{H}^\bullet(\mathrm{Z}_{\bullet1},\mathbf{Z}) \arrow[r]
 & \mathrm{H}^\bullet(\mathrm{Z}_\bullet,\mathbf{Z})
 & \mathrm{H}^\bullet(\mathrm{Z}_{\bullet2},\mathbf{Z}) \arrow[l]
\end{tikzcd}
\]

\textit{Définition} \textbf{(8.3.4)}. — \textit{La} structure de Hodge mixte \textit{de \(\mathrm{H}^\bullet(\mathrm{X}_\bullet,\mathbf{Z})\) est celle définie plus haut.}

Il résulte de (8.3.2) (ii) que cette structure de Hodge mixte est fonctorielle en \(\mathrm{X}_\bullet\).

\textit{Proposition} \textbf{(8.3.5)}. — \textit{Soit \(\mathrm{X}_\bullet\) un schéma simplicial séparé. La suite spectrale}
\[
\mathrm{E}_1^{pq}=\mathrm{H}^q(\mathrm{X}_p,\mathbf{Z})
\Rightarrow\mathrm{H}^{p+q}(\mathrm{X}_\bullet,\mathbf{Z})
\]
\textit{est une suite spectrale de structures de Hodge mixtes.}

Pour un diagramme comme en (8.3.1), les applications \(p_n:\mathrm{Z}_n\to\mathrm{X}_n\) sont propres et surjectives. On peut alors appliquer (6.2.4) à la catégorie des objets simpliciaux de la catégorie \(\mathcal{C}\) (8.3.1) pour obtenir comme en (6.2.5) le lemme suivant :

\textit{Lemme} \textbf{(8.3.6)}. — \textit{Soit \(\mathrm{X}_\bullet\) un schéma simplicial séparé. Il existe un schéma bisimplicial \(\mathrm{Z}_{\bullet\bullet}\) augmenté vers \(\mathrm{X}_\bullet\) :}
\[
\varepsilon_{n,m}:\mathrm{Z}_{n,m}\to\mathrm{X}_n
\]
\textit{et \(i:\mathrm{Z}_{\bullet\bullet}\hookrightarrow\overline{\mathrm{Z}}_{\bullet\bullet}\) tels que :}

\textit{a) \(\overline{\mathrm{Z}}_{\bullet\bullet}\) est un schéma bisimplicial propre et lisse, et \(\mathrm{Z}_{\bullet\bullet}\) le complément d’un diviseur à croisements normaux de \(\overline{\mathrm{Z}}_{\bullet\bullet}\);}

\textit{b) pour tout \(n\), le schéma simplicial augmenté \(\mathrm{Z}_{n\bullet}\to\mathrm{X}_n\) est un hyperrecouvrement propre de \(\mathrm{X}_n\).}

En (8.1.21), nous avons muni \(\mathrm{H}^\bullet(\mathrm{Z}_{\bullet\bullet},\mathbf{Z})\) d’une structure de Hodge mixte. Si \(\delta\mathrm{Z}_{\bullet\bullet}\) est le schéma simplicial diagonal de \(\mathrm{Z}_{\bullet\bullet}\), on a un diagramme commutatif de morphismes
\[
\begin{tikzcd}[column sep=large,row sep=huge]
\mathrm{H}^\bullet(\mathrm{X}_\bullet,\mathbf{Z})
  \arrow[r,"\gamma"] \arrow[dr,"\alpha"]
 & \mathrm{H}^\bullet(\mathrm{Z}_{\bullet\bullet},\mathbf{Z})
   \arrow[d,equal,"\beta"]\\
 & \mathrm{H}^\bullet(\delta\mathrm{Z}_{\bullet\bullet},\mathbf{Z})
\end{tikzcd}
\]
\Page{43}
dans lequel \(\beta\) et \(\gamma\) sont des isomorphismes, et \(\alpha\) et \(\beta\) des morphismes de structures de Hodge mixtes; \(\gamma\) est donc un isomorphisme de structures de Hodge mixtes.

On a un morphisme de suites spectrales
\[
\begin{tikzcd}[column sep=large,row sep=huge]
\mathrm{H}^q(\mathrm{X}_p,\mathbf{Z})
  \arrow[r,Rightarrow] \arrow[d]
 & \mathrm{H}^{p+q}(\mathrm{X}_\bullet,\mathbf{Z}) \arrow[d]
 & \text{\textbf{(5.2.3)}}\\
\mathrm{H}^q(\mathrm{Z}_{p\bullet},\mathbf{Z})
  \arrow[r,Rightarrow]
 & \mathrm{H}^{p+q}(\mathrm{Z}_{\bullet\bullet},\mathbf{Z})
 & \text{\textbf{(8.1.23.1)}}
\end{tikzcd}
\]
ce morphisme est un isomorphisme compatible aux structures de Hodge mixtes, de sorte que (8.3.5) résulte de (8.1.23).

\textit{Exemple} \textbf{(8.3.7)}. — Soit \(\mathrm{X}_\bullet\) un schéma simplicial non nécessairement séparé. Il existe alors un schéma simplicial séparé \(\mathrm{Y}_\bullet\) et \(u:\mathrm{Y}_\bullet\to\mathrm{X}_\bullet\) tel que pour \(n\geq0\), l’application \(u'_n:\mathrm{Y}_n\to(\operatorname{cosq}_{n-1}^{\mathrm{X}_\bullet}\operatorname{sq}_{n-1}\mathrm{Y}_\bullet)_n\) soit étale et surjective (par exemple la somme des ouverts d’un recouvrement ouvert).

L’application \(u^\bullet:\mathrm{H}^\bullet(\mathrm{X}_\bullet,\mathbf{Z})\to\mathrm{H}^\bullet(\mathrm{Y}_\bullet,\mathbf{Z})\) est un isomorphisme ((5.3.5) (V) et (III)). Par définition, la \textit{structure de Hodge mixte} de \(\mathrm{H}^\bullet(\mathrm{X}_\bullet,\mathbf{Z})\) se déduit de celle de \(\mathrm{H}^\bullet(\mathrm{Y}_\bullet,\mathbf{Z})\) \textit{via} \(u^\bullet\). Elle ne dépend pas du choix de \(\mathrm{Y}_\bullet\).

Pour \(\mathrm{X}\) un schéma non séparé, la \textit{structure de Hodge mixte} de \(\mathrm{H}^\bullet(\mathrm{X},\mathbf{Z})\) est celle de \(\mathrm{H}^\bullet(\mathrm{X}_\bullet,\mathbf{Z})\), pour \(\mathrm{X}_\bullet\) le schéma simplicial constant défini par \(\mathrm{X}\). Il est vraisemblable, mais non démontré, que (8.2.4) (iv) reste valable pour \(\mathrm{X}\) lisse non séparé.

\textit{Exemple} \textbf{(8.3.8)}. — Soit \(u:\mathrm{X}_\bullet\to\mathrm{Y}_\bullet\) un morphisme de schémas simpliciaux. Les groupes de cohomologie relative \(\mathrm{H}^n(\mathrm{Y}_\bullet\operatorname{mod}\mathrm{X}_\bullet,\mathbf{Z})\) sont par définition (6.3.2) les groupes de cohomologie du schéma simplicial \(\mathrm{C}(u)\). Ils sont donc munis d’une structure de Hodge mixte.

\textit{Proposition} \textbf{(8.3.9)}. — \textit{La suite exacte longue de cohomologie relative}
\[
\mathrm{H}^n(\mathrm{Y}_\bullet,\mathbf{Z})
\to\mathrm{H}^n(\mathrm{X}_\bullet,\mathbf{Z})
\xrightarrow{j}
\mathrm{H}^{n+1}(\mathrm{Y}_\bullet\operatorname{mod}\mathrm{X}_\bullet,\mathbf{Z})
\to\mathrm{H}^{n+1}(\mathrm{Y}_\bullet,\mathbf{Z})
\]
\textit{est une suite exacte longue de structures de Hodge mixtes.}

On se ramène au cas où \(\mathrm{X}_\bullet\) et \(\mathrm{Y}_\bullet\) sont lisses, et où, de plus, il existe un diagramme commutatif
\[
\begin{tikzcd}[column sep=large,row sep=large]
\mathrm{X}_\bullet \arrow[r,"u"] \arrow[d,hook]
 & \mathrm{Y}_\bullet \arrow[d,hook]\\
\overline{\mathrm{X}}_\bullet \arrow[r,"\bar{u}"]
 & \overline{\mathrm{Y}}_\bullet
\end{tikzcd}
\]
avec \(\overline{\mathrm{X}}_\bullet\) et \(\overline{\mathrm{Y}}_\bullet\) propres et lisses, et \(\mathrm{X}_\bullet\) et \(\mathrm{Y}_\bullet\) compléments de diviseurs à croisements normaux \(\mathrm{D}_\bullet\) et \(\mathrm{E}_\bullet\) dans \(\overline{\mathrm{X}}_\bullet\) et \(\overline{\mathrm{Y}}_\bullet\).

\Page{44}
Il existe alors des résolutions bifiltrées (de gradué à composantes acycliques) \(a:\Omega_{\overline{\mathrm{X}}_\bullet}^\bullet(\log\mathrm{D}_\bullet)\to\mathrm{K}\) et \(b:\Omega_{\overline{\mathrm{Y}}_\bullet}^\bullet(\log\mathrm{E}_\bullet)\to\mathrm{L}\) qui s’insèrent dans un diagramme commutatif de complexes bifiltrés
\[
\begin{tikzcd}[column sep=large,row sep=huge]
\Omega_{\overline{\mathrm{X}}_\bullet}^\bullet(\log\mathrm{D}_\bullet)
  \arrow[d,"a"]
 & \Omega_{\overline{\mathrm{Y}}_\bullet}^\bullet(\log\mathrm{E}_\bullet)
   \arrow[l,"\bar{u}^*"'] \arrow[d,"b"]\\
\mathrm{K}
 & \mathrm{L} \arrow[l,"v"']
\end{tikzcd}
\]
(\(\bar{u}^*\) et \(v\) sont des \(\bar{u}\)-morphismes). La suite exacte longue (8.3.9) s’identifie alors à la suite exacte déduite du triangle distingué
\[
\begin{tikzcd}[column sep=small,row sep=large]
 & \mathrm{C}(\Gamma v) \arrow[dl] &\\
\Gamma\mathrm{L} \arrow[rr,"v"]
 & & \Gamma\mathrm{K} \arrow[ul]
\end{tikzcd}
\]
Le complexe \(\mathrm{C}(v)\) sur \(\mathrm{C}(\bar{u})\) est une résolution bifiltrée de \(\mathrm{C}(u^*:\Omega_{\overline{\mathrm{Y}}_\bullet}^\bullet(\log\mathrm{E}_\bullet)\to\Omega_{\overline{\mathrm{X}}_\bullet}^\bullet(\log\mathrm{D}_\bullet))\). De là résulte que les filtrations de Hodge et par le poids de \(\mathrm{H}^\bullet(\mathrm{Y}_\bullet\operatorname{mod}\mathrm{X}_\bullet,\mathbf{C})\) se déduisent des filtrations du cône \(\mathrm{C}(\Gamma v)\), et (8.3.9) en résulte formellement.

\textbf{(8.3.10)} \textit{Démonstration de} (8.2.4) (ii).

La démonstration procède par récurrence sur \(\mathrm{N}=\dim(\mathrm{X})\). On peut supposer \(\mathrm{X}\) réduit. Soit \(\mathrm{U}\) un ouvert de Zariski dense lisse de \(\mathrm{X}\) et \(\mathrm{Y}=\mathrm{X}-\mathrm{U}\). D’après la résolution des singularités, il existe un morphisme propre birationnel \(p:\mathrm{X}'\to\mathrm{X}\) de source lisse qui induise un isomorphisme de \(p^{-1}(\mathrm{U})\) avec \(\mathrm{U}\). Soit \(\mathrm{Y}'=p^{-1}(\mathrm{Y})\). Considérons pour \(n\geq\mathrm{N}\), le diagramme de structures de Hodge mixtes
\[
\begin{tikzcd}[column sep=small,row sep=huge]
 & \mathrm{H}^n(\mathrm{X}\operatorname{mod}\mathrm{Y},\mathbf{Z})
   \arrow[r] \arrow[d,"p^*"]
 & \mathrm{H}^n(\mathrm{X},\mathbf{Z}) \arrow[r] \arrow[d]
 & \mathrm{H}^n(\mathrm{Y},\mathbf{Z}) \arrow[d]\\
\mathrm{H}^{n-1}(\mathrm{Y}',\mathbf{Z}) \arrow[r]
 & \mathrm{H}^n(\mathrm{X}'\operatorname{mod}\mathrm{Y}',\mathbf{Z})
   \arrow[r]
 & \mathrm{H}^n(\mathrm{X}',\mathbf{Z}) \arrow[r]
 & \mathrm{H}^n(\mathrm{Y}',\mathbf{Z})
\end{tikzcd}
\]

Soit \(\Phi\) la famille de supports sur \(\mathrm{U}\) formée des \(\mathrm{A}\subset\mathrm{U}\) tels que, dans \(\mathrm{X}\) (ou \(\mathrm{X}'\), cela revient au même), \(\overline{\mathrm{A}}\subset\mathrm{U}\). On a
\[
\mathrm{H}^n(\mathrm{X}\operatorname{mod}\mathrm{Y},\mathbf{Z})
=\mathrm{H}^n_\Phi(\mathrm{U},\mathbf{Z})
=\mathrm{H}^n(\mathrm{X}'\operatorname{mod}\mathrm{Y}',\mathbf{Z})
\]
de sorte que la flèche marquée \(p^*\) est un isomorphisme.

Pour \(h^{pq}\) un nombre de Hodge non nul d’une des structures de Hodge mixtes en jeu, on a :

% ---- 045-052 ----
\Page{45}
\noindent\textit{a)} pour \(\mathrm{H}^{n}(\mathrm{X}',\mathbf{Z})\) : \((p,q)\in[n-\mathrm{N},\mathrm{N}]\times[n-\mathrm{N},\mathrm{N}]\)

\noindent\textit{b)} pour \(\mathrm{H}^{n-1}(\mathrm{Y}',\mathbf{Z})\) : \((p,q)\in[n-\mathrm{N},\mathrm{N}-1]\times[n-\mathrm{N},\mathrm{N}-1]\)

\noindent\textit{c)} pour \(\mathrm{H}^{n}(\mathrm{X}'\text{ mod }\mathrm{Y}',\mathbf{Z})=\mathrm{H}^{n}(\mathrm{X}\text{ mod }\mathrm{Y},\mathbf{Z})\) : \((p,q)\in[n-\mathrm{N},\mathrm{N}]\times[n-\mathrm{N},\mathrm{N}]\)

\noindent\textit{d)} pour \(\mathrm{H}^{n}(\mathrm{Y},\mathbf{Z})\) : \((p,q)\in[n-\mathrm{N}+1,\mathrm{N}-1]\times[n-\mathrm{N}+1,\mathrm{N}-1]\)

\noindent\textit{e)} pour \(\mathrm{H}^{n}(\mathrm{X},\mathbf{Z})\) : \((p,q)\in[n-\mathrm{N},\mathrm{N}]\times[n-\mathrm{N},\mathrm{N}]\).

\textit{a)} résulte de ce que \(\mathrm{X}'\) est lisse, et que l’assertion est vraie pour les termes \({}_{\mathrm{W}}\mathrm{E}_{1}\) de la suite spectrale (3.2.4.1). \textit{b)} et \textit{d)} résultent de l’hypothèse de récurrence : on a \(\dim\mathrm{Y},\dim\mathrm{Y}'\leq\mathrm{N}-1\). \textit{c)} résulte de \textit{a)} et \textit{b)} et \textit{e)} résulte de \textit{c)} et \textit{d)}. Ceci achève la démonstration.

\section*{9. EXEMPLES ET APPLICATIONS}

\section*{9.1. Cohomologie des groupes et des classifiants.}

Soit \(\mathrm{G}\) un groupe algébrique linéaire connexe complexe. On se propose de calculer la structure de Hodge mixte de \(\mathrm{H}^{\bullet}(\mathrm{G})\). On commencera par calculer celle de \(\mathrm{H}^{\bullet}(\mathrm{B}_{\bullet\mathrm{G}})\), par réduction au cas des tores; on calculera ensuite celle de \(\mathrm{H}^{\bullet}(\mathrm{G})\) \textit{via} la suite spectrale d’Eilenberg-Moore (6.1.5). En prime, on obtient l’étrange théorème (9.1.7).

\textit{Théorème} \textbf{(9.1.1)}. — \textit{Soit \(\mathrm{G}\) un groupe algébrique linéaire. On a \(\mathrm{H}^{2n+1}(\mathrm{B}_{\bullet\mathrm{G}},\mathbf{Q})=0\), et \(\mathrm{H}^{2n}(\mathrm{B}_{\bullet\mathrm{G}},\mathbf{Q})\otimes\mathbf{C}\) est purement de type \((n,n)\) \textup{(III.0.5)}.}

Soit \(\mathrm{G}^{0}\) la composante neutre de \(\mathrm{G}\) et soit \(\mathrm{T}\subset\mathrm{G}^{0}\) un tore maximal. Puisque \(\mathrm{H}^{\bullet}(\mathrm{B}_{\bullet\mathrm{G}},\mathbf{Q})=\mathrm{H}^{\bullet}(\mathrm{B}_{\bullet\mathrm{G}^{0}},\mathbf{Q})^{\mathrm{G}/\mathrm{G}^{0}}\), on a d’après le \textit{splitting principle} (voir (6.1.6))
\[
\mathrm{H}^{\bullet}(\mathrm{B}_{\bullet\mathrm{G}},\mathbf{Q})
\hookrightarrow
\mathrm{H}^{\bullet}(\mathrm{B}_{\bullet\mathrm{T}},\mathbf{Q}).
\]
Cette inclusion est un morphisme de structures de Hodge mixtes : il suffit de traiter le cas d’un tore, voire même, d’après la formule de Künneth, le cas où \(\mathrm{G}=\mathbf{G}_{m}\). Voici deux démonstrations de (9.1.1) dans ce cas.

\textit{1\textsuperscript{re} démonstration.} — La cohomologie de \(\mathbf{G}_{m}\) (\(=\)la droite projective moins deux points) est : \(\mathrm{H}^{0}=\mathbf{Z}\), \(\mathrm{H}^{1}=\mathbf{Z}\). Les seuls nombres de Hodge \(h^{pq}\) de \(\mathrm{H}^{1}(\mathbf{G}_{m})\) qui pourraient être non nuls sont \(h^{01}=h^{10}\) et \(h^{11}\) ((8.2.4) (i) (iv)). Puisque \(h^{01}+h^{10}+h^{11}=2h^{01}+h^{11}=1\) on a \(h^{01}=0\) et \(\mathrm{H}^{1}(\mathbf{G}_{m})\) est purement de type \((1,1)\) (pour une autre démonstration, voir (10.3.8)). Le terme \(\mathrm{E}_{2}^{p,q}\) de la suite spectrale (6.1.5), nul pour \(p+q\) impair, est donc purement de type \(\left(\dfrac{p+q}{2},\dfrac{p+q}{2}\right)\). Son aboutissement présente la même homogénéité.

\textit{2\textsuperscript{e} démonstration.} — On sait que \(\mathbf{P}^{r}(\mathbf{C})\), muni de \(\mathscr{O}(1)\), est une « approximation » de \(\mathrm{B}_{\mathbf{G}_{m}}\) : l’homomorphisme \([\mathscr{O}(1)]:\mathrm{H}^{i}(\mathrm{B}_{\bullet\mathbf{G}_{m}})\to\mathrm{H}^{i}(\mathbf{P}^{r}(\mathbf{C}))\) est bijectif pour \(i\leq2r\). Le groupe de cohomologie \(\mathrm{H}^{2i}(\mathbf{P}^{r}(\mathbf{C}))\) est engendré par la classe de cohomologie d’un cycle algébrique (à savoir, un sous-espace linéaire), donc est de type \((i,i)\). La proposition suivante permet d’en déduire que \(\mathrm{H}^{2i}(\mathrm{B}_{\mathbf{G}_{m}})\) est de type \((i,i)\).

\Page{46}
\textit{Proposition} \textbf{(9.1.2)}. — \textit{Soient \(\mathrm{G}\) un groupe algébrique (sur \(\mathbf{C}\)) et \(\mathrm{P}\) un espace principal homogène (algébrique) de groupe \(\mathrm{G}\) sur un schéma \(\mathrm{X}\). Le morphisme}
\[
[\mathrm{P}]:\mathrm{H}^{\bullet}(\mathrm{B}_{\mathrm{G}})
=\mathrm{H}^{\bullet}(\mathrm{B}_{\bullet\mathrm{G}})
\to\mathrm{H}^{\bullet}(\mathrm{X})
\]
\textit{est un morphisme de structures de Hodge mixtes.}

Résulte de la définition (6.1.3.1) de \([\mathrm{P}]\) et de la fonctorialité (8.3.4).

\textit{Corollaire} \textbf{(9.1.3)}. — \textit{Les classes caractéristiques, dans \(\mathrm{H}^{2n}(\mathrm{S})\), d’un espace principal homogène (algébrique) \(\mathrm{P}\) sur \(\mathrm{S}\), de groupe structural un groupe linéaire \(\mathrm{G}\), sont purement de type \((n,n)\) : elles proviennent de morphismes \(\mathbf{Q}(-n)\to\mathrm{H}^{2n}(\mathrm{S},\mathbf{Q})\).}

Dans le cas particulier où \(\mathrm{S}\) est quasi-projectif et lisse, on sait que les classes caractéristiques de \(\mathrm{P}\) sont même algébriques, i.e. dans l’image de l’anneau de Chow \((\otimes\mathbf{Q})\).

\textbf{(9.1.4)} Soit \(\mathrm{G}\) un groupe algébrique linéaire sur un corps fini \(\mathbf{F}_{q}\), et soit \(\overline{\mathrm{G}}/\overline{\mathbf{F}}_{q}\) le groupe qui s’en déduit par extension des scalaires à une clôture algébrique de \(\mathbf{F}_{q}\). Une démonstration parallèle à (9.1.1) montre que, en cohomologie \(\ell\)-adique, les valeurs propres du Frobenius (géométrique) agissant sur \(\mathrm{H}^{2n}(\mathrm{B}_{\bullet\overline{\mathrm{G}}},\mathbf{Q}_{\ell})\) sont de la forme \(q^{n}.\zeta\) (\(\zeta\) racine de l’unité). Ceci est en accord avec la philosophie (I.3). De même, l’énoncé (9.1.5) ci-dessous est inspiré des formules donnant les nombres de points des groupes réductifs sur les corps finis.

\textit{Théorème} \textbf{(9.1.5)}. — \textit{Soit \(\mathrm{G}\) un groupe algébrique linéaire connexe. Soit \(\mathrm{P}^{\bullet}(\mathrm{G})\subset\mathrm{H}^{\bullet}(\mathrm{G},\mathbf{Q})\) la partie primitive de l’algèbre de Hopf \(\mathrm{H}^{\bullet}(\mathrm{G},\mathbf{Q})\). C’est une sous-structure de Hodge mixte. On a \(\mathrm{P}^{2i}(\mathrm{G})=0\), et \(\mathrm{P}^{2i-1}(\mathrm{G})\) est purement de type \((i,i)\). Enfin, l’isomorphisme}
\[
\bigwedge\mathrm{P}^{\bullet}(\mathrm{G})
\to\mathrm{H}^{\bullet}(\mathrm{G},\mathbf{Q})
\tag*{\textbf{(9.1.5.1)}}
\]
\textit{est un isomorphisme de structures de Hodge mixtes.}

Dans la suite spectrale (6.1.5), \(\mathrm{E}^{0,q}=0\) pour \(q>0\). Pour \(n>0\), le \textit{edge homomorphism}
\[
\mathrm{H}^{2n}(\mathrm{B}_{\mathrm{G}},\mathbf{Q})
\to\mathrm{E}_{1}^{1,2n}
=\mathrm{H}^{2n-1}(\mathrm{G},\mathbf{Q})
\]
est un morphisme de structures de Hodge mixtes (8.3.5); son image \(\mathrm{P}^{2n-1}(\mathrm{G})\) est donc une sous-structure de Hodge mixte, elle est purement de type \((n,n)\) d’après (9.1.1). L’isomorphisme (9.1.5.1) est donné par le cup-produit, et on conclut par (8.2.11).

\textit{Corollaire} \textbf{(9.1.6)}. — \textit{Soit \(\mathrm{G}\) un groupe algébrique linéaire. Pour \(n>0\), on a}
\[
\mathrm{W}_{n}(\mathrm{H}^{n}(\mathrm{G},\mathbf{Q}))=0.
\]
On se ramène à supposer \(\mathrm{G}\) connexe, et on applique (9.1.5).

\textit{Théorème} \textbf{(9.1.7)}. — \textit{Soit \(\mathrm{X}\) un espace algébrique propre et lisse sur lequel agit un groupe algébrique linéaire connexe \(\mathrm{G}\). Soit \(\rho:\mathrm{G}\times\mathrm{X}\to\mathrm{X}\). Le diagramme}
\Page{47}
\[
\begin{tikzcd}[ampersand replacement=\&, column sep=large, row sep=3.5em]
\mathrm{H}^{\bullet}(\mathrm{X},\mathbf{Q})
  \arrow[r, "\rho^{*}"] \arrow[d, equal]
\&
\mathrm{H}^{\bullet}(\mathrm{X},\mathbf{Q})\otimes
\mathrm{H}^{\bullet}(\mathrm{G},\mathbf{Q})
\\
\mathrm{H}^{\bullet}(\mathrm{X},\mathbf{Q})\otimes_{\mathbf{Q}}\mathbf{Q}
  \arrow[r, equal]
\&
\mathrm{H}^{\bullet}(\mathrm{X},\mathbf{Q})\otimes
\mathrm{H}^{0}(\mathrm{G},\mathbf{Q})
  \arrow[u]
\end{tikzcd}
\]
\textit{est commutatif.}

Puisque l’application composée
\[
\mathrm{H}^{\bullet}(\mathrm{X},\mathbf{Q})
\xrightarrow{\rho^{*}}
\mathrm{H}^{\bullet}(\mathrm{X},\mathbf{Q})\otimes\mathrm{H}^{\bullet}(\mathrm{G},\mathbf{Q})
\to
\mathrm{H}^{\bullet}(\mathrm{X},\mathbf{Q})\otimes\mathrm{H}^{0}(\mathrm{G},\mathbf{Q})
=\mathrm{H}^{\bullet}(\mathrm{X},\mathbf{Q})
\]
est l’identité (le composé \(\mathrm{X}=\{e\}\times\mathrm{X}\to\mathrm{G}\times\mathrm{X}\to\mathrm{X}\) est l’identité), il suffit de prouver que pour \(i>0\), la \(i\)-ème composante de Künneth de \(\rho^{*}\)
\[
\rho_{i}^{*}:\mathrm{H}^{n}(\mathrm{X},\mathbf{Q})
\to
\mathrm{H}^{n-i}(\mathrm{X},\mathbf{Q})\otimes\mathrm{H}^{i}(\mathrm{G},\mathbf{Q})
\]
est nulle. Puisque \(\mathrm{H}^{n}(\mathrm{X},\mathbf{Q})\) est purement de poids \(n\) et que \(\mathrm{H}^{n-i}(\mathrm{X},\mathbf{Q})\) est purement de poids \(n-i\), on a
\[
\operatorname{Im}(\mathrm{H}^{n}(\mathrm{X},\mathbf{Q}))
\subset
\mathrm{W}_{n}(\mathrm{H}^{n-i}(\mathrm{X},\mathbf{Q})\otimes\mathrm{H}^{i}(\mathrm{G},\mathbf{Q}))
=
\mathrm{H}^{n-i}(\mathrm{X},\mathbf{Q})\otimes\mathrm{W}_{i}\mathrm{H}^{i}(\mathrm{G},\mathbf{Q}).
\]
D’après (9.1.6), cette image est donc nulle.

\textit{Variante} \textbf{(9.1.8)}. — \textit{Soit \(\mathrm{X}\) un espace algébrique séparé sur lequel agit un groupe algébrique linéaire connexe \(\mathrm{G}\).}

(i) \textit{Si \(\mathrm{X}\) est lisse et que \(\overline{\mathrm{X}}\) est une compactification de \(\mathrm{X}\), le diagramme}
\[
\begin{tikzcd}[ampersand replacement=\&, column sep=1.8em, row sep=3.5em]
\mathrm{H}^{\bullet}(\overline{\mathrm{X}},\mathbf{Q})
  \arrow[r] \arrow[d]
\&
\mathrm{H}^{\bullet}(\mathrm{X},\mathbf{Q})
  \arrow[r, "\rho^{*}"]
\&
\mathrm{H}^{\bullet}(\mathrm{X},\mathbf{Q})\otimes
\mathrm{H}^{\bullet}(\mathrm{G},\mathbf{Q})
\\
\mathrm{H}^{\bullet}(\mathrm{X},\mathbf{Q})
  \arrow[rr, equal]
\& \&
\mathrm{H}^{\bullet}(\mathrm{X},\mathbf{Q})\otimes
\mathrm{H}^{0}(\mathrm{G},\mathbf{Q})
  \arrow[u, hook']
\end{tikzcd}
\]
\textit{est commutatif.}

(ii) \textit{Si \(\mathrm{X}\) est propre et que \(\widetilde{\mathrm{X}}\) est une résolution des singularités de \(\mathrm{X}\), le diagramme}
\[
\begin{tikzcd}[ampersand replacement=\&, column sep=1.5em, row sep=3.5em]
\mathrm{H}^{\bullet}(\mathrm{X},\mathbf{Q})
  \arrow[r, "\rho^{*}"] \arrow[d]
\&
\mathrm{H}^{\bullet}(\mathrm{X},\mathbf{Q})\otimes
\mathrm{H}^{\bullet}(\mathrm{G},\mathbf{Q})
  \arrow[r]
\&
\mathrm{H}^{\bullet}(\widetilde{\mathrm{X}},\mathbf{Q})\otimes
\mathrm{H}^{\bullet}(\mathrm{G},\mathbf{Q})
\\
\mathrm{H}^{\bullet}(\widetilde{\mathrm{X}},\mathbf{Q})
  \arrow[rr, equal]
\& \&
\mathrm{H}^{\bullet}(\widetilde{\mathrm{X}},\mathbf{Q})\otimes
\mathrm{H}^{0}(\mathrm{G},\mathbf{Q})
  \arrow[u, hook']
\end{tikzcd}
\]
\textit{est commutatif.}

\Page{48}
Il suffit de prouver que pour \(i>0\), les composantes de Künneth
\[
\begin{aligned}
&\mathrm{H}^{n}(\overline{\mathrm{X}},\mathbf{Q})
\to\mathrm{H}^{n}(\mathrm{X},\mathbf{Q})
\to\mathrm{H}^{n-i}(\mathrm{X},\mathbf{Q})\otimes\mathrm{H}^{i}(\mathrm{G},\mathbf{Q})
\\
&\mathrm{H}^{n}(\mathrm{X},\mathbf{Q})
\to\mathrm{H}^{n-i}(\mathrm{X},\mathbf{Q})\otimes\mathrm{H}^{i}(\mathrm{G},\mathbf{Q})
\to\mathrm{H}^{n-i}(\widetilde{\mathrm{X}},\mathbf{Q})\otimes\mathrm{H}^{i}(\mathrm{G},\mathbf{Q})
\end{aligned}
\]
sont respectivement nulles. Dans le cas (i) (resp. (ii)), on a (8.2.4)
\[
\mathrm{W}_{n}\mathrm{H}^{n}(\overline{\mathrm{X}},\mathbf{Q})
=\mathrm{H}^{n}(\overline{\mathrm{X}},\mathbf{Q})
\quad
(\text{resp. }\mathrm{W}_{n}\mathrm{H}^{n}(\mathrm{X},\mathbf{Q})
=\mathrm{H}^{n}(\mathrm{X},\mathbf{Q})),
\]
tandis que, par (8.2.4) et (9.1.6)
\[
\mathrm{W}_{n}(\mathrm{H}^{n-i}(\mathrm{X},\mathbf{Q})\otimes\mathrm{H}^{i}(\mathrm{G},\mathbf{Q}))=0
\quad
(\text{resp. }\mathrm{W}_{n}(\mathrm{H}^{n-i}(\widetilde{\mathrm{X}},\mathbf{Q})\otimes\mathrm{H}^{i}(\mathrm{G},\mathbf{Q}))=0)
\]
d’où l’assertion.

\section*{9.2. Théorie de Hodge des hypersurfaces lisses, d’après Griffiths.}

Ce numéro ne dépend que des §§ 1 à 3. On y explique le théorème (8.6) de [7].

\textbf{(9.2.1)} Soient \(\mathrm{X}\) une variété algébrique propre et lisse, et \(\mathrm{Y}\) un diviseur lisse dans \(\mathrm{X}\). Soit \(j\) l’inclusion de \(\mathrm{U}=\mathrm{X}-\mathrm{Y}\) dans \(\mathrm{X}\). On a \(\mathrm{R}^{i}j_{*}\mathbf{Z}=0\) pour \(i\neq0,\ 1\) et la suite spectrale de Leray pour \(j\) s’identifie à la suite exacte longue de cohomologie
\begin{align*}
\cdots\to\mathrm{H}_{\mathrm{Y}}^{n}(\mathrm{X})
&\to\mathrm{H}^{n}(\mathrm{X})
\to\mathrm{H}^{n}(\mathrm{U})
\xrightarrow{\partial}\cdots
\tag*{\textbf{(9.2.1.1)}}\\
\mathrm{H}_{\mathrm{Y}}^{n}(\mathrm{X})
&=\mathrm{H}^{n-2}(\mathrm{Y},\mathscr{H}_{\mathrm{Y}}^{2}).
\end{align*}
Pour une raison rendue claire en (3.1.9), nous identifierons \(\mathscr{H}_{\mathrm{Y}}^{2}\) au faisceau constant \(\mathbf{Z}(-1)_{\mathrm{Y}}\). La suite spectrale de Leray de \(j\) convergeant par définition vers la structure de Hodge mixte de \(\mathrm{H}^{n}(\mathrm{U})\), on trouve que la suite exacte
\[
\cdots\to
\mathrm{H}^{n-2}(\mathrm{Y},\mathbf{Z})(-1)
\to\mathrm{H}^{n}(\mathrm{X},\mathbf{Z})
\to\mathrm{H}^{n}(\mathrm{U},\mathbf{Z})
\xrightarrow{\partial}\cdots
\tag*{\textbf{(9.2.1.2)}}
\]
est une suite exacte de structures de Hodge mixtes.

La suite spectrale de Leray de \(j\) est aussi la suite spectrale du complexe filtré \((\Omega_{\mathrm{X}}^{\bullet}(\log\mathrm{Y}),\mathrm{W})\); la suite exacte précédente s’identifie donc encore à celle déduite de la suite exacte courte
\[
0\to\Omega_{\mathrm{X}}^{\bullet}
\to\Omega_{\mathrm{X}}^{\bullet}(\log\mathrm{Y})
\xrightarrow{\operatorname{Res}}
\Omega_{\mathrm{Y}}^{\bullet}[-1]\to0.
\tag*{\textbf{(9.2.1.3)}}
\]
Pour cette raison, l’homomorphisme \(\partial\) dans (9.2.1.2) porte encore le nom d’\textit{opérateur résidu}.

\textbf{(9.2.2)} L’opérateur résidu
\[
\operatorname{Res}:\mathrm{H}^{n+1}(\mathrm{U},\mathbf{Z})
\to\mathrm{H}^{n}(\mathrm{Y},\mathbf{Z})(-1)
\tag*{\textbf{(9.2.2.1)}}
\]
est strictement compatible à la filtration de Hodge. La filtration de Hodge de \(\mathrm{H}^{n+1}(\mathrm{U},\mathbf{C})\) détermine donc celle de \(\operatorname{Ker}(\mathrm{H}^{n}(\mathrm{Y},\mathbf{C})\to\mathrm{H}^{n+2}(\mathrm{X},\mathbf{C}))\).

\Page{49}
Supposons \(\mathrm{Y}\) ample. Si \(\mathrm{X}\) est purement de dimension \(\mathrm{N}+1\), on a alors \(\mathrm{H}^{i}(\mathrm{X},\mathbf{Z})\xrightarrow{\sim}\mathrm{H}^{i}(\mathrm{Y},\mathbf{Z})\) pour \(i<\mathrm{N}\) (Lefschetz). Par dualité de Poincaré, pour \(i\neq\mathrm{N}\), \(\mathrm{H}^{i}(\mathrm{Y},\mathbf{C})\) se calcule donc à partir de \(\mathrm{H}^{\bullet}(\mathrm{X},\mathbf{C})\). Posons
\[
\begin{aligned}
\mathrm{H}_{\text{év}}^{\mathrm{N}}(\mathrm{Y},\mathbf{C})
&=\operatorname{Ker}(\mathrm{H}^{\mathrm{N}}(\mathrm{Y},\mathbf{C})\to\mathrm{H}^{\mathrm{N}+2}(\mathrm{X},\mathbf{C}))\\
&=\text{orthogonal de }\mathrm{H}^{\mathrm{N}}(\mathrm{X},\mathbf{C})
\hookrightarrow\mathrm{H}^{\mathrm{N}}(\mathrm{Y},\mathbf{C})
\end{aligned}
\]
(partie évanescente de la cohomologie). On a
\[
\mathrm{H}^{\mathrm{N}}(\mathrm{Y},\mathbf{C})
=\mathrm{H}_{\text{év}}^{\mathrm{N}}(\mathrm{Y},\mathbf{C})
\oplus\mathrm{H}^{\mathrm{N}}(\mathrm{X},\mathbf{C})
\]
et on obtient donc beaucoup d’information sur la structure de Hodge de \(\mathrm{H}^{\bullet}(\mathrm{Y},\mathbf{C})\) quand on connaît celle de \(\mathrm{H}^{\bullet}(\mathrm{X},\mathbf{C})\) et la filtration de Hodge de \(\mathrm{H}^{\mathrm{N}+1}(\mathrm{U},\mathbf{C})\).

\textbf{(9.2.3)} Par définition, la filtration de Hodge de \(\mathrm{H}^{\bullet}(\mathrm{U},\mathbf{C})\) est l’aboutissement de la suite spectrale \textit{dégénérée} en \(\mathrm{E}_{1}\)
\[
\mathrm{E}_{1}^{pq}
=\mathrm{H}^{q}(\mathrm{X},\Omega_{\mathrm{X}}^{p}(\log\mathrm{Y}))
\Rightarrow\mathrm{H}^{p+q}(\mathrm{U},\mathbf{C}).
\]
D’après (3.1.11), cette suite spectrale coïncide avec la suite spectrale du complexe \(j_{*}\Omega_{\mathrm{U}}^{\bullet}\), muni de la filtration par l’ordre du pôle \(\mathrm{P}\).

En particulier, on a
\[
\mathrm{F}^{p}(\mathrm{H}^{\bullet}(\mathrm{U},\mathbf{C}))
=\mathbf{H}^{\bullet}(\mathrm{X},\mathrm{P}^{p}(j_{*}\Omega_{\mathrm{U}}^{\bullet})).
\]
Rappelons que \(\mathrm{P}^{p}(j_{*}\Omega_{\mathrm{U}}^{\bullet})\) est le complexe suivant, de première composante non nulle en degré \(p\)
\[
\Omega_{\mathrm{X}}^{p}(\mathrm{Y})
\to\Omega_{\mathrm{X}}^{p+1}(2\mathrm{Y})
\to\cdots\to
\Omega_{\mathrm{X}}^{i}((i-p+1)\mathrm{Y})
\to\cdots
\]
Si \(\mathrm{H}^{i}(\mathrm{X},\Omega_{\mathrm{X}}^{p}(n\mathrm{Y}))=0\) pour \(i>0\), \(n>0\) et \(p\geq0\), on a simplement
\[
\mathbf{H}^{\bullet}(\mathrm{X},\mathrm{P}^{p}(j_{*}\Omega_{\mathrm{U}}^{\bullet}))
=\mathrm{H}^{\bullet}(\Gamma(\mathrm{X},\mathrm{P}^{p}(j_{*}\Omega_{\mathrm{U}}^{\bullet}))),
\]
d’où le résultat suivant :

\textit{Proposition} \textbf{(9.2.4)}. — \textit{Supposons que \(\mathrm{H}^{i}(\mathrm{X},\Omega_{\mathrm{X}}^{p}(n\mathrm{Y}))=0\) pour \(i>0\), \(n>0\) et \(p\geq0\) (tel est le cas si \(\mathrm{Y}\) est suffisamment ample). Alors une classe de cohomologie \(c\in\mathrm{H}^{d}(\mathrm{U},\mathbf{C})\) est de filtration de Hodge \(\geq p\) si et seulement si elle peut se représenter par une \(d\)-forme fermée \(\alpha\) sur \(\mathrm{U}\), avec un pôle d’ordre \(\leq d-p+1\) le long de \(\mathrm{Y}\). Si \(\alpha\) est une forme fermée sur \(\mathrm{U}\), avec un pôle d’ordre \(k\) le long de \(\mathrm{Y}\), et si la classe de \(\alpha\) dans \(\mathrm{H}^{d}(\mathrm{U},\mathbf{C})\) est nulle, alors \(\alpha=d\beta\), avec \(\beta\) présentant un pôle d’ordre \(\leq k-1\) le long de \(\mathrm{Y}\); pour \(k\leq1\), on a même \(\alpha=0\).}

\textbf{(9.2.5)} Il est bien connu (théorème de Bott) que les hypothèses de (9.2.4) sont vérifiées pour \(\mathrm{Y}\) une hypersurface dans l’espace projectif \(\mathbf{P}^{n}(\mathbf{C})\) \((n\geq1)\), i.e. que
\[
\mathrm{H}^{i}(\mathbf{P}^{n}(\mathbf{C}),\Omega^{j}(m))=0
\qquad\text{pour}\qquad i>0,\ m>0
\]
(une démonstration figure par exemple dans SGA 7 XI). De plus, pour \(n>1\), l’application résidu identifie \(\mathrm{H}^{n}(\mathrm{U},\mathbf{C})\) à la partie primitive de la cohomologie de \(\mathrm{H}^{n-1}(\mathrm{Y},\mathbf{C})\). Dès lors ([7], (8.6)) :

\Page{50}
\textit{Proposition} \textbf{(9.2.6)}. — \textit{Soit \(\mathrm{Y}\) une hypersurface lisse de \(\mathbf{P}^{n}(\mathbf{C})\) \((n>1)\) et \(\mathrm{U}=\mathbf{P}^{n}(\mathbf{C})-\mathrm{Y}\).}

(i) \textit{Pour qu’une classe de cohomologie primitive \(c\in\mathrm{H}^{n-1}(\mathrm{Y},\mathbf{C})\) soit de filtration de Hodge \(\geq p\), il faut et il suffit qu’elle soit le résidu d’une \(n\)-forme différentielle régulière sur \(\mathrm{U}\), avec un pôle d’ordre \(\leq n-p\) le long de \(\mathrm{Y}\).}

(ii) \textit{Soit \(\alpha\) une \(n\)-forme différentielle régulière sur \(\mathrm{U}\), avec un pôle d’ordre \(k\geq1\) le long de \(\mathrm{Y}\). Pour que \(\operatorname{Res}(\alpha)\in\mathrm{H}^{n-1}(\mathrm{Y},\mathbf{C})\) soit nul, il faut et il suffit que \(\alpha=d\beta\), pour une \((n-1)\)-forme sur \(\mathrm{U}\), avec un pôle d’ordre \(\leq k-1\) le long de \(\mathrm{Y}\).}

\section*{9.3. Construction de complexes d’opérateurs différentiels du 1\textsuperscript{er} ordre.}

Ce numéro ne dépend que des §§ 5 et 6. On y prouve le théorème suivant.

\textit{Proposition} \textbf{(9.3.1)}. — \textit{Soit \(\mathrm{X}\) un schéma quasi-projectif (sur \(\mathbf{C}\)). Il existe sur \(\mathrm{X}\) un complexe de faisceaux \(\mathrm{K}\) ayant les propriétés suivantes.}

\textit{a) Les \(\mathrm{K}^{n}\) sont des faisceaux algébriques cohérents et \(\mathrm{K}^{n}=0\) pour \(n<0\).}

\textit{b) Les différentielles \(d:\mathrm{K}^{n}\to\mathrm{K}^{n+1}\) sont des opérateurs différentiels algébriques du 1\textsuperscript{er} ordre.}

\textit{c) Le complexe de faisceaux analytiques cohérents \(\mathrm{K}^{\mathrm{an}}\) sur \(\mathrm{X}^{\mathrm{an}}\) est une résolution du faisceau constant \(\underline{\mathbf{C}}\). Plus précisément, \(\mathscr{H}^{i}(\mathrm{K})=0\) pour \(i>0\), et il existe}
\[
f:\Omega_{\mathrm{X}}^{\bullet}\to\mathrm{K}
\]
\textit{avec \(f^{n}\) algébrique linéaire, tel que l’application composée \(\underline{\mathbf{C}}\to\mathscr{O}_{\mathrm{X}}\xrightarrow{f^{0}}\mathrm{K}^{0}\) identifie \(\underline{\mathbf{C}}\) à \(\mathscr{H}^{0}(\mathrm{K})\).}

\textit{d) L’application naturelle}
\[
\mathbf{H}^{i}(\mathrm{X},\mathrm{K})
\to\mathbf{H}^{i}(\mathrm{X}^{\mathrm{an}},\mathrm{K}^{\mathrm{an}})
\xleftarrow[(c)]{\sim}
\mathbf{H}^{i}(\mathrm{X}^{\mathrm{an}},\mathbf{C})
\]
\textit{est un isomorphisme.}

Comme corollaire, on retrouve un cas particulier d’un théorème de Bloom et Herrera [2].

\textit{Corollaire} \textbf{(9.3.2)}. — \textit{L’application déduite de \(\underline{\mathbf{C}}\to\Omega_{\mathrm{X}}^{\bullet\mathrm{an}}\) :}
\[
\mathrm{H}^{\bullet}(\mathrm{X}^{\mathrm{an}},\mathbf{C})
\to\mathbf{H}^{\bullet}(\mathrm{X}^{\mathrm{an}},\Omega_{\mathrm{X}}^{\bullet\mathrm{an}})
\]
\textit{identifie \(\mathrm{H}^{\bullet}(\mathrm{X}^{\mathrm{an}},\mathbf{C})\) à un facteur direct de \(\mathbf{H}^{\bullet}(\mathrm{X}^{\mathrm{an}},\Omega_{\mathrm{X}}^{\bullet\mathrm{an}})\).}

Nous prouverons d’abord le corollaire. La proposition s’obtiendra en reprenant la démonstration au niveau des complexes. Soit \(\varepsilon:\mathrm{X}_{\bullet}\to\mathrm{X}\) un schéma simplicial augmenté vers \(\mathrm{X}\). On suppose que \(\varepsilon\) est un hyperrecouvrement propre et que les \(\mathrm{X}_{n}\) sont lisses. On dispose d’un diagramme commutatif
\[
\begin{tikzcd}[ampersand replacement=\&, column sep=large, row sep=3em]
\mathrm{H}^{\bullet}(\mathrm{X}^{\mathrm{an}},\mathbf{C})
  \arrow[r, "\sim", "\text{\textcircled{1}}"']
  \arrow[d]
\&
\mathrm{H}^{\bullet}(\mathrm{X}_{\bullet}^{\mathrm{an}},\mathbf{C})
  \arrow[d, "\wr", "\text{\textcircled{2}}"']
\\
\mathbf{H}^{\bullet}(\mathrm{X}^{\mathrm{an}},\Omega_{\mathrm{X}}^{\bullet})
  \arrow[r]
\&
\mathbf{H}^{\bullet}(\mathrm{X}_{\bullet}^{\mathrm{an}},\Omega_{\mathrm{X}_{\bullet}}^{\bullet})
\end{tikzcd}
\tag*{\textbf{(9.3.2.1)}}
\]

\Page{51}
La flèche 1 est bijective d’après (5.3.5) (II) (IV). La flèche 2 est bijective, car \(\Omega_{\mathrm{X}_{\bullet}}^{\bullet}\) est une résolution de \(\mathbf{C}\) (lemme de Poincaré sur les \(\mathrm{X}_{n}\)). Le diagramme (9.3.2.1) fournit une rétraction, compatible au cup-produit, de \(\mathbf{H}^{\bullet}(\mathrm{X}^{\mathrm{an}},\Omega_{\mathrm{X}}^{\bullet})\) sur \(\mathrm{H}^{\bullet}(\mathrm{X}^{\mathrm{an}},\mathbf{C})\), et ceci prouve (9.3.2).

\textbf{(9.3.3)} Le complexe \(\mathrm{K}\) que nous voulons construire ne sera autre, dans la catégorie dérivée, que \(\mathbf{s}\mathrm{R}\varepsilon_{\bullet *}\Omega_{\mathrm{X}_{\bullet}}^{\bullet}\). Nous ne traiterons que le cas où l’espace simplicial augmenté \(\mathrm{X}_{\bullet}\) est \(s\)-scindé, défini comme en (6.2.5) par des applications propres et surjectives, de sources quasi-projectives et lisses
\[
f'_{k}:\mathrm{N}_{k}\to
(\operatorname{cosq}\operatorname{sq}_{k-1}(\mathrm{X}_{\bullet}))_{k}
\qquad(k\geq0).
\]
Nous allons construire un « système compatible » de recouvrements ouverts affines des \(\mathrm{X}_{n}\). Il y a intérêt, pour cette construction, à regarder un recouvrement ouvert d’un espace \(\mathrm{Y}\) comme un diagramme
\[
\begin{tikzcd}[ampersand replacement=\&, column sep=1.4em]
\mathrm{Y}
\&
\mathrm{U}\arrow[l, "u"']\arrow[r, "\varphi"]
\&
\mathrm{I}
\end{tikzcd}
\]
avec \(\mathrm{I}\) discret, \(u\) surjectif et \(u_{i}:\mathrm{U}_{i}\underset{\mathrm{dfn}}{=}\varphi^{-1}(i)\hookrightarrow\mathrm{Y}\) un plongement ouvert. Procédant comme en (6.2.5), on construit un diagramme
\[
\begin{tikzcd}[ampersand replacement=\&, column sep=1.4em]
\mathrm{X}_{\bullet}
\&
\mathrm{U}_{\bullet}\arrow[l, "u"']\arrow[r, "\varphi"]
\&
\mathrm{I}_{\bullet}
\end{tikzcd}
\tag*{\textbf{(9.3.3.1)}}
\]
du type suivant.

\textit{a)} \(\mathrm{I}_{\bullet}\) est un ensemble simplicial discret. Les \(\mathrm{I}_{n}\) sont finis.

\textit{b)} Pour tout \(k\geq0\), (9.3.3.1) induit un diagramme
\[
\begin{tikzcd}[ampersand replacement=\&, column sep=1.2em, row sep=3em]
(\operatorname{cosq}\operatorname{sq}_{k-1}(\mathrm{X}_{\bullet}))_{k}
\&
(\operatorname{cosq}\operatorname{sq}_{k-1}(\mathrm{U}_{\bullet}))_{k}
  \arrow[l]\arrow[r, "\varphi"]
\&
(\operatorname{cosq}\operatorname{sq}_{k-1}(\mathrm{I}_{\bullet}))_{k}
\\
\mathrm{N}_{k}\arrow[u, "f'_{k}"']
\&
\mathrm{N}_{k}(\mathrm{U}_{\bullet})
  \arrow[l]\arrow[u]\arrow[r, "\varphi"]
\&
\mathrm{N}_{k}(\mathrm{I}_{\bullet})\arrow[u, "a'_{k}"']
\end{tikzcd}
\]

\textit{c)} Dans le diagramme \textit{b)}, pour \(i\in\operatorname{cosq}\operatorname{sq}_{k-1}(\mathrm{I}_{\bullet})_{k}\), les \(\varphi^{-1}(j)\) pour \(a'_{k}(j)=i\) forment un recouvrement ouvert affine de
\[
\mathrm{N}_{k}
\times_{(\operatorname{cosq}\operatorname{sq}_{k-1}(\mathrm{X}_{\bullet}))_{k}}
\varphi^{-1}(i).
\]

\textit{d)} (résulte par récurrence de \textit{c)}). La première (et donc la deuxième) ligne du diagramme \textit{b)} définit un recouvrement ouvert. Les \(\mathrm{X}_{k}\leftarrow\mathrm{U}_{k}\to\mathrm{I}_{k}\) forment donc aussi des recouvrements ouverts.

\textit{e)} (résulte d’un choix convenable des recouvrements affines en \textit{c)}). \(\mathrm{U}_{n_i}=\varphi_{n}^{-1}(i)\) est le complément d’un sous-schéma \(\mathrm{D}_{n,i}\) de \(\mathrm{X}_{n}\), réunion de quelques composantes connexes de \(\mathrm{X}_{n}\) et d’un diviseur. Pour tout faisceau cohérent \(\mathscr{F}\) sur \(\mathrm{X}_{n}\) et des entiers \((a_{i})_{i\in\mathrm{I}_{n}}\) \((a_{i}\geq0)\), on pose
\[
\mathscr{F}\Bigl(\sum_{i}a_{i}\mathrm{D}_{n,i}\Bigr)
=
\begin{cases}
0\text{ sur les composantes de }\mathrm{X}_{n}
 \text{ contenues dans un }\mathrm{D}_{n,i}\text{ avec }a_{i}>0\\[3pt]
\displaystyle
\mathscr{F}\otimes\mathscr{O}\Bigl(\sum_{a_{i}>0}a_{i}\mathrm{D}_{n,i}\Bigr)
 \text{ sur les autres.}
\end{cases}
\]

\Page{52}
Dans la construction inductive de (9.3.3.1), les recouvrements ouverts (\textit{c}) sont choisis arbitrairement. Ceci permet de supposer

\textit{f)} Quels que soient les \(a_{i}>0\) \((i\in\mathrm{I}_{k})\)
\[
\mathrm{R}^{\ell}\varepsilon_{k*}
\Bigl(\Omega_{\mathrm{X}_{k}}^{m}\Bigl(\sum_{i}a_{i}\mathrm{D}_{k,i}\Bigr)\Bigr)=0
\qquad\text{pour}\qquad\ell>0.
\]

\textit{g)} Quel que soit \(\alpha:\Delta_{m}\to\Delta_{n}\) et \(i\in\mathrm{I}_{n}\), on a
\[
\mathrm{X}_{\bullet}(\alpha)^{-1}(\mathrm{D}_{m,\alpha(i)})
\subset\mathrm{D}_{n,i}
\qquad\text{(comme schéma).}
\]
Pour chaque \(k\), le complexe de faisceaux sur \(\mathrm{X}\)
\[
\mathrm{R}\varepsilon_{k*}(\Omega_{\mathrm{X}_{k}}^{\bullet})
\]
peut se calculer de façon čechiste. Si, pour \(\mathrm{P}\subset\mathrm{I}_{k}\), \(j_{\mathrm{P}}\) est l’inclusion de \(\mathrm{U}_{\mathrm{P}}=\bigcap_{i\in\mathrm{P}}\mathrm{U}_{k_i}\) dans \(\mathrm{X}_{k}\), c’est le complexe simple associé au double complexe des
\[
\check{\mathrm{K}}^{k,p,q}
=\sum_{\#\mathrm{P}=p+1}
(\varepsilon_{k}j_{\mathrm{P}})_{*}j_{\mathrm{P}}^{*}\Omega_{\mathrm{X}_{k}}^{q}.
\]
Ce complexe double admet comme sous-complexe double le complexe de composantes les faisceaux cohérents
\[
\mathrm{K}^{k,p,q}
=\sum_{\#\mathrm{P}=p+1}
\varepsilon_{k*}\Omega_{\mathrm{X}_{k}}^{q}
\Bigl((q+1)\sum_{i\in\mathrm{P}}\mathrm{D}_{k,i}\Bigr).
\]
Pour \(k\) et \(q\) fixes, il résulte de \textit{f)} que l’inclusion de \(\mathrm{K}^{k,\bullet,q}\) dans \(\check{\mathrm{K}}^{k,\bullet,q}\) est un quasi-isomorphisme. Pour \(k\) fixe, l’inclusion de \(\mathrm{K}^{k,\bullet,\bullet}\) dans \(\check{\mathrm{K}}^{k,\bullet,\bullet}\) elle aussi est un quasi-isomorphisme. Enfin, pour \(k\) variable, les \(\mathrm{K}^{k,\bullet,\bullet}\) forment, d’après \textit{g)}, un complexe double simplicial. Le complexe simple associé \(\mathrm{K}\) coïncide avec \(\mathrm{R}\varepsilon_{\bullet *}\Omega_{\mathrm{X}_{\bullet}}^{\bullet}\), dans la catégorie dérivée.

Prouvons que \(\mathrm{K}\) vérifie (9.3.1). Les propriétés \textit{a)} et \textit{b)} de \(\mathrm{K}\) sont claires; on définit \(f:\Omega_{\mathrm{X}}^{\bullet}\to\mathrm{K}\) par les applications évidentes des \(\Omega_{\mathrm{X}}^{q}\) dans les \(\mathrm{K}^{0,0,q}\).

Dans la catégorie analytique, le complexe de faisceaux analytiques cohérents \(\mathrm{K}^{\mathrm{an}}\) est donné par les mêmes formules que \(\mathrm{K}\). On a un diagramme commutatif
\[
\begin{tikzcd}[ampersand replacement=\&, column sep=large, row sep=3em]
\mathbf{H}^{i}(\mathrm{X},\mathrm{K})
  \arrow[r, dash, "\sim"]\arrow[d]
\&
\mathbf{H}^{i}(\mathrm{X}_{\bullet},\Omega_{\mathrm{X}_{\bullet}}^{\bullet})
  \arrow[d, "\text{\textcircled{1}}"]
\\
\mathbf{H}^{i}(\mathrm{X}^{\mathrm{an}},\mathrm{K}^{\mathrm{an}})
  \arrow[r, dash, "\sim"]
\&
\mathbf{H}^{i}(\mathrm{X}_{\bullet}^{\mathrm{an}},
              \Omega_{\mathrm{X}_{\bullet}^{\mathrm{an}}}^{\bullet})
\end{tikzcd}
\]
et 1 est l’aboutissement d’un morphisme de suites spectrales
\[
\begin{tikzcd}[ampersand replacement=\&, column sep=large, row sep=3em]
\mathbf{H}^{a}(\mathrm{X}_{b},\Omega_{\mathrm{X}_{\bullet}}^{\bullet})
  \arrow[r, Rightarrow]\arrow[d]
\&
\mathbf{H}^{a+b}(\mathrm{X}_{\bullet},\Omega_{\mathrm{X}_{\bullet}}^{\bullet})
  \arrow[d]
\\
\mathbf{H}^{a}(\mathrm{X}_{b}^{\mathrm{an}},
              \Omega_{\mathrm{X}_{\bullet}^{\mathrm{an}}}^{\bullet})
  \arrow[r, Rightarrow]
\&
\mathbf{H}^{a+b}(\mathrm{X}_{\bullet}^{\mathrm{an}},
                \Omega_{\mathrm{X}_{\bullet}^{\mathrm{an}}}^{\bullet}).
\end{tikzcd}
\]
Ce morphisme est un isomorphisme d’après [9], d’où \textit{d)}.

% ---- 053-060 ----
\Page{53}
Pour prouver \textit{c)}, on considère le diagramme commutatif
\[
\begin{tikzcd}[column sep=2.4em, row sep=3.5em]
\underline{\mathbf{C}}
  \arrow[rr, dash, "\approx", "{\text{\textcircled{1}}}"'{pos=.2}]
  \arrow[d]
& &
\mathrm{R}\varepsilon_{\bullet *}\underline{\mathbf{C}}
  \arrow[d, "\approx", "{\text{\textcircled{2}}}"'{pos=.2}]
\\
\Omega_{\mathrm{X}_{\bullet}}^{\bullet\mathrm{an}}
  \arrow[r, "f"]
&
\mathrm{K}^{\mathrm{an}}
  \arrow[r, dash, "\approx"]
&
\mathrm{R}\varepsilon_{\bullet *}
  \Omega_{\mathrm{X}_{\bullet}}^{\bullet\mathrm{an}}
\end{tikzcd}
\]
Que 1 soit un quasi-isomorphisme résulte de (5.3.5).\par
Que 2 en soit un résulte du lemme de Poincaré sur les \(\mathrm{X}_n\).
L’assertion \textit{c)} en résulte.

\textbf{(9.3.4)} La même construction fournit des résolutions vérifiant
\textit{a)}, \textit{b)} pour d’autres faisceaux \(\mathscr{F}\) sur
\(\mathrm{X}^{\mathrm{an}}\) que le faisceau constant \(\mathbf{C}\)
(on peut sans doute prendre pour \(\mathscr{F}\) n’importe quel faisceau
« algébriquement constructible »). Malheureusement, je ne puis pas
prouver que les résolutions obtenues sont essentiellement uniques
(dans une catégorie dérivée convenable). En l’absence d’une telle
unicité, (9.3.1) est sans grand intérêt.

\section*{10. THÉORIE DE HODGE EN NIVEAU \(\leq\mathbf{1}\)}

\section*{10.1. 1-motifs}

\textit{Définition} \textbf{(10.1.1)}. — (i)
\textit{Une structure de Hodge mixte \(\mathrm{H}\) est
\textup{sans torsion} si \(\mathrm{H}_{\mathbf{Z}}\) est sans torsion.}

(ii) \textit{Si \(\mathrm{H}\) est sans torsion, on désigne encore par
\(\mathrm{W}\) la filtration de
\(\mathrm{H}_{\mathbf{Z}}\subset\mathrm{H}_{\mathbf{Q}}\) induite par
\(\mathrm{W}\). On désigne par
\(\operatorname{Gr}^{\mathrm{W}}_n(\mathrm{H})\) le
\(\mathbf{Z}\)-module sans torsion
\(\operatorname{Gr}^{\mathrm{W}}_n(\mathrm{H}_{\mathbf{Z}})\), muni de
sa structure de Hodge de poids \(n\).}

\textit{Définition} \textbf{(10.1.2)}. —
\textit{Un \textup{1-motif} sur un corps algébriquement clos \(k\)
consiste en :}

\noindent a) \textit{un \(\mathbf{Z}\)-module libre de type fini
\(\mathrm{X}\), une variété abélienne \(\mathrm{A}\) et un tore
\(\mathrm{T}\);}\par
\noindent b) \textit{une extension \(\mathrm{G}\) de \(\mathrm{A}\)
par \(\mathrm{T}\);}\par
\noindent c) \textit{un homomorphisme
\(u:\mathrm{X}\to\mathrm{G}(k)\).}

Il est souvent utile de regarder un 1-motif \(\mathrm{M}\) comme étant
un complexe de schémas en groupes, concentré en degrés 0 et 1
\[
\mathrm{M}=[\mathrm{X}\xrightarrow{u}\mathrm{G}].
\]

\textit{Construction} \textbf{(10.1.3)}. —
\textit{On construira une équivalence
\(\mathrm{M}\mapsto\mathrm{T}(\mathrm{M})\) de la catégorie des
1-motifs sur \(\mathbf{C}\) avec la catégorie des structures de Hodge
mixtes sans torsion \(\mathrm{H}\) de type}
\[
\{(0,0),(0,-1),(-1,0),(-1,-1)\}
\]
\Page{54}
\textit{telles que \(\operatorname{Gr}^{\mathrm{W}}_{-1}(\mathrm{H})\)
soit polarisable. De plus, on construira des isomorphismes naturels en
\(\mathrm{M}=(\mathrm{X},\mathrm{A},\mathrm{T},\mathrm{G},u)\) :}
\[
\begin{aligned}
\mathrm{H}_1(\mathrm{T},\mathbf{Z})
  &\simeq \operatorname{Gr}^{\mathrm{W}}_{-2}
          \mathrm{T}(\mathrm{M})_{\mathbf{Z}}
\\
\mathrm{H}_1(\mathrm{A},\mathbf{Z})
  &\simeq \operatorname{Gr}^{\mathrm{W}}_{-1}\mathrm{T}(\mathrm{M})
  \qquad \text{\textit{(isomorphisme de structures de Hodge)}}
\\
\mathrm{X}
  &\simeq \operatorname{Gr}^{\mathrm{W}}_0
          \mathrm{T}(\mathrm{M})_{\mathbf{Z}}
\end{aligned}
\]

Soit \(\mathrm{M}\) un 1-motif. L’application exponentielle
\(\exp:\operatorname{Lie}(\mathrm{G})\to\mathrm{G}\) a pour noyau
\(\mathrm{H}_1(\mathrm{G})\). On définit
\(\mathrm{T}_{\mathbf{Z}}(\mathrm{M})\) (que nous noterons plutôt
\(\mathrm{T}_{\mathrm{Z}}(\mathrm{M})\)) comme étant le produit fibré
de \(\operatorname{Lie}(\mathrm{G})\) et \(\mathrm{X}\) au-dessus de
\(\mathrm{G}\)
\[
\tag*{\textbf{(10.1.3.1)}}
\begin{tikzcd}[column sep=1.5em, row sep=3.5em]
0 \arrow[r]
& \mathrm{H}_1(\mathrm{G}) \arrow[r]
& \operatorname{Lie}(\mathrm{G}) \arrow[r, "\exp"]
& \mathrm{G} \arrow[r]
& 0
\\
0 \arrow[r]
& \mathrm{H}_1(\mathrm{G}) \arrow[r] \arrow[u, equal]
& \mathrm{T}_{\mathrm{Z}}(\mathrm{M})
    \arrow[r, "\beta"] \arrow[u, "\alpha"']
& \mathrm{X} \arrow[r] \arrow[u, "u"']
& 0
\end{tikzcd}
\]
On pose
\[
\begin{aligned}
\mathrm{W}_{-1}(\mathrm{T}_{\mathrm{Z}}(\mathrm{M}))
  &=\operatorname{Ker}(\beta)=\mathrm{H}_1(\mathrm{G})
\\
\mathrm{W}_{-2}(\mathrm{T}_{\mathrm{Z}}(\mathrm{M}))
  &=\mathrm{H}_1(\mathrm{T})
    =\operatorname{Ker}(\mathrm{H}_1(\mathrm{G})\to\mathrm{H}_1(\mathrm{A})).
\end{aligned}
\]
Ceci définit la \textit{filtration par le poids}. Le morphisme
\(\alpha\) se prolonge en
\[
\alpha_{\mathbf{C}}:
\mathrm{T}_{\mathrm{Z}}(\mathrm{M})\otimes\mathbf{C}
\to\operatorname{Lie}(\mathrm{G}).
\]
On pose
\(\mathrm{F}^0(\mathrm{T}_{\mathrm{Z}}(\mathrm{M})\otimes\mathbf{C})
=\operatorname{Ker}(\alpha_{\mathbf{C}})\), et ceci définit la
\textit{filtration de Hodge} de
\[
\mathrm{T}_{\mathbf{C}}(\mathrm{M})
=\mathrm{T}_{\mathrm{Z}}(\mathrm{M})\otimes\mathbf{C}.
\]

\textit{Lemme} \textbf{(10.1.3.2)}. —
\textit{Le triplet
\(\mathrm{T}(\mathrm{M})=(\mathrm{T}_{\mathrm{Z}}(\mathrm{M}),
\mathrm{W},\mathrm{F})\) est une structure de Hodge mixte sans torsion
de type \(\{(0,0)(0,-1)(-1,0)(-1,-1)\}\) et
\(\operatorname{Gr}^{\mathrm{W}}_{-1}(\mathrm{T}(\mathrm{M}))\)
est polarisable.}
\[
\begin{tikzcd}[column sep=1.1em, row sep=2.3em]
& & 0 \arrow[d] & 0 \arrow[d] &
\\
& &
\mathrm{W}_{-1}(\mathrm{T}_{\mathbf{C}}(\mathrm{M}))\cap\mathrm{F}^0
  \arrow[r, "\sim", "{\text{\textcircled{2}}}"']
  \arrow[d]
&
\operatorname{Ker}(\alpha_{\mathrm{A},\mathbf{C}})
  \arrow[d]
&
\\
0 \arrow[r]
&
\mathrm{H}_1(\mathrm{T},\mathbf{Z})\otimes\mathbf{C}
  \arrow[r]
  \arrow[d, "\approx", "{\text{\textcircled{1}}}"']
&
\mathrm{H}_1(\mathrm{G},\mathbf{Z})\otimes\mathbf{C}
  \arrow[r]
  \arrow[d, "{\alpha_{\mathrm{G},\mathbf{C}}}"]
&
\mathrm{H}_1(\mathrm{A},\mathbf{Z})\otimes\mathbf{C}
  \arrow[r]
  \arrow[d, "{\alpha_{\mathrm{A},\mathbf{C}}}"]
&
0
\\
0 \arrow[r]
&
\operatorname{Lie}(\mathrm{T}) \arrow[r]
&
\operatorname{Lie}(\mathrm{G}) \arrow[r] \arrow[d]
&
\operatorname{Lie}(\mathrm{A}) \arrow[r] \arrow[d]
&
0
\\
& & 0 & 0 &
\end{tikzcd}
\]

\Page{55}
\textit{a)} On a
\(\mathrm{W}_{-2}(\mathrm{T}_{\mathbf{C}}(\mathrm{M}))
\cap\mathrm{F}^0=0\), car 1 est injectif : dès lors,
\((\mathrm{W}_{-2}(\mathrm{T}_{\mathbf{C}}(\mathrm{M})),\mathrm{F})\)
est une structure de Hodge de type \((-1,-1)\).

\textit{b)} Puisque 2 est surjectif, \(\mathrm{F}\) induit sur
\(\operatorname{Gr}^{\mathrm{W}}_{-1}
(\mathrm{T}_{\mathrm{Z}}(\mathrm{M}))\otimes\mathbf{C}
\simeq\mathrm{H}_1(\mathrm{A},\mathbf{Z})\otimes\mathbf{C}\)
la filtration de Hodge de
\(\mathrm{H}_1(\mathrm{A},\mathbf{Z})\otimes\mathbf{C}\) :
c’est une structure de Hodge polarisable de type
\((-1,0)+(0,-1)\) (cf. (4.4.3)).

\textit{c)}
\[
\begin{tikzcd}[column sep=1.2em, row sep=2.8em]
& 0 \arrow[d] & 0 \arrow[d] & &
\\
0 \arrow[r]
&
\mathrm{W}_{-1}(\mathrm{T}_{\mathbf{C}}(\mathrm{M}))\cap\mathrm{F}^0
  \arrow[r] \arrow[d]
&
\mathrm{F}^0 \arrow[r] \arrow[d]
&
\mathrm{V} \arrow[r] \arrow[d, "\wr"]
&
0
\\
0 \arrow[r]
&
\mathrm{H}_1(\mathrm{G},\mathbf{Z})\otimes\mathbf{C}
  \arrow[r] \arrow[d]
&
\mathrm{T}_{\mathrm{Z}}(\mathrm{M})\otimes\mathbf{C}
  \arrow[r] \arrow[d]
&
\mathrm{X}\otimes\mathbf{C} \arrow[r]
&
0
\\
&
\operatorname{Lie}(\mathrm{G}) \arrow[r, equal] \arrow[d]
&
\operatorname{Lie}(\mathrm{G}) \arrow[d]
& &
\\
& 0 & 0 & &
\end{tikzcd}
\]
\(\mathrm{F}^0\) s’envoie sur
\(\operatorname{Gr}^{\mathrm{W}}_0(\mathrm{T}_{\mathbf{C}}(\mathrm{M}))\),
qui est donc de type \((0,0)\).\par
Ceci termine la construction de \(\mathrm{T}(\mathrm{M})\).
Il est clair que \(\mathrm{T}(\mathrm{M})\) est fonctoriel en
\(\mathrm{M}\).

Soit \(\mathrm{H}\) une structure de Hodge mixte sans torsion de type
\[
\{(0,0)\ (0,-1)\ (-1,0)\ (-1,-1)\}.
\]
Supposons \(\operatorname{Gr}^{\mathrm{W}}_{-1}(\mathrm{H})\)
polarisable; le tore complexe
\[
\mathrm{A}
=\operatorname{Gr}^{\mathrm{W}}_{-1}(\mathrm{H}_{\mathbf{Z}})\backslash
\operatorname{Gr}^{\mathrm{W}}_{-1}(\mathrm{H}_{\mathbf{C}})
/
\mathrm{F}^0\operatorname{Gr}^{\mathrm{W}}_{-1}
(\mathrm{H}_{\mathbf{C}})
\]
est alors une variété abélienne (cf. (4.4.3)). Soit \(\mathrm{T}\) le
tore de groupe de caractères le dual de
\(\operatorname{Gr}^{\mathrm{W}}_{-2}(\mathrm{H}_{\mathbf{Z}})\) :
\[
\mathrm{H}_1(\mathrm{T})
=\operatorname{Gr}^{\mathrm{W}}_{-2}(\mathrm{H}_{\mathbf{Z}}).
\]
Le groupe analytique complexe
\[
\mathrm{G}
=\mathrm{W}_{-1}(\mathrm{H}_{\mathbf{Z}})\backslash
\mathrm{W}_{-1}(\mathrm{H}_{\mathbf{C}})
/
\mathrm{F}^0\cap\mathrm{W}_{-1}(\mathrm{H}_{\mathbf{C}})
\]
est une extension de \(\mathrm{A}\) par \(\mathrm{T}\).

\textit{Lemme} \textbf{(10.1.3.3)}. —
\textit{Le foncteur \(\mathrm{E}\mapsto\mathrm{E}^{\mathrm{an}}\)
est une équivalence de la catégorie des groupes algébriques extensions
de \(\mathrm{A}\) par \(\mathrm{T}\) avec celle des groupes analytiques
extensions de \(\mathrm{A}\) par \(\mathrm{T}\).}

\Page{56}
On se ramène au cas où \(\mathrm{T}=\mathbf{G}_m\).
D’après G.A.G.A., \(\mathrm{E}\mapsto\mathrm{E}^{\mathrm{an}}\) est
alors plus généralement une équivalence de la catégorie des torseurs
sous \(\mathbf{G}_m\) sur \(\mathrm{A}\) avec celle des torseurs sous
\(\mathbf{G}_m\) sur \(\mathrm{A}^{\mathrm{an}}\).

On dispose donc de \(\mathrm{G}\), extension de \(\mathrm{A}\) par
\(\mathrm{T}\). Soit
\(\mathrm{X}=\operatorname{Gr}^{\mathrm{W}}_0(\mathrm{H}_{\mathbf{Z}})\).
On définit ainsi \(u:\mathrm{X}\to\mathrm{G}\) (cf. (2.2.1)) :
\[
\begin{tikzcd}[column sep=1.1em, row sep=1.5em]
0 \arrow[r]
&
\mathrm{W}_{-1}(\mathrm{H}_{\mathbf{Z}}) \arrow[r]
&
\mathrm{W}_{-1}(\mathrm{H}_{\mathbf{C}})
/
\mathrm{W}_{-1}(\mathrm{H}_{\mathbf{C}})\cap\mathrm{F}^0
  \arrow[r] \arrow[d, "\wr"]
&
\mathrm{G} \arrow[r]
&
0
\\
& &
\mathrm{H}_{\mathbf{C}}/\mathrm{F}^0
& &
\\
0 \arrow[r]
&
\mathrm{W}_{-1}(\mathrm{H}_{\mathbf{Z}})
  \arrow[r] \arrow[uu, equal]
&
\mathrm{H}_{\mathbf{Z}} \arrow[r] \arrow[u]
&
\mathrm{X} \arrow[r] \arrow[uu]
&
0
\end{tikzcd}
\]
Il est clair que la construction précédente, qui à \(\mathrm{H}\)
associe un 1-motif sans torsion, est inverse de \(\mathrm{T}\), et est
fonctorielle.

Si \(\mathrm{G}_i\) \((i=1,2)\) est extension d’une variété abélienne
\(\mathrm{A}_i\) par un tore \(\mathrm{T}_i\), pour tout morphisme
\(u:\mathrm{G}_1\to\mathrm{G}_2\),
\(u(\mathrm{T}_1)\subset\mathrm{T}_2\) et
\(\operatorname{Ker}(u:\mathrm{A}_1\to\mathrm{A}_2)^0
=\operatorname{Im}(\operatorname{Ker}(u:\mathrm{G}_1\to\mathrm{G}_2)
\to\mathrm{A}_1)^0\).
Ceci correspond au fait qu’un morphisme de \(\mathbf{Q}\)-structures
de Hodge mixtes est strictement compatible à la filtration par le poids.

\textbf{(10.1.4)} Soit \(\mathrm{H}\) une structure de Hodge mixte
sans torsion de type
\[
\{(0,0)\ (0,-1)\ (-1,0)\ (-1,-1)\}.
\]
La filtration \(\mathrm{W}\) de \(\mathrm{H}_{\mathbf{Z}}\) définit
alors une filtration de \(\mathrm{H}\) par des sous-structures de
Hodge mixtes. Parallèlement, si
\(\mathrm{M}=(\mathrm{X},\mathrm{A},\mathrm{T},\mathrm{G},u)\)
est un 1-motif, on désigne par \(\mathrm{W}\) la filtration croissante
suivante de \(\mathrm{M}\) :
\[
\begin{aligned}
\mathrm{W}_i(\mathrm{M})
  &=0 \quad \text{pour }i<-2
    \quad\text{et}\quad
    \mathrm{W}_i(\mathrm{M})=\mathrm{M}
    \quad\text{pour }i\geq0;
\\
\mathrm{W}_{-1}(\mathrm{M})
  &=\mathrm{G}
    \quad\text{(i.e. }(\{0\},\mathrm{A},\mathrm{T},\mathrm{G},0)\text{)};
\\
\mathrm{W}_{-2}(\mathrm{M})
  &=\mathrm{T}
    \quad\text{(i.e. }(\{0\},0,\mathrm{T},\mathrm{T},0)\text{)}.
\end{aligned}
\]
En un sens évident, les
\(\operatorname{Gr}^{\mathrm{W}}_i(\mathrm{M})\) successifs pour
\(i=0,-1,-2\) sont \(\mathrm{X}\), \(\mathrm{A}\) et \(\mathrm{T}\).

\textbf{(10.1.5)} La construction (10.1.3) qui à un 1-motif
\(\mathrm{M}\) sur \(\mathbf{C}\) associe
\(\mathrm{T}_{\mathrm{Z}}(\mathrm{M})\) est transcendante.
Nous allons montrer que
\[
\widehat{\mathrm{T}}(\mathrm{M})
=\mathrm{T}_{\mathrm{Z}}(\mathrm{M})\otimes\widehat{\mathbf{Z}}
=\prod_{\ell}\mathrm{T}_{\mathrm{Z}}(\mathrm{M})\otimes\mathbf{Z}_{\ell}
\]
peut se définir de façon purement algébrique.

Soit \(\mathrm{M}\) un 1-motif sur un corps \(k\) algébriquement clos
de caractéristique 0. On identifie \(\mathrm{M}\) à un complexe
\([\mathrm{X}\xrightarrow{u}\mathrm{G}]\) (degré 0 et 1).
Pour tout entier \(n>0\), soit le complexe
\([\mathbf{Z}\xrightarrow{n}\mathbf{Z}]\) (degré \(-1\) et 0),
et soit \(\mathrm{T}_{\mathbf{Z}/n\mathbf{Z}}(\mathrm{M})\) le
\(\mathrm{H}^0\) du complexe
\([\mathrm{X}\xrightarrow{u}\mathrm{G}]
\otimes[\mathbf{Z}\xrightarrow{n}\mathbf{Z}]\) :
\Page{57}
\[
\begin{tikzcd}[column sep=2em, row sep=2.5em]
\mathrm{X} \arrow[r, "u"] & \mathrm{G}
\\
\mathrm{X} \arrow[r, "u"] \arrow[u, "n"']
& \mathrm{G} \arrow[u, "-n"']
\end{tikzcd}
\]
On a
\[
\mathrm{T}_{\mathbf{Z}/n\mathbf{Z}}(\mathrm{M})
=\{(x,g)\mid u(x)=ng\}/\{(nx,u(x))\mid x\in\mathrm{X}\}.
\]
Dans la catégorie dérivée, on a
\[
\mathrm{T}_{\mathbf{Z}/n\mathbf{Z}}(\mathrm{M})
=\mathrm{H}^0
\bigl(\mathrm{M}\overset{\mathrm{L}}{\otimes}\mathbf{Z}/n\mathbf{Z}\bigr).
\]
Pour \(n=md\), on définit des morphismes de transition
\[
\varphi_{m,n}:
\mathrm{T}_{\mathbf{Z}/n\mathbf{Z}}(\mathrm{M})
\to\mathrm{T}_{\mathbf{Z}/m\mathbf{Z}}(\mathrm{M})
\qquad\text{par}\qquad
\varphi_{m,n}((x,g))=(x,dg).
\]
On pose
\[
\widehat{\mathrm{T}}(\mathrm{M})
=\varprojlim\mathrm{T}_{\mathbf{Z}/n\mathbf{Z}}(\mathrm{M}).
\]
La filtration \(\mathrm{W}\) de \(\mathrm{M}\) induit une filtration
\(\mathrm{W}\) des
\(\mathrm{T}_{\mathbf{Z}/n\mathbf{Z}}(\mathrm{M})\) et de
\(\widehat{\mathrm{T}}(\mathrm{M})\). On a
\[
\begin{aligned}
&\operatorname{Gr}^{\mathrm{W}}_0(\widehat{\mathrm{T}}(\mathrm{M}))
  =\mathrm{X}\otimes\widehat{\mathbf{Z}}
\\
&\operatorname{Gr}^{\mathrm{W}}_{-1}(\widehat{\mathrm{T}}(\mathrm{M}))
  =\varprojlim\mathrm{A}_n=\widehat{\mathrm{T}}(\mathrm{A})
\\
&\operatorname{Gr}^{\mathrm{W}}_{-2}(\widehat{\mathrm{T}}(\mathrm{M}))
  =\varprojlim\mathrm{T}_n=\mathrm{Y}\otimes\widehat{\mathbf{Z}}(1),
\end{aligned}
\]
pour \(\mathrm{Y}\) le dual du groupe des caractères de
\(\mathrm{T}=\mathrm{W}_{-2}(\mathrm{M})\).

\textit{Construction} \textbf{(10.1.6)}. —
\textit{Soit \(\mathrm{M}=[\mathrm{X}\xrightarrow{u}\mathrm{G}]\)
un 1-motif sur \(\mathbf{C}\). On a}
\[
\tag*{\textbf{(10.1.6.1)}}
\widehat{\mathrm{T}}(\mathrm{M})
\simeq\mathrm{T}_{\mathrm{Z}}(\mathrm{M})\otimes\widehat{\mathbf{Z}}.
\]
Le morphisme naturel
\([\mathrm{T}_{\mathrm{Z}}(\mathrm{M})\to\operatorname{Lie}(\mathrm{G})]
\to[\mathrm{X}\to\mathrm{G}]\) est un quasi-isomorphisme.
Les quasi-isomorphismes
\[
{
\left[
\begin{tikzcd}[ampersand replacement=\&, column sep=1.5em, row sep=2.6em]
\mathrm{X} \arrow[r, "u"] \& \mathrm{G}
\\
\mathrm{X} \arrow[r, "u"] \arrow[u, "n"']
\& \mathrm{G} \arrow[u, "-n"']
\end{tikzcd}
\right]
}
\longleftarrow
{
\left[
\begin{tikzcd}[ampersand replacement=\&, column sep=1.5em, row sep=2.6em]
\mathrm{T}_{\mathrm{Z}} \arrow[r, "u"]
\& \operatorname{Lie}(\mathrm{G})
\\
\mathrm{T}_{\mathrm{Z}} \arrow[r, "u"] \arrow[u, "n"']
\& \operatorname{Lie}(\mathrm{G}) \arrow[u, "\wr", "-n"']
\end{tikzcd}
\right]
}
\longrightarrow
{
\left[
\begin{tikzcd}[ampersand replacement=\&, column sep=1.5em, row sep=2.6em]
\mathrm{T}_{\mathrm{Z}} \arrow[r] \& 0
\\
\mathrm{T}_{\mathrm{Z}} \arrow[r, "c"] \arrow[u, "n"']
\& 0 \arrow[u]
\end{tikzcd}
\right]
}
\]
fournissent des isomorphismes
\[
\tag*{\textbf{(10.1.6.2)}}
\mathrm{T}_{\mathbf{Z}/n\mathbf{Z}}(\mathrm{M})
\simeq\mathrm{T}_{\mathrm{Z}}/n\mathrm{T}_{\mathrm{Z}}.
\]
Par cet isomorphisme, à \(t\in\mathrm{W}_{-1}(\mathrm{T}_{\mathrm{Z}})\)
on associe \(\exp(t/n)\in\mathrm{G}_n\). Par passage à la limite,
on déduit (10.1.6.1) de (10.1.6.2).

Des techniques de Grothendieck permettent de transposer en cohomologie
de De Rham ce que nous venons de faire en cohomologie \(\ell\)-adique.
Soit \(\mathrm{M}\) un 1-motif sur un corps algébriquement clos \(k\).

\Page{58}
\textit{Construction} \textbf{(10.1.7)}. —
\textit{Construisons un espace vectoriel
\(\mathrm{T}_{\mathrm{DR}}(\mathrm{M})\) sur \(k\), muni de filtrations
\(\mathrm{W}\) et \(\mathrm{F}\).}

Soit \(\mathrm{M}=(\mathrm{X},\mathrm{A},\mathrm{T},\mathrm{G},u)\),
et considérons \(\mathrm{M}\) comme un complexe concentré en degrés
\(-1\) et 0. Une \textit{extension} de \(\mathrm{M}\) par
\(\mathbf{G}_a\) est une extension de \(\mathrm{M}\) par le complexe
réduit à \(\mathbf{G}_a\) placé en degré 0. On a
\(\operatorname{Hom}(\mathrm{G},\mathbf{G}_a)
=\operatorname{Ext}^1(\mathrm{X},\mathbf{G}_a)=0\).
Dès lors :

\noindent \textit{a)} les extensions de \(\mathrm{M}\) par
\(\mathbf{G}_a\) n’ont pas d’automorphismes;\par
\noindent \textit{b)} la suite
\[
0\to\operatorname{Hom}(\mathrm{X},\mathbf{G}_a)
\to\operatorname{Ext}^1(\mathrm{M},\mathbf{G}_a)
\to\operatorname{Ext}^1(\mathrm{G},\mathbf{G}_a)\to0
\]
est exacte.

On a
\(\operatorname{Hom}(\mathrm{T},\mathbf{G}_a)
=\operatorname{Ext}^1(\mathrm{T},\mathbf{G}_a)=0\), de sorte que
\[
\tag*{\textit{c)}}
\operatorname{Ext}^1(\mathrm{A},\mathbf{G}_a)
\xrightarrow{\sim}
\operatorname{Ext}^1(\mathrm{G},\mathbf{G}_a).
\]
L’espace vectoriel \(\operatorname{Ext}^1(\mathrm{M},\mathbf{G}_a)\)
est donc de dimension finie et, d’après \textit{a)}, il existe donc une
extension universelle de \(\mathrm{M}\) par un groupe vectoriel;
on la note \(\mathrm{M}^{\natural}:\mathrm{X}\to\mathrm{G}^{\natural}\).
\[
\begin{tikzcd}[column sep=1.8em, row sep=3.5em]
&
0 \arrow[r] \arrow[d]
&
\mathrm{X} \arrow[r, equal] \arrow[d]
&
\mathrm{X} \arrow[d]
&
\\
0 \arrow[r]
&
\operatorname{Ext}^1(\mathrm{M},\mathbf{G}_a)^* \arrow[r]
&
\mathrm{G}^{\natural} \arrow[r]
&
\mathrm{G} \arrow[r]
&
0
\end{tikzcd}
\]
On pose
\[
\begin{aligned}
\mathrm{T}_{\mathrm{DR}}(\mathrm{M})
  &=\operatorname{Lie}(\mathrm{G}^{\natural})
\\
\mathrm{F}^0\mathrm{T}_{\mathrm{DR}}(\mathrm{M})
  &=\operatorname{Ker}
    (\operatorname{Lie}\mathrm{G}^{\natural}\to\operatorname{Lie}\mathrm{G})
    \simeq(\operatorname{Ext}^1(\mathrm{M},\mathbf{G}_a))^*.
\end{aligned}
\]
\(\mathrm{T}_{\mathrm{DR}}\) est fonctoriel en \(\mathrm{M}\)
(ainsi que \(\mathrm{M}^{\natural}\)), et on définit la filtration
\(\mathrm{W}\) de \(\mathrm{T}_{\mathrm{DR}}(\mathrm{M})\) comme
provenant de la filtration \(\mathrm{W}\) de \(\mathrm{M}\).

\textit{Construction} \textbf{(10.1.8)}. —
\textit{Soit \(\mathrm{M}\) un 1-motif sur \(\mathbf{C}\). On a}
\[
(\mathrm{T}_{\mathrm{DR}}(\mathrm{M}),\mathrm{F},\mathrm{W})
\simeq
(\mathrm{T}_{\mathbf{C}}(\mathrm{M}),\mathrm{F},\mathrm{W}).
\]
Les \(\operatorname{Ext}^i(\mathrm{X},\mathbf{G}_a)\),
\(\operatorname{Ext}^i(\mathrm{A},\mathbf{G}_a)\) et
\(\operatorname{Ext}^i(\mathrm{T},\mathbf{G}_a)\) \((i=0,1)\)
sont les mêmes dans les catégories algébrique et analytique.
Soit l’application
\(\mathrm{T}_{\mathbf{C}}(\mathrm{M})
=\mathrm{T}_{\mathrm{Z}}(\mathrm{M})\otimes\mathbf{C}
\to\operatorname{Lie}(\mathrm{G})\to\mathrm{G}\).
Le diagramme
\[
\begin{tikzcd}[column sep=2em, row sep=3em]
\mathrm{X} \arrow[r] \arrow[d, equal]
&
\mathrm{T}_{\mathbf{C}}(\mathrm{M})/\mathrm{H}_1(\mathrm{G})
  \arrow[d]
\\
\mathrm{X} \arrow[r]
&
\mathrm{G}
\end{tikzcd}
\]
\Page{59}
définit une extension de \(\mathrm{M}^{\mathrm{an}}\) par le groupe
vectoriel \(\mathrm{F}^0\mathrm{T}_{\mathbf{C}}\mathrm{M}\), donc une
extension de \(\mathrm{M}\) par ce groupe. Il nous faut prouver que
c’est l’extension universelle. Pour cela, il suffit de vérifier que
la catégorie des extensions de
\([\mathrm{X}\to\mathrm{T}_{\mathbf{C}}(\mathrm{M})/\mathrm{H}_1(\mathrm{G})]\)
par \(\mathbf{G}_a\) est triviale (un seul objet, pas d’automorphisme).
C’est aussi la catégorie des extensions de
\(\mathrm{T}_{\mathbf{C}}(\mathrm{M})/\mathrm{T}_{\mathrm{Z}}(\mathrm{M})\)
par \(\mathbf{G}_a\), et on a en effet
\[
\operatorname{Ext}^i
(\mathrm{T}_{\mathbf{C}}(\mathrm{M})/\mathrm{T}_{\mathrm{Z}}(\mathrm{M}),
\mathbf{G}_a)=0
\qquad\text{pour}\qquad i=0,1.
\]

\textit{Remarque} \textbf{(10.1.9)}. —
\(\mathrm{M}^{\natural}\) est caractérisé par le diagramme commutatif
\[
\begin{tikzcd}[column sep=2em, row sep=3.2em]
\mathrm{X} \arrow[r]
&
\mathrm{G}^{\natural} \arrow[r]
&
\mathrm{G}
\\
\mathrm{T}_{\mathrm{Z}}(\mathrm{M}) \arrow[r] \arrow[u]
&
\operatorname{Lie}(\mathrm{G}^{\natural}) \arrow[r] \arrow[u]
&
\operatorname{Lie}(\mathrm{G}) \arrow[u]
\end{tikzcd}
\]
de carré extérieur (10.1.3.1), et induisant un isomorphisme
\(\mathrm{T}_{\mathbf{C}}(\mathrm{M})
\xrightarrow{\sim}\operatorname{Lie}(\mathrm{G}^{\natural})\).

\textit{Variante} \textbf{(10.1.10)}. —
Un \textit{1-motif lisse} \(\mathrm{M}\) sur un schéma \(\mathrm{S}\)
consiste en :

\noindent \textit{a)} un schéma en groupes \(\mathrm{X}\) sur
\(\mathrm{S}\), qui localement pour la topologie étale est un schéma en
groupes constants défini par un \(\mathbf{Z}\)-module libre de type
fini; un schéma abélien \(\mathrm{A}\) sur \(\mathrm{S}\), et un tore
\(\mathrm{T}\) sur \(\mathrm{S}\);\par
\noindent \textit{b)} une extension \(\mathrm{G}\) de \(\mathrm{A}\)
par \(\mathrm{T}\);\par
\noindent \textit{c)} un morphisme \(u:\mathrm{X}\to\mathrm{G}\).

Pour tout entier \(n\), une construction analogue à (10.1.5) associe
à un 1-motif lisse \(\mathrm{M}\) sur \(\mathrm{S}\) un schéma en
groupes \(\mathrm{T}_{\mathbf{Z}/n\mathbf{Z}}(\mathrm{M})\) fini et
plat sur \(\mathrm{S}\). Pour \(\ell\) inversible sur \(\mathrm{S}\),
le système projectif
\[
\mathrm{T}_{\ell}(\mathrm{M})
=\varprojlim\mathrm{T}_{\mathbf{Z}/\ell^n\mathbf{Z}}(\mathrm{M})
\]
est un \(\mathbf{Z}_{\ell}\)-faisceau sur \(\mathrm{S}\). En général,
\(\mathrm{T}_{\ell}(\mathrm{M})\) correspond à un groupe
\(\ell\)-divisible (\(=\) de Barsotti-Tate) sur \(\mathrm{S}\).

La filtration \(\mathrm{W}\) de \(\mathrm{M}\), définie comme en
(10.1.4), définit une filtration \(\mathrm{W}\) de
\(\mathrm{T}_{\ell}(\mathrm{M})\). On a encore
\[
\begin{aligned}
&\operatorname{Gr}^{\mathrm{W}}_0\mathrm{T}_{\ell}(\mathrm{M})
  =\mathrm{X}\otimes\mathbf{Z}_{\ell}
\\
&\operatorname{Gr}^{\mathrm{W}}_{-1}\mathrm{T}_{\ell}(\mathrm{M})
  =\mathrm{T}_{\ell}(\mathrm{A})
\\
&\operatorname{Gr}^{\mathrm{W}}_{-2}\mathrm{T}_{\ell}(\mathrm{M})
  =\mathrm{Y}\otimes\mathbf{Z}_{\ell}(1)
\end{aligned}
\]
pour \(\mathrm{Y}=\mathit{Hom}(\mathrm{T},\mathbf{G}_m)^{\vee}\).

De même, une construction analogue à (10.1.7) associe à \(\mathrm{M}\)
un fibré vectoriel \(\mathrm{T}_{\mathrm{DR}}(\mathrm{M})\).

\textit{Terminologie} \textbf{(10.1.11)}. —
Soit \(\mathrm{M}\) un 1-motif. On appelle \(\mathrm{T}(\mathrm{M})\),
\(\mathrm{T}_{\ell}(\mathrm{M})\) et
\(\mathrm{T}_{\mathrm{DR}}(\mathrm{M})\) les
\textit{réalisations de Hodge, \(\ell\)-adique} et
\textit{de De Rham} du 1-motif \(\mathrm{M}\).

\Page{60}
\section*{10.2. 1-motifs et biextensions}

Les résultats de ce numéro ne seront pas utilisés dans la suite de ce
paragraphe. Nous ferons usage de la notion de biextension
(S.G.A. 7, VII (2.1)) et de la généralisation suivante.

\textbf{(10.2.1)} Soient
\(\mathrm{K}_i:[\mathrm{A}_i\to\mathrm{B}_i]\) \((i=1,2)\)
deux complexes de faisceaux, concentrés en degré 0 et \(-1\)
(sur un topos quelconque). Une \textit{biextension} de
\(\mathrm{K}_1\) et \(\mathrm{K}_2\) par un faisceau abélien
\(\mathrm{H}\) consiste en :

\noindent \textit{a)} une biextension \(\mathrm{P}\) de
\(\mathrm{B}_1\) et \(\mathrm{B}_2\) par \(\mathrm{H}\), i.e. un
\(\mathrm{H}\)-torseur \(\mathrm{P}\) sur
\(\mathrm{B}_1\times\mathrm{B}_2\), avec
\(\mathrm{P}_{b_1b_2}\) dépendant biadditivement de \(b_1\) et \(b_2\);\par
\noindent \textit{b)} une trivialisation (\(=\) section biadditive)
de la biextension de \(\mathrm{B}_1\) et \(\mathrm{A}_2\) par
\(\mathrm{H}\), image réciproque de \(\mathrm{P}\);\par
\noindent \textit{c)} une trivialisation de la biextension de
\(\mathrm{A}_1\) et \(\mathrm{B}_2\) par \(\mathrm{H}\), image
réciproque de \(\mathrm{P}\); on exige que les trivialisations
\textit{b)} et \textit{c)} coïncident sur
\(\mathrm{A}_1\times\mathrm{A}_2\).

On désigne par
\(\symbfit{Biext}(\mathrm{K}_1,\mathrm{K}_2;\mathrm{H})\)
la catégorie des biextensions de \(\mathrm{K}_1\) et \(\mathrm{K}_2\)
par \(\mathrm{H}\). C’est une catégorie de Picard (les objets
s’additionnent). Le groupe
\(\operatorname{Biext}^0(\mathrm{K}_1,\mathrm{K}_2;\mathrm{H})\)
des automorphismes d’une quelconque biextension de \(\mathrm{K}_1\)
et \(\mathrm{K}_2\) par \(\mathrm{H}\) est
\[
\operatorname{Biext}^0(\mathrm{K}_1,\mathrm{K}_2;\mathrm{H})
=\operatorname{Hom}
(\mathrm{H}^0(\mathrm{K}_1)\otimes\mathrm{H}^0(\mathrm{K}_2),\mathrm{H}).
\]
On note
\(\operatorname{Biext}^1(\mathrm{K}_1,\mathrm{K}_2;\mathrm{H})\)
le groupe des classes d’isomorphismes de biextensions. On vérifie comme
dans (S.G.A. 7, VII (3.6.5)) que
\[
\operatorname{Biext}^i(\mathrm{K}_1,\mathrm{K}_2;\mathrm{H})
=\operatorname{Ext}^i
(\mathrm{K}_1\overset{\mathrm{L}}{\otimes}\mathrm{K}_2;\mathrm{H})
\qquad(i=0,1).
\]
Sur cette formule, ou de façon élémentaire, on vérifie que si
\(\mathrm{K}'_i\to\mathrm{K}_i\) est un quasi-isomorphisme
\((i=1,2)\), il revient au même de se donner une biextension de
\(\mathrm{K}_1\) et \(\mathrm{K}_2\) par \(\mathrm{H}\) ou de
\(\mathrm{K}'_1\) et \(\mathrm{K}'_2\) par \(\mathrm{H}\).

Ceci s’applique encore en remplaçant « faisceaux » par « groupes
algébriques » (resp. « groupes analytiques complexes ») : on peut
interpréter ceux-ci comme des faisceaux sur un site convenable.

\textbf{(10.2.2)} Dans ce paragraphe, nous identifierons un 1-motif
\(\mathrm{M}=(\mathrm{X},\mathrm{A},\mathrm{T},\mathrm{G},u)\)
à un complexe concentré en \textit{degrés \(0\) et \(-1\)}
\[
\mathrm{M}:[\mathrm{X}\to\mathrm{G}].
\]
Sur \(\mathbf{C}\), nous désignerons par \(\mathrm{M}^{\mathrm{an}}\)
le complexe de groupes analytiques complexes
\([\mathrm{X}\to\mathrm{G}^{\mathrm{an}}]\).
On a le résultat de rigidité :

\textit{Lemme} \textbf{(10.2.2.1)}. —
\textit{Soient
\(\mathrm{M}_i=(\mathrm{X}_i,\mathrm{A}_i,\mathrm{T}_i,\mathrm{G}_i,u_i)\)
\((i=1,2)\) deux 1-motifs. On a}
\[
\operatorname{Biext}^0(\mathrm{M}_1,\mathrm{M}_2;\mathbf{G}_m)=0.
\]
\textit{Sur \(\mathbf{C}\), on a de même
\(\operatorname{Biext}^0
(\mathrm{M}_1^{\mathrm{an}},\mathrm{M}_2^{\mathrm{an}};\mathbf{G}_m)=0\).}

Tout morphisme biadditif
\(\mathrm{G}_1\times\mathrm{G}_2\to\mathbf{G}_m\)
est en effet trivial.

% ---- 061-068 ----
\Page{61}
\textit{Construction} \textbf{(10.2.3)}. — \textit{Soient
\(\mathrm{M}_i=(\mathrm{X}_i,\mathrm{A}_i,\mathrm{T}_i,\mathrm{G}_i,u_i)\)
\((i=1,2)\) deux \(1\)-motifs sur \(\mathbf{C}\). On a}
\[
\operatorname{Biext}^{1}(\mathrm{M}_1,\mathrm{M}_2;\mathbf{G}_m)
=\operatorname{Hom}(\mathrm{T}(\mathrm{M}_1)\otimes
\mathrm{T}(\mathrm{M}_2),\mathbf{Z}(1)).
\]

Nous allons d’abord calculer
\(\operatorname{Biext}^{1}(\mathrm{M}_1^{\mathrm{an}},
\mathrm{M}_2^{\mathrm{an}};\mathbf{G}_m)\).

Si \(\mathrm{V}\) et \(\mathrm{W}\) sont deux espaces vectoriels, on a
\textit{dans la catégorie analytique}
\[
\operatorname{Biext}^{0}(\mathrm{V},\mathrm{W};\mathbf{G}_m)
\xleftarrow[\exp]{\sim}
\operatorname{Hom}(\mathrm{V}\otimes\mathrm{W},\mathbf{C})
\tag*{\textbf{(10.2.3.1)}\textsubscript{an}}
\]
\[
\operatorname{Biext}^{1}(\mathrm{V},\mathrm{W};\mathbf{G}_m)=0.
\tag*{\textbf{(10.2.3.2)}\textsubscript{an}}
\]
De ce qu’une extension de \(\mathrm{W}\) par \(\mathbf{G}_m\), sur une base
quelconque, est toujours localement triviale, il résulte en effet
formellement que
\[
\begin{aligned}
\operatorname{Biext}^{i}(\mathrm{V},\mathrm{W};\mathbf{G}_m)
&=\operatorname{Ext}^{i}(\mathrm{V},
  \mathit{Hom}(\mathrm{W},\mathbf{G}_m))\\
&=\operatorname{Ext}^{i}(\mathrm{V},\mathrm{W}^{*})
\qquad(i=1,2).
\end{aligned}
\]

Soient \(\mathrm{V}_{\mathbf{Z}}\) et \(\mathrm{W}_{\mathbf{Z}}\) deux
\(\mathbf{Z}\)-modules libres de rang fini. Puisque
\[
[\mathrm{V}_{\mathbf{Z}}\to\mathrm{V}_{\mathbf{C}}]
\to[0\to\mathrm{V}_{\mathbf{C}}/\mathrm{V}_{\mathbf{Z}}]
\]
est un quasi-isomorphisme, le foncteur « image réciproque » est une
équivalence.
\[
\symbfit{Biext}(\mathrm{V}_{\mathbf{C}}/\mathrm{V}_{\mathbf{Z}},
\mathrm{W}_{\mathbf{C}}/\mathrm{W}_{\mathbf{Z}};\mathbf{G}_m)
\xrightarrow{\approx}
\symbfit{Biext}([\mathrm{V}_{\mathbf{Z}}\to\mathrm{V}_{\mathbf{C}}],
[\mathrm{W}_{\mathbf{Z}}\to\mathrm{W}_{\mathbf{C}}];\mathbf{G}_m).
\tag*{\textbf{(10.2.3.3)}\textsubscript{an}}
\]
Tout morphisme biadditif de
\(\mathrm{V}_{\mathbf{C}}/\mathrm{V}_{\mathbf{Z}}\times
\mathrm{W}_{\mathbf{C}}/\mathrm{W}_{\mathbf{Z}}\)
dans \(\mathbf{G}_m\) est trivial, d’où \textit{a)} ci-dessous :

\textit{Lemme} \textbf{(10.2.3.4)}
\[
\begin{aligned}
\text{\textit{a)}}\quad&
\operatorname{Biext}^{0}(\mathrm{V}_{\mathbf{C}}/\mathrm{V}_{\mathbf{Z}},
\mathrm{W}_{\mathbf{C}}/\mathrm{W}_{\mathbf{Z}};\mathbf{G}_m)=0.\\
\text{\textit{b)}}\quad&
\operatorname{Biext}^{1}(\mathrm{V}_{\mathbf{C}}/\mathrm{V}_{\mathbf{Z}},
\mathrm{W}_{\mathbf{C}}/\mathrm{W}_{\mathbf{Z}};\mathbf{G}_m)
=\operatorname{Hom}(\mathrm{V}_{\mathbf{Z}}\otimes
\mathrm{W}_{\mathbf{Z}},\mathbf{Z}(1)).
\end{aligned}
\]

Soient \(\psi_i:\mathrm{V}_{\mathbf{C}}\otimes
\mathrm{W}_{\mathbf{C}}\to\mathbf{C}\) deux applications bilinéaires,
avec \(\psi=\psi_1-\psi_2\) à valeurs dans
\(\mathbf{Z}(1)=2\pi i\mathbf{Z}\) sur
\(\mathrm{V}_{\mathbf{Z}}\otimes\mathrm{W}_{\mathbf{Z}}\).
Soit \(\mathrm{P}(\psi_1,\psi_2)\) la biextension suivante de
\([\mathrm{V}_{\mathbf{Z}}\to\mathrm{V}_{\mathbf{C}}]\) et
\([\mathrm{W}_{\mathbf{Z}}\to\mathrm{W}_{\mathbf{C}}]\) par
\(\mathbf{G}_m\) : la biextension triviale de
\(\mathrm{V}_{\mathbf{C}}\) et \(\mathrm{W}_{\mathbf{C}}\) par
\(\mathbf{G}_m\), munie des trivialisations
\(\varphi_1=\exp\psi_1:\mathrm{V}_{\mathbf{Z}}\otimes
\mathrm{W}_{\mathbf{C}}\to\mathbf{G}_m\) et
\(\varphi_2=\exp\psi_2:\mathrm{V}_{\mathbf{C}}\otimes
\mathrm{W}_{\mathbf{Z}}\to\mathbf{G}_m\).
D’après (10.2.3.2), toute biextension est de cette forme; d’après
(10.2.3.1)
\(\mathrm{P}(\psi_1,\psi_2)\simeq\mathrm{P}(\psi_1',\psi_2')\)
si et seulement si \(\psi_1'-\psi_2'=\psi_1-\psi_2\).
Les \(\mathrm{P}(\psi,0)\) représentent donc chaque biextension une
fois et une seule.

\textbf{(10.2.3.5)} Sous les hypothèses générales de (10.2.1), soient
des complexes
\(\mathrm{K}_i=[\mathrm{A}_i+\mathrm{F}_i\to\mathrm{B}_i]\)
\((i=1,2)\). Soit \(\mathrm{K}_i'\) le sous-complexe
\([\mathrm{A}_i\to\mathrm{B}_i]\). On vérifie qu’il revient au même
de se donner soit une biextension de \(\mathrm{K}_1\) et
\(\mathrm{K}_2\) par \(\mathrm{H}\), soit :

\noindent\textit{a)} une biextension \(\mathrm{P}\) de
\(\mathrm{K}_1'\) et \(\mathrm{K}_2'\) par \(\mathrm{H}\); et

\noindent\textit{b)} des trivialisations des biextensions de
\([0\to\mathrm{F}_1]\) et \(\mathrm{K}_2'\)
(resp. \(\mathrm{K}_1'\) et \([0\to\mathrm{F}_2]\)) par
\(\mathrm{H}\), images réciproques de \(\mathrm{P}\); ces
trivialisations doivent coïncider sur
\(\mathrm{F}_1\times\mathrm{F}_2\).

Dans le cas particulier où
\begin{align*}
&\operatorname{Biext}^{i}(\mathrm{F}_1,\mathrm{K}_2';\mathrm{H})=0
\qquad(i=1,2)\\
\text{et}\qquad&
\operatorname{Biext}^{i}(\mathrm{K}_1',\mathrm{F}_2;\mathrm{H})=0
\qquad(i=1,2),
\end{align*}
\Page{62}
il existe une et une seule trivialisation \(\varphi_1\), section de
\(\mathrm{P}\) sur \(\mathrm{F}_1\times\mathrm{B}_2\)
(resp. \(\varphi_2\), section de \(\mathrm{P}\) sur
\(\mathrm{B}_1\times\mathrm{F}_2\)) comme en \textit{b)}.
La catégorie
\(\symbfit{Biext}(\mathrm{K}_1,\mathrm{K}_2;\mathrm{H})\)
s’identifie alors à la sous-catégorie de
\(\symbfit{Biext}(\mathrm{K}_1',\mathrm{K}_2';\mathrm{H})\)
formée des \(\mathrm{P}\) pour lesquels \(\varphi_1=\varphi_2\)
sur \(\mathrm{F}_1\times\mathrm{F}_2\).

\textbf{(10.2.3.6)} Soient \(\mathrm{V}_{\mathbf{Z}}\) et
\(\mathrm{W}_{\mathbf{Z}}\) deux \(\mathbf{Z}\)-modules libres de
type fini, et \(\mathrm{F},\mathrm{G}\) des sous-espaces vectoriels
de \(\mathrm{V}_{\mathbf{C}}\) et \(\mathrm{W}_{\mathbf{C}}\). On a
\[
\operatorname{Biext}^{i}(\mathrm{F},
[\mathrm{W}_{\mathbf{Z}}\to\mathrm{W}_{\mathbf{C}}];\mathbf{G}_m)=0.
\]
Pour \(i=0\), cela résulte de (10.2.3.1) et de ce qu’une application
bilinéaire
\[
\mathrm{B}:\mathrm{F}\times\mathrm{W}_{\mathbf{C}}\to\mathbf{C}
\qquad\text{avec}\qquad
\mathrm{B}(\mathrm{F},\mathrm{W}_{\mathbf{Z}})\subset\mathbf{Z}(1)
\]
est nulle. Pour \(i=1\), cela résulte de ce que tout morphisme
biadditif
\(\mathrm{F}\times\mathrm{W}_{\mathbf{Z}}\to\mathbf{G}_m\)
provient d’un morphisme biadditif
\(\mathrm{F}\times\mathrm{W}_{\mathbf{C}}\to\mathbf{G}_m\).
Autrement dit, on utilise la suite exacte longue
\[
\begin{aligned}
0&\to\operatorname{Biext}^{0}(\mathrm{F},
[\mathrm{W}_{\mathbf{Z}}\to\mathrm{W}_{\mathbf{C}}];\mathbf{G}_m)
\to\operatorname{Biext}^{0}(\mathrm{F},\mathrm{W}_{\mathbf{C}};\mathbf{G}_m)
\xrightarrow{\sim}
\operatorname{Biext}^{0}(\mathrm{F},\mathrm{W}_{\mathbf{Z}};\mathbf{G}_m)\\
&\to\operatorname{Biext}^{1}(\mathrm{F},
[\mathrm{W}_{\mathbf{Z}}\to\mathrm{W}_{\mathbf{C}}];\mathbf{G}_m)
\to\operatorname{Biext}^{1}(\mathrm{F},\mathrm{W}_{\mathbf{C}};\mathbf{G}_m).
\end{aligned}
\]

Nous pouvons donc appliquer (10.2.3.5) aux complexes
\(\mathrm{K}_1=[\mathrm{V}_{\mathbf{Z}}\oplus\mathrm{F}
\to\mathrm{V}_{\mathbf{C}}]\) et
\(\mathrm{K}_2=[\mathrm{W}_{\mathbf{Z}}\oplus\mathrm{G}
\to\mathrm{W}_{\mathbf{C}}]\).
On trouve que les biextensions de \(\mathrm{K}_1\) et
\(\mathrm{K}_2\) par \(\mathbf{G}_m\) s’identifient à certaines
biextensions de
\(\mathrm{V}_{\mathbf{C}}/\mathrm{V}_{\mathbf{Z}}\) et
\(\mathrm{W}_{\mathbf{C}}/\mathrm{W}_{\mathbf{Z}}\) par
\(\mathbf{G}_m\), i.e. à certaines formes
\(\psi:\mathrm{V}_{\mathbf{Z}}\otimes\mathrm{W}_{\mathbf{Z}}
\to\mathbf{Z}(1)\) (10.2.3.4).
Pour que \(\mathrm{P}(\psi_1,\psi_2)\) (10.2.3.4) provienne d’une
biextension de \(\mathrm{K}_1\) et \(\mathrm{K}_2\) par
\(\mathbf{G}_m\), il faut et suffit que
\(\exp\psi_2:\mathrm{F}\times\mathrm{W}_{\mathbf{C}}\to\mathbf{G}_m\)
et
\(\exp\psi_1:\mathrm{V}_{\mathbf{C}}\times\mathrm{G}\to\mathbf{G}_m\)
coïncident sur \(\mathrm{F}\times\mathrm{G}\), c’est-à-dire que
\(\psi(\mathrm{F},\mathrm{G})=0\).

\textit{Lemme} \textbf{(10.2.3.7)}. — \textit{Avec les notations
précédentes, on a
\(\operatorname{Biext}^{0}(\mathrm{K}_1,\mathrm{K}_2;\mathbf{G}_m)=0\)
et
\(\operatorname{Biext}^{1}(\mathrm{K}_1,\mathrm{K}_2;\mathbf{G}_m)\)
s’identifie à l’ensemble des formes
\(\psi:\mathrm{V}_{\mathbf{Z}}\otimes\mathrm{W}_{\mathbf{Z}}
\to\mathbf{Z}(1)\) de complexifiées vérifiant
\(\psi(\mathrm{F},\mathrm{G})=0\).}

Pour \(\psi\) comme en (10.2.3.7), avec \(\psi=\psi_1-\psi_2\),
la biextension correspondant à \(\psi\) est donnée par les
trivialisations suivantes de la biextension triviale de
\(\mathrm{V}_{\mathbf{C}}\) et \(\mathrm{W}_{\mathbf{C}}\) par
\(\mathbf{G}_m\) :
\[
\begin{aligned}
&\text{sur }(\mathrm{V}_{\mathbf{Z}}\oplus\mathrm{F})
\times\mathrm{W}_{\mathbf{C}}
&&:\quad \exp\psi_1(v_{\mathbf{Z}},w_{\mathbf{C}})
.\exp\psi_2(f,w_{\mathbf{C}})\\
&\text{sur }\mathrm{V}_{\mathbf{C}}\times
(\mathrm{W}_{\mathbf{Z}}\oplus\mathrm{G})
&&:\quad \exp\psi_2(v_{\mathbf{C}},w_{\mathbf{Z}})
.\exp\psi_1(v_{\mathbf{C}},g).
\end{aligned}
\]

\textbf{(10.2.3.8)} Pour les 1-motifs \(\mathrm{M}_i\), on a des
quasi-isomorphismes
\[
\begin{tikzcd}[row sep=3.5em, column sep=2em]
\mathrm{T}_{\mathbf{Z}}(\mathrm{M}_i)\oplus
\mathrm{F}^{0}\mathrm{T}_{\mathbf{C}}(\mathrm{M}_i)
  \arrow[r] \arrow[d]
&
\mathrm{T}_{\mathbf{Z}}(\mathrm{M}_i)
  \arrow[r] \arrow[d]
&
\mathrm{X}_i \arrow[d]
\\
\mathrm{T}_{\mathbf{C}}(\mathrm{M}_i) \arrow[r]
&
\operatorname{Lie}(\mathrm{G}_i) \arrow[r]
&
\mathrm{G}_i
\end{tikzcd}
\]
\Page{63}
Le groupe
\(\operatorname{Biext}^{1}(\mathrm{M}_1^{\mathrm{an}},
\mathrm{M}_2^{\mathrm{an}};\mathbf{G}_m)\)
s’identifie donc au groupe des applications
\[
\psi:\mathrm{T}_{\mathbf{Z}}(\mathrm{M}_1)\otimes
\mathrm{T}_{\mathbf{Z}}(\mathrm{M}_2)\to\mathbf{Z}(1)
\]
de complexifié compatible à la filtration de Hodge.

\textit{Lemme} \textbf{(10.2.3.9)}. — \textit{L’application
\(\operatorname{Biext}^{1}(\mathrm{M}_1,\mathrm{M}_2;\mathbf{G}_m)
\to\operatorname{Biext}^{1}(\mathrm{M}_1^{\mathrm{an}},
\mathrm{M}_2^{\mathrm{an}};\mathbf{G}_m)\)
est injective. Son image est formée des biextensions \(\mathrm{P}\)
dont la restriction à \(\mathrm{G}_1\times\mathrm{G}_2\) est image
réciproque d’une biextension de \(\mathrm{A}_1\) et
\(\mathrm{A}_2\) par \(\mathbf{G}_m\).}

On a
\begin{align*}
&\operatorname{Biext}^{1}(\mathrm{A}_1,\mathrm{A}_2;\mathbf{G}_m)
\xrightarrow{\sim}
\operatorname{Biext}^{1}(\mathrm{G}_1,\mathrm{G}_2;\mathbf{G}_m)
\qquad\text{(S.G.A. 7, VIII (3.5))}\\
\text{et}\qquad&
\operatorname{Biext}^{1}(\mathrm{A}_1,\mathrm{A}_2;\mathbf{G}_m)
\xrightarrow{\sim}
\operatorname{Biext}^{1}(\mathrm{A}_1^{\mathrm{an}},
\mathrm{A}_2^{\mathrm{an}};\mathbf{G}_m)
\qquad\text{(G.A.G.A.).}
\end{align*}

Soit \(\mathrm{P}\) une biextension de \(\mathrm{M}_1\) et
\(\mathrm{M}_2\) par \(\mathbf{G}_m\). Si
\(\mathrm{P}^{\mathrm{an}}\) est trivial, alors :

\textit{a)} la restriction de \(\mathrm{P}\) à
\(\mathrm{G}_1\times\mathrm{G}_2\) est image d’une biextension
\(\mathrm{P}_0\) de \(\mathrm{A}_1\) et \(\mathrm{A}_2\) par
\(\mathbf{G}_m\); cette dernière est triviale. Sinon elle serait
analytiquement non triviale, donnerait lieu à une forme bilinéaire
non nulle \(\psi_0\), et la forme \(\psi\) correspondant à
\(\mathrm{P}\) serait non nulle \textit{a fortiori}. La restriction
de \(\mathrm{P}\) à \(\mathrm{G}_1\times\mathrm{G}_2\) est donc
triviale (et uniquement trivialisable, même analytiquement).

\textit{b)} \(\mathrm{P}\) est donc défini par des applications
biadditives \(\mathrm{X}_1\times\mathrm{G}_2\to\mathbf{G}_m\) et
\(\mathrm{G}_1\times\mathrm{X}_2\to\mathbf{G}_m\). Celles-ci sont
nulles dans la catégorie analytique, donc sont nulles, et
\(\mathrm{P}\) est trivial.

Soit \(\mathrm{P}\) une biextension de
\(\mathrm{M}_1^{\mathrm{an}}\) et \(\mathrm{M}_2^{\mathrm{an}}\)
par \(\mathbf{G}_m\) qui vérifie la condition (10.2.3.9).
Il reste à prouver que \(\mathrm{P}\) est algébrisable. Par
hypothèse et G.A.G.A., sa restriction \(\mathrm{P}_1\) à
\(\mathrm{G}_1\times\mathrm{G}_2\) l’est. On conclut en montrant
qu’une trivialisation de \(\mathrm{P}_1\) sur
\(\mathrm{X}_1\times\mathrm{G}_2\)
(resp. \(\mathrm{G}_1\times\mathrm{X}_2\)) est automatiquement
algébrique : \(\mathrm{P}_1\) s’interprète comme une extension de
\(\mathrm{X}_1\otimes\mathrm{G}_2\) par \(\mathbf{G}_m\), ce qui
ramène à montrer que
\[
\operatorname{Ext}^{i}(\mathrm{G}_1,\mathbf{G}_m)
\xrightarrow{\sim}
\operatorname{Ext}^{i}(\mathrm{G}_1^{\mathrm{an}},\mathbf{G}_m)
\qquad(i=0,1).
\]
Cela résulte de la même assertion pour une variété abélienne
(G.A.G.A.) et pour \(\mathbf{G}_m\).

D’après (10.2.3.9),
\(\operatorname{Biext}^{1}(\mathrm{M}_1,\mathrm{M}_2;\mathbf{G}_m)\)
s’identifie à l’ensemble des
\[
\psi:\mathrm{T}_{\mathbf{Z}}(\mathrm{M}_1)\otimes
\mathrm{T}_{\mathbf{Z}}(\mathrm{M}_2)\to\mathbf{Z}(1)
\]
tels que :

\noindent\textit{a)} \(\psi\) est compatible à la filtration de
Hodge (10.2.3.8);

\noindent\textit{b)} la restriction de \(\psi\) à
\(\mathrm{W}_{-1}\mathrm{T}_{\mathbf{Z}}(\mathrm{M}_1)\otimes
\mathrm{W}_{-1}\mathrm{T}_{\mathbf{Z}}(\mathrm{M}_2)\)
provient d’une forme sur
\[
\operatorname{Gr}^{\mathrm{W}}_{-1}
\mathrm{T}_{\mathbf{Z}}(\mathrm{M}_1)\otimes
\operatorname{Gr}^{\mathrm{W}}_{-1}
\mathrm{T}_{\mathbf{Z}}(\mathrm{M}_2).
\]
La condition \textit{b)} signifie que \(\psi\) est compatible à
\(\mathrm{W}\), et ceci achève la construction (10.2.3).

\textit{Remarque} \textbf{(10.2.4)}. — Une biextension
\(\mathrm{P}\) de \(\mathrm{M}_1\) et \(\mathrm{M}_2\) par
\(\mathbf{G}_m\) définit de façon évidente une biextension
\(\mathrm{P}^{0}\) de \(\mathrm{M}_2\) et \(\mathrm{M}_1\) par
\(\mathbf{G}_m\). Si \(\psi\) (resp. \(\psi^{0}\)) est la forme
bilinéaire
\Page{64}
définie par \(\mathrm{P}\) (resp. \(\mathrm{P}^{0}\)), on a
\(\psi(x,y)=-\psi^{0}(y,x)\). C’est clair sur les formules de
(10.2.3.7).

La construction (10.2.3) a des analogues purement algébriques en
cohomologie \(\ell\)-adique ou de De Rham.

\textit{Construction} \textbf{(10.2.5)}. — \textit{Soient
\(\mathrm{M}_i=(\mathrm{X}_i,\mathrm{A}_i,\mathrm{T}_i,\mathrm{G}_i,u_i)\)
\((i=1,2)\) deux \(1\)-motifs sur un corps algébriquement clos
\(k\) de caractéristique \(0\). Soit \(\mathrm{P}\) une biextension
de \(\mathrm{M}_1\) et \(\mathrm{M}_2\) par \(\mathbf{G}_m\).
On lui associe un morphisme}
\[
\mathrm{T}_{\mathbf{Z}/n\mathbf{Z}}(\mathrm{M}_1)\otimes
\mathrm{T}_{\mathbf{Z}/n\mathbf{Z}}(\mathrm{M}_2)
\to\mathbf{Z}/n\mathbf{Z}(1).
\]

Soit \(m_i\in\mathrm{T}_{\mathbf{Z}/n\mathbf{Z}}(\mathrm{M}_i)\),
représenté par \((x_i,g_i)\) avec \(u_i(x_i)=ng_i\).
Construisons deux isomorphismes \(a_1\) et \(a_2\) de
\(\mathrm{P}_{g_1g_2}^{\otimes n}\) avec le torseur trivial
\(\mathbf{G}_m\) :
\[
\begin{aligned}
a_1:\quad
\mathrm{P}_{g_1,g_2}^{\otimes n}
&\xrightarrow{\sim}\mathrm{P}_{ng_1,g_2}
=\mathrm{P}_{u(x_1),g_2}\xrightarrow{\sim}\mathbf{G}_m\\
a_2:\quad
\mathrm{P}_{g_1,g_2}^{\otimes n}
&\xrightarrow{\sim}\mathrm{P}_{g_1,ng_2}
=\mathrm{P}_{g_1,u(x_2)}\xrightarrow{\sim}\mathbf{G}_m.
\end{aligned}
\]
On pose \(a_2=\Phi(m_1,m_2)a_1\). On vérifie que
\(\Phi(m_1,m_2)\) ne dépend pas du choix des \((x_i,g_i)\) et
est une racine \(n\)-ième de l’unité. C’est la forme cherchée.

\textit{Proposition} \textbf{(10.2.6)}. — \textit{Sur
\(\mathbf{C}\), la forme \((10.2.5)\) se déduit de celle construite
en \((10.2.3)\) par réduction modulo \(n\).}

Cela se vérifie à l’aide des formules données en (10.2.3.7).

\textit{Construction} \textbf{(10.2.7)}. — \textit{Soient
\(\mathrm{M}_i=(\mathrm{X}_i,\mathrm{A}_i,\mathrm{T}_i,\mathrm{G}_i,u_i)\)
deux \(1\)-motifs sur un corps algébriquement clos \(k\).
Soit \(\mathrm{P}\) une biextension de \(\mathrm{M}_1\) et
\(\mathrm{M}_2\) par \(\mathbf{G}_m\).
On lui associe un morphisme}
\[
\mathrm{T}_{\mathrm{DR}}(\mathrm{M}_1)\otimes_k
\mathrm{T}_{\mathrm{DR}}(\mathrm{M}_2)\to k.
\]

Rappelons les définitions (dues à Grothendieck) des
\(\natural\)-extensions et biextensions.

\textbf{(10.2.7.1)} Une \(\natural\)-\textit{extension} d’un groupe
commutatif lisse \(\mathrm{G}\) par un groupe commutatif lisse
\(\mathrm{H}\) consiste en :

\noindent\textit{a)} une extension \(\mathrm{E}\) de \(\mathrm{G}\)
par \(\mathrm{H}\); si
\(\mu:\mathrm{G}\times\mathrm{G}\to\mathrm{G}\) est l’addition,
on regarde \(\mathrm{E}\) comme un \(\mathrm{H}\)-torseur sur
\(\mathrm{G}\) muni d’un morphisme de \(\mathrm{H}\)-torseurs sur
\(\mathrm{G}\times\mathrm{G}\) (l’addition)
\[
\nu:\operatorname{pr}_1^{*}\mathrm{E}
+\operatorname{pr}_2^{*}\mathrm{E}\to\mu^{*}\mathrm{E};
\]

\noindent\textit{b)} une connexion sur le \(\mathrm{H}\)-torseur
\(\mathrm{E}\), telle que l’application \(\nu\) soit horizontale.
Cette connexion est automatiquement à courbure nulle.

Si \(\mathrm{G}\) est extension d’une variété abélienne par un
tore, toute extension de \(\mathrm{G}\) par \(\mathbf{G}_m\)
admet une \(\natural\)-structure; deux \(\natural\)-structures
diffèrent par une forme invariante sur \(\mathrm{G}\) (à valeurs
dans l’algèbre de Lie de \(\mathbf{G}_m\)). On définit de façon
évidente la somme de Baer de deux \(\natural\)-extensions.

\Page{65}
\textbf{(10.2.7.2)} Soit \(\mathrm{P}\) une biextension de
\(\mathrm{G}_1\) et \(\mathrm{G}_2\) par \(\mathrm{H}\).
Regardons \(\mathrm{P}\) comme un \(\mathrm{H}\)-torseur
\(\mathrm{P}\) sur \(\mathrm{G}_1\times\mathrm{G}_2\), muni de
morphismes
\[
\begin{aligned}
\nu_1:\quad
\operatorname{pr}_{13}^{*}\mathrm{P}
+\operatorname{pr}_{23}^{*}\mathrm{P}
&\to(\mu_1\times\mathrm{Id})^{*}\mathrm{P}
&&\quad\text{sur}\quad
\mathrm{G}_1\times\mathrm{G}_1\times\mathrm{G}_2\\
\nu_2:\quad
\operatorname{pr}_{12}^{*}\mathrm{P}
+\operatorname{pr}_{13}^{*}\mathrm{P}
&\to(\mathrm{Id}\times\mu_2)^{*}\mathrm{P}
&&\quad\text{sur}\quad
\mathrm{G}_1\times\mathrm{G}_2\times\mathrm{G}_2.
\end{aligned}
\]
Une \(\natural\)-1-\textit{structure} sur \(\mathrm{P}\) est une
connexion sur le \(\mathrm{H}\)-torseur \(\mathrm{P}\), relative à
\(\mathrm{G}_1\times\mathrm{G}_2\to\mathrm{G}_2\) (la dérivée
covariante n’est définie que pour des champs de vecteurs parallèles
à \(\mathrm{G}_1\)), telle que \(\nu_1\) et \(\nu_2\) soient
horizontaux. On définit de même les
\(\natural\)-2-\textit{structures}. Une
\(\natural\)-\textit{structure} est la donnée d’une
\(\natural\)-1-structure et d’une \(\natural\)-2-structure.
C’est aussi une connexion sur le \(\mathrm{H}\)-torseur
\(\mathrm{P}\), telle que \(\nu_1\) et \(\nu_2\) soient horizontaux.

La \textit{courbure} \(\mathrm{R}\) d’une \(\natural\)-biextension
\(\mathrm{P}\) est la courbure de la connexion de \(\mathrm{P}\).
C’est une 2-forme invariante par translation sur
\(\mathrm{G}_1\times\mathrm{G}_2\), à valeurs dans
\(\operatorname{Lie}(\mathrm{H})\), i.e. c’est une forme alternée
sur
\(\operatorname{Lie}(\mathrm{G}_1)\times
\operatorname{Lie}(\mathrm{G}_2)\), à valeurs dans
\(\operatorname{Lie}(\mathrm{H})\).
Sa restriction à \(\operatorname{Lie}(\mathrm{G}_1)\) et à
\(\operatorname{Lie}(\mathrm{G}_2)\) est nulle, de sorte que
\(\mathrm{R}\) définit un accouplement (encore appelé
\textit{courbure} de \(\mathrm{P}\))
\[
\Phi:\operatorname{Lie}(\mathrm{G}_1)\otimes
\operatorname{Lie}(\mathrm{G}_2)\to\operatorname{Lie}(\mathrm{H}),
\tag*{\textbf{(10.2.7.3)}}
\]
avec
\[
\mathrm{R}(g_1+g_2,g_1'+g_2')
=\Phi(g_1,g_2')-\Phi(g_1',g_2).
\]

Pour des complexes de groupes lisses
\(\mathrm{K}_i:[\mathrm{A}_i\to\mathrm{B}_i]\), une
\(\natural\)-\textit{biextension}
(resp. \(\natural\)-\(i\)-\textit{biextension}) de
\(\mathrm{K}_1\) et \(\mathrm{K}_2\) par \(\mathrm{H}\) est une
\(\natural\)-biextension (resp. \(\natural\)-\(i\)-biextension)
de \(\mathrm{B}_1\) et \(\mathrm{B}_2\) par \(\mathrm{H}\),
munie de trivialisations (en tant que \(\natural\)-biextension,
resp. \(\natural\)-\(i\)-biextension) comme en (10.2.1).

\textit{Proposition} \textbf{(10.2.7.4)}. — \textit{Soit
\(\mathrm{P}^{\natural}\) la biextension de
\(\mathrm{M}_1^{\natural}\) et \(\mathrm{M}_2^{\natural}\) par
\(\mathbf{G}_m\) image réciproque de \(\mathrm{P}\). Il existe
sur \(\mathrm{P}^{\natural}\) une et une seule
\(\natural\)-structure.}

Prouvons l’unicité : il faut vérifier que sur la biextension
triviale de \(\mathrm{M}_1^{\natural}\) et
\(\mathrm{M}_2^{\natural}\) par \(\mathbf{G}_m\), toute
\(\natural\)-structure est triviale. Une \(\natural\)-structure
est d’abord une connexion sur le \(\mathbf{G}_m\)-torseur trivial
sur \(\mathrm{G}_1^{\natural}\times\mathrm{G}_2^{\natural}\),
i.e. un champ de formes sur
\(\mathrm{G}_1^{\natural}\times\mathrm{G}_2^{\natural}\),
soit \(\Gamma_{g_1,g_2}(t_1+t_2)\).
L’axiome des \(\natural\)-structures donne que
\[
\Gamma_{g_1g_2}(t_1+t_2)
=\Phi_1(g_1,t_2)+\Phi_2(t_1,g_2)
\]
avec \(\Phi_1\) et \(\Phi_2\) biadditifs.
\(\mathrm{G}_i^{\natural}\) est extension de
\(\mathrm{X}_i\otimes k\) par l’extension universelle
\(\mathrm{G}_i'\) de \(\mathrm{G}_i\) par un groupe additif.
Tout morphisme \(\mathrm{G}_i'\to\mathbf{G}_a\) est trivial.
Dès lors, \(\Phi_1\) se factorise par
\((\mathrm{X}_1\otimes k)\otimes
\operatorname{Lie}(\mathrm{G}_2^{\natural})\).
Enfin, la restriction de \(\Gamma\) à
\(\mathrm{X}_1\times\mathrm{G}_2^{\natural}\) doit être triviale,
donc aussi la restriction de \(\Phi_1\) à
\(\mathrm{X}_1\times\operatorname{Lie}(\mathrm{G}_2^{\natural})\).
On a donc \(\Phi_1=0\). De même, \(\Phi_2=0\), et \(\Gamma=0\).

Pour prouver l’existence, il suffit de construire une
\(\natural\)-2-structure sur la biextension de
\(\mathrm{M}_1^{\natural}\) et \(\mathrm{M}_2\) par
\(\mathbf{G}_m\) image réciproque de \(\mathrm{P}\) : par image
réciproque, on obtiendra une \(\natural\)-2-structure sur
\(\mathrm{P}^{\natural}\) et, symétriquement, une
\(\natural\)-1-structure.

Pour \(g\in\mathrm{G}_1\), \(\mathrm{P}_g\) est une extension
de \(\mathrm{G}_2\) par \(\mathbf{G}_m\).
Soit \(\mathrm{C}_g\) l’ensemble des \(\natural\)-structures sur
\(\mathrm{P}_g\). C’est un torseur sous
\(\operatorname{Lie}(\mathrm{G}_2)^{*}\).
Pour \(g\in\mathrm{G}_1\), les \(\mathrm{C}_g\) sont les fibres
d’un torseur \(\mathrm{C}\) sur \(\mathrm{G}_1\); l’addition
de Baer des \(\natural\)-extensions fait même de \(\mathrm{C}\)
une extension de \(\mathrm{G}_1\) par
\Page{66}
\(\operatorname{Lie}(\mathrm{G}_2)^{*}\).
Relevons \(u_1:\mathrm{X}\to\mathrm{G}_1\) en
\(u_1':\mathrm{X}\to\mathrm{C}\) en associant à
\(x\in\mathrm{X}\) la connexion triviale de l’extension
trivialisée \(\mathrm{P}_{u_1(x)}\).
On obtient ainsi une \(\natural\)-2-structure sur la biextension
de \([\mathrm{X}\to\mathrm{C}]\) et \(\mathrm{M}_2\) par
\(\mathbf{G}_m\) image réciproque de \(\mathrm{P}\).

Vu la propriété universelle de \(\mathrm{M}_1^{\natural}\),
il existe un unique diagramme commutatif
\[
\begin{tikzcd}[row sep=3em, column sep=1.5em]
& {[\mathrm{X}\longrightarrow\mathrm{G}]} &\\
\mathrm{M}_1^{\natural}
  \arrow[ur] \arrow[rr, "v"]
&& {[\mathrm{X}\longrightarrow\mathrm{C}]} \arrow[ul]
\end{tikzcd}
\]
prenant une image réciproque par \(v\), on trouve la
\(\natural\)-2-structure cherchée.

Par définition, l’accouplement (10.2.7) est l’opposé de la
courbure de la \(\natural\)-biextension (10.2.7.4).

\textit{Proposition} \textbf{(10.2.8)}. — \textit{Sur
\(\mathbf{C}\), l’accouplement \((10.2.7)\) est le complexifié
de l’accouplement \((10.2.3)\).}

La \(\natural\)-biextension \(\mathrm{P}^{\natural}\) définit,
dans la catégorie analytique, une \(\natural\)-biextension de
\([\mathrm{T}_{\mathbf{Z}}(\mathrm{M}_1)\to
\mathrm{T}_{\mathbf{C}}(\mathrm{M}_1)]\) et
\([\mathrm{T}_{\mathbf{Z}}(\mathrm{M}_2)\to
\mathrm{T}_{\mathbf{C}}(\mathrm{M}_2)]\) par \(\mathbf{G}_m\).
La proposition (10.2.8) résulte alors du

\textit{Lemme} \textbf{(10.2.9)}. — \textit{Soient
\(\mathrm{V}_{\mathbf{Z}}\) et \(\mathrm{W}_{\mathbf{Z}}\)
deux \(\mathbf{Z}\)-modules libres de rang fini, et \(\mathrm{P}^{\natural}\)
une \(\natural\)-biextension de
\([\mathrm{V}_{\mathbf{Z}}\to\mathrm{V}_{\mathbf{C}}]\) et
\([\mathrm{W}_{\mathbf{Z}}\to\mathrm{W}_{\mathbf{C}}]\) par
\(\mathbf{G}_m\), de biextension sous-jacente \(\mathrm{P}\).}

\textit{Alors, la courbure
\(\psi:\mathrm{V}_{\mathbf{C}}\otimes\mathrm{W}_{\mathbf{C}}
\to\mathbf{C}\) de \(\mathrm{P}^{\natural}\) coïncide avec
l’opposé de la complexifiée de la forme correspondant à
\(\mathrm{P}\) par \((10.2.3.4)\).}

Prenons \(\mathrm{P}=\mathrm{P}(\psi_1,\psi_2)\) (10.2.3.4).
Une \(\natural\)-structure sur \(\mathrm{P}\) est d’abord une
connexion sur le \(\mathbf{G}_m\)-torseur trivial sur
\(\mathrm{V}_{\mathbf{C}}\times\mathrm{W}_{\mathbf{C}}\),
i.e. un champ de formes \(\Gamma_{v,w}(v'+w')\) sur
\(\mathrm{V}_{\mathbf{C}}\times\mathrm{W}_{\mathbf{C}}\).
Cette connexion définit une \(\natural\)-structure sur ce
torseur si et seulement si
\[
\Gamma_{v,w}(v'+w')=\Phi_1(v,w')+\Phi_2(v',w)
\]
avec \(\Phi_1\) et \(\Phi_2\) bilinéaires.
Les trivialisations \(\exp\psi_1(v_{\mathbf{Z}},w_{\mathbf{C}})\)
et \(\exp\psi_2(v_{\mathbf{C}},w_{\mathbf{Z}})\) sont
horizontales si et seulement si \(\Phi_i=-\psi_i\).
La courbure de la connexion est le champ de 2-formes \(d\Gamma\) :
\[
\mathrm{R}_{v,w}(v'+w',v''+w'')
=\Phi_1(v',w'')+\Phi_2(v'',w')
-\Phi_1(v'',w')-\Phi_2(v',w'').
\]
La courbure de \(\mathrm{P}\) est donc la forme
\[
\Phi_1(v',w'')-\Phi_2(v',w'')
=-(\psi_1(v',w'')-\psi_2(v',w''))
=-\psi(v',w'').
\]

\textbf{(10.2.10)} Soit \(\mathrm{M}\) un 1-motif sur
\(\mathbf{C}\). A \(\mathrm{M}\) correspond une structure de
Hodge mixte \(\mathrm{T}(\mathrm{M})\) de type
\(\{(0,0)\ (0,-1)\ (-1,0)\ (-1,-1)\}\).
La structure de Hodge mixte
\(\operatorname{Hom}(\mathrm{T}(\mathrm{M}),\mathbf{Z}(1))\)
est encore de ce type, donc définit un 1-motif \(\mathrm{M}^{*}\).
L’application évidente
\(\mathrm{T}(\mathrm{M})\otimes\mathrm{T}(\mathrm{M}^{*})
\to\mathbf{Z}(1)\) définit une biextension \(\mathrm{P}\) de
\(\mathrm{M}\) et \(\mathrm{M}^{*}\) par \(\mathbf{G}_m\).
Le 1-motif \(\mathrm{M}^{*}\), muni de la biextension
\(\mathrm{P}\), est le \textit{dual de Cartier} de \(\mathrm{M}\).
On appelle \(\mathrm{P}\) la \textit{biextension de Poincaré}.

\Page{67}
\textbf{(10.2.11)} La construction (10.2.10) peut être rendue
purement algébrique. A un 1-motif
\(\mathrm{M}=(\mathrm{X},\mathrm{A},\mathrm{T},\mathrm{G},u)\)
sur un corps algébriquement clos \(k\), nous allons associer
un \textit{dual de Cartier}
\(\mathrm{M}^{*}=(\mathrm{X}',\mathrm{A}',\mathrm{T}',
\mathrm{G}',u')\) :

\textit{a)} \(\mathrm{X}'\) est le groupe des caractères de
\(\mathrm{T}\), \(\mathrm{A}'\) la variété abélienne duale de
\(\mathrm{A}\) et \(\mathrm{T}'\) le tore de groupe de
caractères \(\mathrm{X}\).

\textit{b)} Traitons tout d’abord le cas où \(\mathrm{T}=0\).
Une extension de \(\mathrm{M}\) par \(\mathbf{G}_m\) consiste
alors en une extension \(\mathrm{E}\) de \(\mathrm{A}\) par
\(\mathbf{G}_m\), plus une trivialisation \(\tilde{u}\) de
l’extension \(u^{*}\mathrm{E}\) de \(\mathrm{X}\) par
\(\mathbf{G}_m\)
\[
\begin{tikzcd}[row sep=2.5em, column sep=2em]
&&& \mathrm{X} \arrow[dl, "\tilde{u}"'] \arrow[d, "u"] &\\
0 \arrow[r] & \mathbf{G}_m \arrow[r] &
\mathrm{E} \arrow[r] & \mathrm{A} \arrow[r] & 0
\end{tikzcd}
\]
On dispose d’une suite exacte
\[
0\to\operatorname{Hom}(\mathrm{X},\mathbf{G}_m)
\to\operatorname{Ext}^{1}(\mathrm{M},\mathbf{G}_m)
\to\operatorname{Ext}^{1}(\mathrm{A},\mathbf{G}_m)\to0.
\]
Les extensions de \(\mathrm{M}\) par \(\mathbf{G}_m\) n’ont
pas d’automorphisme non trivial. Le foncteur des classes
d’isomorphie d’extensions est représentable; on définit
\(\mathrm{G}'\) comme étant le groupe algébrique qui représente
ce foncteur; c’est une extension de \(\mathrm{A}'\) par
\(\mathrm{T}'\). On dispose d’une biextension \(\mathrm{P}\)
de \(\mathrm{M}\) et \(\mathrm{G}'\) par \(\mathbf{G}_m\).

\textit{c)} Dans le cas général, on pose
\(\mathrm{G}'=\operatorname{Ext}^{1}
(\mathrm{M}/\mathrm{W}_{-2}\mathrm{M},\mathbf{G}_m)\).
La biextension naturelle \(\mathrm{P}''\) de
\(\mathrm{M}/\mathrm{W}_{-2}\mathrm{M}\) et \(\mathrm{G}'\)
par \(\mathbf{G}_m\) induit une biextension \(\mathrm{P}'\)
de \(\mathrm{M}\) et \(\mathrm{G}'\) par \(\mathbf{G}_m\).
Pour chaque \(\chi\in\mathrm{X}'\), l’extension \(\mathrm{M}\)
de \(\mathrm{M}/\mathrm{W}_{-2}\mathrm{M}\) par \(\mathrm{T}\)
définit une extension de
\(\mathrm{M}/\mathrm{W}_{-2}\mathrm{M}\) par \(\mathbf{G}_m\),
d’où \(u'(\chi)\in\mathrm{G}'\).
\[
\begin{tikzcd}[row sep=3em, column sep=2em]
0 \arrow[r] &
\mathrm{T} \arrow[r] \arrow[d] &
\mathrm{M} \arrow[r] \arrow[d, "v"] &
\mathrm{M}/\mathrm{W}_{-2}\mathrm{M}
  \arrow[r] \arrow[d, equal] &
0\\
0 \arrow[r] &
\mathbf{G}_m \arrow[r] &
\mathrm{P}''_{u'(\chi)} \arrow[r] &
\mathrm{M}/\mathrm{W}_{-2}\mathrm{M} \arrow[r] &
0
\end{tikzcd}
\]
L’application \(v\) trivialise l’extension
\(\mathrm{P}'_{u'(\chi)}\) de \(\mathrm{M}\) par \(\mathbf{G}_m\)
\[
\begin{tikzcd}[row sep=3em, column sep=2em]
0 \arrow[r] &
\mathbf{G}_m \arrow[r] \arrow[d, equal] &
\mathrm{P}'_{u'(\chi)} \arrow[r] \arrow[d] &
\mathrm{M} \arrow[r] \arrow[d] \arrow[dl, "v"'] &
0\\
& \mathbf{G}_m \arrow[r] &
\mathrm{P}''_{u'(\chi)} \arrow[r] &
\mathrm{M}/\mathrm{W}_{-2}\mathrm{M} &
\end{tikzcd}
\]
Ces trivialisations font de \(\mathrm{P}'\) une biextension de
\(\mathrm{M}\) et
\(\mathrm{M}^{*}=[\mathrm{X}'\xrightarrow{u'}\mathrm{G}']\)
par \(\mathbf{G}_m\). Le 1-motif \(\mathrm{M}^{*}\), muni de
\(\mathrm{P}'\), est le dual de Cartier cherché.

\Page{68}
\textbf{(10.2.12)} Cette construction permet de donner des
1-motifs une description plus symétrique.

Soit \(\mathrm{M}=(\mathrm{X},\mathrm{A},\mathrm{T},\mathrm{G},u)\)
un 1-motif, de dual de Cartier
\(\mathrm{M}^{*}=(\mathrm{X}',\mathrm{A}',\mathrm{T}',
\mathrm{G}',u')\), et soit \(\mathrm{P}_0\) la biextension de
Poincaré. Les morphismes \(u\) et \(u'\) définissent par passage
au quotient des morphismes
\[
\bar{u}:\mathrm{X}\to\mathrm{A}
\qquad\text{et}\qquad
\bar{u}':\mathrm{X}'\to\mathrm{A}'.
\tag*{\textbf{(10.2.12.1)}}
\]
La biextension de \(\mathrm{G}\) et \(\mathrm{G}'\) par
\(\mathbf{G}_m\) déduite de \(\mathrm{P}_0\) est image réciproque
de la biextension de Poincaré \(\mathrm{P}\) de \(\mathrm{A}\)
et \(\mathrm{A}'\) par \(\mathbf{G}_m\). Par hypothèse, les
trivialisations (des images réciproques) de \(\mathrm{P}\) sur
\(\mathrm{X}\times\mathrm{G}'\) et
\(\mathrm{G}\times\mathrm{X}'\) coïncident sur
\(\mathrm{X}\times\mathrm{X}'\), d’où
\[
\text{une trivialisation }\psi\text{ de }
(\bar{u}\times\bar{u}')^{*}\mathrm{P}.
\tag*{\textbf{(10.2.12.2)}}
\]
Un morphisme \(f:\mathrm{M}_1\to\mathrm{M}_2\), de transposé
(en un sens évident) \(\mathrm{M}_2^{*}\to\mathrm{M}_1^{*}\)
donne lieu à des diagrammes commutatifs
\[
\begin{tikzcd}[row sep=2.5em, column sep=2em]
\mathrm{X}_1 \arrow[r] \arrow[d, "f_{\mathrm{X}}"] &
\mathrm{A}_1 \arrow[d, "f_{\mathrm{A}}"]\\
\mathrm{X}_2 \arrow[r] & \mathrm{A}_2
\end{tikzcd}
\qquad
\begin{tikzcd}[row sep=2.5em, column sep=2em]
\mathrm{X}_1' \arrow[r] & \mathrm{A}_1'\\
\mathrm{X}_2' \arrow[r] \arrow[u, "f_{\mathrm{X}}'"'] &
\mathrm{A}_2' \arrow[u, "f_{\mathrm{A}}'"']
\end{tikzcd}
\tag*{\textbf{(10.2.12.3)}}
\]
De plus, \(f_{\mathrm{A}}\) et \(f_{\mathrm{A}}'\) sont transposés
et, \textit{via} l’isomorphisme
\((1,f_{\mathrm{A}}')^{*}\mathrm{P}_1
\simeq(f_{\mathrm{A}},1)^{*}\mathrm{P}_2\) on a
\[
\psi_1(x_1,f_{\mathrm{X}}'(x_2'))
=\psi_2(f_{\mathrm{X}}(x_1),x_2')
\tag*{\textbf{(10.2.12.4)}}
\]
\[
\begin{tikzcd}[row sep=4em, column sep=2em]
\mathrm{X}_1\times\mathrm{X}_1' \arrow[r, "\psi_1"] &
(\mathrm{P}_1\text{ sur }\mathrm{A}_1\times\mathrm{A}_1')
\\
\mathrm{X}_1\times\mathrm{X}_2'
  \arrow[u] \arrow[d] \arrow[r, dotted] &
(1,f_{\mathrm{A}}')^{*}\mathrm{P}_1
=(f_{\mathrm{A}},1)^{*}\mathrm{P}_2
\text{ sur }\mathrm{A}_1\times\mathrm{A}_2'
  \arrow[u] \arrow[d]
\\
\mathrm{X}_2\times\mathrm{X}_2' \arrow[r, "\psi_2"] &
(\mathrm{P}_2\text{ sur }\mathrm{A}_2\times\mathrm{A}_2').
\end{tikzcd}
\]

\textbf{(10.2.13)} On obtient ainsi un foncteur de la catégorie
des 1-motifs dans la catégorie suivante. Les objets consistent en :

\noindent\textit{a)} deux variétés abéliennes en dualité
\(\mathrm{A}\) et \(\mathrm{A}'\);

\noindent\textit{b)} deux \(\mathbf{Z}\)-modules libres de rang
fini \(\mathrm{X}\) et \(\mathrm{X}'\);

\noindent\textit{c)} des morphismes
\(v:\mathrm{X}\to\mathrm{A}\) et
\(v':\mathrm{X}'\to\mathrm{A}'\);

\noindent\textit{d)} une trivialisation \(\psi\) de l’image
réciproque par \((v,v')\) de la biextension de Poincaré de
\(\mathrm{A}\) et \(\mathrm{A}'\) par \(\mathbf{G}_m\).

% ---- 069-077 ----
\Page{69}
Les flèches consistent en les diagrammes (10.2.12.3) vérifiant les conditions de (10.2.12).

\textit{Proposition} \textbf{(10.2.14)}. — \textit{Le foncteur \textup{(10.2.13)} est une équivalence de catégories.}

La vérification est laissée en exercice au lecteur.

\section*{10.3. Interprétation algébrique du \(\mathbf{H}^{1}\) mixte : cas des courbes.}

\textbf{(10.3.1)} Soit \(\mathrm{X}_0\) une courbe algébrique sur un corps algébriquement clos \(k\). Soit \(p:\mathrm{X}'\to\mathrm{X}_0\) la normalisée de \(\mathrm{X}_0\), et soit \(\overline{\mathrm{X}}'\) la courbe projective non singulière dont \(\mathrm{X}'\) est un ouvert dense. On désignera par \(\mathrm{X}\) (resp. par \(\overline{\mathrm{X}}\)) la courbe déduite de \(\mathrm{X}'\) (resp. de \(\overline{\mathrm{X}}'\)) en contractant en un point chacun des ensembles finis \(p^{-1}(s)\) pour \(s\in\mathrm{X}_0\).
\[
\begin{tikzcd}[column sep=2em,row sep=1.5em]
\mathrm{X}' \arrow[r,hook,"j'"'] \arrow[d,"r"']
  & \overline{\mathrm{X}}' \arrow[d,"\bar r"'] \\
\mathrm{X} \arrow[r,hook,"j"'] \arrow[d,"p_0"']
  & \overline{\mathrm{X}} \\
\mathrm{X}_0 &
\end{tikzcd}
\]

On peut caractériser \(\mathrm{X}\) et \(\overline{\mathrm{X}}\) par les propriétés suivantes.

\textit{a)} \(p_0\) est fini et radiciel, donc induit des isomorphismes
\[
\mathrm{H}^{\bullet}(\mathrm{X}_0,\mathbf{Z}_{\ell})
\longrightarrow
\mathrm{H}^{\bullet}(\mathrm{X},\mathbf{Z}_{\ell})
\qquad
(\ell\text{ premier à la caractéristique}).
\]

\textit{b)} Les singularités de \(\mathrm{X}\) sont analytiquement isomorphes à celles que présente une réunion d’axes de coordonnées dans l’espace affine.

\textit{c)} \(\overline{\mathrm{X}}\) est projective et \(\overline{\mathrm{X}}-\mathrm{X}\) est un ensemble fini de points où \(\overline{\mathrm{X}}\) est lisse.

\textbf{(10.3.2)} Le groupe de Picard \(\operatorname{Pic}(\overline{\mathrm{X}})\) se dévisse comme suit.

\textit{a)} Soit \(\mathrm{I}\) l’ensemble des composantes irréductibles de \(\overline{\mathrm{X}}\) (= l’ensemble des composantes irréductibles de \(\mathrm{X}_0\)). Pour chaque composante irréductible \(\mathrm{Y}\) de \(\overline{\mathrm{X}}\), le degré de la restriction à \(\mathrm{Y}\) d’un faisceau inversible est un homomorphisme
\[
\deg_{\mathrm{Y}}:\operatorname{Pic}(\overline{\mathrm{X}})
\longrightarrow\mathbf{Z}:
\mathcal{L}\longrightarrow\deg_{\mathrm{Y}}(\mathcal{L}).
\]
On en déduit
\[
\mathit{deg}:\operatorname{Pic}(\overline{\mathrm{X}})
\longrightarrow\mathbf{Z}^{\mathrm{I}}.
\]
Cet homomorphisme est surjectif, de noyau \(\operatorname{Pic}^{0}(\overline{\mathrm{X}})\).

\textit{b)} Le morphisme de \(\operatorname{Pic}^{0}(\overline{\mathrm{X}})\) dans la variété abélienne \(\operatorname{Pic}^{0}(\overline{\mathrm{X}}')\) est surjectif. Son noyau est un tore.

\Page{70}
\textbf{(10.3.3)} Soient \(\mathrm{S}\) l’ensemble fini \(\overline{\mathrm{X}}-\mathrm{X}\), \(\alpha_0:\mathrm{S}\to\mathrm{I}\) l’application qui à \(s\in\mathrm{S}\) associe l’unique composante irréductible de \(\overline{\mathrm{X}}\) telle que \(s\in\alpha_0(s)\) et soit \(\alpha:\mathbf{Z}^{\mathrm{S}}\to\mathbf{Z}^{\mathrm{I}}\) induit par \(\alpha_0\).

Chaque \(s\in\mathrm{S}\) définit un faisceau inversible \(\mathcal{O}(s)\), d’où \(u_0:\mathbf{Z}^{\mathrm{S}}\to\operatorname{Pic}(\overline{\mathrm{X}})\). On a \(\mathit{deg}\,u_0=\alpha\).

\textit{Définition} \textbf{(10.3.4)}. — \textit{Le \(\mathrm{H}^{1}\) \textup{motivique} de \(\mathrm{X}_0\), noté \(\mathrm{H}^{1}_{m}(\mathrm{X}_0)(1)\) est le \(1\)-motif}
\[
\operatorname{Ker}(\alpha)\longrightarrow
\operatorname{Pic}^{0}(\overline{\mathrm{X}}).
\]

Par définition, on a donc \(\mathrm{H}^{1}_{m}(\mathrm{X}_0)(1)=\mathrm{H}^{1}_{m}(\mathrm{X})(1)\).

\textit{Remarques} \textbf{(10.3.5)}. — (i) On a
\[
\begin{aligned}
\operatorname{Gr}^{\mathrm{W}}_0(\mathrm{H}^{1}_{m}(\mathrm{X})(1))
  &=\operatorname{Ker}(\alpha:\mathbf{Z}^{\mathrm{S}}\to\mathbf{Z}^{\mathrm{I}})\\
\operatorname{Gr}^{\mathrm{W}}_{-1}(\mathrm{H}^{1}_{m}(\mathrm{X})(1))
  &=\operatorname{Pic}^{0}(\overline{\mathrm{X}}').
\end{aligned}
\]

(ii) Le conoyau de \(\alpha\) est sans torsion : il s’identifie au \(\mathbf{Z}\)-module libre de base \(\mathrm{I}-\alpha_0(\mathrm{S})\). On dispose d’une suite exacte courte de complexes et d’un quasi-isomorphisme
\[
\begin{tikzcd}[column sep=small,row sep=3em]
0 \arrow[r]
 & {[\operatorname{Ker}(\alpha)\longrightarrow
       \operatorname{Pic}^{0}(\overline{\mathrm{X}})]} \arrow[r]
 & {[\mathbf{Z}^{\mathrm{S}}\longrightarrow
       \operatorname{Pic}(\overline{\mathrm{X}})]} \arrow[r]
 & {[\operatorname{Im}(\alpha)\longrightarrow\mathbf{Z}^{\mathrm{I}}]}
       \arrow[r] \arrow[d,"\wr"]
 & 0 \\
 & & & {[0\longrightarrow\operatorname{coker}(\alpha)]} &
\end{tikzcd}
\]

\textit{Construction} \textbf{(10.3.6)}. — \textit{Pour \(n\) un entier \(>0\) premier à la caractéristique de \(k\), nous construisons un isomorphisme}
\[
\mathrm{H}^{1}(\mathrm{X},\mathbf{Z}/n\mathbf{Z})(1)
=\mathrm{H}^{1}(\mathrm{X},\mu_n)
\simeq
\mathrm{T}_{\mathbf{Z}/n\mathbf{Z}}(\mathrm{H}^{1}_{m}(\mathrm{X})(1)).
\]

On sait que \(\mathrm{H}^{1}(\mathrm{X},\mu_n)\) s’identifie à l’ensemble des classes d’isomorphie de faisceaux inversibles \(\mathcal{L}\) sur \(\mathrm{X}\), munis d’un isomorphisme \(\mathcal{L}^{\otimes n}\simeq\mathcal{O}\). Puisque \(p_0:\mathrm{X}\to\mathrm{X}_0\) est fini et radiciel, on a \(\mathrm{H}^{1}(\mathrm{X},\mathbf{Z}/n)\xrightarrow{\sim}\mathrm{H}^{1}(\mathrm{X}_0,\mathbf{Z}/n)\), et \(\mathrm{H}^{1}(\mathrm{X},\mu_n)\) est encore l’ensemble des classes d’isomorphie de couples \((\mathcal{L},\alpha)\) pour \(\mathcal{L}\) un faisceau inversible sur \(\mathrm{X}_0\), et \(\alpha\) un isomorphisme \(\alpha:\mathcal{L}^{\otimes n}\simeq\mathcal{O}\).

Soit un tel couple \((\mathcal{L},\alpha)\). Le faisceau inversible \(\mathcal{L}\) se prolonge en \(\overline{\mathcal{L}}\) sur \(\overline{\mathrm{X}}\), et il existe un diviseur \(\mathrm{D}\) de support \(\mathrm{S}\) tel que \(\alpha\) se prolonge en \(\bar\alpha:\overline{\mathcal{L}}^{\otimes n}\xrightarrow{\sim}\mathcal{O}(\mathrm{D})\). S’il existe un isomorphisme \(\beta\) de \(\overline{\mathcal{L}}^{\otimes n}\) avec \(\mathcal{O}(\mathrm{E})\), cet isomorphisme est uniquement déterminé à multiplication près par un élément du groupe divisible \(\mathrm{H}^{0}(\overline{\mathrm{X}},\mathcal{O}^{*})\). On en déduit que \((\overline{\mathcal{L}},\mathrm{D})\) détermine \((\mathcal{L},\alpha)\) à isomorphisme près. Pour qu’un couple \((\overline{\mathcal{L}},\mathrm{D})\) provienne d’un couple convenable \((\mathcal{L},\alpha)\), il faut et il suffit que \(n[\overline{\mathcal{L}}]=[\mathcal{O}(\mathrm{D})]\) dans \(\operatorname{Pic}(\overline{\mathrm{X}})\). Il provient de \((\mathcal{O}_{\overline{\mathrm{X}}},0)\) si et seulement s’il est de la forme \((\mathcal{O}_{\overline{\mathrm{X}}}(\mathrm{D}),n\mathrm{D})\).

\Page{71}
Ceci identifie \(\mathrm{H}^{1}(\mathrm{X},\mu_n)\) au \(\mathrm{H}^{0}\) du complexe produit tensoriel de \([\mathbf{Z}^{\mathrm{S}}\to\operatorname{Pic}(\overline{\mathrm{X}})]\) (degrés \(0\) et \(1\)) et \([\mathbf{Z}\xrightarrow{n}\mathbf{Z}]\) (degrés \(-1\) et \(0\)) :
\[
\begin{tikzcd}[column sep=2em,row sep=2.5em]
\mathbf{Z}^{\mathrm{S}} \arrow[r,"u_0"]
 & \operatorname{Pic}(\overline{\mathrm{X}})\\
\mathbf{Z}^{\mathrm{S}} \arrow[r,"u_0"] \arrow[u,"n"']
 & \operatorname{Pic}(\overline{\mathrm{X}}) \arrow[u,"-n"']
\end{tikzcd}
\]
\[
\mathrm{H}^{1}(\mathrm{X},\mu_n)
\simeq
\mathrm{H}^{0}\bigl(
[\mathbf{Z}^{\mathrm{S}}\to\operatorname{Pic}(\overline{\mathrm{X}})]
\overset{\mathrm{L}}{\otimes}\mathbf{Z}/n\mathbf{Z}\bigr).
\]

Puisque \(\operatorname{coker}(\alpha)\) est sans torsion, on lit sur (10.3.5) (ii) que
\[
\begin{aligned}
\mathrm{H}^{1}(\mathrm{X},\mu_n)
&\simeq
\mathrm{H}^{0}\bigl(
[\mathbf{Z}^{\mathrm{S}}\to\operatorname{Pic}(\overline{\mathrm{X}})]
\overset{\mathrm{L}}{\otimes}\mathbf{Z}/n\mathbf{Z}\bigr)
\xleftarrow{\sim}
\mathrm{H}^{0}\bigl(
\mathrm{H}^{1}_{m}(\mathrm{X})(1)
\overset{\mathrm{L}}{\otimes}\mathbf{Z}/n\mathbf{Z}\bigr)\\
&\simeq
\mathrm{T}_{\mathbf{Z}/n\mathbf{Z}}(\mathrm{H}^{1}_{m}(\mathrm{X})(1)).
\end{aligned}
\]

\textit{Variante} \textbf{(10.3.7)}. — \textit{Pour \(\ell\) premier à la caractéristique \(p\) de \(k\), on a}
\[
\mathrm{H}^{1}(\mathrm{X}_0,\mathbf{Z}_{\ell}(1))
\simeq
\mathrm{T}_{\ell}(\mathrm{H}^{1}_{m}(\mathrm{X}_0)(1)).
\]

\textit{Construction} \textbf{(10.3.8)}. — \textit{Soit \(\mathrm{X}_0\) une courbe sur \(\mathbf{C}\). On construira un isomorphisme de structures de Hodge mixtes}
\[
\mathrm{H}^{1}(\mathrm{X}_0)(1)
\simeq\mathrm{T}(\mathrm{H}^{1}_{m}(\mathrm{X}_0)(1)).
\]
Nous pouvons supposer et supposerons que \(\mathrm{X}=\mathrm{X}_0\).

\textit{Construction} \textbf{(10.3.9)}. — (i) \textit{Soit \(j_*^{m}\mathcal{O}_{\mathrm{X}}^{*}\) le sous-faisceau des fonctions méromorphes dans \(j_*\mathcal{O}_{\mathrm{X}}^{*}\). On a}
\[
\mathrm{H}^{1}(\mathrm{X},\mathbf{Z}(1))
=\mathbf{H}^{1}\bigl(\overline{\mathrm{X}},
[\mathcal{O}_{\overline{\mathrm{X}}}\xrightarrow{\exp}
j_*^{m}\mathcal{O}_{\mathrm{X}}^{*}]\bigr).
\]
\[
\mathrm{H}^{1}(\mathrm{X},\mathbf{C})
=\mathbf{H}^{1}\bigl(\overline{\mathrm{X}},
[\mathcal{O}_{\overline{\mathrm{X}}}\xrightarrow{d}
r_*\Omega^{1}_{\overline{\mathrm{X}}'}(\log\mathrm{S})]\bigr).
\tag*{(ii)}
\]

(iii) \textit{L’inclusion de \(\mathrm{H}^{1}(\mathrm{X},\mathbf{Z})\) dans \(\mathrm{H}^{1}(\mathrm{X},\mathbf{C})\) est définie par le morphisme de complexes}
\[
\begin{tikzcd}[column sep=2em,row sep=2.5em]
\mathcal{O}_{\overline{\mathrm{X}}}
  \arrow[r,"\exp"] \arrow[d,"\wr"']
 & j_*^{m}\mathcal{O}_{\mathrm{X}}^{*} \arrow[d,"df/f"]\\
\mathcal{O}_{\overline{\mathrm{X}}} \arrow[r]
 & r_*\Omega^{1}_{\overline{\mathrm{X}}'}(\log\mathrm{S})
\end{tikzcd}
\]

Soit le diagramme commutatif
\[
\begin{tikzcd}[column sep=0.7em,row sep=3.5em]
\mathrm{H}^{1}(\mathrm{X},\mathbf{Z}(1))
  \arrow[r,"\sim"] \arrow[d]
  \arrow[dr,phantom,"\text{\textcircled{1}}" description]
&
\mathbf{H}^{1}(\mathrm{X},
 [\mathcal{O}_{\mathrm{X}}\xrightarrow{\exp}\mathcal{O}_{\mathrm{X}}^{*}])
  \arrow[d]
  \arrow[dr,phantom,"\text{\textcircled{2}}" description]
&
\mathbf{H}^{1}(\overline{\mathrm{X}},
 [j_*\mathcal{O}_{\mathrm{X}}\xrightarrow{\exp}j_*\mathcal{O}_{\mathrm{X}}^{*}])
  \arrow[l,"\sim"'] \arrow[d]
  \arrow[dr,phantom,"\text{\textcircled{3}}" description]
&
\mathbf{H}^{1}(\overline{\mathrm{X}},
 [\mathcal{O}_{\overline{\mathrm{X}}}\xrightarrow{\exp}
 j_*^{m}\mathcal{O}_{\mathrm{X}}^{*}])
  \arrow[l,"\sim"'] \arrow[d]
\\
\mathrm{H}^{1}(\mathrm{X},\mathbf{C}) \arrow[r]
&
\mathbf{H}^{1}(\mathrm{X},
 [\mathcal{O}_{\mathrm{X}}\xrightarrow{d}
 r_*\Omega^{1}_{\overline{\mathrm{X}}'}])
&
\mathbf{H}^{1}(\overline{\mathrm{X}},
 [j_*\mathcal{O}_{\mathrm{X}}\longrightarrow
 j_*r_*\Omega^{1}_{\mathrm{X}'}])
 \arrow[l]
&
\mathbf{H}^{1}(\overline{\mathrm{X}},
 [\mathcal{O}_{\overline{\mathrm{X}}}\longrightarrow
 r_*\Omega^{1}_{\overline{\mathrm{X}}'}(\log\mathrm{S})])
 \arrow[l,"\sim"']
\end{tikzcd}
\]

\Page{72}
Les flèches horizontales du carré 1 sont des isomorphismes, car les suites
\[
0\longrightarrow\mathbf{Z}(1)
\longrightarrow\mathcal{O}_{\mathrm{X}}
\xrightarrow{\exp}\mathcal{O}_{\mathrm{X}}^{*}
\longrightarrow0
\]
\[
0\longrightarrow\mathbf{C}
\longrightarrow\mathcal{O}_{\mathrm{X}}
\xrightarrow{d}r_*\Omega^{1}_{\mathrm{X}'}
\longrightarrow0
\]
sont exactes. Celles de 2 le sont car \(\mathrm{R}^{1}j_*\mathcal{O}_{\mathrm{X}}=0\). Celles de 3 le sont car les morphismes de complexes qui les définissent sont des quasi-isomorphismes. Ce diagramme définit (10.3.9).

\textbf{(10.3.10)} Rappelons que si \(d:\mathrm{F}\to\mathrm{G}\) est un morphisme de faisceaux abéliens sur un espace \(\mathrm{Z}\), alors \(\mathbf{H}^{1}(\mathrm{Z},[\mathrm{F}\to\mathrm{G}])\) s’identifie à l’ensemble des classes d’isomorphismes de couples \((\mathrm{P},\alpha)\), où \(\mathrm{P}\) est un \(\mathrm{F}\)-torseur (= espace principal homogène sous \(\mathrm{F}\)) et où \(\alpha\) est une trivialisation du \(\mathrm{G}\)-torseur \(d\mathrm{P}\).

En particulier :

\textit{a)} \(\mathbf{H}^{1}(\overline{\mathrm{X}},[\mathcal{O}_{\overline{\mathrm{X}}}^{*}\to r_*\Omega^{1}_{\overline{\mathrm{X}}'}(\log\mathrm{S})])\) s’identifie à l’ensemble des classes d’isomorphie de couples \((\mathcal{L},\alpha)\), pour \(\mathcal{L}\) un faisceau inversible sur \(\overline{\mathrm{X}}\), et \(\alpha\) une connexion sur \(\bar r^{*}\mathcal{L}\), holomorphe sur \(\mathrm{X}'\) et pouvant présenter un pôle simple en les points de \(\mathrm{S}\). Soit \(a\) la projection de \(\overline{\mathrm{X}}\) sur \(\operatorname{Spec}(\mathbf{C})\) (= un point). En topologie \textit{fppf}, le faisceau \(\mathrm{R}^{1}a_*([\mathcal{O}_{\overline{\mathrm{X}}}^{*}\to r_*\Omega^{1}_{\overline{\mathrm{X}}'}(\log\mathrm{S})])\) est représentable par un schéma en groupes \(\operatorname{Pic}^{\natural}(\mathrm{X})\), dont \(\mathbf{H}^{1}(\overline{\mathrm{X}},[\mathcal{O}_{\overline{\mathrm{X}}}^{*}\to r_*\Omega^{1}_{\overline{\mathrm{X}}'}(\log\mathrm{S})])\) est l’ensemble des points. \(\operatorname{Pic}^{\natural}(\mathrm{X})\) est une extension d’un sous-groupe de \(\operatorname{Pic}(\overline{\mathrm{X}})\), contenant \(\operatorname{Pic}^{0}(\overline{\mathrm{X}})\), par \(\mathrm{H}^{0}(\overline{\mathrm{X}}',\Omega^{1}_{\overline{\mathrm{X}}'}(\log\mathrm{S}))\).

On définit une application
\[
\mathbf{Z}^{\mathrm{S}}\longrightarrow\operatorname{Pic}^{\natural}(\mathrm{X})
\tag*{\textbf{(10.3.10.1)}}
\]
en associant à chaque diviseur \(\mathrm{D}\) de \(\overline{\mathrm{X}}\), concentré en \(\mathrm{S}\), le faisceau inversible \(\mathcal{O}(\mathrm{D})\).

\textit{b)} \(\mathrm{H}^{1}(\mathrm{X},\mathbf{C})=\mathbf{H}^{1}(\overline{\mathrm{X}},[\mathcal{O}_{\overline{\mathrm{X}}}\to r_*\Omega^{1}_{\overline{\mathrm{X}}'}(\log\mathrm{S})])\) est l’algèbre de Lie de \(\operatorname{Pic}^{\natural}(\mathrm{X})\). C’est l’ensemble des classes d’isomorphie de couples \((\mathrm{P},\alpha)\), pour \(\mathrm{P}\) un \(\mathcal{O}_{\overline{\mathrm{X}}}\)-torseur et \(\alpha\) une connexion comme en \textit{a)} sur \(\mathrm{P}\).

\textit{c)} \(\mathrm{H}^{1}(\mathrm{X},\mathbf{Z}(1))=\mathbf{H}^{1}(\overline{\mathrm{X}},[\mathcal{O}_{\overline{\mathrm{X}}}\to j_*^{m}\mathcal{O}_{\mathrm{X}}^{*}])\) est l’ensemble des classes d’isomorphie de couples \((\mathrm{P},\alpha)\), pour \(\mathrm{P}\) un \(\mathcal{O}_{\overline{\mathrm{X}}}\)-torseur, et \(\alpha\) un isomorphisme du faisceau inversible \(\exp(\mathrm{P})\) avec un faisceau inversible \(\mathcal{O}_{\overline{\mathrm{X}}}(\mathrm{D})\) (\(\mathrm{D}\) concentré en \(\mathrm{S}\)). L’application
\[
\operatorname{Aut}(\mathrm{P})=\mathbf{C}
\longrightarrow\operatorname{Aut}(\exp(\mathrm{P}))=\mathbf{C}^{*}
\]
est surjective. Ceci permet encore d’identifier \(\mathrm{H}^{1}(\mathrm{X},\mathbf{Z}(1))\) à l’ensemble des couples \((p,d)\), pour \(p\) une classe d’isomorphie de \(\mathcal{O}_{\overline{\mathrm{X}}}\)-torseurs, i.e. un élément de \(\operatorname{Lie}(\operatorname{Pic}\overline{\mathrm{X}})\), et \(d\in\operatorname{Ker}(\alpha)\), définissant un diviseur concentré en \(\mathrm{S}\) ayant \(\exp(p)\) pour image dans \(\operatorname{Pic}^{0}(\overline{\mathrm{X}})\). En d’autres termes, on a défini un isomorphisme
\[
\mathrm{H}^{1}(\mathrm{X},\mathbf{Z}(1))
\xrightarrow{\sim}
\mathrm{T}_{\mathbf{Z}}(\mathrm{H}^{1}_{m}(\mathrm{X})(1)).
\tag*{\textbf{(10.3.10.2)}}
\]

De la partie (i) du lemme suivant, on déduit que cet isomorphisme est compatible à la filtration par le poids.

\Page{73}
\textit{Lemme} \textbf{(10.3.11)}. — (i) \textit{On a}
\[
\mathrm{W}^{1}(\mathrm{H}^{1}(\mathrm{X},\mathbf{Z}))
=\operatorname{Im}\bigl(
\mathrm{H}^{1}(\overline{\mathrm{X}},\mathbf{Z})
\to\mathrm{H}^{1}(\mathrm{X},\mathbf{Z})\bigr)
\]
\[
\mathrm{W}^{0}(\mathrm{H}^{1}(\mathrm{X},\mathbf{Z}))
=\operatorname{Ker}\bigl(
\mathrm{H}^{1}(\overline{\mathrm{X}},\mathbf{Z})
\to\mathrm{H}^{1}(\overline{\mathrm{X}}',\mathbf{Z})\bigr).
\]

(ii) \textit{La suite spectrale définie par la filtration bête de
\([\mathcal{O}_{\overline{\mathrm{X}}}\to
\bar r_*\Omega^{1}_{\overline{\mathrm{X}}'}(\log\mathrm{S})]\)
dégénère et aboutit à la filtration de Hodge de
\(\mathbf{H}^{*}(\mathrm{X},\mathbf{C})\).}

Prouvons (i). Il suffit de prouver la première assertion : la seconde résulte de (8.2.5). Comparons les suites exactes
\[
\begin{tikzcd}[column sep=1.5em,row sep=3em]
\mathrm{H}^{1}(\overline{\mathrm{X}},\mathbf{Z})
 \arrow[r] \arrow[d]
&
\mathrm{H}^{1}(\mathrm{X},\mathbf{Z})
 \arrow[r,"\partial"] \arrow[d]
&
\mathrm{H}^{2}(\overline{\mathrm{X}}\bmod\mathrm{X},\mathbf{Z})
 \arrow[d,"\wr"]
&
\\
\mathrm{H}^{1}(\overline{\mathrm{X}}',\mathbf{Z})
 \arrow[r]
&
\mathrm{H}^{1}(\mathrm{X}',\mathbf{Z})
 \arrow[r,"\partial"]
&
\mathrm{H}^{2}(\overline{\mathrm{X}}'\bmod\mathrm{X}',\mathbf{Z})
 \arrow[r]
&
\mathrm{H}^{2}(\overline{\mathrm{X}}',\mathbf{Z})
\end{tikzcd}
\]

Il résulte de (8.3.7) et (3.2.17) que \(\mathrm{H}^{2}(\overline{\mathrm{X}}'\bmod\mathrm{X}',\mathbf{Z})\), donc \(\mathrm{H}^{2}(\overline{\mathrm{X}}\bmod\mathrm{X},\mathbf{Z})\) est purement de poids 2. Compte tenu de (8.2.4), ceci prouve (i).

Pour prouver (ii), il nous faudra revenir à la définition (8.2.1) et utiliser que le schéma simplicial \(((\mathrm{X}'/\mathrm{X})^{\Delta_n})_{n\geq0}\) est lisse (le produit fibré itéré \((n+1)\)-uple \((\mathrm{X}'/\mathrm{X})^{\Delta_n}\) est somme disjointe de \(\mathrm{X}'\), diagonal, et d’un nombre fini de points) et qu’il admet \(((\overline{\mathrm{X}}'/\overline{\mathrm{X}})^{\Delta_n})_{n\geq0}\) pour compactification lisse. Soit \(\varepsilon:((\overline{\mathrm{X}}'/\overline{\mathrm{X}})^{\Delta_n})_{n\geq0}\to\overline{\mathrm{X}}\) le morphisme d’augmentation. Soit \((\mathrm{K},\mathrm{F})\) le complexe \(\mathcal{O}\to\Omega^{1}(\log\mathrm{S})\) sur \(((\overline{\mathrm{X}}'/\overline{\mathrm{X}})^{\Delta_n})_{n\geq0}\), muni de la filtration bête. On dispose d’un morphisme d’augmentation
\[
\alpha:
[\mathcal{O}_{\overline{\mathrm{X}}}\to
 \bar r_*\Omega^{1}_{\overline{\mathrm{X}}'}(\log\mathrm{S})]
\longrightarrow s\varepsilon_{\bullet *}\mathrm{K}.
\]

Le diagramme
\[
\begin{tikzcd}[column sep=0.7em,row sep=3em]
\mathrm{H}^{1}(\mathrm{X},\mathbf{C})
 \arrow[r,phantom,"{\textstyle=}" description]
 \arrow[d,equal]
&
\mathbf{H}^{1}\bigl(
 ((\overline{\mathrm{X}}'/\overline{\mathrm{X}})^{\Delta_n})_{\bullet},
 \mathrm{K}\bigr)
&
\mathbf{H}^{1}(\overline{\mathrm{X}},s\varepsilon_{\bullet *}\mathrm{K})
 \arrow[l,"\text{\textcircled{1}}"']
\\
\mathrm{H}^{1}(\mathrm{X},\mathbf{C}) \arrow[rr,equal]
&&
\mathbf{H}^{1}\bigl(\overline{\mathrm{X}},
 [\mathcal{O}_{\overline{\mathrm{X}}}\longrightarrow
 \bar r_*\Omega^{1}_{\overline{\mathrm{X}}'}(\log\mathrm{S})]\bigr)
 \arrow[u]
\end{tikzcd}
\]
est commutatif, et 1 est un isomorphisme car les \(\varepsilon_n\) sont finis, de sorte que \(\mathrm{R}^{i}\varepsilon_{n*}=0\) pour \(i>0\). L’assertion résulte alors de ce que \(\alpha\) est un quasi-isomorphisme filtré.

\Page{74}
\textbf{(10.3.12)} Achevons la construction (10.3.8). Le diagramme
\[
\begin{tikzcd}[column sep=1.5em,row sep=3em]
\mathrm{H}^{1}(\mathrm{X},\mathbf{Z}(1))
 \arrow[r,"\sim"] \arrow[d]
&
\mathrm{T}_{\mathbf{Z}}(\mathrm{H}^{1}_{m}(\mathrm{X})(1))
 \arrow[r]
&
\operatorname{Lie}\operatorname{Pic}(\overline{\mathrm{X}})
 =\mathrm{H}^{1}(\overline{\mathrm{X}},\mathcal{O}_{\overline{\mathrm{X}}})
 \arrow[d,equal]
\\
\mathrm{H}^{1}(\mathrm{X},\mathbf{C})
 \arrow[r,"\sim"]
&
\operatorname{Lie}\operatorname{Pic}^{\natural}(\mathrm{X})
 \arrow[r]
&
\operatorname{Lie}\operatorname{Pic}(\overline{\mathrm{X}})
 =\mathrm{H}^{1}(\overline{\mathrm{X}},\mathcal{O}_{\overline{\mathrm{X}}})
\end{tikzcd}
\]
est commutatif. D’après (10.3.11) (ii) et la définition des filtrations de Hodge, (10.3.10.2) est compatible à la filtration de Hodge.

\textit{Corollaire} \textbf{(10.3.13)}. — \textit{\(\mathrm{H}^{1}_{m}(\mathrm{X})(1)^{\natural}\) est l’extension}
\[
\operatorname{Ker}(\alpha)\longrightarrow\operatorname{Pic}^{\natural}(\mathrm{X})^{0}
\]
\textit{de \(\mathrm{H}^{1}_{m}(\mathrm{X})(1)\) déduite de \textup{(10.3.10.1)}.}

On a en effet un diagramme commutatif
\[
\begin{tikzcd}[column sep=2em,row sep=3em]
\operatorname{Ker}(\alpha) \arrow[r]
&
\operatorname{Pic}^{\natural}(\mathrm{X})^{0} \arrow[r]
&
\operatorname{Pic}^{0}(\overline{\mathrm{X}})
\\
\mathrm{H}^{1}(\mathrm{X},\mathbf{Z}) \arrow[r] \arrow[u]
&
\operatorname{Lie}\operatorname{Pic}^{\natural}(\mathrm{X})
 \arrow[r] \arrow[u]
&
\operatorname{Lie}\operatorname{Pic}(\overline{\mathrm{X}}), \arrow[u]
\end{tikzcd}
\]
et on applique (10.1.9).

\textit{Exercice} \textbf{(10.3.14)}. — Vérifier que le diagramme
\[
\begin{tikzcd}[column sep=2em,row sep=3em]
\mathrm{H}^{1}(\mathrm{X},\mathbf{Z}(1))
 \arrow[r,equal] \arrow[d]
&
\mathrm{T}_{\mathbf{Z}}(\mathrm{H}^{1}_{m}(\mathrm{X})(1))
 \arrow[d]
\\
\mathrm{H}^{1}(\mathrm{X},\mu_n) \arrow[r,equal]
&
\mathrm{T}_{\mathbf{Z}/n\mathbf{Z}}(\mathrm{H}^{1}_{m}(\mathrm{X})(1))
\end{tikzcd}
\]
est commutatif.

\textit{Remarque} \textbf{(10.3.15)}. — L’isomorphisme déduit de (10.3.8) et (10.3.9) (ii), entre \(\mathbf{H}(\overline{\mathrm{X}},\mathcal{O}_{\overline{\mathrm{X}}}\to r_*\Omega^{1}_{\overline{\mathrm{X}}'}(\log\mathrm{S}))\) (qu’on peut appeler \(\mathrm{H}^{1}_{\mathrm{DR}}(\mathrm{X})\)) et \(\mathrm{T}_{\mathrm{DR}}(\mathrm{H}^{1}_{m}(\mathrm{X})(1))\), égal à \(\operatorname{Lie}\operatorname{Pic}^{\natural}(\mathrm{X})\) d’après (10.3.13), a été défini de façon purement algébrique.

\Page{75}
\section*{10.4. Traduction d’un théorème de Picard.}

\textbf{(10.4.1)} Soit \(\mathcal{H}\) une structure de Hodge mixte telle que \(h^{pq}=0\) pour \(p\) ou \(q<0\) (resp. pour \(p\) ou \(q>n\)). Je noterai provisoirement \(\mathrm{I}(\mathcal{H})\) (resp. \(\mathrm{II}_n(\mathcal{H})\)) le 1-motif défini par la plus grande sous-structure de Hodge mixte de \((\mathcal{H}_{\mathbf{Z}}/\mathrm{torsion})(1)\) (resp. par le plus grand quotient de \((\mathcal{H}_{\mathbf{Z}}/\mathrm{torsion})(n)\)) qui soit purement de type
\[
\{(-1,-1),(-1,0),(0,-1),(0,0)\}.
\]
Si \(\mathrm{X}\) est une variété algébrique complexe de dimension \(\leq\mathrm{N}\), je conjecture que les 1-motifs \(\mathrm{I}(\mathrm{H}^{n}(\mathrm{X},\mathbf{Z}))\), \(\mathrm{II}_n(\mathrm{H}^{n}(\mathrm{X},\mathbf{Z}))\) (pour \(n\leq\mathrm{N}\)) et \(\mathrm{II}_{\mathrm{N}}(\mathrm{H}^{n}(\mathrm{X},\mathbf{Z}))\) (pour \(n\geq\mathrm{N}\)) admettent une définition purement algébrique. Les morphismes
\[
\begin{aligned}
\mathrm{T}_{\ell}(\mathrm{I}(\mathrm{H}^{n}(\mathrm{X},\mathbf{Z})))
 &\longrightarrow
 (\mathrm{H}^{n}(\mathrm{X},\mathbf{Z}_{\ell})/\mathrm{torsion})(1)\\
\mathrm{T}_{\ell}(\mathrm{II}_n(\mathrm{H}^{n}(\mathrm{X},\mathbf{Z})))
 &\longleftarrow
 \mathrm{H}^{n}(\mathrm{X},\mathbf{Z}_{\ell})(n)
 &&(\text{pour }n\leq\mathrm{N})\\
\mathrm{T}_{\ell}(\mathrm{II}_{\mathrm{N}}(\mathrm{H}^{n}(\mathrm{X},\mathbf{Z})))
 &\longleftarrow
 \mathrm{H}^{n}(\mathrm{X},\mathbf{Z}_{\ell})(\mathrm{N})
 &&(\text{pour }n\geq\mathrm{N})
\end{aligned}
\]
et leurs analogues en cohomologie de De Rham devraient aussi admettre une définition purement algébrique. C’est ce que nous avons vérifié pour \(\dim(\mathrm{X})=1\). Voici un autre exemple.

\textbf{(10.4.2)} Soient \(\mathrm{S}\) une surface projective et lisse, \(\mathrm{U}\) un ouvert (de Zariski) dense de \(\mathrm{S}\) et \(\mathrm{C}=\mathrm{S}-\mathrm{U}\). Posons
\[
\mathrm{H}^{\log}_{2}(\mathrm{U},\mathbf{Z})
=\operatorname{Ker}\bigl(
\mathrm{H}_{2}(\mathrm{U},\mathbf{Z})
\to\mathrm{H}_{2}(\mathrm{S},\mathbf{Z})\bigr).
\]
D’après [8] \S~9, ce groupe ne dépend pas de la compactification lisse choisie \(\mathrm{S}\) de \(\mathrm{U}\). Modulo groupes finis, ce n’est autre (3.2.17) que \(\mathrm{W}_{-3}\mathrm{H}_{2}(\mathrm{U},\mathbf{Z})\) (\(\mathrm{H}_{2}(\mathrm{U},\mathbf{Q})\), dual de \(\mathrm{H}^{2}(\mathrm{U},\mathbf{Q})\), est muni de la structure de Hodge mixte duale). On peut vérifier que le 1-motif défini par \((\mathrm{H}^{\log}_{2}(\mathrm{U},\mathbf{Z})/\mathrm{torsion})(-1)\) admet la description suivante :

\textit{a)} Par dualité de Poincaré, on a \(\mathrm{H}_{2}(\mathrm{U},\mathbf{Z})=\mathrm{H}^{2}_{c}(\mathrm{U},\mathbf{Z})\). La suite exacte longue de cohomologie de \((\mathrm{S},\mathrm{C})\) fournit une suite exacte
\[
\mathrm{H}^{1}(\mathrm{S},\mathbf{Z})
\longrightarrow\mathrm{H}^{1}(\mathrm{C},\mathbf{Z})
\longrightarrow\mathrm{H}^{\log}_{2}(\mathrm{U},\mathbf{Z})
\longrightarrow0.
\]

\textit{b)} Soit \(\mathrm{H}^{1}_{m}(\mathrm{C})\) le quotient de \(\operatorname{Pic}^{0}(\mathrm{C})\) par son radical unipotent. On sait que \(\mathrm{H}^{1}(\mathrm{C},\mathbf{Z})=\mathrm{T}_{\mathbf{Z}}(\mathrm{H}^{1}_{m}(\mathrm{C}))\) et \(\mathrm{H}^{1}(\mathrm{S},\mathbf{Z})=\mathrm{T}_{\mathbf{Z}}(\operatorname{Pic}^{0}(\mathrm{S}))\). On dispose donc d’une application
\[
\mathrm{H}^{\log}_{2}(\mathrm{U},\mathbf{Z})
\longrightarrow
\mathrm{T}_{\mathbf{Z}}(\mathrm{H}^{1}_{m}(\mathrm{C})/\operatorname{Pic}^{0}(\mathrm{S})),
\]
qui est un isomorphisme modulo groupes finis.

\textit{c)} L’application déduite de \textit{b)}
\[
(\mathrm{H}^{\log}_{2}(\mathrm{U},\mathbf{Z})/\mathrm{torsion})(-1)
\longrightarrow
\mathrm{T}_{\mathbf{Z}}(\mathrm{H}^{1}_{m}(\mathrm{C})/\operatorname{Pic}^{0}(\mathrm{S}))
\tag*{\textbf{(10.4.2.1)}}
\]
est un morphisme de structures de Hodge mixte (et un isomorphisme modulo groupes finis).

\Page{76}
On peut définir de façon purement algébrique les isomorphismes (modulo groupes finis) déduits de (10.4.2.1) entre les analogues de \(\mathrm{H}^{\log}_{2}(\mathrm{U},\mathbf{Z})(-1)\) en cohomologie \(\ell\)-adique ou de De Rham, et les réalisations \(\ell\)-adiques ou de De Rham de \(\mathrm{H}^{1}_{m}(\mathrm{C})/\operatorname{Pic}^{0}(\mathrm{S})\).

\textbf{(10.4.3)} Nous ne démontrerons pas les constructions et compatibilités précédentes. Le résultat est en germe dans le théorème suivant de E. Picard ([11], tome I, p. 62). \(\mathrm{P}\) et \(\mathrm{A}\) y désignent des polynômes.

\textit{Tout résidu de l’intégrale double}
\[
\iint\frac{\mathrm{P}(x,y)\,dx\,dy}{\mathrm{A}(x,y)}
\]
\textit{peut être regardé comme une période logarithmique ou cyclique d’une intégrale abélienne.}

Je rappelle que :

\textit{a)} Une \textit{intégrale \(k\)-uple} sur une variété projective non singulière \(\mathrm{V}\) de dimension \(d\) est une \(k\)-forme rationnelle fermée sur \(\mathrm{V}\) : c’est un élément \(\alpha\in\mathrm{H}^{0}(\mathrm{U},\Omega^{k})\), pour \(\mathrm{U}\) un ouvert dense de \(\mathrm{V}\), qui vérifie \(d\alpha=0\). Ici, \(k=2\) et \(\mathrm{V}=\mathbf{P}^{2}\), de sorte que \(k=\dim\mathrm{V}\) et que la condition \(d\alpha=0\) est vide.

\textit{b)} Une intégrale \(k\)-uple \(\alpha\) définit un élément encore noté \(\alpha\) de \(\mathrm{H}^{k}(\mathrm{U},\mathbf{C})\). Les \textit{périodes} de \(\alpha\) sont les nombres \(\langle\alpha,c\rangle\) pour \(c\in\mathrm{H}_{k}(\mathrm{U},\mathbf{Z})\). Si \(c\) est défini par une chaîne singulière \(\mathrm{C}\), on a
\[
\langle\alpha,c\rangle=\int_{\mathrm{C}}\alpha.
\]

\textit{c)} Les \textit{résidus} d’une intégrale \(k\)-uple \(\alpha\in\mathrm{H}^{0}(\mathrm{U},\Omega^{k})\) sont les nombres \(\langle\alpha,c\rangle\) pour
\[
c\in\operatorname{Ker}\bigl(
\mathrm{H}_{k}(\mathrm{U},\mathbf{Z})
\to\mathrm{H}_{k}(\mathrm{V},\mathbf{Z})\bigr).
\]

Dans le cas considéré par Picard, l’application \(\mathrm{H}_{2}(\mathrm{U},\mathbf{Z})\to\mathrm{H}_{2}(\mathbf{P}^{2},\mathbf{Z})\) est nulle et il n’y a pas lieu de distinguer entre périodes et résidus. Pour le cas général, cf. Lefschetz, [10], note I.

\textit{d)} Par « période logarithmique ou cyclique d’une intégrale abélienne », Picard entend simplement « période d’une intégrale simple \(\beta\) sur une courbe \(\mathrm{Y}\) ». Il est sous-entendu que \(\beta\) et \(\mathrm{Y}\) sont déterminés rationnellement par \(\mathrm{P}\) et \(\mathrm{A}\).

\textbf{(10.4.4)} Traduisons ce théorème. Soient donc \(\mathrm{S}\) et \(\mathrm{U}\) comme en (10.4.2) (par exemple \(\mathrm{S}=\mathbf{P}^{2}(\mathbf{C})\)), une intégrale double \(\beta\in\mathrm{H}^{0}(\mathrm{U},\Omega^{2}_{\mathrm{U}})\) et un cycle \(c\in\mathrm{H}^{\log}_{2}(\mathrm{U},\mathbf{Z})\). Par définition
\[
\langle\beta,c\rangle=\int_c\beta
\]
est un résidu de \(\beta\). Soit \(\mathrm{A}=\mathrm{H}^{1}_{m}(\mathrm{C})(1)/\operatorname{Pic}^{0}(\mathrm{S})\) (une extension d’une variété abélienne par un tore). La 2-forme \(\beta\) définit une classe de cohomologie \(\operatorname{cl}(\beta)\in\mathrm{H}^{2}_{\mathrm{DR}}(\mathrm{U})\) et induit
\Page{77}
une forme linéaire \(\beta'\) sur \(\mathrm{H}^{\log}_{\mathrm{DR}2}(\mathrm{U})=\mathrm{T}_{\mathrm{DR}}(\mathrm{A})\). La correspondance \(\beta\mapsto\beta'\) est rationnelle. La classe \(c\) définit \(c'\in\mathrm{T}_{\mathbf{Z}}(\mathrm{A})\otimes\mathbf{Q}\), et on a
\[
\langle\beta,c\rangle=2\pi i\langle\beta',c'\rangle.
\tag*{\textbf{(10.4.4.1)}}
\]
On peut représenter \(c'\) par un cycle sur \(\mathrm{C}\), et \(\beta'\) comme une intégrale simple sur \(\mathrm{C}\).

{\small
\begin{center}
\textbf{BIBLIOGRAPHIE}
\end{center}

\noindent [1] A. \textsc{Blanchard}, Sur les variétés analytiques complexes, \textit{Ann. Sc. E.N.S.}, \textbf{73} (1956).

\noindent [2] T. \textsc{Bloom} and M. \textsc{Herrera}, De Rham cohomology of an analytic space, \textit{Inv. Math.}, \textbf{7} (1969), 275-296.

\noindent [3] P. \textsc{Deligne}, Théorème de Lefschetz et critères de dégénérescence de suites spectrales, \textit{Publ. Math. I.H.E.S.}, \textbf{35} (1968), 107-126.

\noindent [4] P. \textsc{Deligne}, Théorie de Hodge I, \textit{Actes du Congrès international des Mathématiciens (Nice, 1970)}, Gauthier-Villars, 1971, \textbf{1}, 425-430.

\noindent [5] P. \textsc{Deligne}, Théorie de Hodge II, \textit{Publ. Math. I.H.E.S.}, \textbf{40} (1971), 5-58.

\noindent [6] P. A. \textsc{Griffiths}, Periods of integrals on algebraic manifolds. Summary of main results and discussion of open problems, \textit{Bull. Am. Math. Soc.}, \textbf{76} (1970), 228-296.

\noindent [7] P. A. \textsc{Griffiths}, On the periods of certain rational integrals : I, \textit{Annals of Math.}, \textbf{90} (1969), 460-495.

\noindent [8] A. \textsc{Grothendieck}, Le groupe de Brauer III : Exemples et compléments, dans \textit{Dix exposés sur la cohomologie des schémas}, North Holland Publ. Co., 88-188.

\noindent [9] A. \textsc{Grothendieck}, On the De Rham cohomology of algebraic varieties, \textit{Publ. Math. I.H.E.S.}, \textbf{29} (1966), 96-103.

\noindent [10] S. \textsc{Lefschetz}, \textit{L’analysis situs et la géométrie algébrique}, Gauthier-Villars, 1924 (reproduit dans \textit{Selected Papers}, Chelsea Publ. Co., 1971).

\noindent [11] E. \textsc{Picard} et G. \textsc{Simart}, \textit{Théorie des fonctions algébriques de deux variables indépendantes}, vol. I et II, Paris, Gauthier-Villars, 1895, 1906.

\noindent [12] D. G. \textsc{Quillen}, \textit{Notes on the homology of commutative rings}, M.I.T., 1968.

\noindent [13] B. \textsc{Saint-Donat} (d’après des notes de P. \textsc{Deligne}), Techniques de descente cohomologique, S.G.A. \textbf{4} V \textit{bis} (à paraître avec S.G.A. \textbf{4}, dans les \textit{Lecture Notes}, Springer-Verlag).

\noindent [14] G. \textsc{Segal}, Classifying spaces and spectral sequences, \textit{Publ. Math. I.H.E.S.}, \textbf{34} (1968), 105-112.

\noindent [15] J.-P. \textsc{Serre}, Géométrie algébrique et géométrie analytique, \textit{Ann. Inst. Fourier}, \textbf{6} (1956), cité G.A.G.A.

\begin{center}
\textbf{SIGLES}
\end{center}

\noindent S.G.A., Séminaire de Géométrie algébrique du Bois-Marie, \textbf{1} à \textbf{7}, diffusé par l’I.H.E.S.; \textbf{1}, \textbf{3}, \textbf{4}, \textbf{6}, \textbf{7} I sont parus aux \textit{Springer Lecture Notes}, n° \textbf{224}; \textbf{151}, \textbf{152} et \textbf{153}; \textbf{269}, \textbf{270}, \textbf{305}; \textbf{225}, \textbf{288} — S.G.A. \textbf{2} a été publié par North Holland Publ. Co.
\par
}

\begin{flushright}
\textit{Manuscrit reçu le 2 mars 1972.}
\end{flushright}

\end{document}
